% \iffalse
%
% Trademarks appear throughout this documentation without any trademark
% symbol, so you can't assume that a name is free. There is no intention
% of infringement; the usage is to the benefit of the trademark owner.
%
%
%  S O F T W A R E   L I C E N S E
% =================================
%
% The files  listings.dtx  and  listings.ins  and all files generated
% from only these two files are referred to as `the listings package'
% or simply `the package'. lstdrvrs.dtx  and the files generated from
% that file are `drivers'.
%
% The listings package is copyright 1996--2004 Carsten Heinz, and
% continued maintenance on the package is copyright 2006--2007 Brooks
% Moses. From 2013 on copyright is Jobst Hoffmann, who is the maintainer
% since july 2013. The drivers are copyright 1997/1998/1999/2000/2001/
% 2002/2003/2004/2006/2007/2013 any individual author listed in the
% driver files.
%
% The listings package and its drivers may be distributed and/or modified
% under the conditions of the LaTeX Project Public License, either version
% 1.3c of this license or (at your option) any later version.
% The latest version of this license is in
%   http://www.latex-project.org/lppl.txt
% and version 1.3c or later is part of all distributions of LaTeX
% version 2003/12/01 or later.
%
% The package has the LPPL maintenance status "maintained".
%
% The Current Maintainer is Jobst Hoffmann <j.hoffmann(at)fh-aachen.de>.
%
% end of software license
%
%
%<*driver>
\documentclass[a4paper]{ltxdoc}
% \DisableCrossrefs
% \OnlyDescription

\usepackage{lstdoc}             % the basic support package for generating
                                % the documentation of this package
\usepackage[UTF8]{ctex}
\usepackage{textcomp}           % needed for upquote apostrophes
\usepackage{mdframed}           % frames for external files
\usepackage{moreverb}           % writing external files
\usepackage{xcolor}             % because of coloring the background
\usepackage{tablefootnote}      % because of \tablefootnote

\makeindex
\PageIndex
% 补丁：给 lstdoc 直写 \index 的键名/外观条目补 hdpindex 封装，使这些条目页码也可点击
\makeatletter
\def\SpecialMainAspectIndex#1{%
    \@bsphack
    \index{aspects:\levelchar\protect\aspectname{#1}\encapchar hdpindex{main}}%
    \@esphack}
\def\SpecialMainKeyIndex#1{%
    \@bsphack
    \index{keys\levelchar\protect\keyname{#1}\encapchar hdpindex{main}}%
    \@esphack}
\def\lstaspectindex#1#2{%
    \global\@namedef{lstkandc@\string#1}{}%
    \@ifundefined{lstisaspect@\string#1}
        {\index{unknown\levelchar
                \protect\texttt{\protect\string\string#1}#2\encapchar hdpindex{main}}}%
        {\index{\@nameuse{lstisaspect@\string#1}\levelchar
                \protect\texttt{\protect\string\string#1}#2\encapchar hdpindex{main}}}%
}
\makeatother
\renewcommand{\indexname}{索引}
\renewcommand{\lstlistingname}{列表}
\renewcommand{\lstlistlistingname}{列表清单}

\begin{document}
    \DocInput{listings-zh.dtx}
\end{document}
%</driver>
% \fi
%
%^^A
%^^A  Command/key to aspect relation
%^^A ================================
%^^A
%\lstisaspect[strings]{string,morestring,deletestring,stringstyle,showstringspaces}
%\lstisaspect[comments]{comment,morecomment,deletecomment,commentstyle}
%\lstisaspect[comment styles]{b,d,l,n,s,ib,id,il,in,is}
%\lstisaspect[pod]{printpod,podcomment}
%\lstisaspect[escape]{texcl,escapebegin,escapeend,escapechar,escapeinside,mathescape}
%\lstisaspect[keywords]{sensitive,classoffset,keywords,morekeywords,^^A
%   deletekeywords,keywordstyle,ndkeywords,morendkeywords,deletendkeywords,^^A
%   ndkeywordstyle,keywordsprefix,otherkeywords}
%\lstisaspect[emph]{emph,moreemph,deleteemph,emphstyle}
%\lstisaspect[tex]{texcs,moretexcs,deletetexcs,texcsstyle}
%\lstisaspect[directives]{directives,moredirectives,deletedirectives,directivestyle}
%\lstisaspect[html]{tag,usekeywordsintag,tagstyle,markfirstintag}
%\lstisaspect[keywordcomments]{keywordcomment,morekeywordcomment,^^A
%   deletekeywordcomment,keywordcommentsemicolon}
%\lstisaspect[index]{index,moreindex,deleteindex,indexstyle,\string\lstindexmacro}
%\lstisaspect[procnames]{procnamestyle,indexprocnames,procnamekeys,^^A
%   moreprocnamekeys,deleteprocnamekeys}
%\lstisaspect[style]{style,\string\lstdefinestyle,\string\lst@definestyle,^^A
%   \string\lststylefiles}
%\lstisaspect[language]{language,alsolanguage,defaultdialect,^^A
%   \string\lstalias,\string\lstdefinelanguage,\string\lst@definelanguage,^^A
%   \string\lstloadlanguages,\string\lstlanguagefiles}
%\lstisaspect[formats]{format,fmtindent,\string\lstdefineformat,^^A
%   \string\lst@defineformat,\string\lstformatfiles}
%\lstisaspect[labels]{numbers,numberstyle,numbersep,stepnumber,^^A
%   numberblanklines,firstnumber,\string\thelstnumber,numberfirstline}
%\lstisaspect[lineshape]{xleftmargin,xrightmargin,resetmargins,linewidth,^^A
%   lineskip,breaklines,breakindent,breakautoindent,prebreak,postbreak,^^A
%   breakatwhitespace}
%\lstisaspect[frames]{framexleftmargin,framexrightmargin,framextopmargin,^^A
%   framexbottommargin,backgroundcolor,fillcolor,rulecolor,rulesepcolor,^^A
%   rulesep,framerule,framesep,frameshape,frameround,frame}
%\lstisaspect[make]{makemacrouse}
%\lstisaspect[fancyvrb]{fancyvrb,fvcmdparams,morefvcmdparams}
%\lstisaspect[lgrind]{lgrindef,\string\lstlgrindeffile}
%\lstisaspect[hyper]{hyperref,morehyperref,deletehyperref,hyperanchor,hyperlink}
%\lstisaspect[kernel]{basewidth,fontadjust,columns,flexiblecolumns,identifierstyle,^^A
%   tabsize,showtabs,tab,showspaces,keepspaces,formfeed,SelectCharTable,^^A
%   MoreSelectCharTable,extendedchars,alsoletter,alsodigit,alsoother,excludedelims,^^A
%   literate,basicstyle,print,firstline,lastline,linerange,^^A
%   consecutivenumbers,nolol,captionpos,abovecaptionskip,^^A
%   belowcaptionskip,label,title,caption,\string\lstlistingname,^^A
%   \string\lstlistingnamestyle,boxpos,float,^^A
%   floatplacement,aboveskip,belowskip,everydisplay,showlines,emptylines,gobble,name,^^A
%   \string\lstname,\string\lstlistlistingname,\string\lstlistoflistings,^^A
%   \string\lstnewenvironment,\string\lstinline,\string\lstinputlisting,lstlisting,^^A
%   \string\lstloadaspects,\string\lstset,\string\thelstlisting,\string\lstaspectfiles,^^A
%   inputencoding,inputpath,delim,moredelim,deletedelim,upquote,numberbychapter,^^A
%   \string\lstMakeShortInline,\string\lstDeleteShortInline}
%\lstisaspect[doc]{lstsample,lstxsample}^^A environment
%\lstisaspect[experimental]{includerangemarker,rangebeginprefix,rangebeginsuffix,^^A
%   rangeendprefix,rangeendsuffix,rangeprefix,rangesuffix}
%
%^^A
%^^A  The long awaited beginning of documentation
%^^A =============================================
%^^A
%\newbox\abstractbox
%\setbox\abstractbox=\vbox{%
%	\begin{abstract}
%	\packagename{listings} 宏包是一个用于 \LaTeX 的源代码打印工具。
%	你可以排版独立文件，也可以用类似于 \texttt{verbatim} 的环境排版列表，
%	还可以用类似于 |\verb| 的命令打印代码片段。
%	有许多参数控制输出，如果你喜欢的编程语言尚未受支持，
%	你可以自己定义。
%   你可以自己定义。
%	\end{abstract}}
%
% \title{\vspace*{-2\baselineskip}\packagename{Listings} 宏包}
% \author{Copyright 1996--2004, Carsten Heinz%
%    \\ Copyright 2006--2007, Brooks Moses
%    \\ Copyright 2013--, Jobst Hoffmann
%    \\ 维护者：Jobst Hoffmann\thanks{Jobst %
%       Hoffmann 自 2013 年起担任 \packagename{listings}
%       宏包的维护者；详情见前言。}~ %
%    \textless\lstemail\textgreater}
% \date{2025/11/14\enspace\enspace Version 1.11b\ \box\abstractbox}
% \def\lstemail{\href{mailto:j.hoffmann@fh-aachen.de}{\texttt{j.hoffmann(at)fh-aachen.de}}}
% \ifhyper
%   \hypersetup{pdftitle=listings 宏包 v1.11,
%               pdfsubject=宏包指南,
%               pdfauthor=Jobst Hoffmann <j.hoffmann(at)fh-aachen.de>,%
%               pdfkeywords={源代码格式化工具, 编程语言}}
% \fi
%
% \csname @twocolumntrue\endcsname
% \maketitle
%^^A \enlargethispage{2\baselineskip}
% \csname @starttoc\endcsname{toc}
% \onecolumn
%
%
% \section*{前言}
%
% \paragraph{宏包维护权的移交}
% 2004 年末，Carsten Heinz 在发布 \packagename{listings} 宏包 1.3b 版后不久便
% 与 \TeX\ 界失去了联系。在多次尝试联系他未果后，Hendri Adriaens 按照 LPPL
% 关于无人维护宏包的程序接管了该宏包的维护工作。随后他将维护权移交给了
% Brooks Moses，后者在该程序进行期间主动请缨。其结果就是 listings 1.4 版。
%
% 本次发布的 1.5 版是我接任维护后的一次小型维护版本。我要感谢 Stephan
% Hennig，他为 Lua 语言定义提供了支持。正是他要求集成一门新语言，并推动我
% 成为本宏包的维护者。
%
%
% \paragraph{新闻与变更}
% 1.5 版是第五个错误修复版本。本版本没有其他变更，只有两项扩展：
% 对现代 Fortran（2003、2008）和 Lua 的支持。
%
%
% \vfill
% \paragraph{致谢}
% 有许多人我要感谢，他们与我进行了富有成效的交流、提出想法、提交错误报告、
% 向 \texttt{lstdrvrs.dtx} 添加编程语言等等。他们的名字列在
% 第~\ref{uClosingAndCredits} 节。
%
% \paragraph{商标}
% 本文档中随处出现的商标均未标注商标符号；它们归各自的商标所有者所有。
% 无意侵权；此处使用系为商标所有者之利益。
%
%
% \clearpage
%
%
% \part{用户指南}
%
%
% \section{入门}\label{uGettingStarted}
%
%
% \subsection{最小文件}\label{uAMinimalFile}
%
% 在使用 \packagename{listings} 宏包之前，你应该熟悉 \LaTeX\ 排版系统。
% 你不需要成为专家。下面是一个使用 \packagename{listings} 的最小文件。
% \begin{verbatim}
%    \documentclass{article}
%    \usepackage{listings}
%
%    \begin{document}
%    \lstset{language=Pascal}
%
%      % Insert Pascal examples here.
%
%    \end{document}\end{verbatim}
% 现在输入这个第一个例子，用 \LaTeX\ 编译看看。
% \begin{advise}
% \item 我真的必须这么做吗？
%       \advisespace
%       既是也不是。有些编程书说这样做有好处。大错特错！输入代码要花时间——
%       如果代码对你来说已经很清楚，那时间就浪费了。而如果你需要这些时间
%       来理解发生了什么，那这本书的作者应该重新考虑其呈现关键内容的方式——
%       你甚至可以这样评价本指南——或者你只是缺乏编程经验。如果只是后一种
%       情况，你应该花更多时间阅读（好的）编程书籍、（好的）文档和别人的
%       （好的）源代码。当然你也应该自己做实验。你会学到很多东西。不过，
%       用 \LaTeX\ 编译这个例子能验证 \packagename{listings} 宏包是否正确安装。
% \item 例子不工作。
%       \advisespace
%       你的系统上是否安装了 \packagename{listings} 和 \packagename{keyval}
%       这两个宏包？请按顺序查阅：你的 \TeX\ 发行版的管理工具、系统管理员、
%       本地 \TeX\ 和 \LaTeX\ 指南、\TeX\ FAQ，以及第~\ref{rInstallation} 节。
%       如果你查遍了\emph{所有}这些来源仍然束手无策，可以考虑在
%       \texttt{comp.text.tex} 之类的 \TeX\ 新闻组发帖求助。
% \item 使用宏包前我应该先阅读软件许可吗？
%       \advisespace
%       应该，但先阅读本\emph{入门}节来决定你是否愿意使用这个宏包。^^A ;-)
% \end{advise}
%
%
% \subsection{排版源代码列表}
%
% 支持三种源代码：代码片段、代码段和独立文件的列表。代码片段放在段落内部，
% 其余两种作为独立段落——区别和行内公式与行间公式的区别相同。
% \begin{advise}
% \item 无论你有哪种源代码，如果列表中包含 \'e、\L、\"a 等国家或地区字符，
%       你必须告诉宏包！第~\lstref{uSpecialCharacters} 节讨论了这个问题。
% \end{advise}
%
% \paragraph{代码片段}
% 众所周知的 \LaTeX\ 命令 |\verb| 逐字排版代码片段。新命令 |\lstinline|
% 则对代码进行美化排版，例如 `\lstinline!var i:integer;!' 排版为
% `{\rstyle|\lstinline|}!var i:integer;!'。感叹号用作代码的定界符，
% 可以替换为代码中不出现的任意字符；|\lstinline$var i:integer;$|
% 给出相同结果。
%
% \paragraph{行间代码}
% \texttt{lstlisting} 环境排版其中包含的源代码。和大多数例子一样，
% 下面这个例子在右侧显示逐字的 \LaTeX\ 代码，左侧显示排版结果。
% 你可以把右侧的内容放到最小文件中，用 \LaTeX\ 编译看看。
% \begin{lstsample}[lstlisting]{}{}
%    \begin{lstlisting}
%    for i:=maxint to 0 do
%    begin
%        { do nothing }
%    end;
%
%    Write('Case insensitive ');
%    WritE('Pascal keywords.');
%    \end{lstlisting}
% \end{lstsample}
% 再简单不过了。
% \begin{advise}
% \item 才不是呢。`\texttt{listing}' 这个名字更短。
%       \advisespace
%       确实。但其他宏包已经定义了同名的环境。为了与这些宏包兼容，
%       \packagename{listings} 宏包的所有命令和环境都使用前缀 `\texttt{lst}'。
% \end{advise}
% 该环境接受一个可选参数。它告诉宏包执行特殊任务，例如只打印第 2--5 行：
% \begin{lstsample}{\lstset{frame=trbl,framesep=0pt}\label{gFirstKey=ValueList}}{}
%    \begin{lstlisting}[firstline=2,
%                       lastline=5]
%    for i:=maxint to 0 do
%    begin
%        { do nothing }
%    end;
%
%    Write('Case insensitive ');
%    WritE('Pascal keywords.');
%    \end{lstlisting}
% \end{lstsample}
% \begin{advise}
% \item 等等！边框从哪来的，有什么用？
%       \advisespace
%       你可以给除代码片段外的所有列表加边框。后面会讲怎么做。边框表明
%       列表末尾的空行不会被打印。本例中就是第 5 行。
% \item 喂，你不能丢掉我的空行！
%       \advisespace
%       你可以告诉宏包不要丢掉它们：键 `\ikeyname{showlines}' 控制这些空行，
%       在第~\ref{rTypesettingListings} 节中描述。警告：请先阅读后面关于
%       如何使用键的一般说明。
% \item 使用 `\ikeyname{firstline}' 时出现莫名其妙的错误信息。
%       \advisespace
%       这不应该发生。请按照第~\lstref{uTroubleshooting} 节的说明提交错误报告。
% \end{advise}
%
% \paragraph{独立文件}
% 最后我们介绍 |\lstinputlisting|，这是用于美化排版独立文件的命令。
% 它有一个可选参数和一个文件名参数。注意你可能需要指定文件的相对路径。
% 这里结果打印在逐字代码下方，因为两者放在一起会超出文本宽度。
% \begin{lstsample}{\lstset{comment=[l]\%,columns=fullflexible}}^^A
%       {\lstset{alsoletter=\\,emph=\\lstinputlisting,emphstyle=\rstyle}^^A
%       \lstaspectindex{\lstinputlisting}{}}
%    \lstinputlisting[lastline=4]{listings.sty}
% \end{lstsample}
% \begin{advise}
% \item 这个例子的间距不一样。
%       \advisespace
%       是的。前两个例子使用对齐列，即相同编号的列具有相同的水平位置——
%       本宏包只做微小调整。本例中的列不对齐。这在
%       第~\ref{uFixedAndFlexibleColumns} 节（关键词：全灵活列格式）中有解释。
% \end{advise}
%
% 现在你已经了解了所有美化排版命令和环境。接下来要学习控制
% \packagename{listings} 宏包工作方式的参数。不过这才是主要任务。
% 下面介绍其中一些参数。
%
%
% \subsection{调整外观}\label{gFigureOutTheAppearance}
%
% 关键词排版为粗体，评论排版为斜体，字符串中的空格显示为 \textvisiblespace。
% 你不喜欢这些设置？看看这个：
%\ifcolor
% \begin{lstxsample}[basicstyle,keywordstyle,identifierstyle,commentstyle,^^A
%       stringstyle,showstringspaces]
%    \lstset{% general command to set parameter(s)
%        basicstyle=\small,          % print whole listing small
%        keywordstyle=\color{black}\bfseries\underbar,
%                                    % underlined bold black keywords
%        identifierstyle=,           % nothing happens
%        commentstyle=\color{white}, % white comments
%        stringstyle=\ttfamily,      % typewriter type for strings
%        showstringspaces=false}     % no special string spaces
% \end{lstxsample}
%\else
% \begin{lstxsample}[basicstyle,keywordstyle,identifierstyle,commentstyle,^^A
%       stringstyle,showstringspaces]
%    \lstset{% general command to set parameter(s)
%        basicstyle=\small,          % print whole listing small
%        keywordstyle=\bfseries\underbar,
%                                    % underlined bold keywords
%        identifierstyle=,           % nothing happens
%        commentstyle=\itshape,      % default
%        stringstyle=\ttfamily,      % typewriter type for strings
%        showstringspaces=false}     % no special string spaces
% \end{lstxsample}
%\fi
% \begin{lstsample}{}{}
%    \begin{lstlisting}
%    for i:=maxint to 0 do
%    begin
%        { do nothing }
%    end;
%
%    Write('Case insensitive ');
%    WritE('Pascal keywords.');
%    \end{lstlisting}
% \end{lstsample}
%\ifcolor
% \begin{advise}
% \item 你要求注释为白色，但我在左侧能看到注释。
%       \advisespace
%       可能有几个原因：
%       (1) 你在非白纸上打印了本文档。
%       (2) 如果你在 \texttt{.dvi} 文件中查看本文档，你的查看器似乎
%           对颜色特殊指令有问题。试着在白纸上打印该页。
%       (3) 如果在白纸上打印仍显示注释，颜色特殊指令不适合你的打印机
%           或打印机驱动。重新生成文档再试一次——并确保
%           \packagename{color} 宏包配置正确。
% \end{advise}
%\fi
% 样式使用两种不同类型的命令。|\ttfamily| 和 |\bfseries| 都不带参数，
% 但 |\underbar| 带参数；它给后面的参数加下划线。一般来说，\emph{最后}
% 一个命令可以恰好读取一个参数，即宏包排版的某段内容。有一个例外。
% \ikeyname{basicstyle} 的最后一个命令\emph{不得}读取任何记号——
% 否则你会陷入大麻烦。
% \begin{advise}
% \item `|basicstyle=\small|' 看起来不错，但比如 `|commentstyle=\tiny|'
%       加上空 basic style 时注释看起来很糟。
%       \advisespace
%       不要在单个列表中使用不同字号。
% \item 但我真的想要那样！
%       \advisespace
%       不，你不想要。
%^^A       宏包在每个列表开始时选择基本样式后会调整内部数据。如果你
%^^A       更改注释或字符串的字号，就会出问题。第~\ref{rColumnAlignment}
%^^A       节展示了如何克服这个问题。但再次强调：除非你确实知道自己在做
%^^A       什么，否则不要在单个列表中使用不同字号。
% \end{advise}
%
% \paragraph{警告}\label{wStrikingStyles}
% 你应该非常谨慎地使用醒目样式；刚才的例子还算是温和的——用过头会很可怕。
% \emph{始终使用得体的高亮。}遗憾的是，很难给出更多建议，因为它们取决于
% 你正在创建的文档类型。例如，幻灯片或其他演示文稿通常比书籍需要更醒目的
% 样式。最终，需要\emph{你自己}找到那个黄金平衡点！
%
%
% \subsection{引得你上手}\label{gSeduceToUse}
%
% 你已经了解了所有美化排版命令和一些主要参数。下面简要介绍其他功能，
% 并不完整。目录和索引也提供信息。
%
% \paragraph{行号}
% 可用于所有行间列表，例如左侧小号数字，每隔一行显示，距列表 5pt：
% \begin{lstxsample}[numbers,numberstyle,stepnumber,numbersep]
%    \lstset{numbers=left, numberstyle=\tiny, stepnumber=2, numbersep=5pt}
% \end{lstxsample}
% \begin{lstsample}{}{}
%    \begin{lstlisting}
%    for i:=maxint to 0 do
%    begin
%        { do nothing }
%    end;
%
%    Write('Case insensitive ');
%    WritE('Pascal keywords.');
%    \end{lstlisting}
% \end{lstsample}
% \begin{advise}
% \item 后续列表中的行号我去不掉了。
%       \advisespace
%       `|numbers=none|' 关闭行号。
% \item 我可以在可选参数中使用这些键吗？
%       \advisespace
%       当然。注意可选参数只修改单个特定列表的值：你为单个列表更改
%       行号的外观、步长或距离。之后恢复之前的值。
% \end{advise}
% 该环境允许你中断列表：你可以结束一个列表，稍后以正确的行号继续它，
% 即使中间有其他列表。请阅读第~\ref{uLineNumbers} 节了解详细讨论。
%
% \paragraph{浮动列表}
% 显示的列表可以浮动：
% \begin{lstsample}{\lstset{frame=tb}}{}
%    \begin{lstlisting}[float,caption=A floating example]
%    for i:=maxint to 0 do
%    begin
%        { do nothing }
%    end;
%
%    Write('Case insensitive ');
%    WritE('Pascal keywords.');
%    \end{lstlisting}
% \end{lstsample}
% 现在不要关心 \ikeyname{caption} 参数。如果你把这个例子放到最小文件中
% 用 \LaTeX\ 编译，请不要惊讶：你会看不到水平线，因为它们在别处描述。
% \begin{advise}
% \item \LaTeX\ 的浮动体机制允许确定浮动体的位置。我如何对这些做同样的操作？
%       \advisespace
%       例如你可以写 `|float=tp|'。
% \end{advise}
%
% \paragraph{其他功能}
% 还有一些功能尚未提及：长行自动断行、在列表中使用 \LaTeX\ 代码的可能性、
% 自动索引，或个人语言定义。
% 再来一点小惊喜？这就给你。但注意结果不是仅由右侧的 \LaTeX\ 代码产生的。
% 主要参数被隐藏了。
% \begin{lstsample}{^^A
%   \lstset{literate={:=}{{$\gets$}}1 {<=}{{$\leq$}}1 {>=}{{$\geq$}}1 ^^A
%     {<>}{{$\neq$}}1}}{}
%    \begin{lstlisting}
%    if (i<=0) then i := 1;
%    if (i>=0) then i := 0;
%    if (i<>0) then i := 0;
%    \end{lstlisting}
% \end{lstsample}
%
% 你不确定是否应该使用 \packagename{listings}？请阅读下一节！
%
%
% \subsection{替代方案}
%
% \begin{advise}
% \item 你为什么要列出替代方案？
%       \advisespace
%       了解竞争对手总是好的。^^A :-)
% \item 我读了下面的描述，\packagename{listings} 宏包似乎包含了所有功能。
%       我为什么要使用其他程序？
%       \advisespace
%       首先，这些描述只给人一点味道，不是完整概述；其次，
%       \packagename{listings} 缺少一些特性；最终，你应该使用最符合你需求的程序。
% \end{advise}
% 本宏包当然不是排版源代码的终极工具。如果你对 \packagename{listings}
% 不满意，其他程序也做得很好。有些独立于 \LaTeX，有些作为独立程序加
% \LaTeX\ 宏包提供，还有一些是不美化排版源代码的宏包。第二类包括转换器、
% 交叉编译器和预处理器。这类程序生成 \LaTeX\ 文件，你可以在文档中使用，
% 或作为独立的可直接运行的 \LaTeX\ 文件。
%
% 注意我这里不涉及任何文学编程工具，它们也可以是替代方案。不过，你应该
% 听说过 \texttt{WEB} 系统，这是 Donald E.~Knuth 教授开发并用于编写和
% 实现 \TeX 的工具。
%
% \paragraph{\href{http://www.infres.enst.fr/~demaille/a2ps}{\packagename{a2ps}}}
% 最初是作为 `ASCII 转 PostScript' 转换器开始的，但今天你可以用
% \texttt{--pretty-print=}\meta{language} 选项调用该程序。如果你喜欢的
% 编程语言尚未被支持，你可以编写自己的所谓样式表。你可以请求行号、边框、
% 页眉、每页多页等。你甚至可以打印 $\forall$ 或 $\alpha$ 等符号而不是
% 它们的冗长形式。如果你只想要程序列表而不是带列表的文档，这是最佳选择。
%
% \paragraph{\href{http://mirror.ctan.org/support/lgrind}{\packagename{LGrind}}}
% 是一个交叉编译器，附带许多预定义编程语言。例如，你可以把右侧的代码
% 放到文档中，用 \texttt{-e} 选项（和文件名）调用 \packagename{LGrind}，
% 然后用 \LaTeX\ 编译生成的文件。你应该得到与左侧类似的结果：
% \begin{center}
% \begin{minipage}{0.45\linewidth}
%\iflgrind
%    \LGindent=0pt
%    \LGinlinefalse\LGbegin\lgrinde
%    \L{\LB{\K{for}_\V{i}:=\V{maxint}_\K{to}_\N{0}_\K{do}}}
%    \L{\LB{\K{begin}}}
%    \L{\LB{____\C{}\{_do_nothing_\}\CE{}}}
%    \L{\LB{\K{end};}}
%    \L{\LB{}}
%    \L{\LB{\V{Write}(\S{}{'}Case_insensitive_{'}\SE{});}}
%    \L{\LB{\V{WritE}(\S{}{'}Pascal_keywords.{'}\SE{});}}
%    \endlgrinde\LGend
%\else
%    \packagename{LGrind} not installed.
%\fi
% \end{minipage}
% \begin{minipage}{0.45\linewidth}
% \begin{verbatim}
% %[
% for i:=maxint to 0 do
% begin
%     { do nothing }
% end;
%
% Write('Case insensitive ');
% WritE('Pascal keywords.');
% %]\end{verbatim}
% \end{minipage}
% \end{center}
% 如果你使用 |%(| 和 |%)| 而不是 |%[| 和 |%]|，你会得到一个代码片段
% 而不是行间列表。此外你还可以在左侧或右侧加行号、在源代码中使用任意
% \LaTeX\ 代码、用符号代替冗长的名称、设置字体等等。你会（不得不）
% 喜欢它（如果你不喜欢 \packagename{listings} 的话）。
%
% 注意 \packagename{LGrind} 包含无销售许可的代码，因此是非自由软件。
%
% \paragraph{\href{ftp://axp3.sv.fh-mannheim.de/cvt2latex}{\packagename{cvt2ltx}}}
% 是一组面向 C、Objective C、\Cpp、IDL 和 Perl 的 `源代码转 \LaTeX' 转换器。
% 可以通过命令行选项选择不同样式、行号和其他限定符。遗憾的是，文档没有
% 说明如何添加其他编程语言。
%
% \paragraph{\href{http://mirror.ctan.org/support/C++2LaTeX-1_1pl1}{\packagename{\Cpp2\LaTeX}}}
% 是一个 C/\Cpp\ 转 \LaTeX\ 转换器。你可以指定注释、指令、关键词和字符串的
% 字体，或制表符大小。但据我所知不能给行编号。
%
% \paragraph{\href{http://mirror.ctan.org/support/slatex}{\packagename{S\LaTeX}}}
% 是一个美化排版的 Scheme 程序（自动调用 \LaTeX），专门为 Scheme 和其他
% Lisp 方言设计。它支持独立文件、文本和行间列表，甚至可以嵌套命令/环境，
% 例如在注释中使用 \LaTeX\ 代码。关键词、常量、变量和符号可定义，
% 可以使用不同样式。无行号。
%
% \paragraph{\href{http://mirror.ctan.org/support/tiny_c2l}^^A
%   {\packagename{tiny\textunderscore c2ltx}}}
% 是一个基于 \packagename{cvt2ltx}（还是反过来？）的 C/\Cpp/Java 转
% \LaTeX\ 转换器。它支持行号、块注释、注释中/作为注释的 \LaTeX\ 代码，
% 以及智能断行。字体选择和制表符是硬编码的，即如果你想更改外观，
% 必须重新编译该程序。
%
% \paragraph{\href{http://mirror.ctan.org/macros/latex/contrib/misc}{\packagename{listing}}}
% ——注意少了一个 \packagename{s}——不是美化排版器，
% \texttt{listing.sty} 末尾关于文档的格言也不正确。\space ^^A :-)
% 它定义了 |\listoflistings| 和一个非浮动的列表环境。所有字体选择和缩进
% 必须手动完成。不过，如果你有另一个工具做这些工作（例如
% \packagename{LGrind}），它就很有用。
%
% \paragraph{\href{http://mirror.ctan.org/macros/latex/contrib/alg}{\packagename{alg}}}
% 提供与 \packagename{algorithms} 基本相同的功能。所以阅读下一段，
% 注意语法会不同。
%
% \paragraph{\href{http://mirror.ctan.org/macros/latex/contrib/algorithms}^^A
%   {\packagename{algorithms}}}
% 走了一条完全不同的路。你描述一个算法，宏包来排版它，例如
% \begin{center}
% \begin{minipage}{0.45\linewidth}
%\ifalgorithmicpkg
%    \begin{algorithmic}
%    \IF {$i\leq0$}
%    \STATE $i\gets1$
%    \ELSE\IF {$i\geq0$}
%    \STATE $i\gets0$
%    \ENDIF\ENDIF
%    \end{algorithmic}
%\else
%    \packagename{algorithms} not installed.
%\fi
% \end{minipage}
% \begin{minipage}{0.45\linewidth}
% \begin{quote}
% |\begin{algorithmic}|\\
% |  \IF{$i\leq0$}|\\
% |  \STATE $i\gets1$|\\
% |  \ELSE\IF{$i\geq0$}|\\
% |  \STATE $i\gets0$|\\
% |  \ENDIF\ENDIF|\\
% |\end{algorithmic}|
%\end{quote}
% \end{minipage}
% \end{center}
% 如本例所示，即使输入不好看，你也能得到好看的算法。该宏包提供更多结构，
% 如 |for| 循环、|while| 循环或注释。你可以请求行号、`ruled'、`boxed'
% 和浮动算法、算法列表，还可以自定义 \textbf{if}、\textbf{then} 等术语。
%
% \paragraph{\href{http://www.mimuw.edu.pl/~wolinski/pretprin.html}^^A
%   {\packagename{pretprin}}}
% 是一个用于美化排版形式语言文本的宏包——正如 TUGboat 第 19 卷
% （1998）第 3 期标题所述。它提供美化排版\emph{并}格式化源代码的环境。
% 已定义 Pascal 和 Prolog 的分析器；添加其他语言很容易——如果你熟悉或
% 稍微了解自动机和形式语言。
%
% \paragraph{\packagename{alltt}}
% 定义了一个类似于 \texttt{verbatim} 的环境，不同之处在于 |\|、|{| 和
% |}| 保持其通常含义。这意味着你可以在 verbatim 中使用命令，例如
% 选择不同字体或进入数学模式。
%
% \paragraph{\href{http://mirror.ctan.org/macros/latex/contrib/moreverb}^^A
%   {\packagename{moreverb}}}
% 需要 \packagename{verbatim}，提供将 verbatim 输出到文件、`boxed'
% verbatim 和行号。
%
% \paragraph{\packagename{verbatim}}
% 定义了标准 \texttt{verbatim} 环境的改进版本，以及逐字输入文件的命令。
%
% \paragraph{\href{http://mirror.ctan.org/macros/latex/contrib/fancyvrb}^^A
%   {\packagename{fancyvrb}}}
% 大致说来，它是 \packagename{alltt} 的超集，
% \packagename{moreverb} 和 \packagename{verbatim}，但有更多参数控制输出。
% 宏包提供边框、左侧或右侧行号、自动断行（有难度）等等。例如，
% 存在与 \packagename{listings} 的接口，即你可以自动美化排版源代码。
% \packagename{fvrb-ex} 宏包基于 \packagename{fancyvrb}，定义了一些环境，
% 用于展示与本指南中类似的例子。
%
%
% \section{后续步骤}\label{uTheNextSteps}
%
% 现在，在真正使用 \packagename{listings} 宏包之前，你应该\emph{认真}
% 阅读软件许可。这不花多少时间，而且提供你可能需要知道的信息。
%
%
% \subsection{软件许可}\label{uSoftwareLicense}
%
% 文件 \texttt{listings.dtx} 和 \texttt{listings.ins} 以及仅由这两个文件
% 生成的所有文件称为 `\packagename{listings} 宏包'或简称`宏包'。
% \texttt{lstdrvrs.dtx} 及由该文件生成的文件称为`驱动'。
%
% \paragraph{版权}
%   \packagename{listings} 宏包版权所有 \copyright\ 1996--2004 Carsten Heinz，
%   \copyright\ 2006 Brooks Moses。驱动的版权归驱动文件中列出的各位作者所有。
%
% \paragraph{分发与修改}
%   \packagename{listings} 宏包及其驱动可在 LaTeX Project Public License
%   条件下分发和/或修改，无论是该许可证的 1.3c 版还是（由你选择）任何更高版本。
%   该许可证的最新版本见
%      \href{http://www.latex-project.org/lppl.txt}{http://www.latex-project.org/lppl.txt}
%   且 1.3c 版或更高版本是 LaTeX 2003/12/01 或更高版本所有发行版的一部分。
%
% \paragraph{联系方式}
%   阅读第~\lstref{uTroubleshooting} 节了解如何提交错误报告。
%   将所有其他意见、想法和附加编程语言发送至 \lstemail，
%   主题中包含 \texttt{listings}。
%
%
% \subsection{加载宏包}\label{uPackageLoading}
%
% 与 \LaTeX 中通常一样，宏包通过
%    |\usepackage[|\meta{options}|]{listings}|
% 加载，其中 |[|\meta{options}|]| 是可选的，给出逗号分隔的选项列表。
% 每个选项要么加载额外的 \packagename{listings} aspect，要么更改默认属性。
% 通常你不需要关心这些选项。但在某些情况下可能需要：如果你想编译用本宏包
% 早期版本创建的文档，或者使用特殊功能。下面是可能选项的不完整列表。
% \begin{advise}
% \item 哪里有所有选项的列表？
%       \advisespace
%       在开发者指南中，因为引入它们是为了更容易调试宏包。
%       阅读第~\ref{uHowTos} 节了解如何获取该指南。
% \end{advise}
% \begin{description}
% \item[\normalfont\texttt{0.21}]\leavevmode
%
%       调用兼容模式，用于编译为 \packagename{listings} 0.21 版编写的文档。
%
% \item[\normalfont\texttt{draft}]\leavevmode
%
%       宏包不打印独立文件，但显示标题并定义相应标签。
%       注意全局 |\documentclass| 选项 \texttt{draft} 会被识别，
%       所以你不需要将其重复为宏包选项。
%
% \item[\normalfont\texttt{final}]\leavevmode\label{uoption:final}
%
%       覆盖全局 \texttt{draft} 选项。
%
% \item[\normalfont\texttt{savemem}]\leavevmode
%
%       尝试节省一些 \TeX 的内存。如果你经常在语言之间切换，
%       它也可能减少编译时间。但这一切取决于具体文档及其列表。
% \end{description}
% 注意各种实验功能也需要通过选项显式加载。阅读第~\ref{rExperimentalFeatures}
% 节中的相应行。
%
% \medbreak
% 加载宏包后，建议用以下命令加载所有使用的编程语言方言。用一个命令加载
% 多种语言比按需加载每种语言更快。
% \begin{syntax}
% \item {\rstyle\icmdname\lstloadlanguages}\marg{逗号分隔的语言列表}
%
%       每种语言的形式为 \oarg{dialect}\meta{language}。如果没有可选的
%       \oarg{dialect}，宏包加载默认方言。所以如果你想要 Visual \Cpp，
%       写 `|[Visual]C++|'；想要 ISO \Cpp，写 `|[ISO]C++|'。两者可以用
%       命令 |\lstloadlanguages{[Visual]C++,[ISO]C++}| 一起加载。
%
%       第~\pageref{uPredefinedLanguages} 页的表~\ref{uPredefinedLanguages}
%       显示了所有已定义的语言及其方言。
% \end{syntax}
%^^A 在加载语言之后甚至之前，你可能想定义默认方言——只是为了不依赖配置文件。
%
%
% \subsection{键值接口}\label{uTheKey=ValueInterface}
%
% 本宏包使用 David Carlisle 的 \packagename{graphics} 捆绑中的
% \packagename{keyval} 宏包。每个参数由关联的键和用户提供的值控制。
% 例如，\ikeyname{firstline} 是一个键，|2| 是该键的有效值。
%
% 命令 {\rstyle\icmdname\lstset} 接受逗号分隔的``键|=|值''对列表。
% 第一个包含多个条目的列表在第~\pageref{gFirstKey=ValueList} 页：
% |firstline=2,lastline=5|。
% \begin{advise}
% \item 所以我可以写一次 `|\lstset{firstline=2,lastline=5}|' 就一劳永逸了？
%       \advisespace
%       不行。`\ikeyname{firstline}' 和 `\ikeyname{lastline}' 属于一小部分
%       只在单个列表上使用的键。不过你的命令并不非法——只是没有效果。
%       你必须在环境或输入命令的可选参数内使用这些键。
% \item 有没有更好的键|=|值表示例？
%       \advisespace
%       第~\ref{gFigureOutTheAppearance} 节有一个。
% \item `|language=[77]Fortran|' 在可选参数中不起作用。
%       \advisespace
%       如果在可选参数内使用带可选参数的值，你必须用花括号将值括起来。
% 在本例中写 `|language={[77]Fortran}|' 来选择 Fortran 77。
% \item 如果我在可选参数内使用 `\ikeyname{language}' 键，排版下一个列表时
%       语言不生效。
%       \advisespace
%       通过 `|\lstset|' 设置的所有参数保持其值直到当前环境或组结束。
%       之后恢复之前的值。两个美化排版命令和 `\texttt{lstlisting}' 环境的
%       可选参数仅对特定列表生效，即值立即恢复。例如，你可以选择一种主语言，
%       为特殊列表更改它。
% \item \icmdname\lstinline\ 有可选参数？
%       \advisespace
%       是的。由此产生一个限制：你不能用左方括号 `|[|' 作为定界符，
%       除非像 `|\lstinline[][var i:integer;[|' 那样至少指定一个空可选参数。
%       如果你忘了这一点，要么会从 \TeX 得到``runaway argument''错误，
%       要么从 \packagename{keyval} 宏包得到错误信息。
% \end{advise}
%
%
% \subsection{编程语言}\label{uProgrammingLanguages}
%
% 你已经知道如何激活编程语言——至少 Pascal。可选参数选择语言的特定方言。
% 例如，|language=[77]Fortran| 选择 Fortran 77，|language=[XSC]Pascal|
% 对 Pascal XSC 做同样的事。一般形式为
%    {\rstyle\ikeyname{language}}|=|\oarg{dialect}\meta{language}。
% 如果你想取消关键词、注释和字符串检测，使用 |language={}| 作为
% |\lstset| 的参数或可选参数。
%
% 表~\ref{uPredefinedLanguages} 显示了所有预定义语言和方言。
% 分别使用列出的名称作为 \meta{language} 和 \meta{dialect}。如果表中
% 没有给出方言或给出`空'，就不要指定方言。每个带下划线的方言是默认值；
% 如果你省略可选参数，就会选择它。预定义的默认值是最新语言版本或标准方言。
%^^A
%^^A  Make table of predefined languages.
%^^A
%\let\lstlanguages\empty
%\makeatletter
%\@for\lst@temp:={lstlang1.sty,lstlang2.sty,lstlang3.sty}\do
%    {\IfFileExists\lst@temp{}{\let\lstlanguages\relax}}
%\makeatother
%\ifx\lstlanguages\relax
%    \PackageWarningNoLine{Listings}
%        {Standard drivers not available.\MessageBreak
%         Please check your installation.\MessageBreak
%         Compilation aborted}
%    \csname @@end\expandafter\endcsname
%\fi
%\lstscanlanguages\lstlanguages{lstlang1.sty,lstlang2.sty,lstlang3.sty}{}^^A
%\def\topfigrule{\hrule\kern-0.4pt\relax}^^A
%\let\botfigrule\topfigrule
%\belowcaptionskip=\smallskipamount
% \begin{table}[tbhp]
% \small
% \caption{预定义语言。
%          注意有些定义是初步的，例如 HTML 和 XML。
%          每个带下划线的方言是默认方言。}^^A
%          \label{uPredefinedLanguages}^^A
% \makeatletter
% \setbox\@tempboxa\hbox{^^A
%    \InputIfFileExists{listings.cfg}{\lst@InputCatcodes}{}}^^A
% \lstprintlanguages\lstlanguages
% \end{table}
%^^A
%^^A end of table
%^^A
%\lstset{defaultdialect=[doc]Pascal}^^A restore
% \begin{advise}
% \item 如何定义默认方言？
%       \advisespace
%       请参阅第~\ref{rLanguagesAndStyles} 节关于 \keyname{defaultdialect}' 的说明。
% \item 我有 C 代码和汇编代码混合在一起。\packagename{listings} 能美化排版
%       这样的源代码吗？即同时高亮两种语言的关键词和注释？
%       \advisespace
%       \ikeyname{alsolanguage}|=|\oarg{dialect}\meta{language}' 在当前语言之外
%       再选择一种语言。所以你只需为你的
%       汇编方言写一个语言定义，它不会与 C 的定义冲突。此外你可能想使用
%       第~\ref{rLanguagesAndStyles} 节中描述的键 \keyname{classoffset}'。
% \item 如何定义自己的语言？
%       \advisespace
%       这在第~\ref{rLanguageDefinitions} 节讨论。如果你认为其他人也能从你的
%       定义中受益，你可以把它发送到第~\ref{uSoftwareLicense} 节中的地址。
%       然后它将在 \LaTeX 项目公共许可下发布。
% \end{advise}
% 注意参数 \meta{language} 和 \meta{dialect} 不区分大小写，空格也无效。
%
% 至少有一种语言（VDM，维也纳开发语言，一种 ISO 标准化的形式化建模语言\footnote{参见 \url{http://www.vdmportal.org}。文献中指出，先前提到的该网址已不再指向 VDM 专用页面。}）并不受 \packagename{listings} 宏包直接支持，它需要独立的宏包：
% \packagename{vdmlisting}。另一方面，\packagename{vdmlisting} 本身也使用了 \packagename{listings} 宏包，因此在这里应当一并提及。
%
%
% \subsubsection{偏好设置}\label{uPreferences}
%
% 有时，语言支持文件的作者会提供自己的配置
% 偏好。这些偏好或来自个人经验，或来自
% IDE 中的设置，并且可以定义为一个 \packagename{listings}
% 样式。从 \packagename{listings} 宏包 1.5b 版本起，这些
% 样式以如下文件名提供：
% |listings-|\meta{language}|.prf|，其中 \meta{language} 是
% 所支持编程语言的小写名称。
%
% 所以如果 \packagename{listings} 宏包的用户想使用这些偏好设置，
% 例如在使用 Python 时，可以在她/他的私有/项目级 |lstlocal.cfg|
% 配置文件末尾写
% \begin{quote}
%     |\input{listings-python.prf}|
% \end{quote}
% 只要文件 |listings-python.prf| 位于 \TeX{} 搜索路径中。当然，
% 该文件可以根据用户偏好进行修改。
%
% 目前有八个这样的偏好设置文件：
% \begin{enumerate}
%   \item |listings-acm.prf|
%   \item |listings-bash.prf|
%   \item |listings-fortran.prf|
%   \item |listings-hansl.prf|
%   \item |listings-lua.prf|
%   \item |listings-python.prf|
%   \item |listings-rexx.prf|
%   \item |listings-riscv.prf|
% \end{enumerate}
% 欢迎所有贡献者提供更多个人偏好设置。
%
%
% \subsection{特殊字符}\label{uSpecialCharacters}
%
%
% \paragraph{制表符}
% 如果你的源代码包含制表符，可能会得到意外的输出。宏包假设制表位在
% 第 9、17、25、33 列等。这是通过 |tabsize=8| 预定义的。如果你把 8
% 改成数字 $n$，制表位将在第 $n+1,2n+1,3n+1$ 列等。
% \begin{lstsample}[tabsize]{}{}
%    \lstset{tabsize=2}
%    \begin{lstlisting}
%    123456789
%    	{ one tabulator }
%    		{ two tabs }
%    123		{ 123 + two tabs }
%    \end{lstlisting}
% \end{lstsample}
% 为了更好地说明，左侧使用 |tabsize=2|，但逐字代码使用 |tabsize=4|。
% 注意 |\lstset| 修改同一环境或组中后续所有列表的值。这里没有问题，
% 因为例子在 minipage 内排版。如果你想为单个列表更改设置，使用可选参数。
%
%
% \paragraph{可见制表符和空格}
% 可以让空格和制表符可见：
% \begin{lstsample}[showspaces,showtabs,tab]{}{}
%    \lstset{showspaces=true,
%            showtabs=true,
%            tab=\rightarrowfill}
%    \begin{lstlisting}
%        for i:=maxint to 0 do
%        begin
%    	{ do nothing }
%        end;
%    \end{lstlisting}
% \end{lstsample}
% 如果你请求 \ikeyname{showspaces} 但不请求 \ikeyname{showtabs}，
% 制表符会被转换为可见空格。\ikeyname{tab} 的默认定义产生一个`宽可见空格'
% \lstinline[showtabs]!	!。所以你可能想用 |$\to$|、|$\dashv$| 或其他符号代替。
% \begin{advise}
% \item 一些建议：(1) 你真的应该缩进源代码行，使列表更易读。
%       (2) 不要有些行用空格缩进，另一些行用制表符。更改制表符大小
%       （你的编辑器或美化排版工具的）会完全打乱列对齐。
%       (3) 因此，永远不要和制表符大小不同的人共享你的文件！^^A
%       只有你使用制表符时才成立，只是开个玩笑 :-)
% \item 为了让 \LaTeX 代码更易读，我缩进了环境中的程序列表。
%       如何在输出中去掉那个缩进？
%       \advisespace
%       阅读第~\ref{uHowTos} 节中的`如何吞掉字符'。
% \end{advise}
%
%
% \paragraph{换页符}
% 另一个特殊字符是换页符，默认产生一个空行。
% {\rstyle\ikeyname{formfeed}}|=\newpage| 会在每个换页符处产生新页。
% 请注意这样的定义（即使是默认值）可能与边框冲突。
%
%
% \paragraph{国家或地区字符}
% 如果你直接输入代码 128--255 的字符，并在列表中使用它们，
% 请告诉宏包——否则你会得到非常有趣的结果。
% {\rstyle\ikeyname{extendedchars}}|=true| 允许、|extendedchars=false|
% 禁止 \packagename{listings} 处理列表中的扩展字符。如果你使用它们，
% 应该加载 \packagename{fontenc}、\packagename{inputenc} 和/或
% 任何其他定义这些字符的宏包。
% \begin{advise}
% \item 我在将 \packagename{inputenc} 与 \packagename{listings} 一起使用时
%       遇到问题。
%       \advisespace
%       这可能是兼容性问题。请按照第~\lstref{uTroubleshooting} 节的说明
%       提交错误报告。
% \end{advise}
% 扩展字符不包括阿拉伯语、中文、希伯来语、日语等——具体来说，
% 不包括任何每字符使用多字节的编码。
%
% 因此，如果你使用支持多字节字符的宏包，例如用于中文和 UTF-8 字符的
% \packagename{CJK} 或 \packagename{ucs} 宏包，你必须避免让
% \packagename{listings} 处理扩展字符。通常最好同时指定
% |extendedchars=false|，以避免 \packagename{listings} 卷入另一个宏包的
% 扩展字符处理中。
%
% 如果你确实在 CJK 环境中有一个列表，并想在列表中包含 CJK 字符，
% 可以将它们放在转义到 \LaTeX 的注释中——参见第~\ref{rEscapingToLaTeX} 节
% 了解如何操作。（如果列表不在 CJK 环境中，你可以简单地在转义到 \LaTeX
% 的注释部分放一个小的 CJK 环境。）
%
% 同样，如果你在列表中使用 UTF-8 扩展字符，它们必须放在转义到 \LaTeX
% 的部分中。
%
% 此外，第~\ref{uNationalCharacters} 节有一些关于在 $\Lambda$ 上下文中
% 处理扩展字符的细节。
%
%
% \subsection{行号}\label{uLineNumbers}
%
% 你已经从第~\ref{gSeduceToUse} 节了解了键 \ikeyname{numbers}、
% \ikeyname{numberstyle}、\ikeyname{stepnumber} 和 \ikeyname{numbersep}。
% 现在我们讨论连续列表。你有两个选项来获得跨列表的一致行号。
%
% \begin{lstsample}[firstnumber]{\lstset{numbers=left,numberstyle=\tiny,^^A
%       stepnumber=2,numbersep=5pt}}{}
%    \begin{lstlisting}[firstnumber=100]
%    for i:=maxint to 0 do
%    begin
%        { do nothing }
%    end;
%
%    \end{lstlisting}
%    And we continue the listing:
%    \begin{lstlisting}[firstnumber=last]
%    Write('Case insensitive ');
%    WritE('Pascal keywords.');
%    \end{lstlisting}
% \end{lstsample}
% 在本例中，\ikeyname{firstnumber} 初始设为 100；几行后值为 \texttt{last}，
% 它延续上一个列表的编号。注意第一部分末尾的空行在这里不打印，
% 但它计入行号。你还应该注意，可以写一次 |\lstset{firstnumber=last}|
% 就得到连续编号的代码行——除非你为特定列表指定不同值。
%
% 另一方面，你可以使用 |firstnumber=auto| 并为列表命名。
% 名称相同（区分大小写！）的列表共享行计数器。
% \begin{lstsample}[name]{\lstset{numbers=left,numberstyle=\tiny,stepnumber=2,numbersep=5pt}}{}
%    \begin{lstlisting}[name=Test]
%    for i:=maxint to 0 do
%    begin
%        { do nothing }
%    end;
%
%    \end{lstlisting}
%    And we continue the listing---note
%    the blank line at the end of the
%    previous code:
%    \begin{lstlisting}[name=Test]
%    Write('Case insensitive ');
%    WritE('Pascal keywords.');
%    \end{lstlisting}
% \end{lstsample}
% 下一个 |Test| 列表将从第 {\makeatletter\lstno@Test} 行继续，
% 无论中间是否有其他列表。
%
% 你还可以选择要打印的行，选项 `\ikeyname{linerange}' 和
% `\ikeyname{consecutivenumbers}' 是你的好帮手。例如在演示文稿中，
% 你不需要程序的注释，所以更喜欢行号连续编号，但结果应该反映程序的行为——
% 你省略冗长输出的部分。所以你可能有下面的程序及其结果。
% \begin{lstsample}[name]{\lstset{numbers=left,numberstyle=\tiny,stepnumber=1,numbersep=5pt}}{}
%    \begin{lstlisting}[name=Test,
%      language={[ansi]C},
%      linerange={1-4,6-7,10-14,
%        17-19,21-22},
%      firstnumber=1]
%    #include <stdio.h>
%    #include <stdlib.h>
%
%    int main(int argc,char* argv[]){
%      /* declaring variables */
%      int i;
%      int limit;
%
%      /* checking arguments */
%      if ( argc > 1 ) {
%        limit = atoi(argv[1]);
%      } else {
%        limit = 100;
%      }
%
%      /* counting lines */
%      for (i = 1;i <= limit;i++) {
%        printf("Line no. %3.0d\n", i);
%      }
%
%      return 0;
%    }
%
%    \end{lstlisting}
%    And these are the results:
%    \begin{lstlisting}[language={},
%      linerange={1-2,6-7},
%      consecutivenumbers=false]
%    Line no.   1
%    Line no.   2
%    Line no.   3
%    Line no.   4
%    Line no.   5
%    Line no.   6
%    Line no.   7
%    \end{lstlisting}
% \end{lstsample}
%
% \begin{advise}
% \item 好的。那我怎么得到递减的行号？
%       \advisespace
%       抱歉，你说什么？
%       \advisespace
%       像第~\pageref{rDecreasingLabels} 页那样的递减行号。
%       \advisespace
%       我可以建议你通过其他方式展现个性吗？如果你不同意，
%       可以试试负的 `\ikeyname{stepnumber}'。
%       （与 `\ikeyname{firstnumber}' 一起使用）。
% \end{advise}
%
% 关于如何引用行号，请参阅第~\ref{uHowTos} 节。
%
%
% \subsection{版面元素}
%
% 通过垂直间距、水平线或不同字号和字体来构建版面结构总是好主意。
% 最能强调整个列表的是——不要同时使用——颜色、边框、垂直间距和标题。
% 后者当然也便于引用列表。
%
% \paragraph{垂直间距}
% 键 {\rstyle\ikeyname{aboveskip}} 和 {\rstyle\ikeyname{belowskip}}
% 控制行间列表上下方的垂直间距。两个键接受尺寸或弹性间距作为值，
% 初始化为 |\medskipamount|。
%
% \paragraph{边框}
% 键 \ikeyname{frame} 接受详细值 \keyvalue{none}、\keyvalue{leftline}、
% \keyvalue{topline}、\keyvalue{bottomline}、\keyvalue{lines}（上下线）、
% \keyvalue{single}（单边框）或 \keyvalue{shadowbox}。
% \begin{lstsample}[frame]{}{}
%    \begin{lstlisting}[frame=single]
%    for i:=maxint to 0 do
%    begin
%        { do nothing }
%    end;
%    \end{lstlisting}
% \end{lstsample}
% \begin{advise}
% \item 线不对齐。
%       \advisespace
%       这可能是本宏包的错误，或者是你的 \texttt{.dvi} 驱动的问题。
%       在向宏包作者发送错误报告\emph{之前}，大量修改第~\ref{rFrames} 节
%       描述的参数。并且要一步一步地做！例如，从 `|framerule=10mm|' 开始。
%       如果线仍然以相同的（小）量错位，问题就不出在线宽上。那就继续
%       下一个参数。另外，Adobe Acrobat 有时会有单像素舍入误差，
%       在屏幕上显示 PDF 文件时会导致角落的小错位；遗憾的是这是正常的。
% \end{advise}
% 或者你可以直接通过使用四个首字母控制 \texttt{t}op、\texttt{r}ight、
% \texttt{b}ottom 和 \texttt{l}eft 的线：单字母表示单线，大写表示双线。
% \begin{lstsample}[frame]{}{}
%    \begin{lstlisting}[frame=trBL]
%    for i:=maxint to 0 do
%    begin
%        { do nothing }
%    end;
%    \end{lstlisting}
% \end{lstsample}
% 注意当且仅当相邻两条线都被请求时，才绘制角。你可能认为线应该画到边缘，
% 但圆角呢？键 \ikeyname{frameround} 必须接受恰好四个字符作为值。
% 第一个字符连到右上角，按顺时针方向继续。`\texttt{t}' 字符使相应的角变圆。
% \begin{lstsample}[frameround]{}{}
%    \lstset{frameround=fttt}
%    \begin{lstlisting}[frame=trBL]
%    for i:=maxint to 0 do
%    begin
%        { do nothing }
%    end;
%    \end{lstlisting}
% \end{lstsample}
% 注意 \ikeyname{frameround} 与 |\lstset| 一起使用，因此该值影响同一组或环境中
% 后续所有列表。由于列表在这里位于 \texttt{minipage} 内，这没有问题。
% \begin{advise}
% \item 不要总是使用边框，特别是不要对短列表使用。那样什么也强调不了。
%       只对 $10\%$ 甚至更少的最重要列表使用边框。
% \item 如果你对浮动列表使用边框，你真的想要边框吗？
%       \advisespace
%       不，我想把浮动体和文本分开。
%       \advisespace
%       那最好重新定义 \LaTeX 的 `|\topfigrule|' 和 `|\botfigrule|'。
%       例如，你可以写 `|\renewcommand*\topfigrule{\hrule\kern-0.4pt\relax}|'
%       并对 |\botfigrule| 做同样的定义。
% \end{advise}
%
% \paragraph{标题}
% 现在我们讨论 \ikeyname{caption} 和 \ikeyname{label}。你可能猜对了，
% 它们可以像 \LaTeX 的 |\caption| 和 |\label| 命令一样使用，
% 尽管这里也可以在列表不在浮动体中时使用标题：\footnote{你应该记住，
% 根据 \cite{Mittelbach:companion} \ldots{} 如果标题文本适合一行，
% 文本居中；如果文本不适合单行，它将作为一个段落排版，宽度等于行宽。}
% \begin{lstsample}[caption,label]{\lstset{xleftmargin=.05\linewidth}}{}
%    \begin{lstlisting}[caption={Useless code},label=useless]
%    for i:=maxint to 0 do
%    begin
%        { do nothing }
%    end;
%    \end{lstlisting}
% \end{lstsample}
% 之后你可以通过 |\ref{useless}| 引用该列表。默认情况下，这样的列表会在
% 列表目录中获得一个条目，可以用命令 {\rstyle\icmdname\lstlistoflistings}
% 打印。键 {\rstyle\ikeyname{nolol}} 抑制环境或输入命令的条目。此外，
% 你可以为列表目录指定短标题：
%    \keyname{caption}|={|\oarg{short}\meta{long}|}|。
% 注意整个值用花括号括起来，因为在可选参数中使用了可选值。
%
% 如果你不想要标签 \texttt{\lstlistingname} 加编号，应该使用 \ikeyname{title}：
% \begin{lstsample}[title]{\lstset{xleftmargin=.05\linewidth}}{}
%    \begin{lstlisting}[title={`Caption' without label}]
%    for i:=maxint to 0 do
%    begin
%        { do nothing }
%    end;
%    \end{lstlisting}
% \end{lstsample}
% \begin{advise}
% \item 我的文档中 `\keyname{title}' 出了问题：标题前面
%       出现了一个定界符。
%       \advisespace
%       结果取决于文档类；有些不兼容。请联系宏包作者获取变通方法。
% \end{advise}
%
% \paragraph{颜色}
% 还有一点。你需要 \packagename{color} 宏包，然后可以通过
% \ikeyname{backgroundcolor}|=|\meta{颜色命令} 请求带颜色的背景。
% \begin{advise}
% \item 太好了！我喜欢颜色。
%       \advisespace
%       好的，是的，真的。我想提醒你注意第~\pageref{wStrikingStyles} 页
%       关于过于醒目的样式的警告。
% \end{advise}
%\ifcolor
% \begin{lstxsample}[backgroundcolor]
%    \lstset{backgroundcolor=\color{yellow}}
% \end{lstxsample}
%\else
% \begin{verbatim}
%    color package not installed\end{verbatim}
%\fi
% \begin{lstsample}{}{}
%    \begin{lstlisting}[frame=single,
%                       framerule=0pt]
%    for i:=maxint to 0 do
%    begin
%        j:=square(root(i));
%    end;
%    \end{lstlisting}
% \end{lstsample}
% 该例子还展示了如何在整个列表周围获得带颜色的间距：
% 使用一个边框线宽为零的边框。
%
%
% \subsection{强调标识符}\label{uEmphasizeIdentifiers}
%
% 回忆一下美化排版命令和环境。|\lstinline| 打印代码片段，
% |\lstinputlisting| 打印整个文件，\texttt{lstlisting} 打印位于
% \LaTeX 文件中的代码段。这些不同`类型'的源代码有什么用？
% 嗯，正好一句话中可能包含一个代码片段。整个文件通常作为附录或放在附录中。
% 不过有些编程书也在正文章节中包含这样的列表——为了增加页数。
% 现在源代码应该放在磁盘或 CD-ROM 上，只排版主要头文件或接口文件以供参考。
% 所以，请不要滥用 \packagename{listings} 宏包。但让我们回到主题。
%
% 显然，`\texttt{lstlisting} 源代码'不是用来生成可执行程序的。
% 这种源代码有某种教育目的，甚至是说教性的。
% \begin{advise}
% \item 教育性和说教性有什么区别？
%       \advisespace
%       教育性的东西可以是好的或坏的，对的或错的。
%       说教性的按定义就是对的。^^A :-)
% \end{advise}
% 通常宏包排版源代码片段时会高亮\emph{关键词}。这对熟悉编程语言的读者
% 来说不是必需的。主要关注的是接口、库或其他函数或变量的呈现。
% 如果这是你关心的，下面就是正确的键。比方说，你想强调函数 |square|
% 和 |root|，例如给它们加下划线。你可以这样做：
% \begin{lstxsample}[emph,emphstyle]
%    \lstset{emph={square,root},emphstyle=\underbar}
% \end{lstxsample}
% \begin{lstsample}{}{}
%    \begin{lstlisting}
%    for i:=maxint to 0 do
%    begin
%        j:=square(root(i));
%    end;
%    \end{lstlisting}
% \end{lstsample}
% \begin{advise}
% \item 注意标识符列表 |{square,root}| 用花括号括起来。否则
%       \packagename{keyval} 宏包会抱怨未定义的键 \keyname{root}，
%       因为逗号结束了键值对。还要注意，如果你在美化排版命令的可选参数内
%       使用键的可选参数，你\emph{必须}用花括号括住值。虽然不是必需的，
%       下面的例子使用了这些花括号。当它们变得必要时，通常会被忘记。
% \end{advise}
%
% 两个键都有一个可选的 \meta{class number} 参数，用于多个标识符列表：
%\ifcolor
% \begin{lstxsample}[emph,emphstyle]
%    \lstset{emph={square},      emphstyle=\color{red},
%            emph={[2]root,base},emphstyle={[2]\color{blue}}}
% \end{lstxsample}
%\else
% \begin{lstxsample}[emph,emphstyle]
%    \lstset{emph={square},      emphstyle=\underbar,
%            emph={[2]root,base},emphstyle={[2]\fbox}}
% \end{lstxsample}
%\fi
% \begin{lstsample}{}{}
%    \begin{lstlisting}
%    for i:=maxint to 0 do
%    begin
%        j:=square(root(i));
%    end;
%    \end{lstlisting}
% \end{lstsample}
% \begin{advise}
% \item 最大的 \meta{class number} 是多少？
%       \advisespace
%       $2^{31}-1=2\,147\,483\,647$。但在你定义这么多不同类之前，
%       \TeX 的内存就会耗尽。
% \end{advise}
%
% 最后一个提示：保持标识符列表互不相交。永远不要在`强调'列表中使用关键词，
% 也不要在两个不同列表中使用同一个名称。即使你的源代码按预期高亮显示，
% 也不能保证如果你更改列表顺序或使用本宏包的下一版本时仍然如此。
%
%
%\iffalse
% \subsection{列表对齐}\label{uListingAlignment}
%
% 这些例子使用居中的 \texttt{minipage} 排版。这就是为什么你看不到
% 行号打印在边距中。现在我们用一条竖线分隔 minipage 边距和 minipage：
% \begin{lstsample}{\lstset{frame=l,framesep=0pt,numberstyle=\tiny,stepnumber=2,^^A
%     numbersep=5pt}}{}
%    Some text before
%    \begin{lstlisting}
%    for i:=maxint to 0 do
%    begin
%        { do nothing }
%    end;
%    \end{lstlisting}
% \end{lstsample}
% 列表与正文对齐。参数 \ikeyname{xleftmargin} 将列表向右移动
%（如果尺寸为负则向左移动）。
% \begin{lstsample}{\lstset{frame=l,framesep=0pt,numberstyle=\tiny,stepnumber=2,^^A
%     numbersep=5pt}}{}
%    Some text before
%    \begin{lstlisting}[xleftmargin=15pt]
%    for i:=maxint to 0 do
%    begin
%        { do nothing }
%    end;
%    \end{lstlisting}
%
%    \begin{lstlisting}{ }
%    Write('Insensitive');
%    WritE('keywords.');
%    \end{lstlisting}
% \end{lstsample}
% 再次注意，可选参数只改变单个列表的设置。
%
% 如果你使用 \texttt{itemize} 或 \texttt{enumerate} 之类的环境，就会有
% 这些环境带来的“自然”缩进。默认情况下，
% \packagename{listings} 宏包会遵循这种缩进。不过你也可以用
% |resetmargins=true|（或 |false|）自行决定。当然可以把它与
% |xleftmargin| 一起使用。
% \begin{advise}
% \item 某些列表会产生严重的 |\hbox| 溢出。
%       \advisespace
%       这是因为列表中有过长的行。有几种办法可以
%       消除这种 |\hbox| 溢出。首先，建议用比周围正文
%       更小的字号排版列表，例如
%       `|basicstyle=\small|'。其次，可以使用灵活
%       列格式。第三，可以增大行宽，或者显式
%       设定行宽，参见第~\ref{rMarginsAndLineShape} 节。
%       如果这些都不管用，可以试着修改
%       `\ikeyname{basewidth}'，但要小心！还有两项内容
%       会在下一节解释。
% \end{advise}
%
% 你可能需要用
% \ikeyname{boxpos} 键来控制列表的垂直位置，例如在 \texttt{minipage} 或
% \texttt{tabular} 环境中使用时。这里“列表”指 \texttt{lstlisting} 或
% |\lstinputlisting|。如下例所示，你甚至可以把这样的列表
% 放在段落内部，但必须用
% |\hbox{| 和 |}| 把列表括起来，才能让宏包这样做。
% \begin{advise}
% \item 在 \LaTeX\ 文档中使用 \TeX 原语 `|\hbox|' 是否得体？
%       \advisespace
%       不，并不好。但本例中 \LaTeX 的 `|\mbox|' 不起作用：
% \end{advise}
% \begin{lstsample}{}{}
%    Here are some multi-line listings inside a paragraph.
%    The `boxpos' key controls their vertical alignment:
%    \hbox{\begin{lstlisting}[boxpos=c]
%    center
%    center
%    \end{lstlisting}}
%    \hbox{\begin{lstlisting}[boxpos=b]
%    bottom baseline
%    bottom baseline
%    \end{lstlisting}}
%    \hbox{\begin{lstlisting}[boxpos=t]
%    top baseline
%    top baseline
%    \end{lstlisting}}
% \end{lstsample}
%\fi
%
%
% \subsection{索引}\label{uIndexing}
%
% 索引就像强调标识符——我是说用法：
% \begin{lstxsample}[index]
%    \lstset{index={square},index={[2]root}}
% \end{lstxsample}
% \begin{lstsample}{}{}
%    \begin{lstlisting}
%    for i:=maxint to 0 do
%    begin
%        j:=square(root(i));
%    end;
%    \end{lstlisting}
% \end{lstsample}
% 当然，你在这里什么也看不到。你得去看索引。
% \begin{advise}
% \item 为什么 `\ikeyname{index}' 键能与多个标识符列表一起工作？
%       \advisespace
%       这个问题与 `{\rstyle\ikeyname{indexstyle}}' 键密切相关。有人可能想
%       创建多个索引，或者想插入 `|constants|'、`|functions|'、`|keywords|'
%       等前缀。`\ikeyname{indexstyle}' 键的工作方式与其他样式键相同，
%       只是最后一个记号\emph{必须}接受一个参数，即当前标识符的
%       （可打印形式）。
%
%       你可以定义 `|\newcommand\indexkeywords[1]{\index{keywords, #1}}|'，
%       并为常量或函数名做类似定义。然后 `|indexstyle=[1]\indexkeywords|'
%       可能满足你的目的。如果你想用
%       \href{http://mirror.ctan.org/macros/latex/contrib/camel}
%       {\packagename{index}} 宏包创建多个索引，这会变得更容易。
%       如果你定义了合适的新索引，就可以写 `|indexstyle=\index[keywords]|'。
%
% \item 比方说，我想索引所有关键词。如果常用的编程语言经常变化，
%       再输入所有关键词会很烦人。
%       \advisespace
%       继续往下读。
% \end{advise}
% \ikeyname{index} 键实际上有两个可选参数。第一个是大家熟悉的
% \meta{class number}，第二个是其他关键词类的逗号分隔列表，
% 它们的标识符会被索引。被索引的标识符随后随定义的关键词自动更改——
% 不是自动变魔术，这不是幻觉。^^A :-)
%
% 最终你需要知道关键词类的名称。通常是键名后加类编号，例如 |emph2|、
% |emph3|、\ldots、|keywords2| 或 |index5|。但第一级类
% |keywords|、|emph|、|directives| 等没有编号。
% \begin{advise}
% \item `|index=[keywords]|' 不起作用。
%       \advisespace
%       宏包猜不出你指的是哪个可选参数。因此如果你想使用第二个，
%       必须同时指定两个。你应该试试 `|index=[1][keywords]|'。
% \end{advise}
%
%
% \subsection{固定列与灵活列}\label{uFixedAndFlexibleColumns}
%
% 读者注意到的第一件事——除了关键词等的不同样式外——就是列对齐。
% Arne John Glenstrup 于 1997 年发明了灵活列格式。此后人们做了一些努力
% 来发展这个分支。目前提供四种列格式：固定、灵活、空格灵活和全灵活。
% 仔细看看下面的例子。
% \begin{center}
%     \def\sample{\begin{lstlisting}^^J WOMEN\ \ are^^A
%                                   ^^J \ \ \ \ \ \ \ MEN^^A
%                                   ^^J WOMEN are^^A
%                                   ^^J better MEN^^J \end{lstlisting}}
%     \newdimen\lstdocwidth
%     \lstdocwidth=80pt
%     \lstset{style={},language={}}
%     \hbox{\begin{minipage}[t]{\lstdocwidth}
%         \centering%
%         \texttt{{\rstyle\ikeyname{columns}}=}\\
%         \phantom{(at ?)}\\
%         {\lstset{basicstyle=\ttfamily, basewidth=0.51em}
%         \sample}
%     \end{minipage}
%     \quad
%     \begin{minipage}[t]{\lstdocwidth}
%         \centering%
%         \texttt{fixed}\\
%         (at {\makeatletter\lst@widthfixed})\\
%         {\lstset{columns=fixed}
%         \sample}
%     \end{minipage}
%     \quad
%     \begin{minipage}[t]{\lstdocwidth}
%         \centering%
%         \texttt{flexible}\\
%         (at {\makeatletter\lst@widthflexible})\\
%         {\lstset{columns=flexible}
%         \sample}
%     \end{minipage}
%     \quad
%     \begin{minipage}[t]{\lstdocwidth}
%         \centering%
%         \texttt{fullflexible}\\
%         (at {\makeatletter\lst@widthflexible})\\
%         {\lstset{columns=fullflexible}
%         \sample}
%     \end{minipage}}
% \end{center}
% \begin{advise}
% \item 为什么女人比男人强？
%       \advisespace
%       你想进行哲学思考吗？好吧，我什么时候说过 女人比男人强' 这句话是真的？
%       我甚至不记得说过 女人是男人' \ldots . ^^A ;-)
%       关于“女人是男人”…… ^^A ;-)
% \end{advise}
% 概括来说：固定列格式破坏了字体设计师设计的间距，而灵活格式破坏了
% 程序员（可能）设计的列对齐。它们的共同点是输入字符被转换为一系列
% 基本输出单元，如
% \begingroup \lstset{gobble=6,xleftmargin=\leftmargini}
% \makeatletter
%^^A  Make \fbox around each output unit.
% \fboxsep=0pt
% \def\lst@alloverstyle#1{\fbox{\kern-\fboxrule\strut#1}\kern-\fboxrule}
% \begin{lstlisting}[basewidth=1em]
%     if x=y then write('align')
%            else print('align');
% \end{lstlisting}
% 现在，固定格式把 $n$ 个字符放进宽度为 $n\times{}$
% 的“基本宽度”盒子中，本例中的基本宽度为 {\makeatletter\lst@widthfixed}。
% 该格式会压缩或拉伸字符之间的空隙，
% 使其适应盒子。如例所示，某些字符串看起来像
%    \hbox to 2em{b\hss a\hss d}
% or
%    \hbox to 2em{w\hss o\hss r\hss s\hss e},
% 但输出是垂直对齐的。
% \endgroup
%
% 如果你不需要或不喜欢这种效果，就应当使用灵活格式。所有
% 字符都按其自然宽度排版。特别是它们绝不会
% 重叠。如果一个词需要的空间超过预留值，该行其余部分
% 就直接向右移动。三种格式的区别在于：
% 完全灵活格式不再关心其他任何因素，而普通灵活格式
% 和空格灵活格式则会设法在字符
% 串占用的空间少于“预留”空间时修正列对齐。普通灵活格式会在空格处、
% 标识符前后，以及其他字符序列前后插入补偿空隙
% 以修正对齐；
% 空格灵活格式则只通过拉伸
% 已有的空格来插入补偿空隙。在上面的灵活列示例中，两个 MEN 是垂直
% 对齐的，因为第四行中插入了空格以修正
% 对齐。而在完全灵活格式中，两个 MEN 并不对齐。
%
% 注意两种灵活模式都将第一行中的两个空格打印为单个空格，但原因不同：
% 普通灵活格式修正列对齐（空格灵活格式也是如此），而完全灵活格式
% 不关心第二个空格。
%
%
% \section{高级技术}\label{uAdvancedTechniques}
%
%
% \subsection{样式定义}
%
% 显然，像这样的美化排版工具需要某种语言选择和定义。前者已经描述过，
% 后者由下一节讨论。不过，如果打印样式也能这样就很方便：在文档的一个
% 中心位置可以轻松修改它们，更改对所有列表生效。
%
% 与语言类似，
%    {\rstyle\ikeyname{style}}|=|\meta{style name}
% 激活先前定义的样式。定义同样简单：
%    {\rstyle|\lstdefinestyle|}\marg{style name}\marg{key=value list}。
% 在这样的定义中未使用的键当然不会被相应的样式选择影响。例如，你可以写
% \begin{verbatim}
%   \lstdefinestyle{numbers}
%       {numbers=left, stepnumber=1, numberstyle=\tiny, numbersep=10pt}
%   \lstdefinestyle{nonumbers}
%       {numbers=none}\end{verbatim}
% 然后只需分别用 |style=nonumbers| 和 |style=numbers|，就能在带行号的列表
% 和不带行号的列表之间切换。
% \begin{advise}
% \item 你甚至可以写
%           `|\lstdefinestyle{C++}{language=C++,style=numbers}|'。
%       样式名和语言名彼此独立，因此可能重合。此外还可以激活其他样式。
%
% \item 用样式很容易搞崩宏包。写
%       '|\lstdefinestyle{crash}{style=crash}|' 和 '|\lstset{style=crash}|'。
%       \TeX 的容量会耗尽，抱歉 [参数栈大小]。只有坏男孩才会用这种递归调用，
%       但只有好女孩才用这个宏包。所以这个问题不太重要。^^A :-)
% \end{advise}
%
%
% \subsection{语言定义}\label{uLanguageDefinitions}
%
% 这些类似于样式定义，只是多了一个可选的方言名和一个可选的基础语言——
% 当然，还有不同的命令名和专门的键。简单情况下是
%    {\rstyle|\lstdefinelanguage|}\marg{language name}\marg{key=value list}。
% 对许多编程语言来说，指定关键词和标准函数名、注释和字符串就足够了。
% 我们来看一个例子。
% \begin{lstxsample}[morekeywords,sensitive,morecomment,morestring]
%    \lstdefinelanguage{rock}
%      {morekeywords={one,two,three,four,five,six,seven,eight,
%          nine,ten,eleven,twelve,o,clock,rock,around,the,tonight},
%       sensitive=false,
%       morecomment=[l]{//},
%       morecomment=[s]{/*}{*/},
%       morestring=[b]",
%      }
% \end{lstxsample}
% \begingroup \csname lst@EndWriteFile\endcsname
% \bigbreak
%
% \noindent
% 关于关键词没什么好说的。它们的定义方式和你想强调的标识符一样。
% 此外你需要指定它们是否区分大小写。是的：你可以在关键词
% \texttt{one} 前面插入 |[2]|，把关键词定义为`第二级'，
% 并用 |keywordstyle={[2]...}| 打印它们。
% \begin{advise}
% \item 选择我的语言时出现 `\texttt{Missing = inserted for }|\ifnum|' 错误。
%       \advisespace
%       你是不是忘了在 `|keywords={...}|' 后面加逗号？如果你在选择语言
%       （或样式）后遇到意外字符，你可能忘了另一个逗号，或者给某个键
%       提供了太多参数，例如 |morecomment=[l]{--}{!}|。
% \end{advise}
%
% 那么我们来讨论注释和字符串。每个值以一个\emph{必需的} \oarg{type}
% 参数开始，后跟数量可变的开始和结束分隔符。注意每个分隔符（对）都需要
% 自己的键值，即使类型相同。因此，如果单引号开始和结束字符串或字符，
% 你需要插入 \texttt{morestring=[b]'}。
% 字面量，正如示例中的双引号那样。
%
% 最终你需要了解定界符的各种类型及其编号。
% 参考指南中有完整列表，这里只讨论最常见的几种。
% 对于字符串，它们是 {\rstyle\texttt{b}} 和 {\rstyle\texttt{d}}，各有一个
% 分隔符。这个分隔符开启和关闭字符串，在字符串内部它要么由一个
% \texttt{反斜杠} 转义，要么 \texttt{双写}。注释类型 {\rstyle\texttt{l}} 恰好需要一个
% 分隔符，它在任意列开始注释。该注释一直到行尾。其他两种最常见的注释
% 类型是 {\rstyle\texttt{s}} 和 {\rstyle\texttt{n}}，各有两个分隔符。第一个分隔符
% 开启注释，第二个分隔符终止注释。与 \texttt{s} 型不同，
% \texttt{n} 型注释可以嵌套。
% \begin{lstxsample}[b,d,l,s,n]
%    \lstset{morecomment=[l]{//},
%            morecomment=[s]{/*}{*/},
%            morecomment=[n]{(*}{*)},
%            morestring=[b]",
%            morestring=[d]'}
% \end{lstxsample}
% \begin{lstsample}{}{}
%    \begin{lstlisting}
%    "str\"ing "    not a string
%    'str''ing '    not a string
%    // comment line
%    /* comment/**/ not a comment
%    (* nested (**) still comment
%       comment  *) not a comment
%    \end{lstlisting}
% \end{lstsample}
% \begin{advise}
% \item 真的\emph{这么}容易吗？
%       \advisespace
%       差不多吧。不过也会遇到一些麻烦。例如，如果
%       `\texttt{-*}' 开始一个注释行，而 `\texttt{-*-}' 开始一个字符串
%       （虽不太可能但确实可以），那么必须先定义较短的那个
%       定界符。
%       另一个问题：默认情况下，关键词内不允许出现某些字符，例如
%       `\texttt{-}'、`\texttt{:}'、`\texttt{.}' 等。参考指南
%       通过引入若干额外的键来处理这个问题，让你可以
%       相应地调整标准字符表。但请注意，关键词内部
%       禁止出现空白字符。
% \end{advise}
% 最后请记住，本节只是语言定义
% 的入门介绍，还有更多的键和可能性。
%
%
% \subsection{定界符}\label{uDelimiters}
%
% 你已经知道两种特殊的分隔符类：注释和字符串。然而，它们的完整语法
% 到目前为止还没有描述。例如，\ikeyname{commentstyle} 应用于所有注释——
% 除非你指定不同的样式。\emph{可选} 的 \oarg{style} 参数跟在
% \emph{必需} 的 \oarg{type} 参数之后。
%\ifcolor
% \begin{lstxsample}
%    \lstset{morecomment=[l][keywordstyle]{//},
%            morecomment=[s][\color{white}]{/*}{*/}}
% \end{lstxsample}
%\else
% \begin{lstxsample}
%    \lstset{morecomment=[l][keywordstyle]{//},
%            morecomment=[s][\underbar]{/*}{*/}}
% \end{lstxsample}
%\fi
% \begin{lstsample}{}{}
%    \begin{lstlisting}
%    // bold comment line
%    a single /* comment */
%    \end{lstlisting}
% \end{lstsample}
% 如你所见，你可以选择用 \LaTeX 命令显式指定样式，或用其他样式键
% 隐式指定。但是，你说得对，有些隐式定义的样式没有单独的键，
% 例如第二级关键词样式。这里——永远不要用数字 1——你只需将级别号
% 附加到基础键上：\texttt{keywordstyle2}。
%
% 你问有什么用途？这就给你一个：可以为注释的“子类型”定义不同的排版
% 样式，例如
%\ifcolor
% \begin{lstxsample}
%    \lstset{morecomment=[s][\color{blue}]{/*+}{*/},
%            morecomment=[s][\color{red}]{/*-}{*/}}
% \end{lstxsample}
%\else
% \begin{lstxsample}
%    \lstset{morecomment=[s][\upshape]{/*+}{*/},
%            morecomment=[s][\bfseries]{/*-}{*/}}
% \end{lstxsample}
%\fi
% \begin{lstsample}{\lstset{morecomment=[s]{/*}{*/}}}{}
%    \begin{lstlisting}
%    /*  normal comment */
%    /*+    keep cool   */
%    /*-     danger!    */
%    \end{lstlisting}
% \end{lstsample}
% 这里，注释样式没有应用到第二行和第三行。
% \begin{advise}
% \item 请记住，两个“额外”注释都必须定义在
%       普通注释\emph{之后}，因为定界符 `\texttt{/*}' 是
%       `\texttt{/*+}' 和 `\texttt{/*-}' 的子串。
%
% \item 我还有一个问题。用 `\texttt{language=}\meta{其他
%       语言}' 是去掉这类额外定界符的唯一方法吗？
%       \advisespace
%       调用 {\rstyle\ikeyname{deletecomment}} 和/或
%       {\rstyle\ikeyname{deletestring}}，参数相同，即可去掉
%       这些定界符（但不需要提供可选的样式
%       参数）。
% \end{advise}
% 最后，你可能想在任何注释类型上使用前缀 \texttt{i}。这样注释不仅
% 不可见，而且会从输出中完全丢弃！
% \begin{lstxsample}[is]
%    \lstset{morecomment=[is]{/*}{*/}}
% \end{lstxsample}
% \begin{lstsample}{}{}
%    \begin{lstlisting}
%    begin /* comment */ end
%    begin/* comment */end
%    \end{lstlisting}
% \end{lstsample}
%
% 好，现在来看真正的挑战。更通用的定界符可以用
% 键 {\rstyle\ikeyname{moredelim}} 来定义。合法类型是 {\rstyle\texttt{l}}
% 和 {\rstyle\texttt{s}}。这些类型前面可以加一个 \texttt{i}，但
% 此时\emph{只有定界符}会从输出中删除。这样
% 你就可以通过标记来选择样式。
% \begin{lstxsample}
%    \lstset{moredelim=[is][\ttfamily]{|}{|}}
% \end{lstxsample}
% \begin{lstsample}{}{}
%    \begin{lstlisting}
%    roman |typewriter|
%    \end{lstlisting}
% \end{lstsample}
% 你甚至可以让宏包在内容内部检测关键词、注释、字符串和其他分隔符。
% \begin{lstxsample}
%    \lstset{moredelim=*[s][\itshape]{/*}{*/}}
% \end{lstxsample}
% \begin{lstsample}{}{}
%    \begin{lstlisting}
%    /* begin
%      (* comment *)
%       ' string ' */
%    \end{lstlisting}
% \end{lstsample}
% 此外，你可以强制样式累积应用。
% \begin{lstxsample}
%    \lstset{moredelim=**[is][\ttfamily]{|}{|}, % cumulative
%            moredelim=*[s][\itshape]{/*}{*/}}  % not so
% \end{lstxsample}
% \begin{lstsample}{}{}
%    \begin{lstlisting}
%    /* begin
%       ' string '
%       |typewriter| */
%
%    | begin
%     ' string '
%     /*typewriter*/ |
%    \end{lstlisting}
% \end{lstsample}
% 仔细看输出并注意区别。第二个 \texttt{begin} 没有用粗体打字机字体
% 打印，因为标准 \LaTeX 没有这样的字体。
%
% 介绍就到这里。现在去寻找更多应用吧。
%
%
% \subsection{结语与致谢}\label{uClosingAndCredits}
%
% 你已经看到了很多键，但离知道所有键还差得远。下一步是真正使用
% \packagename{listings} 宏包。请接受以下建议。首先，查阅参考指南中
% 已知的命令和键，了解那里的记号。其次，用这些键摸索，学习其他参数。
% 然后，希望你在遇到问题或需要特殊功能时就有准备了。
%
% \begin{advise}
% \item
% 有一个问题你在前几页一直没问：该怪谁。Carsten Heinz 写了指南，
% 编写了 \packagename{listings} 宏包，并写了一些语言驱动。
% Brooks Moses 接管维护数年，Jobst Hoffmann 当前维护该宏包。
% 其他人定义了更多语言或贡献了想法；还有许多人提交了错误报告，
% 但只列出第一个发现错误的人。
%^^A
%^^A Thanks for error reports (first bug finder only), new programming
%^^A languages, etc.
%^^A Special thanks for communication which lead to kernel extensions, and to
%^^A Hendri Adriaens for reviving maintenance on the package.
%^^A
% 特别感谢（按字母顺序）
% \begin{quote}
% \hyphenpenalty=10000\relax \rightskip=0pt plus \linewidth
%   \lstthanks{Hendri~Adriaens}{-},
%   \lstthanks{Andreas~Bartelt}{Andreas.Bartelt@Informatik.Uni-Oldenburg.DE},
%   \lstthanks{Jan~Braun}{Jan.Braun@tu-bs.de},
%   \lstthanks{Denis~Girou}{Denis.Girou@idris.fr},
%   \lstthanks{Arne~John~Glenstrup}{panic@diku.dk},
%   \lstthanks{Frank~Mittelbach}{frank.mittelbach@latex-project.org},
%   \lstthanks{Rolf~Niepraschk}{niepraschk@PTB.DE},
%   \lstthanks{Rui~Oliveira}{rco@di.uminho.pt},
%   \lstthanks{Jens~Schwarzer}{schwarzer@schwarzer.dk}，以及
%   \lstthanks{Boris~Veytsman}{boris@plmsc.psu.edu}.
% \end{quote}
% 此外我们还要感谢
% \begin{quote}
% \hyphenpenalty=10000\relax \rightskip=0pt plus \linewidth
%   \lstthanks{Nasser~M.~Abbasi}{nma@12000.org},
%   \lstthanks{Bj{\o}rn~{\AA}dlandsvik}{bjorn@imr.no},
%   \lstthanks{Omair-Inam~Abdul-Matin}{-},
%   \lstthanks{Gaurav~Aggarwal}{gaurav@ics.uci.edu},
%   \lstthanks{Jason~Alexander}{jalex@ea.oac.uci.edu},
%   \lstthanks{Andrei~Alexandrescu}{-},
%   \lstthanks{Holger~Arndt}{-},
%   \lstthanks{Donald~Arseneau}{ASND@erich.triumf.ca},
%   \lstthanks{David~Aspinall}{David.Aspinall@ed.ac.uk},
%   \lstthanks{Frank~Atanassow}{-},
%   \lstthanks{Claus~Atzenbeck}{Claus.Atzenbeck@stud.uni-regensburg.de},
%   \lstthanks{Michael~Bachmann}{-},
%   \lstthanks{Luca~Balzerani}{-},
%   \lstthanks{Peter~Bartke}{bartke@inf.fu-berlin.de} （非常感谢）， ^^A beta tester
%   \lstthanks{Jean-Yves~Baudais}{Jean-Yves.Baudais@insa-rennes.fr},
%   \lstthanks{Heiko~Bauke}{-},
%   \lstthanks{Oliver~Baum}{oli.baum@web.de},
%   \lstthanks{Ralph~Becket}{rbeck@microsoft.com},
%   \lstthanks{Andres~Becerra~Sandoval}{abecerra@univalle.edu.co},
%   \lstthanks{Kai~Below}{below@tu-harburg.de},
%   \lstthanks{Matthias~Bethke}{-},
%   \lstthanks{Javier~Bezos}{javier.bezos@bancoval.es},
%   \lstthanks{Olaf~Trygve~Berglihn}{olafb@pvv.org},
%      ^^A {1999/11/29}{3-char comment delimiter don't work (Python)}
%   \lstthanks{Karl~Berry}{karl@freefriends.org},
%   \lstthanks{Geraint~Paul~Bevan}{geraint@users.sf.net},
%   \lstthanks{Peter~Biechele}{peter.biechele@physik.uni-freiburg.de},
%   \lstthanks{Beat~Birkhofer}{beat@birkhofer.ch},
%   \lstthanks{Fr\'ed\'eric~Boulanger}{Frederic.Boulanger@supelec.fr},
%   \lstthanks{Byron~K.~Boulton}{bkboulton@berriehill.com},
%   \lstthanks{Joachim~Breitner}{-},
%   \lstthanks{Martin~Brodbeck}{Martin.Brodbeck@gmx.de},
%   \lstthanks{Walter~E.~Brown}{WB@fnal.gov},
%   \lstthanks{Achim~D.~Brucker}{brucker@informatik.uni-freiburg.de},
%   \lstthanks{J\'an Bu\v{s}a}{-},
%   \lstthanks{Thomas~ten~Cate}{-},
%   \lstthanks{David~Carlisle}{davidc@nag.co.uk},
%   \lstthanks{Bradford~Chamberlain}{brad@cs.washington.edu},
%   \lstthanks{Brian~Christensen}{-},
%   \lstthanks{Neil~Conway}{-},
%   \lstthanks{Patrick~Cousot}{Patrick.Cousot@wanadoo.fr},
%   \lstthanks{Xavier~Cr\'egut}{cregut@enseeiht.fr},
%   \lstthanks{Christopher~Creutzig}{-},
%   \lstthanks{Holger~Danielsson}{dani@fbg.schwerte.de},
%   \lstthanks{Andreas~Deininger}{deininger@uni-kassel.de},
%   \lstthanks{Robert~Denham}{Robert.Denham@dnr.qld.gov.au},
%   \lstthanks{Detlev~Dr\"oge}{droege@informatik.uni-koblenz.de},
%   \lstthanks{Anders~Edenbrandt}{Anders.Edenbrandt@dna.lth.se},
%   \lstthanks{Mark~van~Eijk}{mark@luon.net},
%   \lstthanks{Norbert~Eisinger}{Norbert.Eisinger@informatik.uni-muenchen.de},
%   \lstthanks{Brian~Elmegaard}{-},
%   \lstthanks{Jon~Ericson}{Jon.Ericson@jpl.nasa.gov},
%   \lstthanks{Luc~Van~Eycken}{Luc.VanEycken@esat.kuleuven.be}
%   \lstthanks{Thomas~Esser}{te@dbs.uni-hannover.de},
%   \lstthanks{Chris~Edwards}{edwch00p@infoscience.otago.ac.nz},
%   \lstthanks{David~John~Evans}{Matrix.Software@dial.pipex.com},
%   \lstthanks{Tanguy~Fautr\'e}{tfautre@pandora.be},
%   \lstthanks{Ulrike~Fischer}{news3@nililand.de},
%   \lstthanks{Robert~Frank}{rf7@ukc.ac.uk},
%   \lstthanks{Michael~Franke}{-},
%   \lstthanks{Ignacio~Fern\'andez~Galv\'an}{-},
%   \lstthanks{Martine~Gautier}{-}
%   \lstthanks{Daniel~Gazard}{gazard_d@epita.fr},
%   \lstthanks{Daniel~Gerigk}{Daniel.Gerigk@ePost.de},
%   \lstthanks{Dr.~Christoph~Giess}{-},
%   \lstthanks{KP~Gores}{kp.gores@web.de},
%   \lstthanks{Adam~Grabowski}{adam@mizar.org},
%   \lstthanks{Jean-Philippe~Grivet}{grivet@cnrs-orleans.fr},
%   \lstthanks{Christian~Gudrian}{Christian.Gudrian@kawo1.rwth-aachen.de},
%   \lstthanks{Jonathan~de~Halleux}{dehalleux@auto.ucl.ac.be},
%   \lstthanks{Carsten~Hamm}{carsten.hamm@siemens.com},
%   \lstthanks{Martina~Hansel}{Martina.Hansel@fhtw-berlin.de},
%   \lstthanks{Harald~Harders}{h.harders@tu-bs.de},
%   \lstthanks{Christian~Haul}{haul@dvs1.informatik.tu-darmstadt.de},
%   \lstthanks{Aidan~Philip~Heerdegen}{Aidan.Heerdegen@anu.edu.au},
%   \lstthanks{Jim~Hefferon}{Hefferon9@aol.com},
%   \lstthanks{Heiko~Heil}{info@heiko-heil.de},
%   \lstthanks{J\"urgen~Heim}{heim@astro.uni-tuebingen.de},
%   \lstthanks{Martin~Heller}{-},
%   \lstthanks{Stephan~Hennig}{-},
%   \lstthanks{Alvaro~Herrera}{alvherre@dcc.uchile.cl},
%   \lstthanks{Richard~Hoefter}{hoefter@gmx.de},
%   \lstthanks{Dr.~Jobst~Hoffmann}{HOFFMANN@rz.rwth-aachen.de},
%   \lstthanks{Torben~Hoffmann}{toho@it.dtu.dk},
%   \lstthanks{Morten~H\o gholm}{-},
%   \lstthanks{Berthold~H\"ollmann}{bhoel@starship.python.net},
%   \lstthanks{G\'erard~Huet}{-},
%   \lstthanks{Hermann~H\"uttler}{hermann.huettler@gmx.net},
%   \lstthanks{Ralf~Imh\"auser}{snoopy@tribal.line.org},
%   \lstthanks{R.~Isernhagen}{R.Isernhagen@FH-Wolfenbuettel.DE},
%   \lstthanks{Oldrich~Jedlicka}{ojedlick@students.zcu.cz},
%   \lstthanks{Dirk~Jesko}{jesko@iti.cs.uni-magdeburg.de},
%   \lstthanks{Lo\"\i c~Joly}{-},
%   \lstthanks{Christian~Kaiser}{chk@combit.net},
%   \lstthanks{Bekir~Karaoglu}{karabekirus@yahoo.com},
%   \lstthanks{Marcin~Kasperski}{Marcin.Kasperski@softax.com.pl},
%   \lstthanks{Christian~Kindinger}{chkind@uni-wuppertal.de},
%   \lstthanks{Steffen~Klupsch}{steffen@vlsi.informatik.tu-darmstadt.de},
%   \lstthanks{Markus~Kohm}{-},
%   \lstthanks{Peter~K\"oller}{pkoeller@metaprojekt.de} （非常感谢）， ^^A beta tester
%   \lstthanks{Reinhard~Kotucha}{Reinhard.Kotucha@web.de},
%   \lstthanks{Stefan~Lagotzki}{info@lagotzki.de},
%   \lstthanks{Tino~Langer}{langer@tournex.de},
%   \lstthanks{Rene~H.~Larsen}{rhl@traceroute.dk},
%   \lstthanks{Olivier~Lecarme}{ol@i3s.unice.fr},
%   \lstthanks{Thomas~Leduc}{Thomas.Leduc@lsv.ens-cachan.fr},
%   \lstthanks{Qing~Lee}{-},
%   \lstthanks{Dr.~Peter~Leibner}{Peter.Leibner@sta.siemens.de},
%   \lstthanks{Thomas~Leonhardt}{leonhardt@informatik.tu-darmstadt.de}
%      （非常感谢）， ^^A beta tester
%   \lstthanks{Magnus~Lewis-Smith}{Magnus.Lewis-Smith@pace.co.uk},
%   \lstthanks{Knut~Lickert}{knut.lickert@gmx.de},
%   \lstthanks{Benjamin~Lings}{-},
%   \lstthanks{Dan~Luecking}{luecking@uark.edu},
%   \lstthanks{Peter~L\"offler}{-},
%   \lstthanks{Markus~Luisser}{-},
%   \lstthanks{Kris~Luyten}{no email available},
%   \lstthanks{Jos\'e~Romildo~Malaquias}{romildo@urano.iceb.ufop.br},
%   \lstthanks{Andreas~Matthias}{amat@kabsi.at},
%   \lstthanks{Patrick~TJ~McPhee}{ptjm@interlog.com},
%   ^^A \lstthanks{Brooks~Moses}{-},
%   \lstthanks{Riccardo~Murri}{riccardo.murri@gmx.it},
%   \lstthanks{Knut~M\"uller}{knut@physik3.gwdg.de},
%   \lstthanks{Svend~Tollak~Munkejord}{svendm@efisms.energy.sintef.no},
%   \lstthanks{Gerd~Neugebauer}{gerd.neugebauer@gmx.de},
%   \lstthanks{Torsten~Neuer}{tneuer@inwise.de},
%   \lstthanks{Enzo~Nicosia}{-},
%   \lstthanks{Michael~Niedermair}{m.g.n@gmx.de},
%   \lstthanks{Xavier~Noria}{fxn@hashref.com},
%   \lstthanks{Heiko~Oberdiek}{oberdiek@ruf.uni-freiburg.de},
%   \lstthanks{Xavier~Olive}{-},
%   \lstthanks{Sebastian Ørsted}{sorsted@gmail.com},
%   \lstthanks{Alessio~Pace}{-},
%   \lstthanks{Markus~Pahlow}{pahlowm@mar.dfo-mpo.gc.ca},
%   \lstthanks{Morten~H.~Pedersen}{mhp@dadlnet.dk},
%   \lstthanks{Xiaobo~Peng}{-},
%   \lstthanks{Zvezdan~V.~Petkovic}{zpetkovic@acm.org},
%   \lstthanks{Michael~Piefel}{piefel@informatik.hu-berlin.de},
%   \lstthanks{Michael~Piotrowski}{mxp@linguistik.uni-erlangen.de},
%   \lstthanks{Manfred~Piringer}{sz0490@rrze.uni-erlangen.de},
%   \lstthanks{Ivo~Pletikosi\'c}{ivo.pletikosic@gmail.com},
%   \lstthanks{Vincent~Poirriez}{Vincent.Poirriez@univ-valenciennes.fr},
%   \lstthanks{Adam~Prugel-Bennett}{apb@ecs.soton.ac.uk},
%   \lstthanks{Ralf~Quast}{rquast@hs.uni-hamburg.de},
%   \lstthanks{Aslak~Raanes}{araanes@ifi.ntnu.no},
%   \lstthanks{Venkatesh~Prasad~Ranganath}{vranganath@cox.net},
%   \lstthanks{Tobias~Rapp}{-},
%   \lstthanks{Jeffrey~Ratcliffe}{-},
%   \lstthanks{Georg~Rehm}{Georg.Rehm@germanistik.uni-giessen.de},
%   \lstthanks{Fermin~Reig}{reig@ics.uci.edu},
%   \lstthanks{Detlef~Reimers}{dreimers@aol.com},
%   \lstthanks{Stephen~Reindl}{stephen.reindl@vodafone.com},
%   \lstthanks{Franz~Rinnerthaler}{-},
%   \lstthanks{Peter~Ruckdeschel}{Peter.Ruckdeschel@uni-bayreuth.de},
%   \lstthanks{Magne~Rudshaug}{magne@ife.no},
%   \lstthanks{Jonathan~Sauer}{jonathan.sauer@gmx.de},
%   \lstthanks{Vespe~Savikko}{vespe@cs.tut.fi},
%   \lstthanks{Mark~Schade}{-},
%   \lstthanks{Gunther~Schmidl}{gschmidl@gmx.at},
%   \lstthanks{Andreas~Schmidt}{-},
%   \lstthanks{Walter~Schmidt}{wschmi@arcor.de},
%   \lstthanks{Christian~Schneider}{-},
%   \lstthanks{Jochen~Schneider}{jschneider@ds3.etech.haw-hamburg.de},
%   \lstthanks{Sven~Schreiber}{svetosch@gmx.net},
%   \lstthanks{Benjamin~Schubert}{benjamin.schubert@berlin.de},
%   \lstthanks{Sebastian~Schubert}{-},
%   \lstthanks{Uwe~Siart}{uwe.siart@ei.tum.de},
%   \lstthanks{Axel~Sommerfeldt}{axel@sommerfeldt.net},
%   \lstthanks{Richard~Stallman}{-},
%   \lstthanks{Nigel~Stanger}{nstanger@infoscience.otago.ac.nz},
%   \lstthanks{Martin~Steffen}{ms@informatik.uni-kiel.de},
%   \lstthanks{Andreas~Stephan}{Andreas.Stephan@victoria.de},
%   \lstthanks{Stefan~Stoll}{stoll@phys.chem.ethz.ch},
%   \lstthanks{Enrico~Straube}{no email available},
%   \lstthanks{Werner~Struckmann}{struck@ips.cs.tu-bs.de},
%   \lstthanks{Martin~S\"u\ss kraut}{Edon.Myder@web.de},
%   \lstthanks{Gabriel~Tauro}{gabriel@informatik.uni-jena.de},
%   \lstthanks{Winfried~Theis}{theis@statistik.uni-dortmund.de},
%   \lstthanks{Jens~T.~Berger~Thielemann}{jensthi@ifi.uio.no},
%   \lstthanks{William~Thimbleby}{-},
%   \lstthanks{Arnaud~Tisserand}{arnaud.tisserand@ens-lyon.fr},
%   \lstthanks{Jens~Troeger}{-},
%   \lstthanks{Kalle~Tuulos}{kalle.tuulos@nic.fi},
%   \lstthanks{Gregory~Van~Vooren}{Gregory.VanVooren@rug.ac.be},
%   \lstthanks{Timothy~Van~Zandt}{tvz@econ.insead.edu},
%   \lstthanks{J\"org~Viermann}{-},
%   \lstthanks{Thorsten~Vitt}{vitt@informatik.hu-berlin.de},
%   \lstthanks{Herbert~Voss}{voss@perce.de} （非常感谢）， ^^A beta tester
%   \lstthanks{Edsko~de~Vries}{devriese@tcd.ie},
%   \lstthanks{Herfried~Karl~Wagner}{hirf@gmx.at},
%   \lstthanks{Dominique~de~Waleffe}{ddw@miscrit.be},
%   \lstthanks{Bernhard~Walle}{-},
%   \lstthanks{Jared~Warren}{warren@cs.queensu.ca},
%   \lstthanks{Michael~Weber}{mweber@informatik.hu-berlin.de},
%   \lstthanks{Sonja~Weidmann}{Sonja.Weidmann@gmx.de},
%   \lstthanks{Andreas~Weidner}{-},
%   \lstthanks{Herbert~Weinhandl}{weinhand@grz08u.unileoben.ac.at},
%   \lstthanks{Robert~Wenner}{robert.wenner@gmx.de},
%   \lstthanks{Michael~Wiese}{wiese@itwm.uni-kl.de},
%   \lstthanks{James~Willans}{-},
%   \lstthanks{J\"orn~Wilms}{wilms@rocinante.colorado.edu},
%   \lstthanks{Kai~Wollenweber}{kai@ece.WPI.EDU},
%   \lstthanks{Ulrich~G.~Wortmann}{uliw@erdw.ethz.ch},
%   \lstthanks{Cameron~H.G.~Wright}{-},
%   \lstthanks{Joseph~Wright}{joseph.wright@morningstar2.co.uk},
%   \lstthanks{Andrew~Zabolotny}{-}，以及
%   \lstthanks{Florian~Z\"ahringer}{-}.
% \end{quote}
% 可能还有其他人贡献了这个宏包。如果我漏掉了你的名字，请发邮件告诉我。
% \end{advise}
%
%
% \part{参考指南}
%
%
% \section{主要参考}\label{rMainReference}
%
% 你的初步训练完成了。既然你已经离开了用户指南，告诉你该做什么的朋友
% 已经走了。多练习，成为熟练工吧！^^A :-)
% \begin{advise}
% \item 实际上，朋友还没走。仍然有一些建议，但只是偶尔出现。
% \end{advise}
%
%
% \subsection{如何阅读参考}
%
% 命令、键和环境按如下方式呈现。
% \begin{syntax}
% \item[1.0,default,hints] \texttt{command}、\texttt{environment} 或
%       \keyname{key} 带 \meta{parameters}
%
%       本栏包含说明；这里我们描述其他栏。
%
%       如果存在，左页边的标签提供额外信息：
%       `\textit{addon}' 表示额外引入的功能，
%       `\textit{changed}' 表示修改过的键，
%       `\textit{data}' 表示仅包含数据的命令（因此可以通过
%       |\renewcommand| 调整），等等。有些键和功能标有 `\emph{bug}'
%       或 \dag 符号。这些功能将来可能更改或被移除，所以请谨慎使用。
%
%       如果有触及右页边的逐字文本，那就是预定义值。注意有些键对每个列表
%       都默认此值，即只能在单个列表上使用的键。
% \end{syntax}
% 关于参数，请记住以下几点：
% \begin{enumerate}
% \item 列表总是指逗号分隔的列表。你必须用花括号括住这样的列表。
%       否则你会和 \packagename{keyval} 宏包发生麻烦；它会抱怨未定义的键。
% \item 如果你在可选的 \oarg{key=value list} 内使用键的
%       \oarg{optional argument}，你必须用参数花括号括住键的整个值：
%       |\begin{lstlisting}[caption=|{\rstyle|{|}|[one]two|{\rstyle|}|}|]|。
% \item 方括号 `|[ ]|' 通常括住可选参数，必须逐字输入。
%       普通方括号 `[ ]' 总是表示可选参数，不要输入。因此 |[*]| 必须
%       原样输入，但 [|*|] 如果你使用这个参数就只得到 |*|。
% \item 竖线表示二选一，例如 ^^A
%       \meta{\alternative{true,false}} 允许 \texttt{true} 或
%       \texttt{false} 作为参数。
% \item 如果你想输入特殊字符 |{}#%\| 之一，该字符必须用反斜杠转义。
%       这意味着你必须写 |\}| 来表示单个字符`右花括号'——当然不是用于
%       结束参数字符。
% \end{enumerate}
%
%
% \subsection{排版列表}\label{rTypesettingListings}
%
% \begin{syntax}
%   \item[0.19] \rcmdname\lstset\marg{key=value list}
%
%     设置指定键的值，另见第
%     \ref{uTheKey=ValueInterface} 节。这些参数的值保持到
%     当前组结束为止。相反，下面所有可选的
%     \meta{键=值列表} 只修改单个
%     列表的参数。
%
%   \item[0.18]
%     \rcmdname\lstinline\oarg{key=value list}\meta{character}^^A
%     \meta{source code}\meta{same character}
%
%     作用类似 |\verb|，但会遵循当前活动的语言和样式。
%     这些列表使用灵活列，除非在可选参数中另有
%     要求；它们不支持边框或背景
%     色。你可以写 `|\lstinline!var i:integer;!|'，得到
%     `\lstinline!var i:integer;!'.
%
%     由于该命令会先向后查找可选参数，因此如果你想用 |[| 作为
%     \meta{字符}，就至少得提供一个空的可选参数。
%
%     \dag\ 为支持
%     语法 |\lstinline|\oarg{键=值列表}\marg{源代码}，已经实现了一个实验版本。你可以试用，
%     并反馈成功或失败的情况。一个已知限制是：
%     在另一个参数内部时，源代码的最后一个记号不能
%     是显式的空格记号——当然，在另一个参数内部使用列表本身
%     就是实验性的，参见第
%     \ref{rListingsInsideArguments} 节。
%
%     另一个限制是该功能不能用于
%     |tabular| 环境的单元格。变通方法见第~\ref{uListingsArguments} 节。
%
%     关于创建 |\lstinline| 命令简短别名的命令，另见第~\ref{rShortInline} 节。
%
%   \item[0.15] |\begin{|\texttt{\rstyle lstlisting}|}|\oarg{key=value
%           list}
%
%       \leavevmode\hspace*{-\leftmargini}|\end{|\texttt{\rstyle
%       lstlisting}|}|
%
%       把夹在中间的代码排版为行间列表。
%
%       与 \packagename{verbatim} 宏包的环境不同，
%       同一行上的 \LaTeX 代码以及环境结束之后的代码会被排版，
%       或者说被执行。
%
%     \item[0.1] \rcmdname\lstinputlisting\oarg{key=value list}\marg{file
%       name}
%
%       把独立的源代码文件排版为行间列表。
%   \end{syntax}
%
%
%   \subsection{选项}
%
% 以下各节描述了所有可用于影响列表外观的键。
%
%
% \subsubsection{查找文件}
%
% \begin{syntax}
%   \item[1.8, {}] \rkeyname{inputpath}|=|\meta{path}
%
%     定义 \meta{文件名} 所指文件所在的路径。
%
%     \rkeyname{inputpath} 会覆盖 |TEXINPUTS| 环境变量，
%     也就是说，如果某个文件位于
%     |TEXINPUTS| 给出的某个路径上，而 \meta{路径} 不在
%     |TEXINPUTS| 中，该文件就不会再被找到。
%
%     把 \rkeyname{inputpath} 作为 \rcmdname\lstinputlisting{} 的选项设置时，
%     会覆盖 \rcmdname\lstset 设定的值。
% \end{syntax}
%
%
% \subsubsection{空间与位置}
%
% \begin{syntax}
%   \item[0.20,floatplacement] \rkeyname{float}|=|[|*|]\meta{subset of
%     \textup{\texttt{tbph}}}\syntaxor\rkeyname{float}
%
%     只对单独的显示列表有意义，使它们可以
%     浮动。该参数控制 \LaTeX\ \emph{允许}把浮动体放在
%     哪里：当前页或下一页的顶部或底部，单独
%     的一页，或者就在列表所在处。
%
%     可选星号可用于在双栏文档中得到跨栏的
%     浮动体。
%
%   \item[0.21,tbp] \rkeyname{floatplacement}|=|\meta{place specifiers}
%
%     当 \keyname{float} 不带值时，就用它作为位置说明符。
%
%   \item[0.21,\medskipamount] \rkeyname{aboveskip}|=|\meta{dimension}
%   \item[0.21,\medskipamount] \rkeyname{belowskip}|=|\meta{dimension}
%
%     定义显示列表上下的间距。
%
%   \item[0.17,0pt,\dag] \rkeyname{lineskip}|=|\meta{dimension}
%
%     指定列表内各行之间的附加间距。
%
%   \item[0.18,c,\dag] \rkeyname{boxpos}|=|\meta{\alternative{b,c,t}}
%
%     有时 \packagename{listings} 宏包会给列表加上一个 |\hbox|——否则
%     它就无法正常打印甚至正确处理。这个
%     键决定相对于周围材料的垂直对齐方式：
%     底部基线、居中或顶部基线。
% \end{syntax}
%
%
% \subsubsection{打印范围}
%
% \begin{syntax}
%   \item[0.12,true]
%     \rkeyname{print}|=|\meta{\alternative{true,false}}\syntaxor\rkeyname{print}
%
%     控制是否排版单个显示列表。即使
%     设为 false，相应的标题仍会打印，标签仍会
%     定义。
%
%     注意：如果加载宏包时没有使用 \texttt{draft} 选项，你可以
%     把这个键与 |\lstset| 一起使用。反之，该键
%     可用于在使用 \texttt{draft} 选项时仍然排版特定的
%     列表。
%
%   \item[0.1,1] \rkeyname{firstline}|=|\meta{number}
%   \item[0.1,9999999] \rkeyname{lastline}|=|\meta{number}
%
%     只能用于单个列表。它们决定打印显示列表时所使用的实际
%     输入行。
%
%   \item[1.2;1.9]
%     \rkeyname{linerange}|={|\meta{first1}\texttt-\meta{last1}\texttt,^^A
%         \meta{first2}\texttt-\meta{last2}\texttt, 等等
%         |}|\label{uoption:linerange}
%
%     只能用于单个列表。给定的列表行范围
%     会被显示。各区间必须有序，且不能
%     相交。
%
%     事实上三元组 \meta{first}\texttt-\meta{last} 的每一部分都可以
%     省略。省略 \meta{first} 或 \meta{last} 意味着
%     范围从第 1 行开始、到 \meta{last} 结束，或从
%     \meta{first} 开始、到文件最后一行结束；单个数字
%     则意味着范围从该数字所给的行开始并结束于该行，
%     而省略 \meta{first} 或 \meta{last} 意味着范围
%     从第 1 行开始、到文件最后一行结束。
%
%   \item[1.8, true]
%     \rkeyname{consecutivenumbers}|=|\meta{\alternative{true,false}}^^A
%     \syntaxor\rkeyname{consecutivenumbers}\label{uoption:consecutivenumbers}
%
%     只能用于单个列表。只有同时使用了 \rkeyname{linerange} 时
%     它才有意义。默认值（true）意味着\emph{所有}行范围的
%     行号是连续编号的，
%     例如 1、2、3、……如果设为 false，不同范围会
%     各自编号（见第~\ref{uLineNumbers} 节）。
%
%   \item[0.20,false]
%     \rkeyname{showlines}|=|\meta{\alternative{true,false}}\syntaxor\rkeyname{showlines}
%
%     若为 true，宏包会在列表末尾打印空行。
%     否则这些行会被丢弃（但仍计入
%     行号计数）。
%
%   \item[1.0] \rkeyname{emptylines}|=|[|*|]\meta{number}
%
%     允许的空行最大数。如果空行块超过
%     \meta{number} 行，则只打印 \meta{number} 行。
%     不加可选星号时，省略空行会打乱行号；
%     加了星号，各行则保留其原始行号。
%
%   \item[0.19,0] \rkeyname{gobble}|=|\meta{number}
%
%     吃掉 \emph{环境}代码每行开头的 \meta{number} 个
%     字符。该键对 \cs{lstinline} 或 \cs{lstinputlisting}
%     不起作用。
%
%     制表符在被吃掉之前先展开为 \ikeyname{tabsize} 个空格。
%     字符数少于 \ikeyname{gobble} 的代码行
%     视为空行。切勿让环境结尾的缩进超过
%     该字符数。
% \end{syntax}
%
%
% \subsubsection{语言与样式}\label{rLanguagesAndStyles}
%
% 请注意，参数 \meta{language}、\meta{dialect} 和 \meta{style name}
% 不区分大小写，空格也不起作用。
% \begin{syntax}
% \item[0.18,{{}}] \rkeyname{style}|=|\meta{style name}
%
%       激活由 |\lstdefinestyle| 保存的键=值列表。
%
% \item[0.19] \rcmdname\lstdefinestyle\marg{style name}\marg{key=value list}
%
%       保存该键=值列表。
%
% \item[0.17,{{}}] \rkeyname{language}|=|\oarg{dialect}\meta{language}
%
%      激活一种编程语言（的某个方言）。``空''的默认
%      语言不检测关键词、注释、字符串等；它对
%      排版纯文本可能有用。
%      如果未指定 \meta{dialect}，宏包会选择默认
%      方言；若没有默认方言，则选择空方言。
%
%      第~\pageref{uPredefinedLanguages} 页的表~\ref{uPredefinedLanguages}
%      列出了 \texttt{lstdrvrs.dtx} 提供的所有语言和方言。
%      预定义的默认方言带下划线。
%
% \item[0.21] \rkeyname{alsolanguage}|=|\oarg{dialect}\meta{language}
%
%      在当前活动语言之外，再激活一种编程语言（的某个方言）。
%      注意有些语言定义会相互干扰，彼此完全不兼容；
%      例如一种区分大小写，另一种不区分。
%
%      如果你想以不同方式高亮各语言的关键词，
%      可以看看第~\ref{rFigureOutTheAppearance} 节中的
%      \ikeyname{classoffset} 键。
%
% \item[0.19] \rkeyname{defaultdialect}|=|\oarg{dialect}\meta{language}
%
%      把 \meta{dialect} 定义为 \meta{language} 的默认方言。
%      如果你定义了非空的默认方言，例如
%      |defaultdialect=[iama]fool|，就不能再选择空方言，
%      即使写 |language=[]fool| 也不行。
% \end{syntax}
%
% 最后是一小列与语言相关的键。
% \begin{syntax}
%   \item[0.19,false,optional]
%     \rkeyname{printpod}|=|\meta{\alternative{true,false}}
%
%     打印或丢弃 Perl 中的 POD。
%
%   \item[0.20,true,{renamed,optional}]
%     \rkeyname{usekeywordsintag}|=|\meta{\alternative{true,false}}^^A
%     \label{uoption:usekeywordsintag}
%
%     宏包要么使用标签中的一阶关键词，要么把 |<>| 内的
%     所有标识符按关键词样式打印。
%
%   \item[1.1,{{}},optional]
%     \rkeyname{tagstyle}|=|\meta{style}\label{uoption:tagstyle}
%
%     决定标签及其内容的打印样式。
%
%   \item[1.1,false,optional]
%     \rkeyname{markfirstintag}|=|\meta{style}\label{uoption:markfirstintag}
%
%     按关键词样式打印标签中的第一个名字。
%
%   \item[0.20,true,optional]
%     \rkeyname{makemacrouse}|=|\meta{\alternative{true,false}}
%
%     宏调用特例：对定义为一阶关键词的标识符作宏调用时，
%     还会把其两侧的
%     |$(| 和 |)| 按关键词样式打印，例如可以得到
%     \textbf{\textdollar(}\textbf{strip} \textdollar(BIBS)\textbf{)}。若
%     取消该设置，则得到 \textdollar(\textbf{strip} \textdollar(BIBS))。
% \end{syntax}
%
%
% \subsubsection{弄清外观}\label{rFigureOutTheAppearance}
%
% \begin{syntax}
%   \item[0.18,{{}}] \rkeyname{basicstyle}|=|\meta{basic style}
%
%     在每个列表开头被选用。你可以用
%     |\footnotesize|、|\small|、|\itshape|、|\ttfamily| 之类
%     的设置。\meta{basic style} 的最后一个记号不能读取任何
%     后续字符。
%
%   \item[0.18,{{}}] \rkeyname{identifierstyle}|=|\meta{style}
%   \item[0.11,\itshape] \rkeyname{commentstyle}|=|\meta{style}
%   \item[0.12,{{}}] \rkeyname{stringstyle}|=|\meta{style}
%
%     决定非关键词、注释和字符串的样式。
%     \emph{最后}一个记号可以是带一个参数的命令，例如 |\textbf| 或
%     |\underbar|。
%
%   \item[0.11,\bfseries,addon]
%     \rkeyname{keywordstyle}|=|\oarg{number}[\textasteriskcentered]^^A
%     \meta{style}\label{roption:keywordstyle}
%
%     用于打印关键词。可选的 \meta{number} 参数
%     是指该样式应作用于的类别号。
%
%     附加功能：如果在（可选的）类别号之后使用可选星号，
%     关键词会以大写打印——即使该语言
%     区分大小写、只定义了关键词的小写形式。也许
%     还应该为小写关键词提供一个选项……
%
%   \item[0.19,keywordstyle,deprecated]
%     \rkeyname{ndkeywordstyle}|=|\meta{style}
%
%     等价于 |keywordstyle=2|\meta{style}。
%
%   \item[1.0,0] \rkeyname{classoffset}|=|\meta{number}
%
%     在分配样式、关键词、标识符等之前，会加到所有类别号上。
%     下面的例子直接定义关键词；你也可以通过选择两种
%     不同的语言来间接实现。
% \end{syntax}
%\ifcolor
% \begin{lstxsample}
%    \lstset{classoffset=0,
%            morekeywords={one,three,five},keywordstyle=\color{red},
%            classoffset=1,
%            morekeywords={two,four,six},keywordstyle=\color{blue},
%            classoffset=0}% restore default
% \end{lstxsample}
%\else
% \begin{lstxsample}
%    \lstset{classoffset=0,
%            morekeywords={one,three,five},keywordstyle=\itshape,
%            classoffset=1,
%            morekeywords={two,four,six},keywordstyle=\bfseries},
%            classoffset=0}% restore default
% \end{lstxsample}
%\fi
% \begin{lstsample}{}{}
%    \begin{lstlisting}
%    one two three
%    four five six
%    \end{lstlisting}
% \end{lstsample}
%
% \begin{syntax}
%   \item[0.20,keywordstyle,{addon,bug,optional}]
%     \rkeyname{texcsstyle}|=|[|*|]\oarg{class
%     number}\meta{style}\label{roption:texcsstyle}
%   \item[0.20,keywordstyle,optional]
%     \rkeyname{directivestyle}|=|\meta{style}
%
%     决定 \TeX\ 控制序列和指令的样式。请注意
%     这些键只有在选择了相应
%     语言时才存在。
%
%     |texcsstyle| 的可选星号还会高亮控制序列名前面的
%     反斜杠。注意该选项对所有
%     |texcs| 列表都生效。
%
%     缺陷：\texttt{texcs\ldots} 会干扰其他关键词列表。例如，
%     如果 \texttt{emph} 中含有 \texttt{foo} 一词，那么
%     控制序列 |\foo| 会以 \texttt{emphstyle} 样式显示。
%
%   \item[0.21] \rkeyname{emph}|=|\oarg{number}\marg{identifier list}
%   \item[0.21] \rkeyname{moreemph}|=|\oarg{number}\marg{identifier list}
%   \item[0.21] \rkeyname{deleteemph}|=|\oarg{number}\marg{identifier list}
%   \item[0.21] \rkeyname{emphstyle}|=|\oarg{number}\marg{style}
%
%     分别定义、添加或删除``强调类 \meta{number}''中的
%     \meta{identifier list}，或者定义该类的样式。
%     如果没有给出可选参数，宏包假定
%     \meta{number}$\,=1$。
%
%     这些键在第~\ref{uEmphasizeIdentifiers} 节
%     有更详细的说明。
%
%   \item[1.0]
%     \rkeyname{delim}|=|[\texttt*[\texttt*]]\texttt[\meta{type}\texttt]^^A
%     [\texttt[\meta{style}\texttt]]\meta{delimiter\textup(s\textup)}
%   \item[1.0]
%     \rkeyname{moredelim}|=|[\texttt*[\texttt*]]\texttt[\meta{type}\texttt]^^A
%     [\texttt[\meta{style}\texttt]]\meta{delimiter\textup(s\textup)}
%   \item[1.0]
%     \rkeyname{deletedelim}|=|[\texttt*[\texttt*]]\texttt[\meta{type}\texttt]^^A
%     \meta{delimiter\textup(s\textup)}
%
%     定义、添加或删除用户提供的定界符。（注意这
%     不会影响字符串或注释。）
%
%     前两种情况下，\meta{style} 用来打印被定界的代码（以及定界符）。
%     这里的 \meta{style} 可以是 |\bfseries| 或 |\itshape| 之类，
%     也可以通过 \texttt{keywordstyle}、\texttt{keywordstyle2}、
%     \texttt{emphstyle}
%     等引用其他样式。
%
%     支持的类型是 \texttt{l} 和 \texttt{s}，解释见第~\ref{uLanguageDefinitions}
%     节中的注释键。如果使用前缀 \texttt i，即~\texttt{il} 或 \texttt{is}，
%     则定界符不会被打印，相当于一种
%     不可见效果。
%
%     如果使用一个可选星号，宏包会检测被定界代码中的关键词、
%     注释和字符串。如果使用两个可选
%     星号，样式还会叠加应用；参见第
%     \ref{uDelimiters} 节。
% \end{syntax}
%
%
% \subsubsection{正确处理所有字符}
%
% \begin{syntax}
%   \item[0.18,true]
%     \rkeyname{extendedchars}|=|\meta{\alternative{true,false}}^^A
%     \syntaxor\rkeyname{extendedchars}
%
%     允许或禁止列表中的扩展字符，即编码 128--255 的
%     （各国的）字符。如果使用扩展
%     字符，就应当加载 \packagename{fontenc} 和/或
%     \packagename{inputenc} 等宏包。
%
%   \item[1.0,{{}}] \rkeyname{inputencoding}|=|\meta{encoding}
%
%     决定输入编码。使用该键需要
%     \packagename{inputenc} 宏包；如果没有加载，什么也不会发生。
%
%   \item[1.1,false]
%     \rkeyname{upquote}|=|\meta{\alternative{true,false}}\label{uoption:upquote}
%
%     决定左右引号打印为 |`'|（false）还是
%     \texttt{\textasciigrave\textquotesingle}（true）。为 true 时该键需要
%     \packagename{textcomp} 宏包，更多信息见
%     第~\pageref{sec:upquote} 页。
%
%   \item[0.12,8] \rkeyname{tabsize}|=|\meta{number}
%
%     设置制表位位于第
%     $\meta{number}+1$、$2\cdot\meta{number}+1$、
%     $3\cdot\meta{number}+1$ 等列。列表中的每个制表符都会
%     把当前列移到下一个制表位。
%
%   \item[0.20,false] \rkeyname{showtabs}|=|\meta{\alternative{true,false}}
%
%     使制表符可见或不可见。可见的制表符形如
%             \lstinline[showtabs]!\t!，
%     不过这一点可以修改。如果选择
%     不可见制表符而可见空格，制表符会被转换成相应数量的空格。
%
%   \item[0.20] \rkeyname{tab}|=|\meta{tokens}
%
%     \meta{tokens} 用于打印可见的制表符。你也许想用
%     |$\to$|、|$\mapsto$|、
%     |$\dashv$| 之类代替这个奇怪的默认
%     定义。
%
%   \item[0.20,false]
%     \rkeyname{showspaces}|=|\meta{\alternative{true,false}}
%
%     让所有空白字符显示为 {\textvisiblespace} 或空格。
%
%   \item[0.12,true]
%     \rkeyname{showstringspaces}|=|\meta{\alternative{true,false}}
%
%     让字符串中的空白字符显示为 {\textvisiblespace} 或
%     空格。
%
%   \item[0.19,\bigbreak] \rkeyname{formfeed}|=|\meta{tokens}
%
%     只要列表中出现换页符，就执行 \meta{tokens}。
% \end{syntax}
%
%
% \subsubsection{行号}\label{rLineNumbers}
%
% \begin{syntax}
%   \item[1.0,none]
%     \rkeyname{numbers}|=|\meta{\alternative{none,left,right}}
%
%     使宏包不打印行号，或者把行号放在列表的
%     左侧或右侧。
%
%   \item[0.16,1] \rkeyname{stepnumber}|=|\meta{number}
%
%     所有``行号 $\equiv
%     0$（模 \meta{number}）''的行都会得到行号。如果用 \keyname{numbers} 开关
%     行号，参数 \keyname{stepnumber} 会保持其值。
%     另一种做法是用 |stepnumber=0| 关闭、用非零数字打开，
%     同时保持
%     \keyname{numbers} 的值。
%
%   \item[1.1,false]
%     \rkeyname{numberfirstline}|=|\meta{\alternative{true,false}}^^A
%     \label{uoption:numberfirstline}
%
%     每个列表的第一行都会得到编号（只要行号功能已打开），
%     即使该行号不能被
%     \keyname{stepnumber} 整除。
%
%   \item[0.16,{{}}] \rkeyname{numberstyle}|=|\meta{style}
%
%     决定行号的字体和字号。
%
%   \item[0.19,10pt] \rkeyname{numbersep}|=|\meta{dimension}
%
%     是行号与列表之间的距离。
%
%   \item[1.0,true]
%     \rkeyname{numberblanklines}|=|\meta{\alternative{true,false}}
%
%     若设为 false，空行不会打印行号。
%
%   \item[0.20,auto]
%     \rkeyname{firstnumber}|=|\meta{\alternative{auto,last,\normalfont\meta{number}}}
%
%     \begin{description}
%       \item[\texttt{auto}] 让宏包自行选择起始编号：新列表
%         从 1 开始；有名字的列表延续最近同名的列表（见下文）；
%         独立文件则从与第一输入行
%         对应的编号开始。
%       \item[\texttt{last}] 延续最近一个列表的编号，
%       \item[\meta{number}] 把它设为给定的（逻辑）编号。
%     \end{description}
%
%   \item[1.0,{{}}] \rkeyname{name}|=|\meta{name}
%
%     只能用于单个列表，用于给列表命名。
%     如果选项 |firstnumber=auto| 生效，同名的显示型环境列表
%     共享一个行计数器。
%
%   \item[0.20,\arabic{lstnumber},data] \rcmdname\thelstnumber
%
%     打印各行的行号。
% \end{syntax}
% 我们展示一个如何重定义 |\thelstnumber| 的例子。但如果你测试它，
% 你不会得到左边显示的结果，因为内部已经设置了 |stepnumber=-1|。
% \begin{lstxsample}
%    \renewcommand*\thelstnumber{\oldstylenums{\the\value{lstnumber}}}
% \end{lstxsample}
% \begin{lstsample}{\lstset{stepnumber=-1}\label{rDecreasingLabels}}{}
%    \begin{lstlisting}[numbers=left,
%                       firstnumber=753]
%    begin { empty lines }
%
%
%
%
%
%
%    end; { empty lines }
%    \end{lstlisting}
% \end{lstsample}
%
% \begin{advise}
%   \item 该例展示一个由 8 个三位数组成的序列
%     $n,n+1,\ldots,n+7$，该序列包含每一个数字
%     $0,1,\ldots,9$。但 8 并非具有这一性质的最小值。请找出
%     最小的数并证明它是最小的。这样的最小
%     序列共有多少个？
%
%     现在来看一般化的问题：已知
%     $k\in\{1,\ldots,10\}$。请找出最小的
%     数 $m\in\{1,\ldots,10\}$，使得存在由 $m$ 个
%     $n,{n+1},\ldots,\allowbreak{n+m-1}$ 组成的
%     $k$ 位数序列，包含每一个数字
%     $\{0,\ldots,9\}$。证明这个数是最小的。这样的
%     最小序列共有多少个？
%
%     如果你用计算机解决这个问题，那就写一个 \TeX\ 程序吧！
% \end{advise}
%
%
% \subsubsection{标题}
%
% 与 \LaTeX 标准行为不同，这里标题和浮动是相互独立的；你可以在
% 非浮动列表中使用标题。
% \begin{syntax}
% \item[0.21] \rkeyname{title}|=|\meta{title text}
%
%       用于没有任何编号或标签的标题。
%
% \item[0.20] \rkeyname{caption}|={|\oarg{short}\meta{caption text}|}|
%
%       标题由 \cs{lstlistingname} 加一个流水号、一个分隔符
%       和 \meta{caption text} 组成。列表目录中要么使用标题文字，
%       要么使用 \meta{short}（如果给出了）。
%
% \item[0.21] \rkeyname{label}|=|\meta{name}
%
%       使列表可以通过 |\ref|\marg{name} 引用。
%
% \item[0.16] \rcmdname\lstlistoflistings
%
%       打印列表目录。每个条目按优先级递减的顺序取短标题、标题、
%       文件名或列表名，另见第
%       \ref{rLineNumbers} 节中的 \keyname{name} 键。从本宏包 v1.10
%       版本起，列表目录可以借助 \packagename{tocbasic} 宏包
%       \cite[Chapter 15]{Kohm:scrguide} 生成，更多信息见
%       第~\pageref{sec:tocbasic} 页。
%
% \item[1.0] \rkeyname{nolol}|=|\meta{\alternative{true,false}}\syntaxor\rkeyname{nolol}
%
%       若为 true，该列表不会进入列表目录。
%
% \item[0.16,Listings,data] \rcmdname\lstlistlistingname
%
%       列表目录的标题名。
%
% \item[0.20,Listing,data] \rcmdname\lstlistingname
%
%       列表的标题标签名。
%
% \item[,{{}},data]        \rcmdname\lstlistingnamestyle
%
%       以简单的方式定制程序列表标题标签的样式，例如
%       |\small|、|\bfseries|，或者几个命令的
%       组合。如果需要对标签作完全定制
%       （对齐、字体、页边距等），
%       应使用 A.~Sommerfeldt 的 \packagename{caption} 宏包
%       \cite{Sommerfeldt:caption}。
%
% \item[0.20,\arabic{lstlisting},data] \rcmdname\thelstlisting
%
%       打印标题的流水号。
%
% \item[1.4,true] \rkeyname{numberbychapter}|=|\meta{\alternative{true,false}}
%
%       若为 true 且 |\thechapter| 存在，列表按章编号。
%       否则从文档开头起顺序编号。
%       该键只能在 |\begin{document}| 之前使用。
%
% \item[0.19] \rcmdname\lstname
%
%       打印当前列表的名字，即文件名或
%       由 \keyname{name} 键定义的名字。该命令可用于定义标题
%       或题目模板，例如用
%       |\lstset{caption=\lstname}|。
%
% \item[0.20,t] \rkeyname{captionpos}|=|\meta{subset of \textup{\texttt{tb}}}
%
%       指定标题的位置：列表的顶部和/或底部。
%
% \item[0.20,\smallskipamount] \rkeyname{abovecaptionskip}|=|\meta{dimension}
% \item[0.20,\smallskipamount] \rkeyname{belowcaptionskip}|=|\meta{dimension}
%
%             分别是每个题注上方或下方的垂直间距。
% \end{syntax}
%
%
% \subsubsection{页边距与行形状}\label{rMarginsAndLineShape}
%
% \begin{syntax}
%   \item[0.21,\linewidth] \rkeyname{linewidth}|=|\meta{dimension}
%
%     定义列表的基本行宽。另外还会计入以下三个键。
%
%   \item[0.19,0pt] \rkeyname{xleftmargin}|=|\meta{dimension}
%   \item[1.0,0pt] \rkeyname{xrightmargin}|=|\meta{dimension}
%
%     这些尺寸用作左右两侧的额外页边距。行号
%     和边框也随之移动。
%
%   \item[0.19,false]
%     \rkeyname{resetmargins}|=|\meta{\alternative{true,false}}
%
%     若为 true，来自 \texttt{enumerate} 或
%     \texttt{itemize} 等列表环境的缩进会被重置，即不再使用。
%
%   \item[0.20,false]
%     \rkeyname{breaklines}|=|\meta{\alternative{true,false}}\syntaxor\rkeyname{breaklines}
%
%     激活或关闭长行的自动换行。
%
%   \item[1.2,false]
%     \rkeyname{breakatwhitespace}|=|\meta{\alternative{true,false}}^^A
%     \syntaxor\rkeyname{breakatwhitespace}\label{uoption:breakatwhitespace}
%
%     若为 true，只允许在空白处换行。
%
%   \item[0.20,{{}}] \rkeyname{prebreak}|=|\meta{tokens}
%   \item[0.20,{{}}] \rkeyname{postbreak}|=|\meta{tokens}
%
%     \meta{tokens} 分别出现在当前行末尾，以及下一条
%     （断开的）行的开头。
%
%     不能使用动态间距（尤其是空格），因为
%     内部使用的是 |\discretionary|。不过 |\space| 已被重定义，
%     可在 \meta{tokens} 内使用。
%
%   \item[0.20,20pt] \rkeyname{breakindent}|=|\meta{dimension}
%
%     是断行后第二、第三……行的缩进。
%
%   \item[0.20,true]
%     \rkeyname{breakautoindent}|=|\meta{\alternative{true,false}}^^A
%     \syntaxor\rkeyname{breakautoindent}
%
%     激活或关闭断开行的自动缩进。该缩进是在
%     \ikeyname{breakindent} 之外额外使用的，见下面的
%     例子。可见空格或可见制表符可能把这个
%     自动缩进设为零。
% \end{syntax}
% 在下面的例子中，我们使用制表符创建长行，但 verbatim 部分使用 |tabsize=1|。
% \begin{lstxsample}
%    \lstset{postbreak=\space, breakindent=5pt, breaklines}
% \end{lstxsample}
% \begin{lstsample}{\lstset{string=[d]",tabsize=6}}{\lstset{tabsize=1}\hfuzz=1in}
%    \begin{lstlisting}
%    		"A long string is broken!"
%    			"Another long line."
%    \end{lstlisting}
%
%    \begin{lstlisting}[breakautoindent
%                                 =false]
%    		{ Now auto indention is off. }
%    \end{lstlisting}
% \end{lstsample}
%
%
% \subsubsection{边框}\label{rFrames}
%
% \begin{syntax}
%   \item[1.0,none]
%     \rkeyname{frame}|=|\meta{\alternative{none,leftline,topline,bottomline,^^A
%       lines,single,shadowbox}}
%
%     画出无边框、左侧单线、顶部单线、底部单线、
%     顶部加底部、完整单线框或阴影框。
%
%     注意 \packagename{fancyvrb} 支持相同的边框类型，除了
%     \texttt{shadowbox}。阴影颜色由 \keyname{rulesepcolor} 决定，见
%     下文。
%
%   \item[0.19,{{}}] \rkeyname{frame}|=|\meta{subset of
%     \textup{\texttt{trblTRBL}}}
%
%     字符 \texttt{trblTRBL} 分别指定列表顶部、底部
%     以及右侧、左侧的线。大写字符用于画双线。因此 |frame=tlrb| 画
%     单线框，|frame=TL| 在顶部和左侧画双线。
%
%     注意边框通常位于列表空间之外。
%
%   \item[0.20,ffff]
%     \rkeyname{frameround}|=|\meta{\alternative{t,f}}\meta{\alternative{t,f}}^^A
%     \meta{\alternative{t,f}}\meta{\alternative{t,f}}
%
%     四个字母依次指右上、右下、左下
%     和左上角。\texttt{t} 使相应的
%     角变圆。如果使用圆角，线宽由 |\thinlines| 和 |\thicklines| 控制。
%
%     注意：四分之一圆的大小取决于 \keyname{framesep}，
%     与边框的额外页边距无关。其大小可能会被
%     调整以贴合 \LaTeX 的圆尺寸。
%
%   \item[0.19,3pt] \rkeyname{framesep}|=|\meta{dimension}
%   \item[0.19,2pt] \rkeyname{rulesep}|=|\meta{dimension}
%
%     控制边框与列表之间，以及两条双线之间的间距。
%
%   \item[0.19,0.4pt] \rkeyname{framerule}|=|\meta{dimension}
%
%     控制线的宽度。
%
%   \item[1.0,0pt] \rkeyname{framexleftmargin}|=|\meta{dimension}
%   \item[1.0,0pt] \rkeyname{framexrightmargin}|=|\meta{dimension}
%   \item[1.0,0pt] \rkeyname{framextopmargin}|=|\meta{dimension}
%   \item[1.0,0pt] \rkeyname{framexbottommargin}|=|\meta{dimension}
%
%     是在 \keyname{framesep} 之外用于构成边框
%     页边距的尺寸。
%
%   \item[0.21] \rkeyname{backgroundcolor}|=|\meta{color command}
%   \item[0.21] \rkeyname{rulecolor}|=|\meta{color command}
%   \item[1.0] \rkeyname{fillcolor}|=|\meta{color command}
%   \item[1.0] \rkeyname{rulesepcolor}|=|\meta{color command}
%
%     分别指定背景、线条、``文本框''与第一条线之间，
%     以及两条线之间空白处的颜色。注意该值需要一个 |\color| 命令，例如
%     \keyname{rulecolor}|=\color{blue}|。
% \end{syntax}
% \ikeyname{frame} 在 |fancyvrb=true| 时不起作用，或者当宏包
% 内部在列表外套一个 |\hbox| 时也不起作用！其他命令肯定还有更多
% 问题；请花点时间提交一个（缺陷）报告。
%\ifcolor
% \begin{lstxsample}
%    \lstset{framexleftmargin=5mm, frame=shadowbox, rulesepcolor=\color{blue}}
% \end{lstxsample}
%\else
%    \lstset{framexleftmargin=5mm, frame=shadowbox}
%\fi
% \begin{lstsample}{}{}
%    \begin{lstlisting}[numbers=left]
%    for i:=maxint to 0 do
%    begin
%        { do nothing }
%    end;
%    \end{lstlisting}
% \end{lstsample}
%
% 注意这里使用 |framexleftmargin| 将行号包含在边框内部。
%
% 你想要奇特的边框吗？试试下面的键，例如，
% \begin{lstsample}{\lstset{frameshape={RYRYNYYYY}{yny}{yny}{RYRYNYYYY}}}{}
%    \begin{lstlisting}
%    for i:=maxint to 0 do
%    begin
%        { do nothing }
%    end;
%    \end{lstlisting}
% \end{lstsample}
% \begin{syntax}
%   \item[0.20,,\dag] \rkeyname{frameshape}|=|\marg{top shape}\marg{left
%     shape}\marg{right shape}\marg{bottom shape}
%
%     让你完全控制绘制的边框部分。参数不区分大小写。
%
%     \meta{left shape} 和 \meta{right shape} 都是``从左到右''的
%     \alternative{y,n} 字符序列（或空）。每个 |y| 让宏包
%     画一条线，否则该处留空。这些竖线按指定的形状
%     ``从左到右''绘制。
%     上面的例子使用了 |yny|。
%
%     \meta{top shape} 和 \meta{bottom shape} 是``左—线—右''
%     序列（或空）。第一个``左—线—右''序列
%     对应最内层的线，第二个对应次外层，依此类推。
%     每个序列有三个字符：``线''为 |y| 或 |n|；
%     ``左''和``右''为 |y|、|n| 或 |r|（使该角变圆）。
%     示例对两个形状都使用 |RYRYNYYYY|：|RYR| 描述
%     最内层（顶部和底部）框形，|YNY| 描述中间层，|YYY|
%     描述最外层。
% \end{syntax}
% 总结一下，上面的例子使用了
% \begin{verbatim}
%    \lstset{frameshape={RYRYNYYYY}{yny}{yny}{RYRYNYYYY}}\end{verbatim}
% 注意你不限于两层或三层。但是如果圆角太大，你会遇到麻烦。
%
%
% \subsubsection{索引}
%
% \begin{syntax}
%   \item[0.19] \rkeyname{index}|=|\oarg{number}\oarg{keyword
%     classes}\marg{identifiers}
%   \item[0.21] \rkeyname{moreindex}|=|\oarg{number}\oarg{keyword
%     classes}\marg{identifiers}
%   \item[0.21] \rkeyname{deleteindex}|=|\oarg{number}\oarg{keyword
%     classes}\marg{identifiers}
%
%     从索引类别列表 \meta{number} 中定义、添加和删除
%     \meta{identifiers} 与 \meta{keyword classes}。如果未指定
%     可选编号，宏包假定 \meta{number} $=1$。
%
%     显式给出的标识符的每次出现，以及指定 \meta{keyword classes}
%     中标识符的每次
%     出现，都会被编入索引。例如，可以写 |index=[1][keywords]| 来
%     为所有关键词建立索引。注意这里的 |[1]| 是必需的——否则我们
%     无法使用第二个可选参数。
%
%   \item[0.19,\lstindexmacro]
%     \rkeyname{indexstyle}|=|\oarg{number}\meta{tokens
%     \textup(one-parameter command\textup)}
%
%     \meta{tokens} 实际上负责为列表 \meta{number} 的标识符
%     建立索引。与样式键不同，\meta{tokens}
%     \emph{必须}读取恰好一个参数，即标识符。
%     默认定义为\icmdname{\lstindexmacro}\vspace*{-\itemsep}
% \begin{verbatim}
%    \newcommand\lstindexmacro[1]{\index{{\ttfamily#1}}}\end{verbatim}
%       \vspace*{-\itemsep}你不应该修改它。
%       请定义你自己的索引命令，并把它们作为该键的参数。
%   \end{syntax}
% 第~\ref{uIndexing} 节详细描述了这一功能。
%
%
% \subsubsection{列对齐}\label{rColumnAlignment}
%
% \begin{syntax}
%   \item[1.0,{[c]fixed}]
%     \rkeyname{columns}|=|\oarg{\alternative{c,l,r}}\meta{alignment}
%
%     选择列对齐方式。\meta{alignment} 可以是 |fixed|、
%     |flexible|、|spaceflexible| 或 |fullflexible|；详见第
%     \ref{uFixedAndFlexibleColumns} 节。
%
%     可选的 |c|、|l| 或 |r| 控制最小输出单元
%     （关键词、标识符等）的水平朝向。各参数
%     的作用如下，竖线用来说明效果：
%     $\vert$\lstinline[columns={[c]fixed}]!listing!$\vert$，
%     $\vert$\lstinline[columns={[l]fixed}]!listing!$\vert$，以及
%     $\vert$\lstinline[columns={[r]fixed}]!listing!$\vert$ 为固定列
%     模式，
%     $\vert$\lstinline[columns={[c]flexible}]!listing!$\vert$，
%     $\vert$\lstinline[columns={[l]flexible}]!listing!$\vert$，以及
%     $\vert$\lstinline[columns={[r]flexible}]!listing!$\vert$ 为
%     灵活列，以及
%     $\vert$\lstinline[columns={[c]fullflexible}]!listing!$\vert$，
%     $\vert$\lstinline[columns={[l]fullflexible}]!listing!$\vert$，以及
%     $\vert$\lstinline[columns={[r]fullflexible}]!listing!$\vert$ 为
%     完全灵活列（它们忽略可选
%     参数，因为不会在可打印字符周围
%     添加额外间距）。
%
%   \item[0.18,false]
%     \rkeyname{flexiblecolumns}|=|\meta{\alternative{true,false}}^^A
%     \syntaxor\rkeyname{flexiblecolumns}
%
%     选择最近一次选定的灵活列或固定列格式，
%     参见第~\ref{uFixedAndFlexibleColumns} 节。
%
%   \item[0.21,false,\dag]
%     \rkeyname{keepspaces}|=|\meta{\alternative{true,false}}
%
%     |keepspaces=true| 告诉宏包不要为了修正列
%     对齐而丢弃空格，并总是把制表符转换成空格。
%
%   \item[0.16] \rkeyname{basewidth}|=|\meta{dimension}\syntaxor
%   \item[0.18,{{0.6em,0.45em}}]
%     \rkeyname{basewidth}|={|\meta{fixed}|,|\meta{flexible mode}|}|
%
%     设置固定列模式和灵活列模式下单字符盒的宽度
%     （可以都设为同一个值，也可以分别设置）。
%
%   \item[0.20,false]
%     \rkeyname{fontadjust}|=|\meta{\alternative{true,false}}^^A
%     \syntaxor\rkeyname{fontadjust}
%
%     若为 true，宏包在每次选择字体时都会调整基准宽度。
%     这只有在 \ikeyname{basewidth} 以字体
%     相关的单位给出（如 `em' 或 `ex'）时才有意义——否则
%     该布尔值没有作用。
%
%     宏包加载后不会在每次选择字体时调整宽度：
%     它在每个列表处查看 \ikeyname{basewidth}，并对
%     整个列表使用该值。如果第~\ref{rFigureOutTheAppearance} 节中的样式
%     键导致大幅度的字号
%     变化，这可能不合适，见下面的例子。
%
%     注意该键可能扰乱列对齐，并可能
%     影响关键词的外观！
% \end{syntax}
% \begin{lstsample}{\lstset{basicstyle=\normalsize}}{}
%    \lstset{commentstyle=\scriptsize}
%    \begin{lstlisting}
%    { scriptsize font
%      doesn't look good }
%    for i:=maxint to 0 do
%    begin
%        { do nothing }
%    end;
%    \end{lstlisting}
% \end{lstsample}
% \begin{lstsample}{\lstset{basicstyle=\normalsize,commentstyle=\scriptsize}}{}
%    \begin{lstlisting}[fontadjust]
%    { scriptsize font
%      looks better now }
%    for i:=maxint to 0 do
%    begin
%        { do nothing }
%    end;
%    \end{lstlisting}
% \end{lstsample}
%
%
% \subsubsection{转义到 \LaTeX}\label{rEscapingToLaTeX}
%
% \textbf{注意：} {\itshape 任何到 \LaTeX\ 的转义都可能扰乱列
% 对齐，因为宏包无法控制那里的间距。}
% \begin{syntax}
% \item[0.18,false] \rkeyname{texcl}|=|\meta{\alternative{true,false}}\syntaxor\rkeyname{texcl}
%
%       激活或关闭 \LaTeX\ 注释行。若激活，注释
%       行的定界符照常打印，但注释行的文字（直到
%       行尾）会被当作 \LaTeX\ 代码读取，并以注释样式排版。
% \end{syntax}
% 该例子使用 \Cpp 注释行（但没有说明如何定义它们）。如果没有 |\upshape|，
% 我们会得到 \textit{calculate}，因为注释样式是 |\itshape|。
% \begin{lstsample}{\lstset{morecomment=[l]//}}{}
%    \begin{lstlisting}[texcl]
%    // \upshape calculate $a_{ij}$
%      A[i][j] = A[j][j]/A[i][j];
%    \end{lstlisting}
% \end{lstsample}
%
% \begin{syntax}
%   \item[0.19,false]
%     \rkeyname{mathescape}|=|\meta{\alternative{true,false}}
%
%     激活或关闭美元符号的特殊行为。若
%     激活，美元符号充当 \TeX 的文本数学切换符。
%
%     如果你想在列表中排版公式，这个键很有用。
%
%   \item[0.19,{{}}]
%     \rkeyname{escapechar}|=|\meta{character}\syntaxor\rkeyname{escapechar}|={}|
%
%     若不为空，给定字符会把用户转义到 \LaTeX：两个这样的
%     字符之间的所有代码都被解释为 \LaTeX\ 代码。注意
%     \TeX 的特殊字符必须带前置反斜杠输入，
%     例如 |escapechar=\%|。
%
%   \item[0.20,{{}}]
%     \rkeyname{escapeinside}|=|\meta{character}\meta{character}^^A
%     \syntaxor\rkeyname{escapeinside}|={}|
%
%     是 \ikeyname{escapechar} 的一般化。若值
%     不为空，宏包会在第一个和第二个
%     字符之间转义到 \LaTeX。
%
%   \item[0.20,{{}}] \rkeyname{escapebegin}|=|\meta{tokens}
%   \item[0.20,{{}}] \rkeyname{escapeend}|=|\meta{tokens}
%
%     这些记号分别在每次转义的开头和结尾执行，
%     尤其是对 \ikeyname{texcl}。应用见第
%     \ref{uNationalCharacters} 节。
% \end{syntax}
%
% \begin{lstsample}{\lstset{morecomment=[l]//}}{}
%    \begin{lstlisting}[mathescape]
%    // calculate $a_{ij}$
%      $a_{ij} = a_{jj}/a_{ij}$;
%    \end{lstlisting}
% \end{lstsample}
%
% \begin{lstsample}{\lstset{morecomment=[l]//}}{}
%    \begin{lstlisting}[escapechar=\%]
%    // calc%ulate $a_{ij}$%
%      %$a_{ij} = a_{jj}/a_{ij}$%;
%    \end{lstlisting}
% \end{lstsample}
%
% \begin{lstsample}{\lstset{morecomment=[l]//}}{}
%    \lstset{escapeinside=`'}
%    \begin{lstlisting}
%    // calc`ulate $a_{ij}$'
%      `$a_{ij} = a_{jj}/a_{ij}$';
%    \end{lstlisting}
% \end{lstsample}
% 第一个例子中，直到 $a_{ij}$ 的注释行由
% \packagename{listings} 宏包以注释样式排版。$a_{ij}$ 本身
% 则以``\TeX\ 数学模式''排版，不带注释样式。第二个例子的
% 注释行约有一半由本宏包排版，
% 其余部分处于``\LaTeX\ 模式''。
%
% 为避免与本宏包当前和未来版本的问题：
% \begin{enumerate}
% \item 当你已转义到 \LaTeX 时，不要使用 \packagename{listings} 宏包的任何命令。
% \item 任何环境必须在同一个转义内开始和结束。
% \item 你可以使用 |\def|、|\edef| 等，但不要假定这些
%       定义稍后仍然存在，除非它们是 |\global| 的。
% \item |\if \else \fi|、分组、数学切换符 |$| 和 |$$| 等必须
%       在每个转义内保持平衡。
% \item \ldots
% \end{enumerate}
% 你可以自行扩充这个列表，并把新条目寄给我。
%
%
% \subsection{与 \packagename{fancyvrb} 的接口}
%
% \packagename{fancyvrb} 宏包——精美 verbatim——由 Timothy van Zandt 编写，
% 提供用于读取、写入和排版 verbatim 代码的宏。它有一些 \packagename{listings}
% 宏包没有的显著特性。（有些是可能的，但你必须找到愿意实现它们的人 |;-)|。）
% \begin{syntax}
%   \item[0.19] \rkeyname{fancyvrb}|=|\meta{\alternative{true,false}}
%
%     激活或关闭该接口。若激活，verbatim 代码由
%     \packagename{fancyvrb} 读取，但由 \packagename{listings} 排版，
%     即带高亮的关键词、字符串、注释等。
%     内部我们使用 |\FancyVerbFormatLine| 的一个非常特殊的
%     定义。
%
%     该接口适用于 |Verbatim|、|BVerbatim| 和 |LVerbatim|。
%     但你不应该使用 \packagename{fancyvrb} 的
%     \keyname{defineactive}。（据我所见这无关紧要，因为它
%     根本不起作用，但为了安全起见……）如果
%     \packagename{fancyvrb} 和 \packagename{listings} 提供类似
%     功能，应当使用 \packagename{fancyvrb} 的。
%
%   \item[1.1,{\overlay 1}]
%     \rkeyname{fvcmdparams}|=|\meta{command$_1$}\meta{number$_1$}^^A
%     \ldots\label{uoption:fvcmdparams}
%   \item[1.1]
%     \rkeyname{morefvcmdparams}|=|\meta{command$_1$}\meta{number$_1$}^^A
%     \ldots\label{uoption:morefvcmdparams}
%
%     如果使用 \packagename{fancyvrb} 的 \keyname{commandchars}，就必须
%     告诉 \packagename{listings} 宏包每个命令
%     接受几个参数。如果命令不接参数，则无需
%     处理。
%
%     键的第一个（第三、第五……）参数是命令，
%     第二个（第四、第六……）是该命令接受的参数个数。因此，
%     如果要在 \packagename{fancyvrb}-\packagename{listings} 接口中使用 |\textcolor{red}{keyword}|，
%     应写 |\lstset{morefvcmdparams=\textcolor 2}|。
% \end{syntax}
%
% \iffancyvrb
% \begin{lstsample}{}{}
%    \lstset{morecomment=[l]\ }% :-)
%    \fvset{commandchars=\\\{\}}
%
%    \begin{BVerbatim}
%    First verbatim line.
%    \fbox{Second} verbatim line.
%    \end{BVerbatim}
%
%    \par\vspace{72.27pt}
%
%    \lstset{fancyvrb}
%    \begin{BVerbatim}
%    First verbatim line.
%    \fbox{Second} verbatim line.
%    \end{BVerbatim}
%    \lstset{fancyvrb=false}
% \end{lstsample}
% 由 \packagename{listings} 宏包排版的行更宽，因为默认的 \ikeyname{basewidth}
% 不等于单个打字机字符的宽度。此外，注意第一个空格开始一个注释，
% 如例子开头所定义的。
% \else
% \begin{center}
%    你的平台上似乎没有 \packagename{fancyvrb}，因此该例子无法在此打印。
% \end{center}
% \fi
%
%
% \subsection{环境}\label{rEnvironments}
%
% 如果你想定义自己的美化排版环境，试试下面的命令。语法来自
% \LaTeX 的 |\newenvironment|。
% \begin{syntax}
% \item[0.19] \rcmdname\lstnewenvironment\\
%       \marg{name}\oarg{number}\oarg{opt.~default~arg.}\\
%       |{|\meta{starting code}|}|\\
%       |{|\meta{ending code}|}|
% \end{syntax}
% 作为一个简单例子，我们可以只选择一种特定语言。
% \begin{lstxsample}
%    \lstnewenvironment{pascal}
%        {\lstset{language=pascal}}
%        {}
% \end{lstxsample}
% \begin{lstsample}{}{}
%    \begin{pascal}
%    for i:=maxint to 0 do
%    begin
%        { do nothing }
%    end;
%    \end{pascal}
% \end{lstsample}
% 做其他事情也同样容易，例如，使用更多键并添加一个可选参数来调整
% 每个列表的设置：
% \begin{verbatim}
%\lstnewenvironment{pascalx}[1][]
%    {\lstset{language=pascal,numbers=left,numberstyle=\tiny,float,#1}}
%    {}\end{verbatim}
%
%
% \subsection{短行内列表命令}\label{rShortInline}
%
% 也可以定义 |\lstinline| 的短等效形式，做法类似于
% \packagename{shortvrb} 提供的短 verbatim 宏。
%
% \begin{syntax}
% \item[1.4] \rcmdname\lstMakeShortInline[\oarg{options}]\meta{character}
%
%       把 \meta{character} 定义为
%       |\lstinline|[\oarg{options}]\meta{character} 的等效形式，
%       在大量使用行内列表时提供方便的语法。
%
% \item[1.4] \rcmdname\lstDeleteShortInline\meta{character}
%
%       删除由 |\lstMakeShortInline| 创建的 \meta{character} 定义，
%       使 \meta{character} 恢复其原有含义。
% \end{syntax}
%
%
% \subsection{语言定义}\label{rLanguageDefinitions}
%
% 你应当先阅读第~\ref{uLanguageDefinitions} 节，了解
% 语言定义。否则你可能对 |\lstdefinelanguage| 的完整
% 语法准备不足。
% \begin{syntax}
%   \item[0.19] \rcmdname\lstdefinelanguage\syntaxnewline[\oarg{dialect}]^^A
%     \marg{language}\syntaxnewline[\oarg{base dialect}^^A
%     \marg{and base language}]\syntaxnewline\marg{key=value list}^^A
%     \syntaxnewline[\oarg{list of required aspects ^^A
%     \textup(keywordcomments,texcs,etc.\textup)}]
%
%		定义（给定方言的）编程语言 \meta{language}。
%       如果该语言定义基于另一个定义，就必须
%       给出完整的 \oarg{base dialect}\marg{and base language}。注意
%       空的 \meta{base dialect} 会使用默认方言！
%
%       最后一个可选参数应当列出所有必需的方面。这一点比较
%       微妙，因为这些方面在开发者指南中描述。
%       你可以用现有的语言作为模板。例如，
%       ANSI C 使用 \aspectname{keywords}、\aspectname{comments}、
%       \aspectname{strings} 和 \aspectname{directives}。
%
%       \icmdname{\lst@definelanguage} 语法相同，用于在
%       驱动文件中定义语言。
%
% \begin{advise}
% \item 我应该把语言定义放在哪里？
%       \advisespace
%       如果你只需要在某个特定文档中使用该语言，把它放在该文档的
%       否则创建本地文件 \texttt{lstlang0.sty}，或者把定义添加到
%       该文件中，但使用 |\lst@definelanguage| 而不是 |\lstdefinelanguage|。
%       不过，你可能想把定义发送到第~\ref{uSoftwareLicense} 节中的地址。
%       这样它就会与宏包分发的其余语言一起包含进来，并在
%       \LaTeX 项目公共许可下发布。
% \end{advise}
%
% \item[0.18] \rcmdname\lstalias\marg{alias}\marg{language}
%
%             为某种编程语言定义别名。每个 \meta{alias} 被
%             重定向到 \meta{language} 的同一方言。
%             也可以只为某一种方言定义别名：
%
% \item[0.18] \rcmdname\lstalias\oarg{alias dialect}\marg{alias}\oarg{dialect}\marg{language}
%
%             这里四个参数全都 \emph{非可选}，而方言为空的别名
%             将选用默认方言。注意别名
%             不能串联：两个别名`|\lstalias{foo1}{foo2}|'和
%             `|\lstalias{foo2}{foo3}|'将 \emph{不会}把 |foo1| 重定向到 |foo3|。
% \end{syntax}
% 本节中其余所有键都用于构建语言
% 定义。\emph{在这样的定义中不应使用任何其他键！}
%
%
% \paragraph{关键词}
% 我们先从关键词构建键开始。注意：{\itshape 如果想把
% {\upshape|\|、|{|、|}|、|%|、|#|} 或 {\upshape|&|} 作为（部分）参数
% 传给下面的键，必须带上前置反斜杠！}
% \begin{syntax}
%   \item[1.0,,{\dag bug}] \rkeyname{keywordsprefix}|=|\meta{prefix}
%
%     所有以 \meta{prefix} 开头的标识符都会按一阶
%     关键词打印。
%
%     缺陷：目前有几项限制。(1) 前缀总是区分
%     大小写。(2) 一次只能定义一个前缀。
%     (3) 如果在语言定义之外``单独''使用，该键可能
%     只有在选择非空语言之后才起作用（必要时再切回
%     空语言）。(4) 该键不理会
%     \keyname{classoffset} 的值，也没有可选的类别
%     \meta{number} 参数。
%
%   \item[0.11] \rkeyname{keywords}|=|\oarg{number}\marg{list of keywords}
%   \item[0.11] \rkeyname{morekeywords}|=|\oarg{number}\marg{list of
%     keywords}
%   \item[0.18] \rkeyname{deletekeywords}|=|\oarg{number}\marg{list of
%     keywords}
%
%     定义、添加或删除关键词列表 \meta{number} 中的关键词。
%     不鼓励使用 \keyname{keywords}，因为它
%     会删除该列表中所有先前定义的关键词，因此与
%     \keyname{alsolanguage} 键不兼容。
%
%     如果你在关键词中使用了特殊字符，请注意
%     下面的 \ikeyname{alsoletter} 和 \ikeyname{alsodigit} 键。
%
%   \item[0.19,,deprecated] \rkeyname{ndkeywords}|=|\marg{list of keywords}
%   \item[0.19,,deprecated] \rkeyname{morendkeywords}|=|\marg{list of
%     keywords}
%   \item[0.19,,deprecated] \rkeyname{deletendkeywords}|=|\marg{list of
%     keywords}
%
%     定义、添加或删除关键词列表 2 中的关键词；注意
%     这等价于 |keywords=[2]|\ldots 等。强烈
%     不鼓励使用 \keyname{ndkeywords}。
%
%   \item[0.19,,{addon,optional}] \rkeyname{texcs}|=|\oarg{class
%     number}\marg{list of control sequences \textup(without
%     backslashes\textup)}
%   \item[0.20,,{addon,optional}] \rkeyname{moretexcs}|=|\oarg{class
%     number}\marg{list of control sequences \textup(without
%     backslashes\textup)}
%   \item[0.21,,{addon,optional}] \rkeyname{deletetexcs}|=|\oarg{class
%     number}\marg{list of control sequences \textup(without
%     backslashes\textup)}
%
%         \TeX\ 和 \LaTeX 中的控制序列同理。
%
%   \item[0.18,,optional] \rkeyname{directives}|=|\marg{list of compiler
%     directives}
%   \item[0.21,,optional] \rkeyname{moredirectives}|=|\marg{list of
%     compiler directives}
%   \item[0.21,,optional] \rkeyname{deletedirectives}|=|\marg{list of
%     compiler directives}
%
%     定义 C、\Cpp、Objective-C 和 POV 中的编译器指令。
%
%   \item[0.14] \rkeyname{sensitive}|=|\meta{\alternative{true,false}}
%
%     分别使关键词、控制序列和指令区分大小写
%     或不区分大小写。该键只在处理列表时影响关键词、控制
%     序列和指令。在其他所有情况下它们都区分大小写，例如
%     |deletekeywords={save,Test}| 会删除 `save' 和 `Test'，但两个
%     `SavE' 和 `test' 都不会被删除。
%
%   \item[0.19] \rkeyname{alsoletter}|=|\marg{character sequence}
%   \item[0.19] \rkeyname{alsodigit}|=|\marg{character sequence}
%   \item[0.19] \rkeyname{alsoother}|=|\marg{character sequence}
%
%     所有标识符（关键词、指令等）都由一个字母
%     后接字母数字字符（字母和数字）组成。例如，如果写
%     |keywords={one-two,\#include}|，减号必须变成数字、
%     井号必须变成字母，否则无法识别这些关键词。
%
%     表~\ref{rStdCharTable} 展示了 \packagename{listings} 宏包的
%     标准配置。这三个键会覆盖默认行为。序列中的每个
%     字符分别变成字母、数字和其他。
%
%   \item[0.20] \rkeyname{otherkeywords}|=|\marg{keywords}
%
%     定义含有其他字符、或以数字开头 的关键词。
%     每个给定的``关键词''都按关键词样式打印，但不
%     改变字符的``字母''、``数字''和``其他''状态。该键
%     用来定义 |=>|、|->|、|-->|、|--|、|::| 等关键词。
%     如果一个关键词是另一个的子序列（如 |--| 和 |-->|），
%     必须先给出较短的那个。
%
%   \item[0.20,,{renamed,optional}]
%     \rkeyname{tag}|=|\meta{character}\meta{character}^^A
%     \syntaxor\rkeyname{tag}|={}|\label{uoption:tag}
%
%     一阶关键词只在第一个和第二个字符之间生效。
%     该键用于 HTML。
% \end{syntax}
%
% \begin{table}[tb]
% \caption{标准字符表}\label{rStdCharTable}
% \begin{tabular}{ll}
% 类别 & 字符\\
% \noalign{\smallskip}
% letter & \texttt{A B C D E F G H I J K L M N O P Q R S T U V W X Y Z}\\
%        & \texttt{a b c d e f g h i j k l m n o p q r s t u v w x y z}\\
%        & \texttt{@ \textdollar\ } |_|\\
% digit  & \texttt{0 1 2 3 4 5 6 7 8 9}\\
% other  & \texttt{!\ " \#\ \%\ \&\ ' ( ) * + , - .\ / :\ ; < = > ?}\\
%        & {\catcode`\|=12\texttt{[ \char92\ ] \textasciicircum\ \char123\
%            | \char125\ \textasciitilde}}\\
% space  & \texttt{chr(32)}\footnotemark{}\\
% tabulator & \texttt{chr(9)}\\
% form feed & \texttt{chr(12)}\\
% \noalign{\smallskip}
% \end{tabular}
% \par\noindent
% 注意：编码 128--255 的扩展字符（如果定义了）\emph{目前}
% 是字母。
% \end{table}
% \footnotetext{这里和其他地方 |chr(|\meta{number}|)| 是指一个
% 返回具有 ASCII 值 \meta{number} 的字符的 Pascal 函数。}
%
%
% \paragraph{字符串}\label{uStrings}
% \begin{syntax}
%   \item[0.12]
%     \rkeyname{string}|=|\oarg{\alternative{b,d,m,bd,s}}\marg{delimiter
%     \textup(character\textup)}
%   \item[0.21]
%     \rkeyname{morestring}|=|\oarg{\alternative{b,d,m,bd,s}}\marg{delimiter}
%   \item[0.21]
%     \rkeyname{deletestring}|=|\oarg{\alternative{b,d,m,bd,s}}\marg{delimiter}
%
%     从字符串定界符列表中定义、添加或删除定界符。
%     起始和结束定界符相同，即源代码中定界符必须
%     彼此匹配。
%
%     可选参数是类型，控制定界符本身在字符串或字符
%     字面量中的表示方式：它被 |b| 反斜杠转义、|d| 重复
%     （两者都允许的写法是 |bd|）。
%     此外，类型也可以指一种不寻常的定界符形式：
%     |s| 字符串定界符（类似 |s| 注释类型）或 |m| Matlab 风格的
%     定界符。后者是 Ada 和 Matlab 以及
%     可能其他语言的特殊类型，这些语言中字符串定界符也用于
%     其他用途。它等价于 |d|，区别在于字母、右圆括号、右方括号
%     或某些其他字符之后
%     不会开始字符串。
% \end{syntax}
%
%
% \paragraph{注释}
% \begin{syntax}
%   \item[0.13]
%     \rkeyname{comment}|=|\oarg{type}\meta{delimiter\textup(s\textup)}
%   \item[0.21]
%     \rkeyname{morecomment}|=|\oarg{type}\meta{delimiter\textup(s\textup)}
%   \item[0.21]
%     \rkeyname{deletecomment}|=|\oarg{type}\meta{delimiter\textup(s\textup)}
%
%     注释同理，但有些类型需要不止一个
%     定界符。下面的概述以 \keyname{morecomment} 为例，
%     但例子同样适用于 \keyname{comment} 和
%     \keyname{deletecomment}。
%
%   \item[0.13] \rkeyname{morecomment}|=[l]|\meta{delimiter}
%
%     该定界符开始一个注释行，注释行一般以
%     定界符开头、到行尾结束。如果字符序列 |//|
%     应当开始注释行（如 \Cpp、Comal 80 或 Java 中），
%     |morecomment=[l]//| 就是正确的声明。对 Matlab 则是 |morecomment=[l]\%|——注意
%     前置反斜杠。
%
%   \item[0.13] \rkeyname{morecomment}|=[s]|\marg{delimiter}\marg{delimiter}
%
%     这里有两个定界符。第二个结束以
%     第一个定界符开始的注释。如果需要两个这样的注释，可以
%     把这个类型用两次。C、Java、PL/I、Prolog 和 SQL 都用
%     |morecomment=[s]{/*}{*/}| 定义单行注释，Algol 用
%     |morecomment=[s]{\#}{\#}|，这意味着井号既是单个注释的开始
%     也是结束。
%
%   \item[0.13] \rkeyname{morecomment}|=[n]|\marg{delimiter}\marg{delimiter}
%
%     类似类型 |s|，但注释可以嵌套。不允许
%     相同的参数——想一想为什么！Modula-2 和
%     Oberon-2 使用 |morecomment=[n]{(*}{*)}|。
%
%   \item[0.18] \rkeyname{morecomment}|=[f]|\meta{delimiter}
%   \item[0.18] \rkeyname{morecomment}|=[f][commentstyle]|\oarg{n=preceding
%     columns}\meta{delimiter}
%
%     当且仅当定界符出现在固定的列号上，
%     即位于第 $n$ 列（从零起算）时，它才开始注释行。
%
%   \item[0.17,,optional] \rkeyname{keywordcomment}|=|\marg{keywords}
%   \item[0.21,,optional] \rkeyname{morekeywordcomment}|=|\marg{keywords}
%   \item[0.21,,optional] \rkeyname{deletekeywordcomment}|=|\marg{keywords}
%
%     关键词注释以关键词开始、以同一关键词
%     结束。考虑 |keywordcomment={comment,co}|。此时
%     `\textbf{comment}\allowbreak\ldots\textbf{comment}' 和
%     `\textbf{co}\ldots\textbf{co}' 都是注释。
%
%   \item[0.17,,optional]
%     \rkeyname{keywordcommentsemicolon}|=|\marg{keywords}\marg{keywords}^^A
%     \marg{keywords}
%
%     ``关键词注释分号''的定义需要三个
%     关键词列表，例如 |{end}{else,end}{comment}|。分号总是
%     结束这样的注释。第一个参数中的任何关键词开始
%     注释，第二个参数中的任何关键词结束它（分号也
%     是）；而以第三个参数中任何关键词开始的注释，只由下一个
%     分号终止。例子中所有可能的注释是 `\textbf{end}\ldots\textbf{else}'、
%     `\textbf{end}\ldots\textbf{end}'（不会再次开始注释）以及
%     `\textbf{comment}\ldots;' 和 `\textbf{end}\ldots;'。也许是个奇怪的
%     定义，但 Algol 和 Simula 就用这样的注释。
%
%     注意：这里的关键词不必是已定义关键词的
%     子集。如果不是，它们不会以关键词样式出现。
%
%   \item[0.17,,optional]
%     \rkeyname{podcomment}|=|\meta{\alternative{true,false}}
%
%     激活或关闭 POD——Perl 专用。
% \end{syntax}
%
%
% \subsection{安装}\label{rInstallation}
%
% \paragraph{软件安装}
% \begin{enumerate}
% \item 按照 \TeX\ 目录结构（TDS），应把 \packagename{listings} 宏包
%       的文件按如下方式放入目录：
%       \begin{center}
%       \begin{tabular}{lcl}
%       \texttt{listings.pdf}&$\to$&\texttt{texmf/doc/latex/listings}\\
%       \texttt{listings.dtx}, \texttt{listings.ins},\\
%       \texttt{listings.ind}, \texttt{lstpatch.sty},\\
%       \texttt{lstdrvrs.dtx}&$\to$&\texttt{texmf/source/latex/listings}
%       \end{tabular}
%       \end{center}
%       \item 把工作目录切换到 \texttt{texmf/source/latex/listings}
%             注意你可能没有补丁文件 \texttt{lstpatch.sty}。
%       如果不使用 TDS，只需相应调整下面的目录。
%             \item 创建目录 \texttt{texmf/tex/latex/listings}，如果它已存在，
%       则删除其中的所有文件，但保留 \texttt{lst}\meta{whatever}\texttt{0.sty} 和
%       \texttt{lstlocal.cfg}，并通过 \TeX\ 运行 \texttt{listings.ins}。
%       \item 如果这一步尚未做，就把生成的文件移到 \texttt{texmf/tex/latex/listings}。
%       \begin{center}
%       \begin{tabular}{lcl}
%       \texttt{listings.sty}, \texttt{lstmisc.sty},
%           &&\qquad(内核与附加件)\\
%       \texttt{listings.cfg},
%           &&\qquad(配置文件)\\
%       \texttt{lstlang}\meta{number}\texttt{.sty},
%           &&\qquad(语言驱动)\\
%       \texttt{lstpatch.sty}&$\to$&\texttt{texmf/tex/latex/listings}
%       \end{tabular}
%       \end{center}
% \item 如果你的 \TeX\ 实现使用文件名数据库，请更新它。
% \item 如果你以后收到补丁文件，把它放在
%       \texttt{listings.sty} 所在的位置（并更新文件名数据库）。
% \end{enumerate}
% 注意 \packagename{listings} 至少需要 \packagename{keyval} 宏包的
% 1.10 版，该宏包由 David Carlisle 包含在 \packagename{graphics} 包中。
%
%
% \paragraph{软件配置}
% 只有在遇到标准配置问题时，或者（例如）想让宏包适配外语时，
% 才需要读这一节。
%
% 永远不要修改 \packagename{listings} 宏包的文件，尤其是
% 配置文件。每次新安装或新版本都会覆盖
% 它。软件许可允许修改，但我不建议这样做。
% 更好的做法是创建一个或多个如下文件
% \begin{center}
% \begin{tabular}{lcl}
% \texttt{lstmisc0.sty} & 用于 & 本地附加件
%                                （见开发者指南），\\
% \texttt{lstlang0.sty} & 用于 & 本地语言定义
%                                （见 \ref{rLanguageDefinitions}），以及\\
% \texttt{lstlocal.cfg} & 作为 & 本地配置文件
% \end{tabular}
% \end{center}
% 并把它们放到其他 \packagename{listings} 文件所在的同一目录。
% 除非你删除它们，否则这三个文件不会被新安装改动。
% 如果 \texttt{lstlocal.cfg} 存在，它会在 \texttt{listings.cfg} 之后被加载。
% 你也许想修改下列参数之一。
% \begin{syntax}
% \item[,,data] \rcmdname\lstaspectfiles\quad contains~\rlap{\texttt{\lstaspectfiles}}
% \item[,,data] \rcmdname\lstlanguagefiles\quad contains~\rlap{\texttt{\lstlanguagefiles}}
%
%       宏包使用指定的这些文件来查找附加件和语言
%       定义。
% \end{syntax}
% 此外，你也许想调整
%   \icmdname\lstlistlistingname,
%   \icmdname\lstlistingname,
%   \icmdname\lstlistingnamestyle,
%   \ikeyname{defaultdialect},
%   \icmdname\lstalias, or
%   \icmdname\lstalias
% \ 等，如前文各节所述。
%
%
% \section{实验功能}\label{rExperimentalFeatures}
%
% 本节描述本宏包中或多或少尚未成熟的部分。它们不太可能全部被移除
% （除非明确说明），但它们可能会（大幅）更改和改进。这类功能在前面的
% 节中已用 \dag 标记。所以，如果你在这里发现任何 \dag 标记的内容，
% 你应该非常非常小心。
%
%
% \subsection{参数内的列表}\label{rListingsInsideArguments}
%
% 如果你想在参数内使用 |\lstinline| 或列表环境，有一些事情需要考虑。
% 由于 \TeX 在执行`\lst-宏'之前读取参数，本宏包无法做任何事情来保留输入：
% 空格收缩为一个空格，制表符和行尾被转换为空格，\TeX 的注释符不可打印，
% 等等。因此，\emph{你}必须多做一些工作。你必须在以下四个字符的每一个
% 前面加反斜杠：|\{}%|。此外，如果 (i) 有两个或更多空格连续出现，
% 或 (ii) 空格是行首字符，你必须以同样的方式保护空格。
% 这还不够：每行必须以`换行' |^^J| 结束。而且你不能在这样的列表中转义到
% \LaTeX！
%
% 最简单的例子是用 |\lstinline|，因为我们不需要换行。
% \begin{verbatim}
%\footnote{\lstinline{var i:integer;} and
%          \lstinline!protected\ \ spaces! and
%          \fbox{\lstinline!\\\{\}\%!}}\end{verbatim}
% yields\lstset{language=Pascal}\footnote{\lstinline{var i:integer;} and
%          \lstinline!protected\ \ spaces! and
%          \fbox{\lstinline!\\\{\}\%!}}
% 如果当前语言是 Pascal。注意本例展示了另一个实验功能：使用参数花括号作为
% 定界符。这在第~\ref{rTypesettingListings} 节中描述。
%
% 现在是一个环境例子：
% \begin{lstsample}{\lstset{language={}}}{}
%    \fbox{%
%    \begin{lstlisting}^^J
%    \ !"#$\%&'()*+,-./^^J
%    0123456789:;<=>?^^J
%    @ABCDEFGHIJKLMNO^^J
%    PQRSTUVWXYZ[\\]^_^^J
%    `abcdefghijklmno^^J
%    pqrstuvwxyz\{|\}~^^J
%    \end{lstlisting}}
% \end{lstsample}
% \begin{advise}
% \item 你可能会奇怪这个功能为什么还是实验性的。原因：
%       你不应该在参数内使用列表；这不总是安全的。
% \end{advise}
%
%
% \subsection{\dag\ 标识符的导出}\label{rExportOfIdentifiers}
%
% 如果能导出函数或过程名称就好了。一般来说这还是个梦想。问题在于编程语言
% 对函数和过程声明或定义使用各种语法。通用接口完全超出了本宏包的范围——
% 那是编译器的工作，不是美化排版工具的工作。不过，对特定语言是可能的：
% 例如在 Pascal 中，每个函数或过程定义和变量声明前面都有特定的关键词。
% 注意你必须用 \texttt{procnames} 选项请求以下键：
% |\usepackage[procnames]{listings}|。
% \begin{syntax}
% \item[0.19,{{}},{\dag optional}] \rkeyname{procnamekeys}|=|\marg{keywords}
% \item[0.21,,\dag optional] \rkeyname{moreprocnamekeys}|=|\marg{keywords}
% \item[0.21,,\dag optional] \rkeyname{deleteprocnamekeys}|=|\marg{keywords}
%
%		每个指定的关键词表示函数或过程定义。该关键词后的任何标识符
%		以`procname'样式出现。对于 Pascal，你可能使用\vspace{-.5\baselineskip}
% \begin{verbatim}
%    procnamekeys={program,procedure,function}\end{verbatim}
%
% \item[0.19,keywordstyle,\dag optional] \rkeyname{procnamestyle}|=|\meta{style}
%
%		定义过程名和函数名出现的样式。
%
% \item[0.19,false,\dag optional] \rkeyname{indexprocnames}|=|\meta{\alternative{true,false}}
%
%		若激活，过程名和函数名也会被编入索引。
% \end{syntax}
% \begin{TODO}
% \aspectname{procnames} 方面并不理想（至少自 2000 年以来
% 未作改动）。它到目前为止只标记并索引函数定义，但
% 例如也可以标记随后的函数调用。
% 可以用一个键控制是否把函数名加入某个特殊关键词类，使之以 `procname' 样式出现。
% 但是这些名字应该全局添加吗？两种做法都有充分的理由。
% 当然，我们还需要一个键来重置名字列表。
% \end{TODO}
%
%
% \subsection{\dag\ 超链接引用}\label{rHyperReferences}
%
% 由于它是实验性的，这个非常小的方面必须通过 \texttt{hyper} 选项请求。
% 将来的一种可能是把这个方面
% 与 \aspectname{procnames} 结合。那样例如就可以点击一个
% 函数名跳到它的定义。
% \begin{syntax}
% \item[0.21,,{\dag optional}] \rkeyname{hyperref}|=|\marg{identifiers}
% \item[0.21,,{\dag optional}] \rkeyname{morehyperref}|=|\marg{identifiers}
% \item[0.21,,{\dag optional}] \rkeyname{deletehyperref}|=|\marg{identifiers}
%
%       把指定的标识符做成超链接（通过 \packagename{hyperref}
%       宏包）。在这些标识符上``点击''会跳到它上一次
%       出现的位置。
%
% \item[0.21,\hyper@@anchor,{\dag optional}] \rkeyname{hyperanchor}|=|\meta{two-parameter macro}
% \item[0.21,\hyperlink,{\dag optional}] \rkeyname{hyperlink}|=|\meta{two-parameter macro}
%
%       分别设置超链接锚点和链接。
%       默认设置适用于 \packagename{hyperref} 宏包。
% \end{syntax}
%
%
% \subsection{文学编程}
%
% 我们从一个例子开始，并隐藏关键的键值列表。
% \begin{lstsample}{%
%   \lstset{literate={:=}{{$\gets$}}1 {<=}{{$\leq$}}1 {>=}{{$\geq$}}1 ^^A
%     {<>}{{$\neq$}}1}}{}
%    \begin{lstlisting}
%    var i:integer;
%
%    if (i<=0) i := 1;
%    if (i>=0) i := 0;
%    if (i<>0) i := 0;
%    \end{lstlisting}
% \end{lstsample}
% 有趣，不是吗？我们可以在列表中保留 |i := 0| 而不是
% i| |$\gets$| |0，但那不够文学化！^^A :-)
% 现在你可能想知道这是怎么做到的。请\emph{仔细}看看下面的键。
% \begin{syntax}
%   \item[0.20,,\dag] \rkeyname{literate}|=|[|*|]\meta{replacement
%     item}\ldots\meta{replacement item}
%
%     首先注意各项之间没有逗号。每项由三个参数组成：\marg{replace}
%     \marg{replacement text}\marg{length}。\meta{replace} 是原始字符序列。
%     我们不打印这些字符，而是使用 \meta{replacement text}，
%     它在输出中占据 \meta{length} 个字符的宽度。
%
%     \meta{replacement text} 中的每个`打印单元'\emph{必须}用花括号括起来，
%     除非它是单个字符。例如，你必须用花括号括住 |$\leq$|。
%     如果你想用 |$\leftarrow1\rightarrow$| 替换 |<-1->|，
%     替换项应该是 |{<-1->}{{$\leftarrow$}1{$\rightarrow$}}3|。
%     注意箭头周围的花括号。
%
%     如果一个 \meta{replace} 是另一个 \meta{replace} 的子序列，
%     你必须先定义较短的序列。例如，必须在 |{--}| 之前定义 |{-}|，
%     在 |{-->}| 之前定义 |{--}|。
%
%     可选星号表示不在字符串、注释和其他分隔文本中进行文学替换。
% \end{syntax}
% 在上面的例子中，我使用了
% \begin{verbatim}
%  literate={:=}{{$\gets$}}1 {<=}{{$\leq$}}1 {>=}{{$\geq$}}1 {<>}{{$\neq$}}1\end{verbatim}
% \begin{TODO}
% 当然，有用于添加和删除单个 \meta{replacement item} 的键会很好。
% 也许键应该以与字符串和注释定义相同的方式工作，即每个键值对一项。
% 这样在子序列的情况下更容易提供更好的自动检测。
% \end{TODO}
%
%
% \subsection{\packagename{LGrind} 定义}\label{rLGrindDefinitions}
%
% 是的，从其他程序窃取语言定义是个糟糕的主意。不过，对
% \packagename{LGrind} 定义文件来说这是可能的——至少部分可能。
% 请注意本文件必须被找到
% \TeX.
% \begin{syntax}
% \item[0.21,,{optional}] \rkeyname{lgrindef}|=|\meta{language}
%
%       扫描 \texttt{lgrindef} 语言定义文件，查找
%       \meta{language} 并在存在时激活它。注意并非所有
%       \packagename{LGrind} 功能都有 \packagename{listings}
%       的对应实现。
%
%       注意 `Linda' 语言不能正常工作，因为它把前面带 `|#|' 的
%       编译器指令定义为关键词。
%
% \item[0.21,lgrindef.,{data,optional}] \rcmdname\lstlgrindeffile
%
%       包含定义文件的（路径和）名字。
% \end{syntax}
%
%
% \subsection{\dag\ 自动格式化}
%
% \lstloadaspects{formats}^^A
% 自动源代码格式化远远谈不上好。首先，关于源代码应该如何格式化没有通用规则。
% 所以`格式定义'必须灵活。这种灵活性需要复杂的接口、强大的`格式定义'解析器，
% 以及幕后的大量代码。目前，格式定义还不够灵活（可能不是定义本身而是结果）。
% 单个`格式项'的形式为
% \begin{itemize}\item[]
%     \meta{input chars}|=|\oarg{exceptional chars}\meta{pre}\oarg{^^A
%       \texttt{\string\string}}\meta{post}
% \end{itemize}
% 每当 \meta{input chars} 后面不跟 \meta{exceptional chars} 之一时，
% 就根据值的其余部分进行格式化。如果未指定 |\string|，输入字符不被打印
% （除非它是标识符或关键词）。否则 \meta{pre} 在打印原始字符串之前`执行'，
% \meta{post} 在之后。这两个是以下内容的`子集'：
% \begin{itemize}
% \item |\newline| ---确保新行；
% \item |\space| ---确保空白；
% \item |\indent| ---增加缩进；
% \item |\noindent| ---减少缩进。
% \end{itemize}
% 现在我们可以给一个例子了。\lstaspectindex{\lstdefineformat}{}\lstaspectindex{format}{}
% \begin{lstxsample}
%    \lstdefineformat{C}{%
%        \{=\newline\string\newline\indent,%
%        \}=\newline\noindent\string\newline,%
%        ;=[\ ]\string\space}
% \end{lstxsample}
% \begin{lstsample}{\lstset{language={}}}{}
%    \begin{lstlisting}[format=C]
%    for (int i=0;i<10; i++){/* wait */};
%    \end{lstlisting}
% \end{lstsample}
% 不好。但有一个（太？）简单的解决方法：
% \begin{lstxsample}
%    \lstdefineformat{C}{%
%        \{=\newline\string\newline\indent,%
%        \}=[;]\newline\noindent\string\newline,%
%        \};=\newline\noindent\string\newline,%
%        ;=[\ ]\string\space}
% \end{lstxsample}
% \begin{lstsample}{\lstset{language={}}}{}
%    \begin{lstlisting}[format=C]
%    for (int i=0;i<10; i++){/* wait */};
%    \end{lstlisting}
% \end{lstsample}
% 有时候问题只是找到合适的格式定义。进一步的格式化很复杂。
% 这里只给出三个难度递增的例子。
% \begin{enumerate}
% \item 插入水平间距以分隔函数/过程名和后面的括号，或分隔函数的参数，
%       例如在逗号后加空格（如果在函数调用内）。
% \item 长行的智能断行。考虑长的`and/or'表达式。格式化应该遵循逻辑结构！
% \item 上下文相关的格式化规则。如果空块或小块在输出中占三行或更多行
%       会很烦人——想想一直向下滚动。所以如果块格式化是上下文相关的就好了。
% \end{enumerate}
% 注意这是提供自动格式化的一个非常初步和笨拙的尝试——之所以笨拙，
% 是因为这个问题并不简单。欢迎任何想法。也欢迎实现。最后你应该知道，
% 你必须在加载宏包时请求格式定义，例如通过 |\usepackage[formats]{listings}|。
%
% \subsection{任意行范围标记}\label{rArbitraryLinerangeMarkers}
%
% 除了在 \keyname{linerange} 中使用行号外，还可以使用文本标记。每个这样的
% 标记由 \meta{prefix}、\meta{text} 和 \meta{suffix} 组成。你定义一次（或多次）
% 前缀和后缀，然后用标记文本代替行号。
% \begin{lstxsample}
%    \lstset{rangeprefix=\{\ ,% curly left brace plus space
%            rangesuffix=\ \}}% space plus curly right brace
% \end{lstxsample}
% \begin{lstsample}{}{}
%    \begin{lstlisting}%
%          [linerange=loop\ 2-end]
%    { loop 1 }
%    for i:=maxint to 0 do
%    begin
%        { do nothing }
%    end;
%    { end }
%    { loop 2 }
%    for i:=maxint to 0 do
%    begin
%        { do nothing }
%    end;
%    { end }
%    \end{lstlisting}
% \end{lstsample}
% 注意 \TeX 的特殊字符如花括号、空格、百分号等必须用反斜杠转义。
% \begin{syntax}
% \item[1.2] \rkeyname{rangebeginprefix}|=|\meta{prefix}
% \item[1.2] \rkeyname{rangebeginsuffix}|=|\meta{suffix}
% \item[1.2] \rkeyname{rangeendprefix}|=|\meta{prefix}
% \item[1.2] \rkeyname{rangeendsuffix}|=|\meta{suffix}
%
%       为开始标记和结束标记分别定义前缀和后缀。
%
% \item[1.2] \rkeyname{rangeprefix}|=|\meta{prefix}
% \item[1.2] \rkeyname{rangesuffix}|=|\meta{suffix}
%
%       为开始标记和结束标记定义相同的前缀和后缀。
%
% \item[1.2,true] \rkeyname{includerangemarker}|=|\meta{\alternative{true,false}}
%
%       在输出中显示或隐藏标记。
%
%       \textbf{注意：} 如果设置了 |firstnumber|，它指的是包含标记的行。
%       所以如果想从数字 1 开始一个范围，必须设置
%       |includerangemarker=false, firstnumber=0|。
% \end{syntax}
% \begin{lstsample}{\lstset{rangeprefix=\{\ ,rangesuffix=\ \}}}{}
%    \begin{lstlisting}%
%          [linerange=loop\ 1-end,
%           includerangemarker=false,
%           frame=single]
%    { loop 1 }
%    for i:=maxint to 0 do
%    begin
%        { do nothing }
%    end;
%    { end }
%    \end{lstlisting}
% \end{lstsample}
%
%
% \subsection{多列列表}\label{rMulticolumnListings}
%
% 当加载 \packagename{multicol} 宏包后，可以用它排版
% 多栏列表。多栏列表由 |multicols| 键指定。例如：
% \begin{lstsample}{}{}
%    \begin{lstlisting}[multicols=2]
%    if (i < 0)
%      i = 0
%      j = 1
%    end if
%    if (j < 0)
%      j = 0
%    end if
%    \end{lstlisting}
% \end{lstsample}
%
% 已知多栏选项在某些键下会失败。
%
% \begin{advise}
% \item 哪些键？
%       \advisespace
%       很遗憾，我不知道。Carsten 把这个选项的代码
%       留在 1.3b 版补丁文件中，只留下这条晦涩的
%       说明作为文档。欢迎提交缺陷报告，不过我并不承诺
%       它们可以修好。——Brooks
% \end{advise}
%
%
%\iffalse
% \section{未来发展？}
%
% 本节相当粗略，只是列出一些我不想
% 忘记的事情。
%
% 首先，我想支持更多语言，例如 Maple、
% PostScript 等。幸好我的寿命有限，所以这些
% 工作可以由别人来做。请把你们的语言定义（电子）邮寄给我。
%
% 其次，将来还有一些想法。有些已经作为`待办事项'
% 提出；有些来自他人，列在下面；还有一些
% 离实现还很远，
%   e.g.~\keyname{linerange}|=|\oarg{inter}\marg{line range list}
% 它会打印范围内的所有行，并在省略某些代码行时执行 \meta{inter}。
% 这里的主要问题是边框和背景色；
% 它们应该怎样处理？事实上，问题在于这如何编码。
% 另一个想法是为特定的代码块改变背景色（或基本样式）。
% 这同样不容易。
%
%^^A Auto-detect whether surplus space (from spaces and tabs) isn't needed to fix
%^^A alignment of wide character combinations like |==| or |<>|.
%^^A
%^^A Make package compatible to calc package.
%^^A
%^^A Rewrite \lst@LAS, \lst@DefDriver, \lst@Require to distinguish loading
%^^A of languages (which don't need base languages at once) and aspects
%^^A (which need required aspects to be loaded).
%
% \lsthelper{Vincent~Poirriez}{1999/11/18}{code examples inside caml comments}:
% 在 caml 注释内部，|[| 和 |]| 之间的代码应以
% basicstyle（或另一个新引入的样式）打印。这些
% `代码示例定界符'允许嵌套，例如~|(* [[x;y]] *)|。
%
% \lsthelper{Claus~Atzenbeck}{1999/12/03}{`extendedchars=false' 在使用扩展字符时不发
% 警告}：在 final 模式下，如果
% \ikeyname{extendedchars}|=false| 但使用了扩展字符，就发出警告。
%
% \lsthelper{Andreas~Matthias}{2000/01/04}{定义页眉/页脚以打印
% 列表名}：让页眉/页脚打印列表名。有些
% 人要求连续的标题。
%\fi
%
%
% \part{技巧与窍门}
%
% 注意：这部分文档正在编写中。
% 第~\ref{uHowTos} 节必须按主题排序并以某种方式组织。
% 此外还计划了一个新节`示例'，但尚未写出。
% 主要问题是缺乏时间……
%
%
% \section{故障排查}\label{uTroubleshooting}
%
% 如果你遇到 \packagename{listings} 宏包的问题，在提交错误报告之前，
% 你应该遵循以下步骤。首先查阅参考指南，看看问题是否已知。如果不是，
% 创建一个\emph{最小}文件来重现问题。遵循以下说明：
% \begin{enumerate}
% \item 从第~\ref{uAMinimalFile} 节的最小文件开始。
% \item 添加导致问题的 \LaTeX 代码，但保持简短。特别是，保持附加宏包的数量少。
% \item 从文件中删除一些代码（以及相应的宏包），直到问题消失。
%       这样你就找到了关键部分。
% \item 再次添加这段代码，从步骤 3 重新开始，直到所有代码和所有宏包
%       都是必要的。
% \item 你现在有了一个最小文件。将错误报告发送到本文档首页的地址，
%       并附上最小文件以及生成的 \texttt{.log} 文件。如果你使用非常特殊的
%       宏包（即不在 CTAN 上的），如果其软件许可允许，也请附上该宏包。
% \end{enumerate}
%
%
% \section{已知错误与变通方法}\label{uBugsWorkarounds}
%
% \subsection{参数内的列表}\label{uListingsArguments}
%
% 很长一段时间里，无法在表格的单元格中使用 \verb-\lstinline{...}-
% \makeatletter\@ifundefined{r@uProcessingInline}{}{%
% （更多信息见第~\pageref{uProcessingInline} 页
% 的 \ref{uProcessingInline} 节）}，%
% \makeatother%
% 但当时有一个推荐的变通方法。这个变通方法至今
% 仍然可用，不过现在可以直接使用 \verb-\lstinline{...}-，当然也可以直接
% 使用 \verb-\lstinline|...|-：
% \begin{lstsample}[lstlisting]{}{}
%    \newcommand\foo{\lstinline{t}}
%    \newcommand\foobar[2][]{\lstinline[#1]{#2}}
%
%    \begin{tabular}{ll}
%    \lstinline|r| & a first variable (standard)\\
%    \lstinline[language=java]|int s;| & a standard declaration \\
%    \foo & a second variable (workaround)\\
%    \foobar[language=java]{int u;} & a (workaraond) declaration \\
%    \lstinline{v} & another variable using braces\\
%    \lstinline[language=java]{int w;} & an additional braced declaration
%    \end{tabular}
% \end{lstsample}
%
%
% \subsection{带背景色和 \LaTeX{} 转义
% 公式的列表}
% \label{uListingsBackgroundColour}
%
% 如果有任何文字转义到 \LaTeX{} 并带有某种彩色背景
% 和周围的边框，那么背景以及
% 构成边框的线条中都会出现空隙。
% \begin{lstsample}[lstlisting]{}{}
%    \begin{lstlisting}[language=C, mathescape,
%      backgroundcolor=\color{yellow!10}, frame=tlb]
%    /* the following code computes $\displaystyle\sum_{i=1}^{n}i$ */
%    for (i = 1; i <= limit; i++) {
%      sum += i;
%    }
%    \end{lstlisting}
% \end{lstsample}
%
% 目前只有一个变通方法：
% \begin{itemize}
%     \item 把你的代码写入外部文件 \meta{filename}。
%     \item 用 |\lstinputlisting|\meta{filename} 把代码读入你的
%         文档，并用 |\begin{mdframed}| 生成的边框把它包围起来
%     \ldots{} |\end{mdframed}|.
% \end{itemize}
% \begin{lstsample}[lstlisting]{}{}
%    \begin{verbatimwrite}{temp.c}
%    /* the following code computes $\displaystyle\sum_{i=1}^{n}i$ */
%    for (i = 1; i <= limit; i++) {
%      sum += i;
%    }
%    \end{verbatimwrite}
%    \begin{mdframed}[backgroundcolor=yellow!10, rightline=false]
%      \lstinputlisting[language=C,mathescape,frame={}]{./temp.c}
%    \end{mdframed}
% \end{lstsample}
% 关于 |verbatimwrite| 环境的更多信息，请看
% \cite{Fairbairns:moreverb}，|mdframed| 环境在
% \cite{DanielSchubert:mdframed}.
%
%
% \section{操作指南}\label{uHowTos}
%
%
% \subsubsection*{如何引用行号}
% 也许你想把 |\label{|\meta{whatever}|}| 放进一个 \LaTeX\ 转义中，而这个转义
% 位于定界符不被打印的注释里？如果你这样做，编译器
% 看不到 \LaTeX\ 代码，因为它位于注释内部，而
% \packagename{listings} 宏包也不会打印任何东西，因为定界符会
% 被丢弃，且 |\label| 不产生任何可打印输出，但你仍然
% 可以引用行号。好，你的愿望实现了。
%
% 例如在 Pascal 中，你可以让宏包把`特殊'
% 注释定界符 |(*@| 和 |@*)| 识别为转义开始和结束序列。
% 这样你就可以用这个特殊注释来放 |\label| 和其他东西。
% \begin{lstsample}{\lstset{numberstyle=\tiny,stepnumber=2,numbersep=5pt}}{}
%    \lstset{escapeinside={(*@}{@*)}}
%
%    \begin{lstlisting}
%    for i:=maxint to 0 do
%    begin
%        { comment }(*@\label{comment}@*)
%    end;
%    \end{lstlisting}
%    Line \ref{comment} shows a comment.
% \end{lstsample}
% \begin{advise}
% \item 我可以改用 `|(*@|' 和 `|*)|' 吗？
%       \advisespace
%       可以。
% \item 我可以改用 `|(*|' 和 `|*)|' 吗？
%       \advisespace
%       当然。只要你愿意。
% \item 我可以改用 `|{@|' 和 `|@}|' 吗？
%       \advisespace
%       不行，绝对不行！第二个定界符不允许。字符 `|@|' 被
%       定义为检查转义是否结束。但读取这个孤立的
%       ``结束参数''花括号时，\TeX\ 会遇到错误 `\texttt{Argument of @
%       has an extra \char125}'。抱歉。
% \item 我可以改用 `|{|' 和 `|}|' 吗？
%       \advisespace
%       不行。第二个定界符同样不允许。这里 \TeX\ 会
%       给你一个 `\texttt{Runaway argument}' 错误。因为 `|}|' 被定义为
%       检查转义是否结束，它无法充当``结束参数''花括号。
% \item 那怎么使用注释行呢？
%       \advisespace
%       例如，写 `|escapeinside={//*}{\^^M}|'。这里 |\^^M| 代表
%       行尾字符。
% \end{advise}
%
%
% \subsubsection*{如何吞掉字符}
% 为了让 \LaTeX\ 代码更易读，你也许想缩进你的
% \texttt{lstlisting} 列表。不过这种缩进不应出现在
% 美化排版后的列表中，所以必须去掉。如果每行代码
% 缩进三个字符，可以用 |gobble=3| 去掉它们：
% \begin{lstsample}{}{\lstset{showspaces}}
%    \begin{lstlisting}[gobble=3]
%    1  for i:=maxint to 0 do
%     2 begin
%      3    { do nothing }
%    123end;
%
%       Write('Case insensitive ');
%       WritE('Pascal keywords.');
%    \end{lstlisting}
% \end{lstsample}
% 注意空行以及环境的开头和结尾
% 不必遵守这个缩进。但切勿让结尾的缩进超过
% `\ikeyname{gobble}' 个字符。此外注意，制表符在吃掉之前会先展开为
% |tabsize| 个空格。
% \begin{advise}
% \item 我可以将 \ikeyname{gobble}' 与 |\lstinputlisting|' 一起使用吗？
%       \advisespace
%       可以，但它没有效果。
%
% \item 注意 \ikeyname{gobble}' 也可以通过 |\lstset|' 设置。
% \end{advise}
%
%
% \subsubsection*{如何插入图形}
% \lsthelper{Herbert~Weinhandl}{1999/09/06}{listings + eps} 发现了一种非常简单的
% 在列表中插入图形的方法。感谢贡献这个想法——一个
% 我永远想不出的想法。
%
% 有些编程语言允许美元符号作为标识符的一部分。
% 但除了中间函数名或库函数之外，这个
% 字符大多不被使用。\packagename{listings} 宏包定义了
% \ikeyname{mathescape} 键，它让 `|$|' 转义到 \TeX 的数学模式。
% 这使得美元字符成为我们这里目的的最佳候选：
% 使用一个能插入图形的宏包，把 \ikeyname{mathescape} 设为 true，
% 并把图形放在两个美元符号之间，而这两个符号位于注释里。
%
% 下面的例子最初来自 Herbert 给我的一个头文件。
% 为了在此展示，我使用 \texttt{lstlisting} 环境和一个
% 头文件摘录。|\includegraphics| 命令来自
% David Carlisle 的 \packagename{graphics} 包。
% \begin{verbatim}
%   \begin{lstlisting}[mathescape=true]
%   /*
%    $ \includegraphics[height=1cm]{defs-p1.eps} $
%    */
%   typedef struct {
%     Atom_T          *V_ptr;   /* pointer to Vacancy in grid    */
%     Atom_T          *x_ptr;   /* pointer to (A|B) Atom in grid */
%   } ABV_Pair_T;
%   \end{lstlisting}\end{verbatim}
% 结果看起来相当不错。可惜你看不到，因为排版
% 本手册时那个图形不可用。
%
%
% \subsubsection*{如何让每页都有闭合的边框}
% 本宏包只对不跨页的列表支持封闭边框。
% 如果一个列表分在两页上，那么既没有页面
% 底部的下边框线，也没有下一页的
% 上边框线。如果你坚持要这些线，可以使用 Donald Arseneau 的 \texttt{framed.sty}。
% 那样你可以写
% \begin{verbatim}
%    \begin{framed}
%    \begin{lstlisting}
%      or \lstinputlisting{...}
%    \end{lstlisting}
%    \end{framed}\end{verbatim}
% 本宏包还提供了一个 \texttt{shaded} 环境。如果使用它，
% 不要忘记用 \packagename{color} 宏包
% 定义 \texttt{shadecolor}。
%
%
% \subsubsection*{如何用 $\Lambda$ 打印各国
%   \packagename{listings}}\label{uNationalCharacters}
%
% 除了直接键入各国字符，你还可以使用第~\ref{rEscapingToLaTeX} 节
% 描述的`转义'功能。
% 键 \ikeyname{escapechar}、\ikeyname{escapeinside} 和 \ikeyname{texcl}
% 允许部分使用 \LaTeX\ 代码。
%
% 现在，如果你使用 $\Lambda$（Lambda，面向 Omega 的 \LaTeX\ 变体）并且
% 想要例如阿拉伯语的注释行，就不必为每个转义的注释行
% 都写 |\begin{arab}| \ldots\ |\end{arab}|。这可以自动化：
% \begin{verbatim}
%    \lstset{escapebegin=\begin{arab},escapeend=\end{arab}}
%
%    \begin{lstlisting}[texcl]
%    // Replace text by Arabic comment.
%    for (int i=0; i<1; i++) { };
%    \end{lstlisting}\end{verbatim}
% 如果你的编程语言没有注释行，就必须使用
% \ikeyname{escapechar} 或 \ikeyname{escapeinside}：
% \begin{verbatim}
%    \lstset{escapebegin=\begin{greek},escapeend=\end{greek}}
%
%    \begin{lstlisting}[escapeinside=`']
%    /* `Replace text by Greek comment.' */
%    for (int i=0; i<1; i++) { };
%    \end{lstlisting}\end{verbatim}
% 注意这里的定界符 |`| 和 |'| 是必需的。没有它们
% 例子无法工作。如果编程语言的注释定界符
% 是单个字符（如 Pascal 中的花括号），还有更聪明的做法：
% \begin{verbatim}
%    \lstset{escapebegin=\textbraceleft\begin{arab},
%            escapeend=\end{arab}\textbraceright}
%
%    \begin{lstlisting}[escapeinside=\{\}]
%    for i:=maxint to 0 do
%    begin
%        { Replace text by Arabic comment. }
%    end;
%    \end{lstlisting}\end{verbatim}
% 请注意，到 $\Lambda$ 的`接口'完全没有经过测试。
% 欢迎报告问题！
%
%
% \subsubsection*{如何得到粗体等宽关键词}
% 使用 \href{http://mirror.ctan.org/fonts/luximono}{\packagename{LuxiMono}} 宏包。
%
% \iffalse
% 很多人要求粗体打字机字体，因为 \LaTeX\ 标准
% 发行版中没有包含它们。下面给出一个
% 在那种情况下如何使用它们的答案。
% \begin{advise}
% \item 请注意，我个人并不认为下面是一个好的
%       解决方案。这种粗体打字机字形太重了。更好的
%       做法是用 \texttt{cmtt} 的细体作为基本字体，用 \texttt{cmtt}
%       或\emph{稍}重的字形来排关键词。
%
% \item 你为什么不告诉我们如何使用更好的方案？
%       \advisespace
%       \texttt{cmtt} 的细体并不存在。一旦它可用，
%       你就可以照下面描述的方法做类似的事。
% \end{advise}
% 首先，你需要粗体打字机的 Metafont 源文件，例如
% \texttt{cmbtt8.mf}、\texttt{cmbtt9.mf} 和 \texttt{cmbtt10.mf}，来自
% \href{http://mirror.ctan.org/fonts/cm/mf-extra/bold}
%      {CTAN/fonts/cm/mf-extra/bold}.
% 其次，你必须生成 \texttt{.tfm} 文件，即对这些源文件
% 运行 Metafont 程序。当你在文档中使用
% 这些字体时，这可能会自动完成。最后，你必须告诉 \LaTeX\ 你已安装
% 了粗体打字机字体。只需在文档的导言区使用
% \begin{verbatim}
%    \DeclareFontShape{OT1}{cmtt}{bx}{n}
%         {<5><6><7><8>cmbtt8%
%          <9>cmbtt9%
%          <10><10.95>cmbtt10%
%          <12><14.4><17.28><20.74><24.88>cmbtt10%
%          }{}\end{verbatim}
% 如果你经常使用这些字体，也许想
% 在本地复制一份 \texttt{ot1cmtt.fd} 并在那里替换声明。
% 但注意，你不允许以原文件名分发
% 修改后的文件！
% \fi
%
%
% \subsubsection*{如何使用纯文本}
% 如果你想用 \packagename{listings} 排纯文本（也许带
% 行号，或者像 |verbatim| 但带换行，等等），
% 就用空语言：|\lstset{language=}|。
%
%
% \subsubsection*{如何获得开发者指南}
% 在 listings 宏包的\emph{源目录}中，即
% \texttt{.dtx} 文件所在处，创建文件 \texttt{ltxdoc.cfg}，
% 内容如下。
% \begin{verbatim}
%    \AtBeginDocument{\AlsoImplementation}\end{verbatim}
% 然后把 \texttt{listings.dtx} 用 \LaTeX\ 跑两遍，运行 Makeindex（带
% |-s gind.ist| 选项），最后再对 \texttt{listings.dtx}
% 跑一遍 \LaTeX。这样就会生成完整文档，包括用户
% 指南、参考指南、开发者指南和实现。
%
% 如果你能运行（GNU）make 程序，执行命令
% \begin{verbatim}
%    make all\end{verbatim}
% or
% \begin{verbatim}
%    make listings-devel.pdf\end{verbatim}
% or
% \begin{verbatim}
%    make pdf-devel\end{verbatim}
% 会得到相同结果——它被命名为 \texttt{listings-devel.pdf}。
%
%
% \subsubsection*{如何从文档中复制粘贴}
% \label{sec:upquote}
% 如果你想提供一份由本宏包（\packagename{listings}）排版、
% 可供复制粘贴的列表文档作为源，就应当
% 确保把 |upquote| 键设为 |true|，并在
% 文档导言区加入 |\usepackage{textcomp}|。否则
% 花引号对 |`'| 会关联到 Unicode 位置
% U+2018 和 U+2019 \cite{CovingtonMittelbachKuhn:upquote}。这些
% 字符在现代编程语言中很少用作定界符。
%
%
% \subsubsection*{如何修改列表清单的版式}
% \label{sec:tocbasic}
% 如果你把 |\lstlistoflistings| 命令放进文档，列表
% 目录就会在该处自动生成并打印。如果你
% 使用标准 \LaTeX{} 文档类，该目录由
% |\addtocontents| 命令或其变体准备（如果你还加载了
% \packagename{float} 宏包）。两种情况下的版式都是
% L. Lamport 定义的那种。
%
% 如果你想改变这个目录的版式，可以按 \cite[Chapter
% 15.3]{Kohm:scrguide} 所述使用 \packagename{tocbasic} 宏包。
% \packagename{tocbasic} 会在你使用 KOMA 文档类
% 之一时自动加载。
%
%
% \makeatletter
%^^A \def\index@prologue{\section*{Index}\markboth{Index}{Index}}
% \def\index@prologue{\part{索引}\markboth{索引}{索引}}
% \makeatother
%^^A \StopEventually{\lstcheckreference\setcounter{IndexColumns}{2}\PrintIndex}
% \StopEventually{%
% \nocite{MF23}
% \begin{thebibliography}{MDB01}
%
%     \bibitem[CMK12]{CovingtonMittelbachKuhn:upquote}
%       Michael A. Covington, Frank Mittelbach and Markus G. Kuhn.
%       \newblock{\packagename{upquote} -- upright-quote and grave-accent glyphs
%       in verbatim}, 2012.
%
%     \bibitem[DS13]{DanielSchubert:mdframed}
%       Marco Daniel and Elke Schubert.
%       \newblock{The \packagename{mdframed} package}, 2013.
%
%     \bibitem[Fai11]{Fairbairns:moreverb}
%       Robin Fairbairns.
%       \newblock{The \packagename{moreverb} package}, 2011.
%
%     \bibitem[Koh23]{Kohm:scrguide}
%       Markus Kohm.
%       \newblock{KOMA-Script -- The Guide}, 2023.
%
%     \bibitem[MF23]{MittelbachFischer:companion}
%       Frank Mittelbach and Ulrike Fischer.
%       \newblock{The \LaTeX-Companion} Part I and II, 2023.
%
%     \bibitem[Mi04]{Mittelbach:companion}
%       Frank Mittelbach, Michel Goossens, Johannes Braams,
%                  David P. Carlisle, and Chris Rowley.
%       \newblock{The \LaTeX-Companion}, 2004.
%
%     \bibitem[Som11]{Sommerfeldt:caption}
%       Axel Sommerfeldt.
%       \newblock{Customizing captions of floating environments}, 2011.
% \end{thebibliography}
% \setcounter{IndexColumns}{2}\PrintIndex}
%
%
% \part{开发者指南}
%
% 首先，我必须为这份开发者指南道歉，因为有些部分
% 没能解释得尽可能好。但请注意你的状况相当不错：
% 这份开发者指南存在！^^A :-)
% 在阅读第~\ref{dBasicConcepts} 节之前，你也许想
% 先浏览第~\ref{dPackageExtensions} 节。
%
%
% \section{基本概念}\label{dBasicConcepts}
%
% \packagename{listings} 宏包的功能似乎
% 分为两部分：一部分是实际排版
% 列表的命令，另一部分是通过 |\lstset| 可调的参数。二者都可以
% 用 \lst-方面（aspect）来实现，而这些方面不过是
% 公共键与命令以及内部钩子与定义的集合。宏包
% 定义了几个方面，尤其是内核，即主引擎。
% 其他方面驱动这个引擎，而语言和样式定义告诉
% 各方面如何驱动。车、司机和副
% 司机之间的关系被精确复现——而我将做你的驾驶教练。
%
%
% \subsection{加载宏包}\label{dPackageLoading}
%
% |\usepackage[|\meta{options}|]{listings}| 中的每个选项都会加载一个方面，或者
% 当方面名\emph{前面加感叹号}时\emph{阻止}宏包加载它。这一机制的设计目的
% 是理清不同宏包部件之间的依赖关系并便于
% 调试宏包。出于这个原因，还有一个选项：
% \begin{syntax}
% \item[0.21,,option] \texttt{noaspects}\leavevmode
%
%             清空要加载的方面列表。注意，例如选项
%             列表 |0.21,!labels,noaspects| 和 |noaspects| 实质上
%             是相同的：只加载内核，不加载其他方面。
% \end{syntax}
% 这对方面测试特别有用，因为我们可以只加载
% 所需的部分。但要注意，如果某个预定义
% 编程语言请求某个方面，该方面会在稍后被加载。也可以手动加载方面：
% \begin{syntax}
% \item[0.20] |\lstloadaspects|\marg{comma separated list of aspect names}
%
%             加载指定的方面（如果尚未加载）。
% \end{syntax}
% 下面是所有方面及相关键和命令的列表——希望这个
% 列表是完整的。
% \begin{description}
% \hyphenpenalty=10000\relax \rightskip=0pt plus \linewidth\relax
% \item[\aspectname{strings}]\leavevmode
%
%       \lstprintaspectkeysandcmds{strings}
%
% \item[\aspectname{comments}]\leavevmode
%
%       \lstprintaspectkeysandcmds{comments}
%
% \item[\aspectname{pod}]\leavevmode
%
%       \lstprintaspectkeysandcmds{pod}
%
% \item[\aspectname{escape}]\leavevmode
%
%       \lstprintaspectkeysandcmds{escape}
%
% \item[\aspectname{writefile}] 需要 1 个 |\toks|、1 个 |\write|
%
%       |\lst@BeginWriteFile|, |\lst@BeginAlsoWriteFile|, |\lst@EndWriteFile|
%
% \item[\aspectname{style}]\leavevmode
%
%             空样式，\lstprintaspectkeysandcmds{style}
%
% \item[\aspectname{language}]\leavevmode
%
%             空语言，\lstprintaspectkeysandcmds{language}
%
% \item[\aspectname{keywords}]\leavevmode
%
%       \lstprintaspectkeysandcmds{keywords}
%
% \item[\aspectname{emph}] 需要 \aspectname{keywords}
%
%       \lstprintaspectkeysandcmds{emph}
%
% \item[\aspectname{html}] 需要 \aspectname{keywords}
%
%       \lstprintaspectkeysandcmds{html}
%
% \item[\aspectname{tex}] 需要 \aspectname{keywords}
%
%       \lstprintaspectkeysandcmds{tex}
%
% \item[\aspectname{directives}] 需要 \aspectname{keywords}
%
%       \lstprintaspectkeysandcmds{directives}
%
% \item[\aspectname{index}] 需要 \aspectname{keywords}
%
%       \lstprintaspectkeysandcmds{index}
%
% \item[\aspectname{procnames}] 需要 \aspectname{keywords}
%
%       \lstprintaspectkeysandcmds{procnames}
%
% \item[\aspectname{keywordcomments}]
%             需要 \aspectname{keywords}、\aspectname{comments}
%
%       \lstprintaspectkeysandcmds{keywordcomments}
%
% \item[\aspectname{labels}] 需要 2 个 |\count|
%
%       \lstprintaspectkeysandcmds{labels}
%
% \item[\aspectname{lineshape}] 需要 2 个 |\dimen|
%
%       \lstprintaspectkeysandcmds{lineshape}
%
% \item[\aspectname{frames}] 需要 \aspectname{lineshape}
%
%       \lstprintaspectkeysandcmds{frames}
%
% \item[\aspectname{make}] 需要 \aspectname{keywords}
%
%       \lstprintaspectkeysandcmds{make}
%
% \item[\aspectname{doc}] 需要 \aspectname{writefile} 和 1 个 |\box|
%
%       \lstprintaspectkeysandcmds{doc}
%
% \item[\aspectname{0.21}] 用新键定义旧键。
% \item[\aspectname{fancyvrb}] 需要 1 个 |\box|
%
%       \lstprintaspectkeysandcmds{fancyvrb}
%
% \item[\aspectname{lgrind}]\leavevmode
%
%       \lstprintaspectkeysandcmds{lgrind}
%
% \item[\aspectname{hyper}] 需要 \aspectname{keywords}
%
%       \lstprintaspectkeysandcmds{hyper}
% \end{description}
% 内核分配 6 个 |\count|、4 个 |\dimen| 和 1 个 |\toks|。
% 此外它还定义下列键、命令和环境：
% \begin{itemize}\item[]
% \hyphenpenalty=10000\relax \rightskip=0pt plus \linewidth\relax
%       \lstprintaspectkeysandcmds{kernel}, \keyname{fancyvrb}
% \end{itemize}
%
%
% \subsection{如何定义 \lst-aspect}\label{dHowToDefineLstAspects}
%
% 至少有三种方式可以添加新功能：(a) 你编写一个普遍
% 感兴趣的方面，寄给我，我直接把它粘贴进
% 实现；(b) 你编写一个不普遍感兴趣的`局部'方面；或
% (c) 你有一个方面的想法，让我来写。(a) 和 (b) 是
% 好选择。^^A :-)
%
% 方面定义以 |\lst@BeginAspect| 加参数开始，以下一个 |\lst@EndAspect|
% 结束。特别地，方面定义不能
% 嵌套。
% \begin{syntax}
% \item[0.20] |\lst@BeginAspect|[\oarg{list of required aspects}]\marg{aspect name}
% \item[0.20] |\lst@EndAspect|
% \end{syntax}
% 可选列表是所需方面名的逗号分隔列表。
% 在下列每种情况下，完整方面不会被定义：
% \begin{enumerate}
% \item \meta{aspect name} 为空。
% \item 该方面已经定义。
% \item 某个所需方面既未定义，也无法通过
%       |\lstloadaspects|.
% \end{enumerate}
% 因此你不能先定义方面的一部分，之后再定义另一部分。
% 但可以先定义方面 $A_1$，之后再定义需要 $A_1$ 的
% 方面 $A_2$。
% \begin{advise}
% \item 把本地附加件放入 `\texttt{lstmisc0.sty}'——默认会先搜索
%             这个文件。如果只想为一个特定文档做附加件，
%             就把外层的 `|\lst@BeginAspect|' 和
%             `|\lst@EndAspect|' 换成 `|\makeatletter|' 和 `|\makeatother|'，并在
%             文档导言区使用这些定义。不过，所需方面必须由你自己
%             加载。
% \end{advise}
% 你可以在两条命令之间放任何 \TeX\ 材料，但注意
% 如果稍后还需要这些定义，它们必须是 |\global| 的——\LaTeX 的 |\newcommand|
% 产生局部定义，且前面不能加 |\global|。所以请使用下面
% 列出的命令、|\gdef| 以及后续各节描述的命令。
% \begin{syntax}
% \item[0.20] |\lst@UserCommand|\meta{macro}\meta{parameter text}\marg{replacement text}
%
%             该宏（主要）等价于 |\gdef|。目的是
%             区分用户命令和内部全局定义。
%
% \item[0.19] |\lst@Key|\marg{key name}\marg{init value}[\oarg{default value}]\marg{definition}
% \item[0.19] |\lst@Key|\marg{key name}|\relax|[\oarg{default value}]\marg{definition}
%
%             使用 David Carlisle 的 \packagename{keyval} 宏包定义一个键。
%             \meta{definition} 是一个带一个参数的宏的替换文本。
%             参数要么是 `key=value' 中的值，要么在未给出 `=value' 时
%             是 \meta{default value}。下面的辅助宏
%             |\lstKV@...| 或许能简化 \meta{definition}。
%
%             如果第二个参数是 |\relax|，该键不会被初始化。
%             否则 \meta{init value} 是赋给该键的初始值。
%             注意我们在局部切换到 |\globalsdefs=1|，以确保
%             初始化不受分组影响。
%
% \item[0.19] |\lst@AddToHook|\marg{name of hook}\marg{\TeX\ material}
%
%             在预定义点添加 \TeX\ 材料。第~\ref{dHooks} 节列出
%             所有钩子及其定义与执行的位置。
%             |\lst@AddToHook{A}{\csa}| 在 |\lst@AddToHook{A}{\csb}| 之前
%             \emph{不}保证 |\csa| 在 |\csb| 之前执行。
%
% \item[0.20] |\lst@AddToHookExe|\marg{name of hook}\marg{\TeX\ material}
%
%             还会为初始化执行 \meta{\TeX\ material}。你可以使用
%             局部变量——\TeX\ 意义上的局部和/或通常编程
%             语言意义上的局部——但当代码为初始化执行时，所有
%             赋值都是全局的：我们把 |\globaldefs| 局部设为一。
%
% \item[0.20] |\lst@UseHook|\marg{name of hook}
%
%             执行该钩子。
% \end{syntax}
% \begin{advise}
% \item 我们来看两个例子。第一个通过添加一些钩子材料来
%             扩展宏包。如果你想要状态消息，可以写
% \begin{verbatim}
%    \lst@AddToHook{Init}{\message{\MessageBreak Processing listing ...}}
%    \lst@AddToHook{DeInit}{\message{complete.\MessageBreak}}\end{verbatim}
%             第二个例子引入两个键，让用户控制这些
%             消息。宏 |\lst@AddTo| 在第
%       \ref{dGeneralPurposeMacros}.
% \begin{verbatim}
%   \lst@BeginAspect{message}
%   \lst@Key{message}{Annoying message.}{\gdef\lst@message{#1}}
%   \lst@Key{moremessage}\relax{\lst@AddTo\lst@message{\MessageBreak#1}}
%   \lst@AddToHook{Init}{\typeout{\MessageBreak\lst@message}}
%   \lst@EndAspect\end{verbatim}
%       不过，当然还有更有用的 aspect。
% \end{advise}
% 以下宏可用于 |\lst@Key| 命令的 \meta{definition} 参数中来求值参数。
% 额外的前缀 |KV|
% 指 \packagename{keyval} 宏包。
% \begin{syntax}
% \item[0.19] |\lstKV@SetIf|\marg{value}\meta{if macro}
%
%             \meta{if macro} 在 \meta{value} 的第一个字符等于
%             |t| 或 |T| 时变成 |\iftrue|。否则它变成 |\iffalse|。
%             通常你会用 |#1| 作为 \meta{value}。
%
% \item[1.0] \cs{lstKV@SwitchCases}\marg{value}\\
%   |{|\meta{string 1}|:|\meta{execute 1}|\\|\\
%   | |\meta{string 2}|:|\meta{execute 2}|\\|\\
%   \hbox to 3em{\hfill\vdots}\\
%   | |\meta{string $n$}|;|\meta{execute $n$}|}|\marg{else}
%
%             要么执行 \meta{else}，要么执行与 \meta{value} 匹配的部分。
%
%             本实现基于 \lsthelper{David~Carlisle}{2016/05/25}{%
%             http://tex.stackexchange.com/questions/311383/^^A
%                 adding-frame-none-to-lstlisting-options-in-lstdefinestyle^^A
%                 -gives-error-with-tabu} 的一个提示，以避免 C. Heinz
%             原始实现的问题\label{gAmpersand}:
%       \begin{quote}
%           \cs{lstKV@SwitchCases}\marg{value}\\
%           |{|\meta{string 1}|&|\meta{execute 1}|\\|\\
%           | |\meta{string 2}|&|\meta{execute 2}|\\|\\
%           \hbox to 3em{\hfill\vdots}\\
%           | |\meta{string $n$}|&|\meta{execute $n$}|}|\marg{else}
%       \end{quote}
%             当列表是 tabular 环境的一部分时就会出现这个问题，
%             这是由
%             \lsthelper{Nasser~M.~Abbasi}{2016/05/25}{listings 宏包中
%             可能的缺陷}发现的。
%
% \item[0.20] |\lstKV@TwoArg|\marg{value}\marg{subdefinition}
% \item[0.20] |\lstKV@ThreeArg|\marg{value}\marg{subdefinition}
% \item[0.20] |\lstKV@FourArg|\marg{value}\marg{subdefinition}
%
%             \meta{subdefinition} 是带两个、三个和四个参数的宏的
%             替换文本。我们用 \meta{value} 给出的参数
%             调用这个宏。如有必要会补上空参数。
%
% \item[0.19] |\lstKV@OptArg|\oarg{default arg.}\marg{value}\marg{subdefinition}
%
%             |[|\meta{default arg.}|]| \emph{不是}可选的。\meta{subdefinition}
%             是参数文本为 |[##1]##2| 的宏的替换文本。
%             注意宏参数字符 |#| 需要双写，因为它用在
%             另一个宏内部。\meta{subdefinition} 通过 |##1| 和 |##2| 访问
%             这些参数。
%
%             \meta{value} 通常是 \packagename{keyval} 宏包传入的
%             参数 |#1|。如果 \meta{value} 没有可选参数，
%             \meta{default arg.} 会被插入，以便向
%       \meta{subdefinition}.
%
% \item[0.21] |\lstKV@XOptArg|\oarg{default arg.}\marg{value}\meta{submacro}
%
%             与 |\lstKV@OptArg| 相同，但第三个参数 \meta{submacro}
%             已经是一个定义，而不是替换文本。
%
% \item[0.20] |\lstKV@CSTwoArg|\marg{value}\marg{subdefinition}
%
%                         \meta{value} 是一个由一个或两个参数组成的逗号分隔列表。
%             它们被传给子定义，该子定义是
%             带两个参数的宏的替换文本。如有必要会补上空的第二个
%             参数。
% \end{syntax}
% \begin{advise}
% \item 再看一个例子。键 `\keyname{sensitive}' 属于方面
%             \aspectname{keywords}。因此它定义在
%             `|\lst@BeginAspect{keywords}|' 和 `|\lst@EndAspect|' 之间，这里没有
%             展示出来。
% \begin{verbatim}
%   \lst@Key{sensitive}\relax[t]{\lstKV@SetIf{#1}\lst@ifsensitive}
%   \lst@AddToHookExe{SetLanguage}{\let\lst@ifsensitive\iftrue}\end{verbatim}
%             最后一行等价于
% \begin{verbatim}
%   \lst@AddToHook{SetLanguage}{\let\lst@ifsensitive\iftrue}
%   \global\let\lst@ifsensitive\iftrue\end{verbatim}
%             我们全局初始化这个变量，因为用户可能在分组内
%             请求某个方面。此后该变量局部使用——\meta{\TeX\ material} 中
%             没有 |\global|。注意我们可以如下定义并
%             初始化该键：
% \begin{verbatim}
%   \lst@Key{sensitive}t[t]{\lstKV@SetIf{#1}\lst@ifsensitive}
%   \lst@AddToHook{SetLanguage}{\let\lst@ifsensitive\iftrue}\end{verbatim}
%\end{advise}
%
%
% \subsection{内部模式}\label{dInternalModes}
%
% 你大概知道 \TeX 的条件命令 |\ifhmode|、|\ifvmode|、
% |\ifmmode| 和 |\ifinner|。它们告诉你 \TeX\ 处于（受限）
% 水平或（内部）垂直或（非显示）数学模式。例如，
% |\ifhmode| 为真且 |\ifinner| 为真表示受限水平
% 模式，也就是说你在一个 |\hbox| 里。典型用户不关心
% 这些模式；\TeX/\LaTeX\ 会管理这一切。但既然你在读
% 开发者指南，我们现在就来讨论 \packagename{listings}
% 宏包的类似机制。它用模式来区分注释与字符串、`注释
% 行'与`单行注释'，等等。
%
% 宏包在读入源代码之前处于`无模式'。在
% 初始化阶段它进入`处理模式'。此后模式取决于
% 实际的源代码。例如，考虑下面的行
% \begin{verbatim}
%    "string" // comment\end{verbatim}
% 并假定 \texttt{language=C++}。读到字符串定界符时，宏包
% 进入`字符串模式'并处理该字符串。匹配的结束定界符
% 离开该模式，即切回通用的`处理模式'。遇到
% 这两个斜杠时，宏包检测到注释行；因此它进入
% `注释行模式'并输出这些斜杠。通常这个模式持续到
% 行尾。
%
% 但在 \texttt{textcl=true} 时，\aspectname{escape} 方面会立即
% 离开`注释行模式'，中断当前模式序列，并进入
% `\TeX\ 注释行模式'。在行尾我们重新进入先前的模式
% 序列`无模式' $\to$ `处理模式'。这个到 \LaTeX 的转义可以
% 工作，因为`无模式'意味着 \TeX 的字符和 catcode 在场，
% 而`处理模式'意味着 \packagename{listings} 的字符和
% catcode 处于活动状态。
%
% \begin{table}[htbp]
% \caption{内部模式}\label{dDefinedInternalModes}
% \def\lsttabspace{\hspace*{1em}\hfill}
% \begin{tabular}{@{}lp{0.56\linewidth}@{}}
% 方面\lsttabspace\meta{模式名} & 用途/我们正在处理……\\
% \noalign{\smallskip}
% 内核\lsttabspace |\lst@nomode| &
%             若该模式处于活动状态，\TeX 的`字符表'在场；反向
%             蕴含不成立。任何其他模式\emph{可能}意味着 catcode
%             和/或字符定义被改变。
% \\
%       \lsttabspace |\lst@Pmode| &
%             是一个通用处理模式。若处于活动状态，我们正在处理列表，
%             但尚未进入更特殊的模式。
% \\
%       \lsttabspace |\lst@GPmode| &
%             语言定义的通用模式。
% \\
% \aspectname{pod}\lsttabspace |\lst@PODmode| &
%             ……一个 POD——Perl 专用。
% \\
% \aspectname{escape}\lsttabspace |\lst@TeXLmode| &
%             ……一个注释行，但 \TeX 的字符表在场——除了
%             行尾字符，它用于终止该模式。
% \\
%       \lsttabspace |\lst@TeXmode| &
%             表示 \TeX 的字符表在场（除了一个用户
%             指定的字符，它用于终止该模式）。
% \\
% \aspectname{directives}\lsttabspace |\lst@CDmode| &
%             表示当前行以编译器指令开始。
% \\
% \aspectname{keywordcomments}\lsttabspace |\lst@KCmode| &
%             ……一个关键词注释。
% \\
%       \lsttabspace |\lst@KCSmode| &
%             ……一个只能由分号终止的关键词注释。
% \\
%   \aspectname{html}\lsttabspace |\lst@htmlmode| &
%             若我们处于 \texttt{<} 和 \texttt{>} 之间则为活动。
%
%             以前的文档在此处使用了模式名 |\lst@insidemode|，
%             但该名称没有被任何文档引用，因此从 v1.10c 起
%             被视为一处文档疏漏。
% \\
% \aspectname{make}\lsttabspace |\lst@makemode| &
%             用于表示一个关键词。
% \end{tabular}
% \end{table}
% 表~\ref{dDefinedInternalModes} 列出了所有静态模式及其
% 所属方面。大多数功能使用动态创建的模式号，
% 例如所有字符串和注释。每个方面只需在方面定义内
% 分配，即可定义自己的模式。
% \begin{syntax}
% \item[0.19] |\lst@NewMode|\meta{mode \textup(control sequence\textup)}
%
%             定义一个新的静态模式，即赋给
%             \meta{mode} 的一个非负整数。\meta{mode} 应有前缀 \texttt{lst@} 和
%             后缀 \texttt{mode}。
%
% \item[0.21] |\lst@UseDynamicMode|\marg{token\textup(s\textup)}
%
%             把一个动态模式号作为参数插入记号。
%
%             当方面被加载或定义时，不能用该宏取得模式号。
%             它只能在每个列表的初始化过程中使用，
%             例如在选出字符表时定义注释。
%
% \item[0.19,,changed] |\lst@EnterMode|\meta{mode}\marg{start tokens}
%
%             打开一个分组层次，进入该模式，并执行 \meta{start tokens}。
%
%             在 \meta{start tokens} 中使用 |\lst@modetrue| 可以禁止后续模式
%             改变——当然，离开模式除外。你必须自己测试
%             是否允许进入，见下文。
%
% \item[0.19] |\lst@LeaveMode|
%
%             当且仅当当前模式不是 |\lst@nomode| 时，通过关闭一个分组层次
%             返回先前的模式。你必须自己测试
%             是否允许离开某个模式，见下文。
%
%\iffalse
% \item[0.19] |\lst@LeaveAllModes|
%
%             返回到 |\lst@nomode|。
%             这算是一种应急宏，所以不要使用它！
%\fi
%
% \item[0.19] |\lst@InterruptModes|
% \item[0.19] |\lst@ReenterModes|
%
%             第一条命令返回到 |\lst@nomode|，但把当前模式
%             序列保存在一个特殊的栈上。之后第二个宏返回到
%             先前的模式。在这些命令之间，你可以进入任何想要的
%             模式。特别地，你可以中断模式、进入一些模式，再
%             说一次`中断模式'。那么两次重新进入会把你带回
%             第一次`中断模式'之前的位置。
%
%             记住 |\lst@nomode| 意味着 \TeX 的字符表
%             处于活动状态。
% \end{syntax}
% 有些变量显示处理的内部状态。你可以读取
% 它们，但\emph{禁止直接写入}。注意：|\lst@ifmode|
% 并\emph{没有}过时，因为该布尔值与当前模式之间没有关系。
% 会出现我们进入某个模式却没有把 |\lst@ifmode| 设为 true 的情况，
% 也会出现我们没有分配任何模式却把它设为 true 的情况！
% \begin{syntax}
% \item[0.18,,counter] |\lst@mode|
%
%             保存当前模式号。用 |\ifnum\lst@mode=|\meta{mode name}
%             来测试某个模式。不要直接修改这个计数器！
%
% \item[0.18,,boolean] |\lst@ifmode|
%
%             若该布尔值为 true，则不允许模式改变——离开当前模式
%             除外。用 |\lst@modetrue| 修改这个变量，但只在
%             \meta{start tokens} 中这样做。
%
% \item[1.0,,boolean] |\lst@ifLmode|
%
%             表示当前模式是否在行尾结束。
% \end{syntax}
%
%
% \subsection{钩子}\label{dHooks}
%
% 如果你想定义一个方面，会出现几个问题。
% 你应当和/或必须
%     (a) 找到（普遍感兴趣的）额外功能并实现它，
%     (b) 创建用户接口，并且
%     (c) 与 \packagename{listings} 宏包对接，即找到正确的
%             钩子并插入合适的 \TeX\ 材料。
% (a) 超出本开发者指南的范围。第~\ref{dHowToDefineLstAspects} 节中的命令 |\lstKV@...|
% 可能有助于 (b)。下面我们
% 描述 \packagename{listings} 宏包的所有钩子。
%
% 所有钩子都在一个总分组内执行。这个分组从
% 某处开始，到每个列表的末尾某处结束。不要对分组作
% 任何其他假设。因此如有必要请全局定义变量——
% 如果不使用自己的分组，就要警觉副作用。
% \begin{syntax}
% \item \hookname{AfterBeginComment}
%
%             在宏包进入注释模式之后执行。
%             调用该钩子时，起始定界符通常已被排版。
%
% \item \hookname{BoxUnsafe}
%
%             包含所有用于停用在 |\hbox| 内可能不安全的命令和寄存器的
%             材料。每当宏包在列表外做盒子和为了 \packagename{fancyvrb} 支持时
%             都会用到它。
%
% \item \hookname{DeInit}
%
%             在列表的最末尾调用，但在从 \hookname{BoxUnsafe} 关闭盒子
%             或结束浮动体之前。
%
% \item \hookname{DetectKeywords}
%
%             这个 \hookname{Output} 子钩子当且仅当允许模式改变时执行，
%             即当且仅当宏包没有在处理注释、
%             字符串等时——见第~\ref{dInternalModes} 节。
%
% \item \hookname{DisplayStyle}
%
%             停用/激活 displaystyle 列表的功能。
%
% \item \hookname{EmptyStyle}
%
%             为选`空'样式而执行——除非用户重新定义了
%             该样式。
%
% \item \hookname{EndGroup}
%
%             每当宏包关闭一个分组时执行，例如在注释或
%             字符串末尾。
%
% \item \hookname{EOL}
%
%             在每个\emph{输入}行末尾调用，就在
%       \hookname{InitVarsEOL}.
%
% \item \hookname{EveryLine}
%
%             在每个\emph{输出}行开头执行，即对于断开的行会
%             执行多次。该钩子不得改变水平或
%             垂直位置。
%
% \item \hookname{EveryPar}
%
%             当输出开始时，每个输入行执行一次。该钩子
%             不得改变水平或垂直位置。
%
%^^A \item \hookname{ExcludeDelims}
%^^A
%^^A       Executed by the \keyname{excludedelims} key before the excluded
%^^A       delimiters are determined.
%^^A
% \item \hookname{ExitVars}
%
%             在 \hookname{DeInit} 之前紧接着执行。
%
% \item \hookname{FontAdjust}
%
%             调整字体相关的内部值（目前只有 |\lst@width|）。
%
% \item \hookname{Init}
%
%             每个列表执行一次，在字符表被改变之前初始化。它在
%             \hookname{PreInit} 之后、在
%       \hookname{InitVars}.
%
% \item \hookname{InitVars}
%
%             为每个列表初始化变量而调用。
%
% \item \hookname{InitVarsBOL}
%
%             在每个输入行开头初始化变量。
%
% \item \hookname{InitVarsEOL}
%
%             在每个输入行末尾更新变量。
%
% \item \hookname{ModeTrue}
%
%             当模式改变变为非法时由宏包执行。
%             这里对注释和字符串关闭关键词检测。
%
% \item \hookname{OnEmptyLine}
%
%             在宏包输出空行\emph{之前}执行。
%
% \item \hookname{OnNewLine}
%
%             在宏包开始一个或多个新行之前执行，
%             即在说 |\par\noindent\hbox{}| 之前（大致地说）。
%
% \item \hookname{Output}
%
%             在打印标识符之前调用。
%             如果你想要特殊的打印样式，修改 |\lst@thestyle|。
%
% \item \hookname{OutputBox}
%
%             在每个输出盒内使用。目前它只用于让
%             宏包与 Lambda 协同工作——希望如此。
%
% \item \hookname{OutputOther}
%
%             在打印其他字符串之前调用。
%             如果你想要特殊的打印样式，修改 |\lst@thestyle|。
%
% \item \hookname{PostOutput}
%
%             在打印标识符或任何其他输出单元之后调用。
%
% \item \hookname{PostTrackKeywords}
%
%             是一个非常特殊的 \hookname{Init} 子钩子，用于插入关键词测试
%             并按需定义关键词。
%             该钩子在 \hookname{TrackKeywords} 之后调用。
%
% \item \hookname{PreInit}
%
%             就在 \hookname{Init} 钩子之前调用。
%
% \item \hookname{PreSet}
%
%             每个排版命令/环境都调用该钩子，以在任何用户
%             提供的键被设置之前初始化内部状态。
%
% \item \hookname{SelectCharTable}
%
%             在宏包选出标准字符表之后执行。
%             各方面在这里调整字符表，并定义字符串和
%             注释定界符等。
%
% \item \hookname{SetFormat}
%
%             在设置格式的内部赋值之前调用。
%             该钩子决定每次选择格式时重置哪些参数。
%
% \item \hookname{SetStyle}
%
%             在设置样式的内部赋值之前调用。
%             该钩子决定每次选择样式时重置哪些参数。
%
% \item \hookname{SetLanguage}
%
%             在设置语言的内部赋值之前调用。
%             该钩子决定每次选择语言时重置哪些
%             参数。
%
% \item \hookname{TextStyle}
%
%             停用/激活 textstyle 列表的功能。
%
% \item \hookname{TrackKeywords}
%
%             是一个非常特殊的 \hookname{Init} 子钩子，用于插入关键词测试并
%             按需定义关键词。
%             该钩子在 \hookname{PostTrackKeywords} 之前调用。
% \end{syntax}
%
%
% \subsection{字符表}\label{dCharacterTables}
%
% 现在你知道汽车长什么样了，如果练习
% 一下，你就能拿到驾照。但如果你想对汽车做大
% 改动，就会遇到困难。所以让我们仔细观察，进入最
% 困难的部分：引擎。我们先看大局，然后
% 一步步填补细节。对我们来说，覆盖 \TeX 的
% 字符表是合适的。首先我们定义一个标准字符表，其中包含
% \begin{itemize}
% \item 字母：标识符由这些字符组成，
% \item 数字：用于标识符或数值常量的字符，
% \item 空格：被视为空白的字符，
% \item 制表符：被视为制表符的字符，
% \item 换页符：被视为换页字符的字符，以及
% \item 其他：所有其他字符。
% \end{itemize}
% 这个字符表会随当前编程语言而改变。
% 我们可以定义字符串和注释定界符或其他特殊字符。
% 第~\pageref{rStdCharTable} 页的表~\ref{rStdCharTable} 展示了标准
% 字符表。它可以用键 \keyname{alsoletter}、
% \keyname{alsodigit} 和 \keyname{alsoother} 修改。
%
% 这些`类'如何协同工作？假设当前字符
% 串是 `|tr|'。那么字母 `|y|' 只是追加这个字母，我们得到
% `|try|'。下一个非字母（且非数字）字符会触发
% 这些字符的输出。然后我们收集随后所有非字母字符，直到再次
% 遇到字母。这又触发非字母的输出，依此类推。内部每个
% 字符都按 \TeX 意义变成活动字符，并被定义为做正确
% 的事，例如我们说
% \begin{verbatim}
%    \def A{\lst@ProcessLetter A}\end{verbatim}
% 其中第一个 `|A|' 是活动的，第二个具有字母 catcode 11。
% 宏 |\lst@ProcessLetter| 取一个记号并把它当字母处理。
% 存在下列宏，其中最后三个不取显式参数。
% \begin{syntax}
% \item[0.18] |\lst@ProcessLetter| \meta{spec.\ token}
% \item[0.18] |\lst@ProcessDigit| \meta{spec.\ token}
% \item[0.18] |\lst@ProcessOther| \meta{spec.\ token}
% \item[0.18] |\lst@ProcessTabulator|
% \item[0.18] |\lst@ProcessSpace|
% \item[0.20] |\lst@ProcessFormFeed|
% \end{syntax}
% \meta{spec.\ token} 应当做两件事。通常它展开为
% 字符的可打印版本。但如果 |\lst@UM| 等价于
% |\@empty|，则 \meta{spec.\ token} 必须展开为一个\emph{字符记号}。
% 例如，井号通常展开为 |\#|，它是通过
% |\chardef| 定义，不是字符记号。但如果 |\lst@UM| 等价于
% |\@empty|，井号展开为字符 `|#|'（catcode 12）。注意：
% 对 |\lst@UM| 的\emph{改动必须是局部的}。不过，你应当
% 不需要自己做这类基础工作。\packagename{listings}
% 宏包提供了使用该特性的高级宏，
% 例如第~\ref{dKeywordsAndWorkingIdentifiers} 节中的 |\lst@InstallKeywords|。
%
% \begin{syntax}
% \item[0.18] |\lst@Def|\marg{character code}\meta{parameter text}\marg{definition}
% \item[0.18] |\lst@Let|\marg{character code}\meta{token}
%
%             定义指定的字符，或分别把 \meta{token} 赋给它。
%             不影响 catcode 表。如果你的定义带参数，要小心：
%             向前多读一个字符是不安全的。
%             此外，参数可以是\emph{任意的}；有时它是下一个
%             源代码字符，有时是 \packagename{listings} 宏包的某个代码，
%             例如 |\relax|、|\@empty|、|\else|、
%             |\fi| 等。因此不要对这样的参数使用 \TeX 的 ord 运算符 |`|，
%             例如不要写 |\ifnum`#1=65| 来测试 `|A|'。
%
%             |\lst@Def| 和 |\lst@Let| 相对较慢。标准字符表的真正定义
%             与下面的例子不同，但它
%             可以以如下内容开始
% \begin{verbatim}
%    \lst@Def{9}{\lst@ProcessTabulator}
%    \lst@Def{32}{\lst@ProcessSpace}
%    \lst@Def{48}{\lst@ProcessDigit 0}
%    \lst@Def{65}{\lst@ProcessLetter A}\end{verbatim}
%
%\iffalse
% \item[0.20] |\lst@activecharstrue|
% \item[0.20] |\lst@activecharsfalse|
%
%             控制选择字符表时是否同时让所有字符
%             变活动（标准/扩展）。这通常为真，因此是默认值。
%             只有 \packagename{fancyvrb} 接口会在局部把它设为 false。
%\fi
% \end{syntax}
% 目前这些就够了。第~\ref{dUsefulInternalDefinitions} 节
% 介绍操作字符表的高级定义，特别是
% 如何添加新的注释或字符串类型。
%
%
% \subsection{关于输出}
%
% \packagename{listings} 宏包用一些变量保存输出
% 数据。不建议写入这些变量。先从简单的说起。
% \begin{syntax}
% \item[0.17,,data] |\lst@lastother|
%
%             等于最后处理过的非标识符字符的 \meta{spec.\ token} 版本。
%             由于编程语言会重新定义标准
%             字符表，我们使用原始的 \meta{spec.\ token}。
%             例如，如果最后处理的是双引号，|\lst@lastother|
%             并不等价于进入和离开字符串模式的宏。
%             它等价于 |\lstum@"|，其中 |"| 属于该控制
%             序列。记住 \meta{spec.\ token} 要么展开为
%             可打印内容，要么展开为字符记号。
%
%             如果这样的字符不可用（例如在行首），|\lst@lastother| 等价于 |\@empty|。
%             有时在处理最后一个`其他'字符之后已经打印了标识符，
%             也就是说该字符已经离得很远。这种情况下 |\lst@lastother|
%             等于 |\relax|。
%
% \item[0.17] |\lst@outputspace|
%
%             用这个预定义的 \meta{spec.\ token}（显然对应字符码
%             32）来测试 |\lst@lastother|。
%
% \item[0.20] |\lstum@backslash|
%
%             用这个预定义的 \meta{spec.\ token}（对应字符码 92）来测试
%             |\lst@lastother|。在 |\lst@Def| 的替换文本中可以写
%             |\ifx| |\lst@lastother| |\lstum@backslash| \ldots\ 来测试
%             最后一个字符是否为反斜杠。
%
% \item[0.20] |\lst@SaveOutputDef|\marg{character code}\meta{macro}
%
%             把与 \meta{character code} 对应的 \meta{spec.\ token}
%             存入 \meta{macro}。这是取得可用于测试 |\lst@lastother| 的
%             正确含义的唯一安全方式，例如
%           |\lst@SaveOutputDef{"5C}\lstum@backslash|.
%
%             如果 \meta{character code} 不在 33 到 126（含）之间，
%             你会得到一个``runaway argument''错误。
% \end{syntax}
% 现在来看更多处理输出数据和
% 状态的宏。
% \begin{syntax}
% \item[1.0] |\lst@XPrintToken|
%
%             输出当前字符串并重置它。该宏会跟踪
%             这里描述的所有变量。
%
% \item[0.18,,token] |\lst@token|
%
%             包含当前字符串。每个`字符'通常
%             展开为它的可打印版本，但如果 |\lst@UM| 等价于
%             |\@empty|，它必须展开为字符记号。
%
% \item[0.12,,counter] |\lst@length|
%
%             是当前字符串的长度。
%
% \item[0.12,,dimension] |\lst@width|
%
%             是单个字符盒的宽度。
%
% \item[0.20,,global dimension] |\lst@currlwidth|
%
%             是到目前为止已打印行的宽度。
%
% \item[0.18,,global counter] |\lst@column|
% \item[0.12,,global counter] |\lst@pos|（非正）
%
%             |\lst@column|$-$|\lst@pos| 是到目前为止已打印行的长度。
%             我们用两个计数器，因为这简化了制表符处理：
%             |\lst@pos| 是`已打印行长度'模 \keyname{tabsize} 的一个非正
%             代表值。它通常不是最大的非正代表值。
%
% \item[0.20] |\lst@CalcColumn|
%
%             |\@tempcnta| 得到 |\lst@column| $-$ |\lst@pos| $+$ |\lst@length|。
%             这是当前列号减一，即从零起算的当前列
%             号。
%
% \item[0.18,,global dimension] |\lst@lostspace|
%
%             等于`丢失的'空间：期望的当前行宽减去实际行宽。
%             只要该尺寸为正，灵活列格式就可以使用
%             这个空间来修正列对齐。
% \end{syntax}
%
%
% \section{宏包扩展}\label{dPackageExtensions}
%
%
% \subsection{关键词和工作标识符}\label{dKeywordsAndWorkingIdentifiers}
%
% \aspectname{keywords} 方面定义两个主要宏。它们的
% 语法分别示于左侧。右侧是宏包
% 实际定义某些键的例子。
% \begin{syntax}
%   \item[0.21] \cs{lst@InstallFamily}
%
%     \marg{prefix}\syntaxfill \texttt k\\
%     \marg{name}\syntaxfill |{keywords}|\\
%     \marg{style name}\syntaxfill |{keywordstyle}|\\
%     \marg{style init}\syntaxfill |\bfseries|\\
%     \marg{default style name}\syntaxfill |{keywordstyle}|\\
%     \marg{working procedure}\syntaxfill |{}|\\
%     \meta{\alternative{l,o}}\syntaxfill |l|\\
%     \meta{\alternative{d,o}}\syntaxfill |d|
%
%         根据 \meta{working procedure} 是否为空，安装关键词类或
%         `工作'标识符类。
%
%         会定义三个键 \meta{name}、\keyname{more}\meta{name} 和
%         \keyname{delete}\meta{name}，以及（若非空）\meta{style name}。
%         后者的第一阶成员若不等价于 |\relax|，会用
%         \meta{style init} 初始化。如果用户让某个
%         类别样式未定义，则改用 \meta{default style name}。
%         因此，请确保这个样式总是已定义。在例子中，
%         第一阶 keywordstyle 被设为 |\bfseries|，并作为
%         所有其他类别的默认值。
%
%         如果 \meta{working procedure} 不为空，则在
%         遇到这样的（用户定义的）标识符时执行该代码。\meta{working procedure}
%         恰好接受一个参数，即实际标识符
%         所属的类别号。如果代码使用变量并需要
%         先前调用的值，你必须把这些变量定义为
%         |\global| 的。并不确定工作过程是否在
%         （单独的）分组内执行。
%
%         \texttt l 表示语言键，即每次选择语言时
%         列表都会被重置。\texttt o 表示`其他'键。关键词
%         或工作测试安装在 \hookname{DetectKeyword}
%         或 \hookname{Output} 钩子上，取决于
%     \meta{\alternative{d,o}}.
%
%   \item[0.20] \cs{lst@InstallKeywords}
%
%     \marg{prefix}\syntaxfill \texttt{cs}\\
%     \marg{name}\syntaxfill |{texcs}|\\
%     \marg{style name}\syntaxfill |{texcsstyle}|\\
%     \marg{style init}\syntaxfill |\relax|\\
%     \marg{default style name}\syntaxfill |{keywordstyle}|\\
%     \marg{working procedure}\syntaxfill see below\\
%     \meta{\alternative{l,o}}\syntaxfill |l|\\
%     \meta{\alternative{d,o}}\syntaxfill |d|
%
%         参数相同，功能相同，只有一个例外。该宏
%         只安装一个关键词类，而不是整个族。因此
%         \meta{working procedure} 的参数是固定的（目前
%         为空）。
%
%         例子中的工作过程读作
%         如下。
% \begin{verbatim}
%    {\ifx\lst@lastother\lstum@backslash
%         \let\lst@thestyle\lst@texcsstyle
%     \fi}\end{verbatim}
%    % \vspace*{-.5\baselineskip}
% 这个过程做什么？首先，只有当找到用户
% 提供的列表（或语言定义）中的关键词时它才被调用。该
% 过程现在检查前面是否有反斜杠，并相应地
% 设置输出样式。
%
% \iffalse
% \item[0.20] |\lst@InstallTest|\marg{prefix}\syntaxbreak
%     |\lst@|\meta{name}|@list||\lst@|\meta{name}~|\lst@g|\meta{name}|@list|^^A
%   |\lst@g|\meta{name}\syntaxbreak
%     |\lst@g|\meta{name}|@sty|~\meta{\alternative{w,s}}\meta{\alternative{d,o}}
%
%     根据 \meta{\alternative{w,s}} 安装`工作标识符'测试或关键词样式。
%     |\lst@g|\meta{name}|@sty| 包含工作过程或样式宏。注意
%     这些测试的行为取决于 \texttt{savemem} 选项。
%     关键词或工作测试安装在 \hookname{DetectKeyword}
%     或 \hookname{Output} 钩子上，取决于 \meta{\alternative{d,o}}。
%   \hookname{Output} hook according to \meta{\alternative{d,o}}.
%
%   ^^A Either each call of this macro or each different \meta{prefix}
%     插入 ^^A 另一个测试（取决于 \texttt{savemem} 选项）。
%
%     |\lst@|\meta{name} 包含当前标识符列表，
%     |\lst@|\meta{name}|@list| 是当前使用类别的
%     `|\lst@|\meta{$n_i$}|\lst@g|\meta{$n_i$}' 序列。如果没有使用其他类别，它
%     等于 |\lst@|\meta{name}|\lst@g|\meta{name}。全局版本
%     |\lst@g|\ldots\ 用于跟踪关键词。（这段
%     描述还需要改进。）
% \fi
% \end{syntax}
%
%
% \subsection{定界符}
%
% 我们分两个阶段描述：给现有定界符类添加新定界符类型，
% 以及编写新类。每个类都有自己的名字；目前有
% \texttt{Comment}、\texttt{String} 和 \texttt{Delim}。如你所知，
% 后者和前者都提供类型 \texttt l，但没有以给定定界符开始、到行尾
% 结束的字符串。所以我们现在就来添加它！
%
% 首先，我们用下列内容扩展字符串类型列表
% \begin{verbatim}
%    \lst@AddTo\lst@stringtypes{,l}\end{verbatim}
% 然后我们必须提供接受用户给定定界符并
% 做相应定义的宏。命令名由前缀
% |\lst@|、定界符名、用于动态模式的 |DM| 组成，后面再接
% |@| 和类型。
% \begin{verbatim}
%    \gdef\lst@StringDM@l#1#2\@empty#3#4#5{%
%        \lst@CArg #2\relax\lst@DefDelimB{}{}{}#3{#1}{#5\lst@Lmodetrue}}\end{verbatim}
% 你可以把这三行放进一个 \texttt{.sty} 文件中，或者在文档导言区用
% |\makeatletter| 和 |\makeatother| 包围它们。
% 就是这样！
%{\makeatletter
%\lst@AddTo\lst@stringtypes{,l}
%\gdef\lst@StringDM@l#1#2\@empty#3#4#5{^^A
%   \lst@CArg #2\relax\lst@DefDelimB{}{}{}#3{#1}{#5\lst@Lmodetrue}}
%}
% \begin{lstsample}{}{}
%    \lstset{string=[l]//}
%    \begin{lstlisting}
%    // This is a string.
%    This isn't a string.
%    \end{lstlisting}
% \end{lstsample}
% 当然，你想要更多细节。先从参数说起。
% \begin{itemize}
% \item 第一个参数（\emph{在} |\@empty| 之后）用于开始
%             定界符。它由定界符类提供。
% \item 第二个参数（\emph{在} |\@empty| 之后）用于结束
%             定界符。它同样由定界符类提供。例子中我们
%             不需要它，见下面的说明。
% \item 第三个参数（\emph{在} |\@empty| 之后）是
%       \marg{style}\meta{start tokens}.
%             它前面加上 |\def\lst@currstyle| 后作为参数传给
%             |\lst@EnterMode|。定界符类也会提供它。例子中我们
%             用 |\lst@Lmodetrue|（行模式为真）`扩展'了 |#5|。
%             该模式会在行尾自动结束，所以我们不需要
%             结束定界符参数。
% \end{itemize}
% 现在看其他参数。在动态模式的情况下，第一个参数
% 是模式号。然后是用户提供的定界符，其数量
% 必须与直到 |\@empty| 的其余参数匹配。对于非动态模式，
% 你必须自己分配一个静态模式，或者使用预定义的模式
% 号。此时定界符从第一个参数开始。
%
% 最后看看宏的替换文本。序列
% |\lst@CArg #2\relax| 在 |\lst@DefDelimB| 之后放两个必需参数。
% 后一个宏的语法是
% \begin{syntax}
% \item[1.0] \cs{lst@DefDelimB}
%
%   |{|\meta{1st}\meta{2nd}\marg{rest}|}|\syntaxfill |{//{}}|\\
%   \meta{save 1st}\syntaxfill |\lst@c/0|\\
%   \marg{execute}\syntaxfill|{}|\\
%   \marg{delim~exe~modetrue}\syntaxfill|{}|\\
%   \marg{delim~exe~modefalse}\syntaxfill|{}|\\
%   \meta{start-delimiter macro}\syntaxfill|#3|\\
%   \meta{mode number}\syntaxfill |{#1}|\\
%   |{|\marg{style}\meta{start tokens}|}|\syntaxfill |{#5\lst@Lmodetrue}|
%
%             把 \meta{1st}\meta{2nd}\meta{rest} 定义为起始定界符。
%             当宏包遇到 \meta{1st} 时执行 \meta{execute}。
%             只有当整个定界符 \meta{1st}\meta{2nd}\meta{rest}
%             被找到时才执行 \meta{delim~exe~modetrue} 和 \meta{delim~exe~modefalse}。
%             依据 |\lst@ifmode| 恰好调用其中一个。
%
%             默认情况下，如果找到定界符\emph{并且} |\lst@ifmode| 为假，
%             宏包就进入该模式。内部我们对 |\lst@bnext|
%             做相应的定义，可以通过把
%             |\@gobblethree| 放在 \meta{delim exe modefalse} 的最末尾来吃掉它。
%             也可以提供自己的定义（并吃掉默认的）。
%
%             \meta{save 1st} 必须是一个未定义的宏，内部用于
%             保存 \meta{1st} 先前的含义。如果定界符严格少于
%             三个字符，参数 \meta{2nd} 和/或 \meta{rest} 为空。
%             \meta{1st}\meta{2nd}\meta{rest} 的所有字符
%             必须已经是活动的（若非空）。这不成问题，
%             因为宏 |\lst@CArgX| 负责这件事。
%
% \item[1.0] \cs{lst@DefDelimE}
%
%   |{|\meta{1st}\meta{2nd}\marg{rest}|}|\\
%   \meta{save 1st}\\
%   \marg{execute}\\
%   \marg{delim~exe~modetrue}\\
%   \marg{delim~exe~modefalse}\\
%   \meta{end-delimiter macro}\\
%   \meta{mode number}
%
%             结束定界符同理，但有细微差别：
%             \meta{delim~exe~modetrue} 和 \meta{delim~exe~modefalse} 的执行
%             取决于 |\lst@mode| 是否等于 \meta{mode}。
%
%             宏包在找到定界符且 |\lst@mode|
%             等于 \meta{mode} 时结束该模式。内部我们对
%             |\lst@enext|（不是 |\lst@bnext|）做相应的定义，可以通过把
%             |\@gobblethree| 放在 \meta{delim exe modetrue} 的最末尾来吃掉它。
%
% \item[1.0] \cs{lst@DefDelimBE}
%
%     后面跟与 |\lst@DefDelimB| 相同的八个参数，以及\ldots\\
%   \meta{end-delimiter macro}
%
%             这是 |\lst@DefDelimB| 与 |\lst@DefDelimE| 的组合，用于
%             起始和结束定界符相同的情况。
% \end{syntax}
% 我们通过两个简单例子结束第一阶段。
% \texttt d 类型字符串由下列方式定义
% \begin{verbatim}
%    \gdef\lst@StringDM@d#1#2\@empty#3#4#5{%
%        \lst@CArg #2\relax\lst@DefDelimBE{}{}{}#3{#1}{#5}#4}\end{verbatim}
% （以及字符串类型列表中的一个条目）。
% 没什么难的。\texttt d 类型注释同理：
% \begin{verbatim}
%    \gdef\lst@CommentDM@s#1#2#3\@empty#4#5#6{%
%        \lst@CArg #2\relax\lst@DefDelimB{}{}{}#4{#1}{#6}%
%        \lst@CArg #3\relax\lst@DefDelimE{}{}{}#5{#1}}\end{verbatim}
% 这里我们只需同时使用 |\lst@DefDelimB| 和 |\lst@DefDelimE|。
% \goodbreak
%
% 现在进入第二阶段。为说明起见，下面是
% \texttt{Delim} 类的定义。传给服务宏的第一个参数
% 分别使它删除该类的所有定界符、添加该定界符，或只删除
% 该特定定界符。
% \begin{verbatim}
%    \lst@Key{delim}\relax{\lst@DelimKey\@empty{#1}}
%    \lst@Key{moredelim}\relax{\lst@DelimKey\relax{#1}}
%    \lst@Key{deletedelim}\relax{\lst@DelimKey\@nil{#1}}\end{verbatim}
% 服务宏本身会用相应参数调用另一个宏。
% \begin{verbatim}
%    \gdef\lst@DelimKey#1#2{%
%        \lst@Delim{}#2\relax{Delim}\lst@delimtypes #1%
%                    {\lst@BeginDelim\lst@EndDelim}
%            i\@empty{\lst@BeginIDelim\lst@EndIDelim}}\end{verbatim}
% 我们必须看看这些参数。上面是 \texttt{Delim} 类的实际参数，
% 下面是 \texttt{Comment} 类的参数。
% 注意用户提供的值占了参数的第二行和第三
% 行。
% \begin{syntax}
%     \item[0.21,,changed] \cs{lst@Delim}
%
%     \meta{default style macro}\syntaxfill \cs{lst@commentstyle}\\ \relax
%     [\texttt*[\texttt*]]\texttt[\meta{type}\texttt][\texttt[\meta{style}^^A
%     \texttt][\texttt[\meta{type option}\texttt]]]\\
%     \meta{delimiter\textup(s\textup)}\cs{relax}\syntaxfill|#2|\cs{relax}\\
%     \marg{delimiter name}\syntaxfill|{Comment}|\\
%     \meta{delimiter types macro}\syntaxfill\texttt{\cs{lst@commenttypes}}\\
%     \alternative{\cs{@empty},\cs{@nil},\cs{relax}}\syntaxfill|#1|\\
%     \marg{begin- and end-delim macro}\syntaxfill|{|\cs{lst@BeginComment}^^A
%     \cs{lst@EndComment}|}|\\
%     \meta{extra prefix}\syntaxfill |i|\\
%     \meta{extra conversion}\syntaxfill |\@empty|\\
%     \marg{begin- and end-delim macro}\syntaxfill|{|\cs{lst@BeginIComment}^^A
%     \cs{lst@EndIComment}|}|
%
%         大多数参数应该很清楚。我们讨论后四个。两个
%         \marg{begin- and end-delim macro} 必须恰好包含两个控制
%         序列，它们被传给 |\lst@|\meta{name}[|DM|]|@|\meta{type}
%         以开始和结束一个定界符。它们就是 |\lst@StringDM@l| 第一个例子中的
%         参数 |#3| 和 |#4|。取决于用户选择的类型是否以
%         \meta{extra prefix} 开头，会使用前两个
%         或后两个控制序列。
%
%         默认情况下，宏包取定界符、把字符
%         变成活动字符，并把它们放在
%         |\lst@|\meta{name}[|DM|]|@|\meta{type} 之后。如果用户类型以
%         \meta{extra prefix} 开头，\meta{extra conversion} 可能会改变
%         |\lst@next| 的定义，以选择另一种转换。
%         默认等价于带 |\lst@false| 的 |\lst@XConvert|。
%
%         注意 \meta{type} 从不会以 \meta{extra prefix} 开头，因为它
%         会被丢弃。该功能必须通过选择一对不同的
%         \marg{begin- and end-delim macro} 来完整实现。
% \end{syntax}
% 你可能需要了解 \meta{begin- and end-delim macro} 的语法。
% 它们的调用方式如下。
% \begin{syntax}
% \item[0.21] |\lst@Begin|\meta{whatever}
%
%   \marg{mode}
%   |{|\marg{style}\meta{start tokens}|}|
%   \meta{delimiter}|\@empty|
%
% \item[0.21] |\lst@End|\meta{whatever}
%
%   \marg{mode}
%   \meta{delimiter}|\@empty|
% \end{syntax}
% 现有的宏在内部是由 |\lst@DelimOpen| 和
% |\lst@DelimClose| 定义的，见实现部分。
%
%
% \subsection{让内核运行起来}
%
% 如果你想要新的美化排版环境，应当对
% 第~\ref{rEnvironments} 节感到满意。像 |\lstinline| 或
% |\lstinputlisting| 这样的新命令更难一些。大致说来，你必须遵循
% 这些步骤。
% \begin{enumerate}
% \item 打开一个分组，使所有改动都是局部的。
% \item \meta{Do whatever you want.}
% \item 无论如何都要调用 |\lsthk@PreSet|。
% \item 现在你\emph{可能}想（但不必）用 |\lstset| 设置一些
%             新值。
% \item \meta{Do whatever you want.}
% \item 执行 |\lst@Init\relax| 完成初始化。
% \item \meta{Do whatever you want.}
% \item 最后是源代码，由内核处理。
%             你必须确保字符要么尚未被读入，要么全部
%             是活动的。此外，\emph{你}必须安装一种检测源代码
%             结束的方法。如果已经到达末尾，你必须\ldots
% \item \ldots\ 调用 |\lst@DeInit| 以安全地关闭内核。
% \item \meta{Do whatever you want.}
% \item 关闭最初打开的分组。
% \end{enumerate}
% 例如，考虑不在参数内部时的 |\lstinline| 命令。
% 此时各步骤如下。
% \begin{enumerate}
% \item |\leavevmode\bgroup| 打开一个分组。
% \item |\def\lst@boxpos{b}| 让列表按`基线'对齐。
% \item |\lsthk@PreSet|
% \item |\lstset{flexiblecolumns,#1}|（|#1| 是用户提供的
%             key=value 列表）
% \item |\lsthk@TextStyle| 停用所有在这里不安全的特性。
% \item |\lst@Init\relax|
% \item |\lst@Def{`#1}{\lst@DeInit\egroup}| 安装`结束行内'
%             检测，其中 |#1| 是 |\lstinline| 之后的那个字符。
%             此外 |chr(13)| 被重定义为以同样方式结束片段，但
%             还会发出错误消息。
% \item 现在是源代码，并且\ldots
% \item \ldots\ |\lst@DeInit|（来自上面的 |\lst@Def|）正确地结束代码片段。
% \item 无事可做。
% \item |\egroup|（同样来自 |\lst@Def|）关闭分组。
% \end{enumerate}
% 真正的定义有所不同，因为我们允许参数内有源代码。
% 如果你真想写美化排版命令，也请阅读第~\ref{iTheInputCommand} 节。
%
%
% \section{有用的内部定义}\label{dUsefulInternalDefinitions}
%
% 本节需要更新。
%
%
% \subsection{通用宏}\label{dGeneralPurposeMacros}
%
% \begin{syntax}
%   \item[0.19] |\lst@AddTo|\meta{macro}\marg{\TeX~material}
%
%         全局地把 \meta{\TeX~material} 加到 \meta{macro} 的内容中。
%
%   \item[0.20] |\lst@Extend|\meta{macro}\marg{\TeX~material}
%
%         在 \meta{\TeX~material} 的第一个记号被
%         |\expand| 之|后|调用 |\lst@AddTo|。例如，|\lst@Extend \a \b| 合并
%         两个宏的内容，并把它全局存到 |\a| 中。
%
%   \item[0.19] |\lst@lAddTo|\meta{macro}\marg{\TeX~material}
%   \item[0.20] |\lst@lExtend|\meta{macro}\marg{\TeX~material}
%
%         是 |\lst@AddTo| 和 |\lst@Extend| 的局部版本。
%
%     \item[0.18] |\lst@DeleteKeysIn|\meta{macro}\meta{macro \textup(要删除的键
%         \textup)}
%
%         两个宏都包含逗号分隔的键（或关键词）列表。第二个宏中
%         出现的所有键都会（局部地）从
%         第一个宏中删除。
%
%     \item[0.19] |\lst@ReplaceIn|\meta{macro}\meta{macro \textup(含
%     replacement list\textup)}
%   \item[0.20] |\lst@ReplaceInArg|\meta{macro}\marg{replacement list}
%
%         替换列表的形式为
%         $a_1b_1$\ldots$a_nb_n$，其中每个 $a_i$ 和
%         $b_i$ 都是字符序列（必要时用花括号括起），
%         可以包含宏，但 $b_i$ 的第一个记号
%         不能等价于 |\@empty|。第一个宏内的每个序列
%         $a_i$ 都会被（局部地）替换为
%         $b_i$。后缀 |Arg| 指带花括号的\emph{第二参数}，
%         而不是（不带花括号的）宏。这提示我们取得的是`真正的'
%         参数，而不是指向参数的`指针'。
%
%     \item[0.20] |\lst@IfSubstring|\marg{字符
%     sequence}\meta{macro}\marg{then}\marg{else}
%
%         如果 \meta{character sequence} 是 \meta{macro} 内容的子串，
%         就执行 \meta{then}。否则调用 \meta{else}。
%
%     \item[0.12] |\lst@IfOneOf|\meta{字符
%     sequence}|\relax|\meta{macro}\marg{then}\marg{else}
%
%         这里 |\relax| 终止第一个参数，因为这样比把它
%         括在花括号里更快。\meta{macro} 包含逗号分隔的
%         标识符列表。如果字符序列是这些
%         标识符之一，就执行 \meta{then}，否则执行 \meta{else}。
%
%   \item[0.21] |\lst@Swap|\marg{tok1}\marg{tok2}
%
%         交换后面两个记号或参数的位置
%         \emph{不}插入花括号。例如，|\lst@Swap{abc}{def}|
%         展开为 |defabc|。
%
%   \item[0.18] |\lst@IfNextChars|\meta{macro}\marg{then}\marg{else}
%     \item[0.19] |\lst@IfNextCharsArg|\marg{字符
%         序列}\marg{then}\marg{else}
%
%         两个宏都根据（在三个参数之后）是否找到给定的字符
%         序列或给定宏的内容，来执行 \meta{then} 或 \meta{else}。
%         注意这些宏与 \LaTeX 的 |\@ifnextchar| 的一个重要
%         区别：我们去除参数后面的字符，直到能够
%         判定必须执行哪一部分。不过，我们把这些字符
%         保存在宏 |\lst@eaten| 中，因此可以通过 \meta{then}
%         或 \meta{else} 把它们插回。
%
%   \item[0.19] |\lst@IfNextCharActive|\marg{then}\marg{else}
%
%         若下一个字符是活动的，就执行 \meta{then}，
%         否则执行 \meta{else}。
%
%   \item[0.20] |\lst@DefActive|\meta{macro}\marg{character sequence}
%
%         把字符序列存入 \meta{macro}，但所有字符
%         都变成活动的。该字符串\emph{不得}包含开始分组、结束
%         分组或转义字符（|{}\|）；它可以包含含义不同（即 catcode 不同）的左花括号、右
%         花括号或反斜杠。如果 \meta{character sequence} 尚未被 \TeX 读入，
%         这条命令就完全是多余的，因为这样的 catcode 很容易
%         改动。这里明确允许字符已被读入，例如在
%     |\def\test{\lst@DefActive\temp{ABC}}|!
%
%         注意该宏会改变 |\lccode| 的 0--9 项而不恢复它们。
%
%   \item[0.20] |\lst@DefOther|\meta{macro}\marg{character sequence}
%
%         把 \meta{character sequence} 存入 \meta{macro}，但所有字符
%         的 catcode 都是 12。此外所有空格被移除，控制
%         序列被转换为其名称（不带前置
%         反斜杠）。例如，|\{ Chip \}| 得到 |{Chip}|，其中所有
%         catcode 都是 12——内部使用了原语 |\meaning|。
%
%     \iffalse
%   \item[0.19] |\lst@MakeActive|\marg{character sequence}
%
%         把字符序列存入 |\lst@arg|，限制条件与 |\lst@DefActive|
%         相同。事实上，后者的定义就使用了
%         这里的这个宏。  \fi
% \end{syntax}
%
%
% \subsection{字符表的操作}
%
% \begin{syntax}
%   \item[0.20] |\lst@SaveDef|\marg{character code}\meta{macro}
%
%         把指定字符的当前定义保存在
%         \meta{macro} 中。在重新定义某个字符之前，你应当总是先保存它的
%         定义！并使用保存的版本，而不是直接写
%         |\lst@Process|\ldots——该字符可能已经被
%         重新定义，因此不再等价于它的标准定义。
%
%     \item[0.20] |\lst@DefSaveDef|\marg{字符
%         码}\meta{macro}\meta{parameter text}\marg{definition}
%     \item[0.20] |\lst@LetSaveDef|\marg{字符
%         码}\meta{macro}\meta{token}
%
%         分别把 |\lst@SaveDef| 与 |\lst@Def| 或 |\lst@Let| 结合。
% \end{syntax}
% 当然，我不该忘记说明\emph{在哪里}修改字符
% 表。在 \hookname{SelectCharTable} 处的钩子材料做出永久改动，
% 即影响所有语言。下面两个键可用于任何
% 语言定义，并且只影响该特定语言。
% \begin{syntax}
% \item[0.20] |SelectCharTable=|\meta{\TeX\ code}
% \item[0.20] |MoreSelectCharTable=|\meta{\TeX\ code}
%
%             （额外）使用 \meta{\TeX\ code} 来选择字符表。
%             该代码在标准字符表选出之后执行，
%             但可能在其他方面做出更多改动之前执行。由于先前的
%             含义总是被保存并在新定义内执行，这
%             应该无害。
% \end{syntax}
% 下面是两个相当无用的例子。如果语言 |useless| 处于活动状态，
% 每个句点都会在终端和 |.log| 文件中
% 产生一条消息 `|.|'：
% \begin{verbatim}
%   \lstdefinelanguage{useless}
%       {SelectCharTable=\lst@DefSaveDef{46}% save chr(46) ...
%            \lsts@point             % ... in \lsts@point and ...
%            {\message{.}\lsts@point}% ... use new definition
%       }\end{verbatim}
% 如果你想数点数，可以这样写
% \begin{verbatim}
%   \newcount\lst@points % \global
%   \lst@AddToHook{Init}{\global\lst@points\z@}
%   \lst@AddToHook{DeInit}{\message{Number of points: \the\lst@points}}
%   \lstdefinelanguage[2]{useless}
%       {SelectCharTable=\lst@DefSaveDef{46}\lsts@point
%            {\global\advance\lst@points\@ne \lsts@point}
%       }\end{verbatim}
% |% \global| 表示所分配的计数器被全局使用。我们在每个
% 列表开头把计数器清零，在列表末尾显示
% 当前值的消息，每个处理过的句点都让
% 计数器加一。
%
% \begin{syntax}
%   \item[0.21] |\lst@CArg|\meta{active characters}|\relax|\meta{macro}
%
%         活动字符的字符串被分割为 \meta{1st}、\meta{2nd}
%         和 \marg{rest}。如果某一项不存在，就使用
%         空参数。然后调用 \meta{macro}，参数为
%         |{|\meta{1st}\meta{2nd}\marg{rest}|}|，外加一个尚未定义的控制
%         序列 \meta{save 1st}。该宏意在保存 \meta{1st} 的当前
%         定义，这样 \meta{1st} 就可以在不丢失信息的情况下
%         被重新定义。
%
%   \item[0.19] |\lst@CArgX|\meta{characters}|\relax|\meta{macro}
%
%         在调用 |\lst@CArg| 之前把 \meta{characters} 变成活动的。
%
%     \item[0.21] |\lst@CDef{|\meta{1st}\meta{2nd}\marg{rest}|}|\meta{save
%          1st}\marg{execute}\marg{pre}\marg{post}
%
%         应与 |\lst@CArg| 或 |\lst@CArgX| 配合使用，
%         即在那里作为 \meta{macro}。\meta{1st}、\meta{2nd} 和 \meta{rest}
%         必须是活动字符，\meta{save 1st} 必须是一个未定义的
%         控制序列。
%
%         每当宏包（在列表中）遇到字符 \meta{1st} 时，
%         就执行 \meta{execute}。如果宏包检测到整个字符串
%         \meta{1st}\meta{2nd}\meta{rest}，我们还会额外执行 \meta{pre}，
%         然后是该字符串，最后是 \meta{post}。
%
%     \item[0.21] |\lst@CDefX|\meta{1st}\meta{2nd}\marg{rest}\meta{save
%          1st}\marg{execute}\marg{pre}\marg{post}
%
%         同上，只是当我们遇到 \meta{1st}\meta{2nd}\meta{rest} 时，
%         执行 \meta{pre} 和 \meta{post} 而不执行原字符串。这
%         意味着该字符串被替换为 \meta{pre}\meta{post}（前面带有
%         \meta{execute}）。
% \end{syntax}
% 作为最后一个例子，下面是 |\lst@DefDelimB| 的定义。
% \begin{verbatim}
%    \gdef\lst@DefDelimB#1#2#3#4#5#6#7#8{%
%        \lst@CDef{#1}#2%
%            {#3}%
%            {\let\lst@bnext\lst@CArgEmpty
%             \lst@ifmode #4\else
%                 #5%
%                 \def\lst@bnext{#6{#7}{#8}}%
%             \fi
%             \lst@bnext}%
%            \@empty}\end{verbatim}
% 你明白了吗？
%
%
% \part{实现}
%
%
% \CheckSum{12444}
%^^A
%^^A Don't index TeX-primitives.
%^^A
% \DoNotIndex{\advance,\afterassignment,\aftergroup,\batchmode,\begingroup}
% \DoNotIndex{\box,\catcode,\char,\chardef,\closeout,\copy,\count,\csname,\def}
% \DoNotIndex{\dimen,\discretionary,\divide,\dp,\edef,\else,\end,\endcsname}
% \DoNotIndex{\endgroup,\endinput,\endlinechar,\escapechar,\everypar}
% \DoNotIndex{\expandafter,\fi,\gdef,\global,\globaldefs,\hbadness,\hbox}
% \DoNotIndex{\hrulefill,\hss,\ht}
% \DoNotIndex{\if,\ifdim,\iffalse,\ifhmode,\ifinner,\ifnum,\ifodd,\iftrue}
% \DoNotIndex{\ifvoid,\ifx,\ignorespaces,\immediate,\input,\jobname,\kern}
% \DoNotIndex{\lccode,\leftskip,\let,\long,\lower,\lowercase,\meaning,\message}
% \DoNotIndex{\multiply,\muskip,\noexpand,\noindent,\openout,\par,\parfillskip}
% \DoNotIndex{\parshape,\parskip,\raise,\read,\relax,\rightskip,\setbox,\skip}
% \DoNotIndex{\string,\the,\toks,\uppercase,\vbox,\vcenter,\vrule,\vtop,\wd}
% \DoNotIndex{\write,\xdef}
%
%^^A
%^^A Don't index LaTeX's private definitions.
%^^A
% \DoNotIndex{\@@end,\@@par,\@M,\@arabic,\@circlefnt,\@currentlabel}
% \DoNotIndex{\@currenvir,\@depth,\@doendpe,\@dottedtocline,\@eha,\@ehc}
% \DoNotIndex{\@empty,\@firstofone,\@firstoftwo,\@float,\@for,\@getcirc}
% \DoNotIndex{\@gobble,\@gobbletwo,\@halfwidth,\@height,\@ifnextchar}
% \DoNotIndex{\@ifundefined,\@ignoretrue,\@makecaption,\@makeother,\@namedef}
% \DoNotIndex{\@ne,\@noligs,\@notprerr,\@onlypreamble,\@secondoftwo,\@spaces}
% \DoNotIndex{\@starttoc,\@totalleftmargin,\@undefined,\@whilenum}
% \DoNotIndex{\@wholewidth,\@width}
% \DoNotIndex{\c@chapter,\c@figure,\c@page,\end@float,\f@family,\filename@area}
% \DoNotIndex{\filename@base,\filename@ext,\filename@parse,\if@twoside}
% \DoNotIndex{\l@ngrel@x,\m@ne,\new@command,\nfss@catcodes,\tw@,\thr@@}
% \DoNotIndex{\z@,\zap@space}
%
%^^A
%^^A Don't index LaTeX's package definitions.
%^^A
% \DoNotIndex{\AtEndOfPackage}
% \DoNotIndex{\CurrentOption,\DeclareOption,\IfFileExists,\InputIfFileExists}
% \DoNotIndex{\MessageBreak,\NeedsTeXFormat,\PackageError,\PackageWarning}
% \DoNotIndex{\ProcessOptions,\ProvidesFile,\ProvidesPackage,\RequirePackage}
%
%^^A
%^^A Don't index LaTeX's public definitions.
%^^A
% \DoNotIndex{\abovecaptionskip,\active,\addcontentsline,\addtocounter,\begin}
% \DoNotIndex{\belowcaptionskip,\bfseries,\bgroup,\bigbreak,\chapter}
% \DoNotIndex{\contentsname,\do,\egroup,\footnotesize,\index,\itshape}
% \DoNotIndex{\linewidth,\llap,\makeatletter,\newbox,\newcommand,\newcount}
% \DoNotIndex{\newcounter,\newdimen,\newtoks,\newwrite,\nointerlineskip}
% \DoNotIndex{\normalbaselines,\normalfont,\numberline,\pretolerance,\protect}
% \DoNotIndex{\qquad,\refstepcounter,\removelastskip,\renewcommand,\rlap}
% \DoNotIndex{\small,\smallbreak,\smallskipamount,\smash,\space,\strut}
% \DoNotIndex{\strutbox,\tableofcontents,\textasciicircum,\textasciitilde}
% \DoNotIndex{\textasteriskcentered,\textbackslash,\textbar,\textbraceleft}
% \DoNotIndex{\textbraceright,\textdollar,\textendash,\textgreater,\textless}
% \DoNotIndex{\textunderscore,\textvisiblespace,\thechapter,\ttdefault}
% \DoNotIndex{\ttfamily,\typeout,\value,\vphantom}
%
%^^A
%^^A Don't index definitions from other packages.
%^^A
% \DoNotIndex{\MakePercentComment,\MakePercentIgnore}
%
%^^A
%^^A Don't index 0.19 definitions.
%^^A
% \DoNotIndex{\listlistingsname,\listoflistings,\lstbox,\lstbox@}
% \DoNotIndex{\lstlanguage@}
%
%^^A
%^^A Don't index 0.20 subdefinitions.
%^^A
% \DoNotIndex{\lst@ATH@,\lst@BOLGobble@,\lst@BOLGobble@@,\lst@CArg@,\lst@CArg@@}
% \DoNotIndex{\lst@CBC@,\lst@CBC@@,\lst@CCECUse@,\lst@CCPutMacro@,\lst@DefLang@}
% \DoNotIndex{\lst@DefLang@@,\lst@DefOther@,\lst@DeleteKeysIn@,\lst@Escape@}
% \DoNotIndex{\lstframe@,\lst@frameH@,\lst@frameL@,\lst@frameR@}
% \DoNotIndex{\lst@FillFixed@,\lst@FindAlias@,\lst@FVConvert@}
% \DoNotIndex{\lst@IfNextChars@,\lst@IfNextChars@@,\lst@InsideConvert@}
% \DoNotIndex{\lst@InstallKeywords@,\lst@Key@,\lst@KeywordTestI@}
% \DoNotIndex{\lst@MakeActive@,\lst@MakeMoreKeywords@}
% \DoNotIndex{\lst@RC@,\lst@RC@@,\lst@ReplaceIn@,\lst@ReplaceInput@}
% \DoNotIndex{\lst@ReportAllocs@,\lst@SKS@,\lst@SKS@@,\lst@UBC@}
% \DoNotIndex{\lst@WorkingTestI@,\lstalias@,\lstalias@@,\lstalso@}
% \DoNotIndex{\lstdefinestyle@,\lstenv@BOLGobble@@}
% \DoNotIndex{\lstenv@ProcessJ@,\lstinline@,\lstKV@OptArg@,\lstKV@SetIf@}
% \DoNotIndex{\lstlang@,\lstnewenvironment@,\lst@outputpos,\lstset@}
%
%
% \section{序曲}
%
% \paragraph{寄存器}
% 每个方面所需的寄存器数目列于第
% \lstref{dPackageLoading} 节。此外，\packagename{keyval} 宏包
% 分配一个记号寄存器。宏、盒子和计数器
% |\@temp|\ldots|a|/|b|、尺寸 |\@tempdim|\ldots，以及宏
% |\@gtempa| 也会被使用，见索引。
%
% \paragraph{命名约定}
% 我们先从面向用户的定义开始。所有这些公共宏
% 都使用小写字母，并含有 |lst|。私有宏和变量使用
% 下列前缀（可能不是最新的？）：
% \begin{itemize}
% \item |\lst@| 用于一般宏或变量，
% \item |\lstenv@| 若为列表环境定义，
% \item |\lsts@| 用于保存的字符含义，
% \item |\lsthk@|\meta{钩子名} 保存钩子材料，
% \item |\lst|\meta{prefix}|@| 用于各种关键词和
%             工作标识符。
% \item |\lstlang@|\meta{language}|@|\meta{dialect} 包含一种语言，
% \item |\lststy@|\meta{the style} 包含样式定义，
% \item |\lstpatch@|\meta{aspect} 用于修补某个方面，
%
% \item |\lsta@|\meta{language}|$|\meta{dialect} 包含别名，
% \item |\lsta@|\meta{language} 包含某语言所有方言的别名，
% \item |\lstdd@|\meta{language} 包含某语言的默认方言
%             （如果有）。
% \end{itemize}
% 为了区分过程类宏和数据类宏，过程
% 宏的名字在每个词首使用大写字母，例如 |\lst@AddTo|。
% 后缀为 |@| 的宏是另一个定义的主工作过程，
% 例如 |\lstMakeShortInline@| 为 |\lstMakeShortInline| 做主要工作。
%
% \paragraph{前导部分}
% 由本 \texttt{listings.dtx} 生成的所有文件都会得到一个文件头。
%    \begin{macrocode}
%<*kernel|misc>
%% Please read the software license in listings.dtx or listings.pdf.
%%
%% (w)(c) 1996--2004 Carsten Heinz and/or any other author listed
%% elsewhere in this file.
%% (c) 2006 Brooks Moses
%% (c) 2013- Jobst Hoffmann
%%
%% Send comments and ideas on the package, error reports and additional
%% programming languages to Jobst Hoffmann at <j.hoffmann@fh-aachen.de>.
%%
%</kernel|misc>
%    \end{macrocode}
%
% \paragraph{标识}
% 所有文件将有相同的日期和版本。
%    \begin{macrocode}
%<*kernel|misc|doc>
\def\filedate{2025/11/14}
\def\fileversion{1.11b}
%</kernel|misc|doc>
%    \end{macrocode}
% 我们需要什么，我们是谁。
%    \begin{macrocode}
%<*kernel>
\NeedsTeXFormat{LaTeX2e}
\AtEndOfPackage{\ProvidesPackage{listings}
             [\filedate\space\fileversion\space(Carsten Heinz)]}
%    \end{macrocode}
% \begin{macro}{\lst@CheckVersion}
% 可以被各个驱动文件用来保证正确的版本。
%    \begin{macrocode}
\def\lst@CheckVersion#1{\edef\reserved@a{#1}%
    \ifx\lst@version\reserved@a \expandafter\@gobble
                          \else \expandafter\@firstofone \fi}
%    \end{macrocode}
%    \begin{macrocode}
\let\lst@version\fileversion
%</kernel>
%    \end{macrocode}
% \end{macro}
% 例如由杂项文件
%    \begin{macrocode}
%<*misc>
\ProvidesFile{lstmisc.sty}
             [\filedate\space\fileversion\space(Carsten Heinz)]
\lst@CheckVersion\fileversion%
    {\typeout{^^J%
     ***^^J%
     *** This file requires `listings.sty' version \fileversion.^^J%
     *** You have a serious problem, so I'm exiting...^^J%
     ***^^J}%
     \batchmode\@@end}
%</misc>
%    \end{macrocode}
% 或由哑元补丁。
%    \begin{macrocode}
%<*patch>
\ProvidesFile{lstpatch.sty}
             [\filedate\space\fileversion\space(Carsten Heinz)]
\lst@CheckVersion\lst@version{}
%%% intentionally empty
%</patch>
%    \end{macrocode}
%    \begin{macrocode}
%<*doc>
\ProvidesPackage{lstdoc}
             [\filedate\space\fileversion\space(Carsten Heinz)]
%</doc>
%    \end{macrocode}
%
% \paragraph{类别码}
% 我们定义两个宏，以便在输入 \packagename{listings} 宏包的
% 其他文件时保证正确的 catcode。
%
% \begin{macro}{\lst@InputCatcodes}
% |@| 和 |"| 变成字母。制表符和行尾被忽略；这可以避免
% 多余的空格——以防我忘了写注释字符。
%    \begin{macrocode}
%<*kernel>
\def\lst@InputCatcodes{%
    \makeatletter \catcode`\"12%
    \catcode`\^^@\active
    \catcode`\^^I9%
    \catcode`\^^L9%
    \catcode`\^^M9%
    \catcode`\%14%
    \catcode`\~\active}
%    \end{macrocode}
% \end{macro}
%
% \begin{macro}{\lst@RestoreCatcodes}
%         为了加载内核，我们会改变一些 catcode 和 lccode。在宏包
%         加载结束时我们恢复它们。\lsthelper{Jobst~Hoffmann}{2000/11/17} {与
%         typehtml 宏包不兼容}报告了与
%         \packagename{typehtml} 宏包的不兼容，下面通过 |\lccode`\/`\/| 解决。与
%         \packagename{typehtml} 宏包的不兼容
%         由现任维护者重新测试
%         （\lsthelper{Jobst~Hoffmann}{2024/10/16} {与
%         typehtml 宏包的不兼容已不复存在}），因为收到了
%         \lsthelper{Sebastian Ørsted}{2024/10/03}{listings 中的缺陷}的一封邮件；在
%         从 |\edef| 中删除下面
%         这一行之后，该不兼容不再出现：
% \begin{verbatim}
% \noexpand\lccode`\noexpand\/`\noexpand\/\end{verbatim}
% 所以 |\lst@RestoreCatcodes| 的定义与 v0.21 相同，只做了
% v1.4 引入的一处小改动（见第~\pageref{catcode12} 页）。
%    \begin{macrocode}
\def\lst@RestoreCatcodes#1{%
    \ifx\relax#1\else
        \noexpand\catcode`\noexpand#1\the\catcode`#1\relax
        \expandafter\lst@RestoreCatcodes
    \fi}
\edef\lst@RestoreCatcodes{%
    \lst@RestoreCatcodes\"\^^I\^^M\~\^^@\relax
    \catcode12\active}
%    \end{macrocode}
% 现在我们准备好进行
%    \begin{macrocode}
\lst@InputCatcodes
\AtEndOfPackage{\lst@RestoreCatcodes}
%</kernel>
%    \end{macrocode}
% \end{macro}
%
% \paragraph{统计}
% \begin{macro}{\lst@GetAllocs}
% \begin{macro}{\lst@ReportAllocs}
% 用于显示已分配的寄存器。
%    \begin{macrocode}
%<*info>
\def\lst@GetAllocs{%
    \edef\lst@allocs{%
        0\noexpand\count\the\count10,1\noexpand\dimen\the\count11,%
        2\noexpand\skip\the\count12,3\noexpand\muskip\the\count13,%
        4\noexpand\box\the\count14,5\noexpand\toks\the\count15,%
        6\noexpand\read\the\count16,7\noexpand\write\the\count17}}
\def\lst@ReportAllocs{%
    \message{^^JAllocs:}\def\lst@temp{none}%
    \expandafter\lst@ReportAllocs@\lst@allocs,\z@\relax\z@,}
\def\lst@ReportAllocs@#1#2#3,{%
    \ifx#2\relax \message{\lst@temp^^J}\else
        \@tempcnta\count1#1\relax \advance\@tempcnta -#3\relax
        \ifnum\@tempcnta=\z@\else
            \let\lst@temp\@empty
            \message{\the\@tempcnta \string#2,}%
        \fi
        \expandafter\lst@ReportAllocs@
    \fi}
\lst@GetAllocs
%    \end{macrocode}
% \end{macro}\end{macro}
% \begingroup
%    \begin{macrocode}
%</info>
%    \end{macrocode}
% \endgroup
%
% \paragraph{其他}
% \begin{macro}{\@lst}
% 只是一个节省内存空间的定义。
%    \begin{macrocode}
%<*kernel>
\def\@lst{lst}
%</kernel>
%    \end{macrocode}
% \end{macro}
%
%
% \section{一般问题}
%
% 本节的所有定义都属内核。
%    \begin{macrocode}
%<*kernel>
%    \end{macrocode}
%
%
%^^A \subsection{Quick `if parameter empty'}
%^^A
%^^A There are many situations where you have to look whether a macro parameter is empty.
%^^A We have at least two possibilities to test |#1|, for example:
%^^A \begin{center}
%^^A \begin{minipage}{0.35\linewidth}
%^^A \begin{verbatim}
%^^A\def\test{#1}%
%^^A\ifx \test\empty
%^^A        % #1 is empty
%^^A\else
%^^A        % #1 is not empty
%^^A\fi\end{verbatim}
%^^A \end{minipage}
%^^A \hskip2em\vrule\hskip2em
%^^A \begin{minipage}{0.35\linewidth}
%^^A \begin{verbatim}
%^^A\ifx \empty#1\empty
%^^A        % #1 is empty
%^^A\else
%^^A        % #1 is not empty
%^^A\fi\end{verbatim}
%^^A \end{minipage}
%^^A \end{center}
%^^A where |\empty| is defined by |\def\empty{}|, of course.
%^^A The left variant should be clear and works in any case.
%^^A The right-hand side works correct if and only if the first token of |#1| is
%^^A not equivalent to |\empty|.
%^^A This granted, the |\ifx| is true if and only if |#1| is empty (since |\empty|
%^^A left from |#1| is (not) compared with |\empty| on the right).
%^^A The two |\empty|s might be replaced by any other macro, which is not
%^^A equivalent to the first token of the argument.
%^^A But the definition of that macro shouldn't be too complex since this would
%^^A slow down the |\ifx|.
%^^A The right example needs about $45\%$ of the left's time.
%^^A Note that this \TeX{}nique lost its importance from version 0.18 on and that
%^^A other packages use |!| or |\relax| instead of |\empty|, for example.
%^^A
%^^A This \TeX nique is described in ``The \TeX book'' on page 376.
%
%
% \subsection{子串测试}\label{iSubstringTests}
%
% \lstset{language=TeX,gobble=4,xleftmargin=20pt,columns=[l]fullflexible,^^A
%   mathescape,keywordstyle=\ttfamily,texcsstyle=\bfseries}
% \let\texverb\lstinline
% \lstnewenvironment{texcode}[1][]{\lstset{#1}}{}
% \lstset{keywords={def,key}}
%
% 判断给定字符序列是否是另一个字符串的子串很容易。
% 例如，对于子串 \texverb|def|，我们可以写
% \begin{texcode}
%   \def \lst@temp#1def#2\relax{%
%       \ifx \@empty#2\@empty
%               % "def" is not a substring
%       \else
%               % "def" is a substring
%       \fi}
%
%   \lst@temp $\meta{another\ string}$def\relax
% \end{texcode}
% 当 \TeX\ 传入参数 |#1| 和 |#2| 时，第二个为空当
% 且仅当 \texverb|def| 不是子串。如果没有附加的
% \texverb|def\relax|，当 \meta{another string} 不含 \texverb|def| 时，
% 就会得到``runaway argument''错误。
%
% 我们主要在标识符与逗号分隔的键或关键词列表这一特殊
% 情况下使用子串测试：
% \begin{texcode}[keywords=key]
%   \def \lst@temp#1,key,#2\relax{%
%       \ifx \@empty#2\@empty
%               % `key' is not a keyword
%       \else
%               % `key' is a keyword
%       \fi}
%
%   \lst@temp,$\meta{list\ of\ keywords}$,key,\relax
% \end{texcode}
% 这工作得非常好，而且相当快。但在 \texttt{key} 是关键词的情况下，
% 我们可以进一步减少运行时间。此时 |#2| 取到字符串的其余部分，
% 即 \texttt{key} 之后的所有关键词。
% 由于 \TeX\ 把 |#2| 插入到 \texverb|\@empty| 之间，它必须丢弃
% |#2| 除第一个字符之外的所有内容——第一个字符会与 \texverb|\@empty| 比较。
% 我们可以把这段其余内容重定向到第三个参数：
% \begin{texcode}[keywords=key]
%   \def \lst@temp#1,key,#2#3\relax{%
%       \ifx \@empty#2%
%               % "key" is not a keyword
%       \else
%               % "key" is a keyword
%       \fi}
%
%   \lst@temp,$\meta{list\ of\ keywords}$,key,\@empty\relax
% \end{texcode}
% 这稍快一些，是 0.20 版的改进。
%
% \begin{macro}{\lst@IfSubstring}
% 从上面的讨论应当可以清楚地看出实现方式。
%    \begin{macrocode}
\def\lst@IfSubstring#1#2{%
    \def\lst@temp##1#1##2##3\relax{%
        \ifx \@empty##2\expandafter\@secondoftwo
                 \else \expandafter\@firstoftwo \fi}%
    \expandafter\lst@temp#2#1\@empty\relax}
%    \end{macrocode}
% \end{macro}
%
% \begin{macro}{\lst@IfOneOf}
% 同上。
%    \begin{macrocode}
\def\lst@IfOneOf#1\relax#2{%
    \def\lst@temp##1,#1,##2##3\relax{%
        \ifx \@empty##2\expandafter\@secondoftwo
                 \else \expandafter\@firstoftwo \fi}%
    \expandafter\lst@temp\expandafter,#2,#1,\@empty\relax}
%    \end{macrocode}
% \end{macro}
% \begin{REMOVED}
% 有一天，如果再次需要不区分大小写的键（词）测试，
% 我们可以用两个 |\uppercase| 来规范化第一个参数：
%    \begin{verbatim}
%\def\lst@IfOneOfInsensitive#1\relax#2{%
%    \uppercase{\def\lst@temp##1,#1},##2##3\relax{%
%        \ifx \@empty##2\expandafter\@secondoftwo
%                 \else \expandafter\@firstoftwo \fi}%
%    \uppercase{%
%        \expandafter\lst@temp\expandafter,#2,#1},\@empty\relax}\end{verbatim}
% 这里我们假定宏 |#2| 已经只包含大写字符，见第
% \ref{iMakingTests} 节最末尾 |\lst@MakeMacroUppercase| 的
% 定义。如果我们\emph{不得}作此假定，只需在
% 第二个 |\uppercase| 与后面的花括号之间插入一个 |\expandafter|。
% 但这会拖慢测试！
% \end{REMOVED}
%
% \begin{macro}{\lst@DeleteKeysIn}
% 子宏做主要工作；我们只需展开第二个
% 宏——要删除的键列表——并附加终止符 |\relax|。
%    \begin{macrocode}
\def\lst@DeleteKeysIn#1#2{%
    \expandafter\lst@DeleteKeysIn@\expandafter#1#2,\relax,}
%    \end{macrocode}
% 把最后一个 |\lst@DeleteKeysIn@| 换成 |\lst@RemoveCommas| 就
% 在这里终止循环。注意 |#2| 之后的 |\@empty| 确保这个
% 宏在 |#2| 为空时也能工作。
%    \begin{macrocode}
\def\lst@DeleteKeysIn@#1#2,{%
    \ifx\relax#2\@empty
        \expandafter\@firstoftwo\expandafter\lst@RemoveCommas
    \else
        \ifx\@empty#2\@empty\else
%    \end{macrocode}
% 如果我们尚未到达列表末尾且键不为空，我们就
% 定义一个临时宏来删除所有出现。
%    \begin{macrocode}
            \def\lst@temp##1,#2,##2{%
                ##1%
                \ifx\@empty##2\@empty\else
                    \expandafter\lst@temp\expandafter,%
                \fi ##2}%
            \edef#1{\expandafter\lst@temp\expandafter,#1,#2,\@empty}%
        \fi
    \fi
    \lst@DeleteKeysIn@#1}
%    \end{macrocode}
% \end{macro}
% \begin{OLDDEF}
% 下面的修改大约需要多 $50\%$ 的运行时间。
% 它不使用 |\edef|，因此也能处理 |#1| 内的 |\{|。
% 不过，我们目前不需要这一功能。
%    \begin{verbatim}
%            \def\lst@temp##1,#2,##2{%
%                \ifx\@empty##2%
%                    \lst@lAddTo#1{##1}%
%                \else
%                    \lst@lAddTo#1{,##1}%
%                    \expandafter\lst@temp\expandafter,%
%                \fi ##2}%
%            \let\@tempa#1\let#1\@empty
%            \expandafter\lst@temp\expandafter,\@tempa,#2,\@empty\end{verbatim}
% \end{OLDDEF}
%
% \begin{macro}{\lst@RemoveCommas}
% 该宏删除开头的逗号，并把新值赋给 |#1|。
%    \begin{macrocode}
\def\lst@RemoveCommas#1{\edef#1{\expandafter\lst@RC@#1\@empty}}
\def\lst@RC@#1{\ifx,#1\expandafter\lst@RC@ \else #1\fi}
%    \end{macrocode}
% \end{macro}
% \begin{OLDDEF}
% 下面的版本可以处理宏 |#1| 内的 |\{|。
%    \begin{verbatim}
%\def\lst@RemoveCommas#1{\expandafter\lst@RC@#1\@empty #1}
%\def\lst@RC@#1{%
%    \ifx,#1\expandafter\lst@RC@
%      \else\expandafter\lst@RC@@\expandafter#1\fi}
%\def\lst@RC@@#1\@empty#2{\def#2{#1}}\end{verbatim}
% \end{OLDDEF}
%
% \begin{macro}{\lst@ReplaceIn}
% \begin{macro}{\lst@ReplaceInArg}
% 这些宏与 |\lst@DeleteKeysIn| 类似，只是\ldots
%    \begin{macrocode}
\def\lst@ReplaceIn#1#2{%
    \expandafter\lst@ReplaceIn@\expandafter#1#2\@empty\@empty}
\def\lst@ReplaceInArg#1#2{\lst@ReplaceIn@#1#2\@empty\@empty}
%    \end{macrocode}
% \ldots\ 我们用 |#3| 替换 |#2|，而不是把 |,#2,| 替换为单个逗号
% （上面的做法则删除了键 |#2|）。
%    \begin{macrocode}
\def\lst@ReplaceIn@#1#2#3{%
    \ifx\@empty#3\relax\else
        \def\lst@temp##1#2##2{%
            \ifx\@empty##2%
                \lst@lAddTo#1{##1}%
            \else
                \lst@lAddTo#1{##1#3}\expandafter\lst@temp
            \fi ##2}%
        \let\@tempa#1\let#1\@empty
        \expandafter\lst@temp\@tempa#2\@empty
        \expandafter\lst@ReplaceIn@\expandafter#1%
    \fi}
%    \end{macrocode}
% \end{macro}
% \end{macro}
%
%
% \subsection{控制流}
%
% \begin{macro}{\@gobblethree}
% 当且仅当未定义时才定义。
%    \begin{macrocode}
\providecommand*\@gobblethree[3]{}
%    \end{macrocode}
% \end{macro}
%
% \begin{macro}{\lst@GobbleNil}
%    \begin{macrocode}
\def\lst@GobbleNil#1\@nil{}
%    \end{macrocode}
% \end{macro}
%
% \begin{macro}{\lst@Swap}
% 也就是这样：
%    \begin{macrocode}
\def\lst@Swap#1#2{#2#1}
%    \end{macrocode}
% \end{macro}
%
% \begin{macro}{\lst@if}
% \begin{macro}{\lst@true}
% \begin{macro}{\lst@false}
% 用于临时使用的通用 |\if|。
%    \begin{macrocode}
\def\lst@true{\let\lst@if\iftrue}
\def\lst@false{\let\lst@if\iffalse}
\lst@false
%    \end{macrocode}
% \end{macro}
% \end{macro}
% \end{macro}
%
% \begin{macro}{\lst@IfNextCharsArg}
% 很容易：我们定义一个宏并调用 |\lst@IfNextChars|。
%    \begin{macrocode}
\def\lst@IfNextCharsArg#1{%
    \def\lst@tofind{#1}\lst@IfNextChars\lst@tofind}
%    \end{macrocode}
% \end{macro}
%
% \begin{macro}{\lst@IfNextChars}
% 我们保存参数并开始一个循环。
%    \begin{macrocode}
\def\lst@IfNextChars#1#2#3{%
    \let\lst@tofind#1\def\@tempa{#2}\def\@tempb{#3}%
    \let\lst@eaten\@empty \lst@IfNextChars@}
%    \end{macrocode}
% 展开我们要查找的字符。
%    \begin{macrocode}
\def\lst@IfNextChars@{\expandafter\lst@IfNextChars@@\lst@tofind\relax}
%    \end{macrocode}
% 现在我们可以细化 |\lst@tofind| 并把输入字符 |#3| 追加到
% |\lst@eaten|.
%    \begin{macrocode}
\def\lst@IfNextChars@@#1#2\relax#3{%
    \def\lst@tofind{#2}\lst@lAddTo\lst@eaten{#3}%
    \ifx#1#3%
%    \end{macrocode}
% 如果字符相同，我们就调用 |\@tempa| 或继续测试。
%    \begin{macrocode}
        \ifx\lst@tofind\@empty
            \let\lst@next\@tempa
        \else
            \let\lst@next\lst@IfNextChars@
        \fi
        \expandafter\lst@next
    \else
%    \end{macrocode}
% 如果字符不同，我们就调用 |\@tempb|。
%    \begin{macrocode}
        \expandafter\@tempb
    \fi}
%    \end{macrocode}
% \end{macro}
%
% \begin{macro}{\lst@IfNextCharActive}
% 我们把字符 |#3| 与它的活动版本 |\lowercase{~}| 比较。
% 注意 |\ifx~| 与 |#3| 之间的右花括号结束了 |\lowercase|。
% |\endgroup| 恢复 |\lccode|。
%    \begin{macrocode}
\def\lst@IfNextCharActive#1#2#3{%
    \begingroup \lccode`\~=`#3\lowercase{\endgroup
    \ifx~}#3%
        \def\lst@next{#1}%
    \else
        \def\lst@next{#2}%
    \fi \lst@next #3}
%    \end{macrocode}
% \end{macro}
%
% \begin{macro}{\lst@for}
% 带循环变量展开的 for 循环。这得益于
% \lsthelper{Hendri~Adriaens}{2006/03/31}{加速
% |\lst@for|} 的建议。
%    \begin{macrocode}
\def\lst@for#1\do#2{%
  \def\lst@forbody##1{#2}%
  \def\@tempa{#1}%
  \ifx\@tempa\@empty\else\expandafter\lst@f@r#1,\@nil,\fi
}
\def\lst@f@r#1,{%
  \def\@tempa{#1}%
  \ifx\@tempa\@nnil\else\lst@forbody{#1}\expandafter\lst@f@r\fi
}
%    \end{macrocode}
% \end{macro}
%
%
% \subsection{catcode 的改动}\label{iCatcodeChanges}
%
% 字符在被读入之后立即得到 catcode，而 \TeX\ 没有
% 原语命令来改变已附着的 catcode。不过，我们可以用
% ASCII 码相同但 catcode 不同的字符替换这些字符。
% 这并不等同，但足够了，因为结果相同。
% 这里我们处理所有字符都变成活动的这个非常特殊的情况。
% 如果我们想让 \texverb|\lst@arg| 包含字符
% |#1| 的活动版本，一个原型宏可以是
% \begin{texcode}
%   \def\lst@MakeActive#1{\lccode`\~=`#1\lowercase{\def\lst@arg{~}}}
% \end{texcode}
% |\lowercase| 把 |~| 的 ASCII 码改为 |#1| 的 ASCII 码，因为
% 我们已说明 |#1| 是 |~| 的小写版本。
% 幸运的是 |\lowercase| 不改动 catcode，所以我们得到 |#1| 的一个活动
% 版本。注意 |~| 通常是活动的。
%
% \begin{macro}{\lst@MakeActive}
% 我们不会逐字符做这件事。
% 为了提高速度，我们同时改变九个字符（如果还剩九个
% 字符）。
% \begin{TODO}
% 这是在每次列表都转换定界符时引入的。现在
% 该转换只在每次选择语言时进行。所以我们也许想
% 再次实现逐字符转换以减少内存
% 占用。
% \end{TODO}
% 我们取得参数，把 |\lst@arg| 清空，然后开始一个循环。
%    \begin{macrocode}
\def\lst@MakeActive#1{%
    \let\lst@temp\@empty \lst@MakeActive@#1%
    \relax\relax\relax\relax\relax\relax\relax\relax\relax}
%    \end{macrocode}
% 有九个 |\relax|，因为 |\lst@MakeActive@| 有九个参数，而
% 我们不希望 |#1| 为空时出问题。
% 我们现在需要九个活动字符，而不是单个 |~|。
% 我们把这些 catcode 改动变成局部的，并把接下来的宏定义为 |\global|。
%    \begin{macrocode}
\begingroup
\catcode`\^^@=\active \catcode`\^^A=\active \catcode`\^^B=\active
\catcode`\^^C=\active \catcode`\^^D=\active \catcode`\^^E=\active
\catcode`\^^F=\active \catcode`\^^G=\active \catcode`\^^H=\active
%    \end{macrocode}
% 首先我们用 |\let| 让下一个操作成为 |\relax|。
% 这会中止我们处理所有字符（默认，之后可能被
% 改变）的循环。
% 然后我们检查是否至少还有一个字符。
% 如果不是，循环终止，一切完成。
%    \begin{macrocode}
\gdef\lst@MakeActive@#1#2#3#4#5#6#7#8#9{\let\lst@next\relax
    \ifx#1\relax
    \else \lccode`\^^@=`#1%
%    \end{macrocode}
% 否则我们说 |^^@|=|chr(0)| 是第一个字符的
% 小写版本。
% 然后我们测试第二个字符。
% 如果没有第二个字符，我们把小写的 |^^@| 追加到 |\lst@temp|。
% 否则我们说 |^^A|=|chr(1)| 是第二个
% 字符的小写版本，于是测试下一个参数，依此类推。
%    \begin{macrocode}
    \ifx#2\relax
        \lowercase{\lst@lAddTo\lst@temp{^^@}}%
    \else \lccode`\^^A=`#2%
    \ifx#3\relax
        \lowercase{\lst@lAddTo\lst@temp{^^@^^A}}%
    \else \lccode`\^^B=`#3%
    \ifx#4\relax
        \lowercase{\lst@lAddTo\lst@temp{^^@^^A^^B}}%
    \else \lccode`\^^C=`#4%
    \ifx#5\relax
        \lowercase{\lst@lAddTo\lst@temp{^^@^^A^^B^^C}}%
    \else \lccode`\^^D=`#5%
    \ifx#6\relax
        \lowercase{\lst@lAddTo\lst@temp{^^@^^A^^B^^C^^D}}%
    \else \lccode`\^^E=`#6%
    \ifx#7\relax
        \lowercase{\lst@lAddTo\lst@temp{^^@^^A^^B^^C^^D^^E}}%
    \else \lccode`\^^F=`#7%
    \ifx#8\relax
        \lowercase{\lst@lAddTo\lst@temp{^^@^^A^^B^^C^^D^^E^^F}}%
    \else \lccode`\^^G=`#8%
    \ifx#9\relax
        \lowercase{\lst@lAddTo\lst@temp{^^@^^A^^B^^C^^D^^E^^F^^G}}%
%    \end{macrocode}
% 如果存在九个字符，我们就追加九个
% 活动字符（的小写版本），并通过重新定义 |\lst@next| 再次调用这个宏。
%    \begin{macrocode}
    \else \lccode`\^^H=`#9%
        \lowercase{\lst@lAddTo\lst@temp{^^@^^A^^B^^C^^D^^E^^F^^G^^H}}%
        \let\lst@next\lst@MakeActive@
    \fi \fi \fi \fi \fi \fi \fi \fi \fi
    \lst@next}
\endgroup
%    \end{macrocode}
% 这个 |\endgroup| 恢复 |chr(0)|--|chr(8)| 的 catcode，但不恢复
% |\lst@MakeActive@| 内部字符的
% catcode，因为它们已经被读入了。
%
% 注意：从任意的`catcode--字符码'表转换回
% \TeX 的 catcode 是可能的，如果我们对字符码做测试（或者
% 通过 |\ifnum|，或者通过 |\ifcase|）。
% 但控制序列以及开始和结束分组字符肯定需要
% 一些特殊处理。
% 不过我还没有检查细节。
% 所以请忽略这一点，不要为这条注释来烦我。:\,--\,)
% \end{macro}
%
% \begin{macro}{\lst@DefActive}
% |\lst@MakeActive| 的一个简单应用。
%    \begin{macrocode}
\def\lst@DefActive#1#2{\lst@MakeActive{#2}\let#1\lst@temp}
%    \end{macrocode}
% \end{macro}
%
% \begin{macro}{\lst@DefOther}
% 我们利用这一事实：|\meaning| 产生的字符 catcode 为 12，只有空格
% 仍是空格。|\escapechar| 被局部修改以抑制
% 转义字符的输出。最后我们用 \LaTeX 的 |\zap@space| 移除空格，这个
% 想法由 \lsthelper{Rolf~Niepraschk}{1997/04/24}{使用 \zap@space}提出——不是在
% 这一语境中，但那没关系。
%    \begin{macrocode}
\def\lst@DefOther#1#2{%
    \begingroup \def#1{#2}\escapechar\m@ne \expandafter\endgroup
    \expandafter\lst@DefOther@\meaning#1\relax#1}
\def\lst@DefOther@#1>#2\relax#3{\edef#3{\zap@space#2 \@empty}}
%    \end{macrocode}
% \end{macro}
%
%
%\ifhyper
% \subsection{对第~\ref*{iCatcodeChanges} 节的应用}\label{iApplicationsTo}
%\else
% \subsection{对第~\ref{iCatcodeChanges} 节的应用}\label{iApplicationsTo}
%\fi
%
% 如果在参数内使用环境，列表已经被读入，
% 我们就无法保留 catcode。
% 不过，在某些情况下环境可以在参数
% 内使用——至少我在用户指南中是这样说的。
% 现在我得为让它成真而努力。
% 此外，我们为 \packagename{fancyvrb} 模式定义一个类似的转换宏。
% \begin{syntax}
% \item[0.19] |\lst@InsideConvert{|\meta{\TeX\ material \textup(already read\textup)}|}|
%
%             把参数的逐字版本\emph{追加}到 |\lst@arg|，但所有
%             追加的字符都是活动的。由于这不是字符到字符的
%             转换，`逐字'需要解释。所有字符都可以
%             原样输入，只有 |\|、|{|、|}| 和 |%| 例外。如果你想要其中
%             某一个，必须改写成 |\\|、|\{|、|\}| 和 |\%|。
%             如果两个空格要相邻，第二个（第三、第四、
%             \ldots）空格必须以带前置反斜杠的方式输入。
%
% \item[0.19] |\lst@XConvert{|\meta{\TeX\ material \textup(already read\textup)}|}|
%
%             把参数的`逐字'版本\emph{追加}到 |\lst@arg|。
%             这里允许把 \TeX\ 材料放进参数花括号中，例如
%       |{(*}{*)}|。这些参数的内容会被转换，花括号
%             保持为花括号。
%
%             如果 |\lst@if| 为真，每个第二参数会被区别对待。
%             只有（定界符的）第一个字符变成活动的。
% \end{syntax}
%
% \begin{macro}{\lst@InsideConvert}
% 如果 \texttt{mathescape} 未开启，我们就（在该定义接近末尾处）调用
% 一个类似 |\zap@space| 的子宏，把单个空格换成活动空格。
% 否则我们检查代码中是否含有一对 |$...$|，并调用
% 相应的宏。
%    \begin{macrocode}
\def\lst@InsideConvert#1{%
   \lst@ifmathescape
      \lst@InsideConvert@e#1$\@nil
      \lst@if
         \lst@InsideConvert@ey#1\@nil
      \else
         \lst@InsideConvert@#1 \@empty
         \expandafter\@gobbletwo
      \fi
      \expandafter\lst@next
   \else
      \lst@InsideConvert@#1 \@empty
   \fi}
\begingroup \lccode`\~=`\ \relax \lowercase{%
%    \end{macrocode}
% 我们把 |#1| 变成活动的，并把这些字符（外加一个活动空格）追加到
% |\lst@arg|.
% 如果没有找到输入的结束 |\@empty|，我们就继续处理。
%    \begin{macrocode}
\gdef\lst@InsideConvert@#1 #2{%
    \lst@MakeActive{#1}%
    \ifx\@empty#2%
        \lst@lExtend\lst@arg{\lst@temp}%
    \else
        \lst@lExtend\lst@arg{\lst@temp~}%
        \expandafter\lst@InsideConvert@
    \fi #2}
%    \end{macrocode}
% 最后我们结束 |\lowercase| 并关闭分组。
%    \begin{macrocode}
}\endgroup
%    \end{macrocode}
% 上面的下一个定义用于检查 |$...$|，再下一个
% 则阻止数学内容被转换。该功能由
% \lsthelper{Dr.~Jobst~Hoffmann}{}{}.
%    \begin{macrocode}
\def\lst@InsideConvert@e#1$#2\@nil{%
   \ifx\@empty#2\@empty \lst@false \else \lst@true \fi}
\def\lst@InsideConvert@ey#1$#2$#3\@nil{%
   \lst@InsideConvert@#1 \@empty
   \lst@lAddTo\lst@arg{%
      \lst@ifdropinput\else
         \lst@TrackNewLines\lst@OutputLostSpace \lst@XPrintToken
         \setbox\@tempboxa=\hbox\bgroup$\lst@escapebegin
         #2%
         \lst@escapeend$\egroup \lst@CalcLostSpaceAndOutput
         \lst@whitespacefalse
      \fi}%
   \def\lst@next{\lst@InsideConvert{#3}}%
}
%    \end{macrocode}
% \end{macro}
%
% \begin{macro}{\lst@XConvert}
% 检查参数\ldots
%    \begin{macrocode}
\def\lst@XConvert{\@ifnextchar\bgroup \lst@XConvertArg\lst@XConvert@}
%    \end{macrocode}
% \ldots，转换该参数，把它连同一对分组定界符加入
% |\lst@arg|，然后我们继续转换。
%    \begin{macrocode}
\def\lst@XConvertArg#1{%
    {\lst@false \let\lst@arg\@empty
     \lst@XConvert#1\@nil
     \global\let\@gtempa\lst@arg}%
    \lst@lExtend\lst@arg{\expandafter{\@gtempa}}%
    \lst@XConvertNext}
%    \end{macrocode}
% 若没有 |\bgroup|，我们就检查是否已找到输入的末尾，并
% 转换一个记号（（非）活动字符或控制序列），然后继续。
%    \begin{macrocode}
\def\lst@XConvert@#1{%
    \ifx\@nil#1\else
        \begingroup\lccode`\~=`#1\lowercase{\endgroup
        \lst@lAddTo\lst@arg~}%
        \expandafter\lst@XConvertNext
    \fi}
\def\lst@XConvertNext{%
    \lst@if \expandafter\lst@XConvertX
      \else \expandafter\lst@XConvert \fi}
%    \end{macrocode}
% 现在我们只让第一个字符变成活动的。
%    \begin{macrocode}
\def\lst@XConvertX#1{%
    \ifx\@nil#1\else
        \lst@XConvertX@#1\relax
        \expandafter\lst@XConvert
    \fi}
\def\lst@XConvertX@#1#2\relax{%
    \begingroup\lccode`\~=`#1\lowercase{\endgroup
    \lst@XCConvertX@@~}{#2}}
\def\lst@XCConvertX@@#1#2{\lst@lAddTo\lst@arg{{#1#2}}}
%    \end{macrocode}
% \end{macro}
%
%
% \subsection{驱动文件的处理*}
%
% \packagename{listings} 宏包被拆成若干驱动文件、
% 杂项（= 方面）文件，以及一个内核文件。
% 所有这些文件都可以按需部分加载——只有提供该功能的
% 内核除外。
% \begin{syntax}
%     \item[0.21] |\lst@Require|\marg{name}\marg{prefix}\marg{feature
%         list}\meta{alias}\meta{file list macro}
%
%         尝试从 \meta{file list macro} 中列出的文件中加载 \meta{feature list} 的
%         所有条目。每个条目的形式为
%         [\oarg{sub}]\meta{feature}。|\lst@if| 等于 |\iftrue| 当且仅
%         当所有条目都可加载。
%
%         宏 \meta{alias} 取一个条目作为参数，必须定义
%         |\lst@oalias| 和 |\lst@malias| 的相应版本。实际上，与这些定义
%         相关的特性会被加载。你可以使用
%         \meta{alias}=|\lst@NoAlias| 表示不作替换。
%
%         \meta{prefix} 在内部标识类型，\meta{name} 用于
%         消息。
%
%         例如，|\lstloadaspects| 使用下列参数，其中
%         |#1| 是方面的列表：|{aspects}|\allowbreak|a|\allowbreak
%     |{#1}|\allowbreak|\lst@NoAlias|\allowbreak|\lstaspectfiles|.
%
%     \item[0.20] |\lst@DefDriver|\marg{name}\marg{prefix}\meta{interface
%         macro}|\if|\alternative{true,false}
%
%   \item[0.21]
%     |\lst@IfRequired|\oarg{sub}\marg{feature}\marg{then}\marg{else}
%
%         由方面、语言或其他定义命令在驱动文件
%         内部使用。\meta{then} 被执行当且仅当
%         \oarg{sub}\marg{feature} 已通过 |\lst@Require| 请求过。
%         否则执行 \meta{else}——对相同 \oarg{sub}\marg{feature} 的
%         后续调用也是如此。
%
%         \meta{then} 和 \meta{else} 可以使用 |\lst@prefix|（只读）。
%
%         第~\ref{iAspectCommands} 节中的
%         |\lst@BeginAspect| 与 |\lst@DefDriver| 可作例子。
% \end{syntax}
%
% \begin{macro}{\lst@Require}
% 初始化变量（如果所需条目非空），\ldots
%    \begin{macrocode}
\def\lst@Require#1#2#3#4#5{%
    \begingroup
    \aftergroup\lst@true
    \ifx\@empty#3\@empty\else
        \def\lst@prefix{#2}\let\lst@require\@empty
%    \end{macrocode}
% \ldots\ 并且对每个非空条目：确定别名，若尚未加载则
% 把它加入 |\lst@require|。
%    \begin{macrocode}
        \edef\lst@temp{\expandafter\zap@space#3 \@empty}%
        \lst@for\lst@temp\do{%
          \ifx\@empty##1\@empty\else \lstKV@OptArg[]{##1}{%
            #4[####1]{####2}%
            \@ifundefined{\@lst\lst@prefix @\lst@malias $\lst@oalias}%
            {\edef\lst@require{\lst@require,\lst@malias $\lst@oalias}}%
            {}}%
          \fi}%
%    \end{macrocode}
% 初始化并输入文件，如必要时并持续进行。
%    \begin{macrocode}
        \global\let\lst@loadaspects\@empty
        \lst@InputCatcodes
        \ifx\lst@require\@empty\else
            \lst@for{#5}\do{%
                \ifx\lst@require\@empty\else
                    \InputIfFileExists{##1}{}{}%
                \fi}%
        \fi
%    \end{macrocode}
% 如果有某些条目无法加载，发出错误并调用 |\lst@false|
% （在关闭局部分组之后）。
%    \begin{macrocode}
        \ifx\lst@require\@empty\else
            \PackageError{Listings}{Couldn't load requested #1}%
            {The following #1s weren't loadable:^^J\@spaces
             \lst@require^^JThis may cause errors in the sequel.}%
            \aftergroup\lst@false
        \fi
%    \end{macrocode}
% 请求方面。
%    \begin{macrocode}
        \ifx\lst@loadaspects\@empty\else
            \lst@RequireAspects\lst@loadaspects
        \fi
    \fi
    \endgroup}
%    \end{macrocode}
% \end{macro}
%
% \begin{macro}{\lst@IfRequired}
% 使用 |\lst@IfOneOf| 并向 \meta{then} 部分添加一些代码：
% 从列表中删除现已加载的条目，并定义
% |\lst|\meta{prefix}|@|\meta{feature}|$|\meta{sub}.
%    \begin{macrocode}
\def\lst@IfRequired[#1]#2{%
    \lst@NormedDef\lst@temp{[#1]#2}%
    \expandafter\lst@IfRequired@\lst@temp\relax}
\def\lst@IfRequired@[#1]#2\relax#3{%
    \lst@IfOneOf #2$#1\relax\lst@require
        {\lst@DeleteKeysIn@\lst@require#2$#1,\relax,%
         \global\expandafter\let
             \csname\@lst\lst@prefix @#2$#1\endcsname\@empty
         #3}}
%    \end{macrocode}
% \end{macro}
%
% \begin{macro}{\lst@require}
%    \begin{macrocode}
\let\lst@require\@empty
%    \end{macrocode}
% \end{macro}
%
% \begin{macro}{\lst@NoAlias}
% 只是定义 |\lst@oalias| 和 |\lst@malias|。
%    \begin{macrocode}
\def\lst@NoAlias[#1]#2{%
    \lst@NormedDef\lst@oalias{#1}\lst@NormedDef\lst@malias{#2}}
%    \end{macrocode}
% \end{macro}
%
% \begin{macro}{\lst@LAS}
%    \begin{macrocode}
\gdef\lst@LAS#1#2#3#4#5#6#7{%
    \lst@Require{#1}{#2}{#3}#4#5%
    #4#3%
    \@ifundefined{lst#2@\lst@malias$\lst@oalias}%
        {\PackageError{Listings}%
         {#1 \ifx\@empty\lst@oalias\else \lst@oalias\space of \fi
          \lst@malias\space undefined}%
         {The #1 is not loadable. \@ehc}}%
        {#6\csname\@lst#2@\lst@malias $\lst@oalias\endcsname #7}}
%    \end{macrocode}
% \end{macro}
%
% \begin{macro}{\lst@RequireAspects}
% \begin{macro}{\lstloadaspects}
% 利用刚刚开发的定义。
%    \begin{macrocode}
\def\lst@RequireAspects#1{%
    \lst@Require{aspect}{asp}{#1}\lst@NoAlias\lstaspectfiles}
\let\lstloadaspects\lst@RequireAspects
%    \end{macrocode}
% \end{macro}
% \end{macro}
%
% \begin{macro}{\lstaspectfiles}
% 该宏被定义当且仅当它尚未定义。
%    \begin{macrocode}
\@ifundefined{lstaspectfiles}
    {\newcommand\lstaspectfiles{lstmisc0.sty,lstmisc.sty}}{}
%    \end{macrocode}
% \end{macro}
%
% \begin{macro}{\lst@DefDriver}
% 测试下一个字符并重新插入参数。
%    \begin{macrocode}
\gdef\lst@DefDriver#1#2#3#4{%
    \@ifnextchar[{\lst@DefDriver@{#1}{#2}#3#4}%
                 {\lst@DefDriver@{#1}{#2}#3#4[]}}
%    \end{macrocode}
% 如果条目已被请求，我们在局部把 |\lst@if| 设为真。
%    \begin{macrocode}
\gdef\lst@DefDriver@#1#2#3#4[#5]#6{%
    \def\lst@name{#1}\let\lst@if#4%
    \lst@NormedDef\lst@driver{\@lst#2@#6$#5}%
    \lst@IfRequired[#5]{#6}{\begingroup \lst@true}%
                           {\begingroup}%
    \lst@setcatcodes
    \@ifnextchar[{\lst@XDefDriver{#1}#3}{\lst@DefDriver@@#3}}
%    \end{macrocode}
% 注意 |\lst@XDefDriver| 取可选的`基'参数，但最终
% 调用 |\lst@DefDriver@@|。我们（在需要时）定义该条目，而
% |\endgroup| 重置一些 catcode 和 |\lst@if|，即 |\lst@XXDefDriver| 知道
% 自己是由公共命令还是内部命令调用的。
%    \begin{macrocode}
\gdef\lst@DefDriver@@#1#2{%
    \lst@if
        \global\@namedef{\lst@driver}{#1{#2}}%
    \fi
    \endgroup
    \@ifnextchar[\lst@XXDefDriver\@empty}
%    \end{macrocode}
% 我们取得方面参数，并且（若非空）在由公共命令调用时
% 立即加载这些方面，或扩充所需方面的列表，或
% 在该条目保持未定义时直接忽略该参数。
%    \begin{macrocode}
\gdef\lst@XXDefDriver[#1]{%
    \ifx\@empty#1\@empty\else
        \lst@if
            \lstloadaspects{#1}%
        \else
            \@ifundefined{\lst@driver}{}%
            {\xdef\lst@loadaspects{\lst@loadaspects,#1}}%
        \fi
    \fi}
%    \end{macrocode}
% 我们插入一个额外的`also' key=value 对。
%    \begin{macrocode}
\gdef\lst@XDefDriver#1#2[#3]#4#5{\lst@DefDriver@@#2{also#1=[#3]#4,#5}}
%    \end{macrocode}
% \end{macro}
%
%
% \subsection{aspect 命令}\label{iAspectCommands}
%
% 本节包含用于定义`\lst-方面'的命令。
% \begin{macro}{\lst@UserCommand}
% 主要等价于 |\gdef|。
%    \begin{macrocode}
%<!info>\let\lst@UserCommand\gdef
%<info>\def\lst@UserCommand#1{\message{\string#1,}\gdef#1}
%    \end{macrocode}
% \end{macro}
%
% \begin{macro}{\lst@BeginAspect}
% 一个直接了当的实现：
%    \begin{macrocode}
\newcommand*\lst@BeginAspect[2][]{%
    \def\lst@curraspect{#2}%
    \ifx \lst@curraspect\@empty
        \expandafter\lst@GobbleAspect
    \else
%    \end{macrocode}
% 如果 \meta{aspect name} 非空，则还有一些其他条件不
% 定义该方面（如第~\ref{dHowToDefineLstAspects} 节所述）。
%    \begin{macrocode}
%<!info>        \let\lst@next\@empty
%<info>        \def\lst@next{%
%<info>            \message{^^JDefine lst-aspect `#2':}\lst@GetAllocs}%
        \lst@IfRequired[]{#2}%
            {\lst@RequireAspects{#1}%
             \lst@if\else \let\lst@next\lst@GobbleAspect \fi}%
            {\let\lst@next\lst@GobbleAspect}%
        \expandafter\lst@next
    \fi}
%    \end{macrocode}
% \end{macro}
%
% \begin{macro}{\lst@EndAspect}
% 结束一个方面的定义。
%    \begin{macrocode}
\def\lst@EndAspect{%
    \csname\@lst patch@\lst@curraspect\endcsname
%<info>    \lst@ReportAllocs
    \let\lst@curraspect\@empty}
%    \end{macrocode}
% \end{macro}
%
% \begin{macro}{\lst@GobbleAspect}
% 丢弃直到下一个 |\lst@EndAspect| 的所有代码。
%    \begin{macrocode}
\long\def\lst@GobbleAspect#1\lst@EndAspect{\let\lst@curraspect\@empty}
%    \end{macrocode}
% \end{macro}
%
% \begin{macro}{\lst@Key}
% 该命令只是定义键。但我们必须处理一个可选
% 参数以及初始化参数 |#2|。
%    \begin{macrocode}
\def\lst@Key#1#2{%
%<info>    \message{#1,}%
    \@ifnextchar[{\lstKV@def{#1}{#2}}%
                 {\def\lst@temp{\lst@Key@{#1}{#2}}
                  \afterassignment\lst@temp
                  \global\@namedef{KV@\@lst @#1}####1}}
%    \end{macrocode}
% 现在来一个改名并修改过的、来自 \packagename{keyval} 宏包的副本：
% 我们需要全局的键定义。
%    \begin{macrocode}
\def\lstKV@def#1#2[#3]{%
    \global\@namedef{KV@\@lst @#1@default\expandafter}\expandafter
        {\csname KV@\@lst @#1\endcsname{#3}}%
    \def\lst@temp{\lst@Key@{#1}{#2}}\afterassignment\lst@temp
    \global\@namedef{KV@\@lst @#1}##1}
%    \end{macrocode}
% 如果 |#2| 的第一个记号不是 |\relax|，我们就初始化该键。
%    \begin{macrocode}
\def\lst@Key@#1#2{%
    \ifx\relax#2\@empty\else
        \begingroup \globaldefs\@ne
        \csname KV@\@lst @#1\endcsname{#2}%
        \endgroup
    \fi}
%    \end{macrocode}
% \end{macro}
%
% \begin{macro}{\lst@UseHook}
% 非常、非常、\ldots、非常（几百倍地）容易。
%    \begin{macrocode}
\def\lst@UseHook#1{\csname\@lst hk@#1\endcsname}
%    \end{macrocode}
% \end{macro}
%
% \begin{macro}{\lst@AddToHook}
% \begin{macro}{\lst@AddToHookExe}
% \begin{macro}{\lst@AddToHookAtTop}
% 都使用同一个子宏。
%    \begin{macrocode}
\def\lst@AddToHook{\lst@ATH@\iffalse\lst@AddTo}
\def\lst@AddToHookExe{\lst@ATH@\iftrue\lst@AddTo}
\def\lst@AddToHookAtTop{\lst@ATH@\iffalse\lst@AddToAtTop}
%    \end{macrocode}
% 当且仅当布尔值为真时，钩子材料被
% 全局执行。
%    \begin{macrocode}
\long\def\lst@ATH@#1#2#3#4{%
    \@ifundefined{\@lst hk@#3}{%
%<info>        \message{^^Jnew hook `#3',^^J}%
        \expandafter\gdef\csname\@lst hk@#3\endcsname{}}{}%
    \expandafter#2\csname\@lst hk@#3\endcsname{#4}%
    \def\lst@temp{#4}%
    #1% \iftrue|false
        \begingroup \globaldefs\@ne \lst@temp \endgroup
    \fi}
%    \end{macrocode}
% \end{macro}
% \end{macro}
% \end{macro}
%
% \begin{macro}{\lst@AddTo}
% 注意该定义是全局的！
%    \begin{macrocode}
\long\def\lst@AddTo#1#2{%
    \expandafter\gdef\expandafter#1\expandafter{#1#2}}
%    \end{macrocode}
% \end{macro}
%
% \begin{macro}{\lst@AddToAtTop}
% 现在我们需要若干个 |\expandafter|。只需注意，在
%   {\small\hspace*{2em}|\expandafter\gdef\expandafter#1\expandafter{\lst@temp|
%    $\langle$\textit{contents of }|#1|$\rangle$|}|}\\
% `第一阶段'展开之后，我们有以下内容。
%    \begin{macrocode}
\def\lst@AddToAtTop#1#2{\def\lst@temp{#2}%
    \expandafter\expandafter\expandafter\gdef
    \expandafter\expandafter\expandafter#1%
    \expandafter\expandafter\expandafter{\expandafter\lst@temp#1}}
%    \end{macrocode}
% \end{macro}
%
% \begin{macro}{\lst@lAddTo}
% |\lst@AddTo| 的局部版本\ldots
%    \begin{macrocode}
\def\lst@lAddTo#1#2{\expandafter\def\expandafter#1\expandafter{#1#2}}
%    \end{macrocode}
% \end{macro}
%
% \begin{macro}{\lst@Extend}
% \begin{macro}{\lst@lExtend}
% \ldots\ 这里我们先展开第二个参数的第一个记号。
%    \begin{macrocode}
\def\lst@Extend#1#2{%
    \expandafter\lst@AddTo\expandafter#1\expandafter{#2}}
\def\lst@lExtend#1#2{%
    \expandafter\lst@lAddTo\expandafter#1\expandafter{#2}}
%    \end{macrocode}
% \begin{TODO}
% 这绝不应改成
%    \begin{verbatim}
%    \def\lst@Extend#1{%
%        \expandafter\lst@AddTo\expandafter#1\expandafter}
%    \def\lst@lExtend#1{%
%        \expandafter\lst@lAddTo\expandafter#1}\end{verbatim}
% 前者在第二个参数是单个（= 未加花括号的）控制序列时不等价，
% 后者在第二个参数加了花括号时不等价。
% \end{TODO}
% \end{macro}
% \end{macro}
%
%
% \subsection{与 \packagename{keyval} 的接口}
%
% \packagename{keyval} 宏包通过唯一的一个参数 |#1| 把值传给
% 键宏的定义部分。下面的命令可以用来分析
% 该值。注意我们至少需要 \packagename{keyval} 宏包的 1.10 版。
% 另请注意，该宏包消除了一项与 AMS 文档类的命名冲突——
% 由 \lsthelper{Ralf~Quast}{1998/01/08}{关键词与 AMS 类冲突}。
% \begingroup
%    \begin{macrocode}
\RequirePackage{keyval}[1997/11/10]
%    \end{macrocode}
% \endgroup
%
% \begin{macro}{\lstKV@TwoArg}
% \begin{macro}{\lstKV@ThreeArg}
% \begin{macro}{\lstKV@FourArg}
% 定义临时宏并用给定参数 |#1| 调用。我们补充空的
% 参数，以应对用户提供的参数不足的情况。
%    \begin{macrocode}
\def\lstKV@TwoArg#1#2{\gdef\@gtempa##1##2{#2}\@gtempa#1{}{}}
\def\lstKV@ThreeArg#1#2{\gdef\@gtempa##1##2##3{#2}\@gtempa#1{}{}{}}
\def\lstKV@FourArg#1#2{\gdef\@gtempa##1##2##3##4{#2}\@gtempa#1{}{}{}{}}
%    \end{macrocode}
% 有一个问题：这些全局定义有什么用？|\lst@Key|
% 可能把 |\globaldefs| 设为 1，并且可能调用这个宏。这就是
% 我们在这里和下面使用全局定义的原因。
% \end{macro}
% \end{macro}
% \end{macro}
%
% \begin{macro}{\lstKV@OptArg}
% 我们定义临时宏 |\@gtempa|，并在必要时
% 插入默认参数。
%    \begin{macrocode}
\def\lstKV@OptArg[#1]#2#3{%
    \gdef\@gtempa[##1]##2{#3}\lstKV@OptArg@{#1}#2\@}
\def\lstKV@OptArg@#1{\@ifnextchar[\lstKV@OptArg@@{\lstKV@OptArg@@[#1]}}
\def\lstKV@OptArg@@[#1]#2\@{\@gtempa[#1]{#2}}
%    \end{macrocode}
% \end{macro}
%
% \begin{macro}{\lstKV@XOptArg}
% 这里 |#3| 已经是一个至少有两个参数的定义，其第一个
% 参数括在方括号中。
%    \begin{macrocode}
\def\lstKV@XOptArg[#1]#2#3{%
    \global\let\@gtempa#3\lstKV@OptArg@{#1}#2\@}
%    \end{macrocode}
% \end{macro}
%
% \begin{macro}{\lstKV@CSTwoArg}
% 只定义临时宏并调用它。
%    \begin{macrocode}
\def\lstKV@CSTwoArg#1#2{%
    \gdef\@gtempa##1,##2,##3\relax{#2}%
    \@gtempa#1,,\relax}
%    \end{macrocode}
% \end{macro}
%
% \begin{macro}{\lstKV@SetIf}
% 我们只测试 |#1| 的小写首字符。
%    \begin{macrocode}
\def\lstKV@SetIf#1{\lstKV@SetIf@#1\relax}
\def\lstKV@SetIf@#1#2\relax#3{\lowercase{%
    \expandafter\let\expandafter#3%
        \csname if\ifx #1t}true\else false\fi\endcsname}
%    \end{macrocode}
% \end{macro}
%
% \begin{macro}{\lstKV@SwitchCases}
% 按子串测试实现。原始版本使用了一个 |&|，
% 这引起了缺陷——见第~\pageref{gAmpersand} 页。
%    \begin{macrocode}
\def\lstKV@SwitchCases#1#2#3{%
    \def\lst@temp##1\\#1:##2\\##3##4\@nil{%
        \ifx\@empty##3%
            #3%
        \else
            ##2%
        \fi
    }%
    \lst@temp\\#2\\#1:\\\@empty\@nil}
%    \end{macrocode}
% \end{macro}
%
% \begin{macro}{\lstset}
% 最后是这个主用户接口宏。
% 我们为读取参数而改变 catcode。
%    \begin{macrocode}
\lst@UserCommand\lstset{\begingroup \lst@setcatcodes \lstset@}
\def\lstset@#1{\endgroup \ifx\@empty#1\@empty\else\setkeys{lst}{#1}\fi}
%    \end{macrocode}
% \end{macro}
%
% \begin{macro}{\lst@setcatcodes}
% 包含 |\lstset| 的所有 catcode 改动。等号是在
% \lsthelper{Bekir~Karaoglu}{2003/09/16}{带 [turkish] 的
% {babel} 时的 keyval 问题}报告缺陷之后添加的——babel 的活动等号与 keyval 的
% 用法冲突。|\catcode`\"=12\relax| 是在
% \lsthelper{Heiko~Bauke}{2004/06/27}{listings 与 ngerman}报告缺陷之后移除的
% ——但愿这不会引入其他缺陷。
%    \begin{macrocode}
\def\lst@setcatcodes{\makeatletter \catcode`\==12\relax}
%    \end{macrocode}
% \begin{TODO}
% 还要改变更多 catcode 吗？
% \end{TODO}
% \end{macro}
%
%
% \subsection{内部模式}
%
% \begin{macro}{\lst@NewMode}
% 我们只是用 |\chardef| 定义模式。计数器 |\lst@mode|
% 主要保存当前模式号。但它也用于在宏 |\lst@newmode| 中推进
% 编号——我们不想再浪费一个计数器。
%    \begin{macrocode}
\def\lst@NewMode#1{%
    \ifx\@undefined#1%
        \lst@mode\lst@newmode\relax \advance\lst@mode\@ne
        \xdef\lst@newmode{\the\lst@mode}%
        \global\chardef#1=\lst@mode
        \lst@mode\lst@nomode
    \fi}
%    \end{macrocode}
% \end{macro}
%
% \begin{macro}{\lst@mode}
% \begin{macro}{\lst@nomode}
% 我们分配计数器和第一个模式。
%    \begin{macrocode}
\newcount\lst@mode
\def\lst@newmode{\m@ne}% init
\lst@NewMode\lst@nomode % init (of \lst@mode :-)
%    \end{macrocode}
% \end{macro}
% \end{macro}
%
% \begin{macro}{\lst@UseDynamicMode}
% 对于动态模式，我们绝不能使用计数器 |\lst@mode|（因为它可能
% 已经被赋值）。|\lst@dynamicmode| 替代 |\lst@newmode|，在这里是一个局部
% 定义，\ldots
%    \begin{macrocode}
\def\lst@UseDynamicMode{%
    \@tempcnta\lst@dynamicmode\relax \advance\@tempcnta\@ne
    \edef\lst@dynamicmode{\the\@tempcnta}%
    \expandafter\lst@Swap\expandafter{\expandafter{\lst@dynamicmode}}}
%    \end{macrocode}
% \ldots\ 每次列表时都用 |\lst@newmode| 的当前`值'初始化。
%    \begin{macrocode}
\lst@AddToHook{InitVars}{\let\lst@dynamicmode\lst@newmode}
%    \end{macrocode}
% \end{macro}
%
% \begin{macro}{\lst@EnterMode}
% 每个模式打开一个分组层级，保存模式号，并执行该模式
% 特定的记号。此外我们（局部地）记录所有这些改动，并在
% 用户需要时调整内部变量。
%    \begin{macrocode}
\def\lst@EnterMode#1#2{%
    \bgroup \lst@mode=#1\relax #2%
    \lst@FontAdjust
    \lst@lAddTo\lst@entermodes{\lst@EnterMode{#1}{#2}}}
%    \end{macrocode}
%    \begin{macrocode}
\lst@AddToHook{InitVars}{\let\lst@entermodes\@empty}
\let\lst@entermodes\@empty % init
%    \end{macrocode}
% 该初始化是在 \lsthelper{Herfried~Karl~Wagner}{2002/05/11}{未定义的控制序列
% \lst@entermodes}.
% \end{macro}
%
% \begin{macro}{\lst@LeaveMode}
% 我们只是关闭分组，并当且仅当当前模式不是
% |\lst@nomode| 时调用 |\lsthk@EndGroup|。
%    \begin{macrocode}
\def\lst@LeaveMode{%
    \ifnum\lst@mode=\lst@nomode\else
        \egroup \expandafter\lsthk@EndGroup
    \fi}
%    \end{macrocode}
%    \begin{macrocode}
\lst@AddToHook{EndGroup}{}% init
%    \end{macrocode}
% \end{macro}
%
% \begin{macro}{\lst@InterruptModes}
% 我们把当前模式序列压入一个栈，并离开所有模式。
%    \begin{macrocode}
\def\lst@InterruptModes{%
    \lst@Extend\lst@modestack{\expandafter{\lst@entermodes}}%
    \lst@LeaveAllModes}
%    \end{macrocode}
%    \begin{macrocode}
\lst@AddToHook{InitVars}{\global\let\lst@modestack\@empty}
%    \end{macrocode}
% \end{macro}
%
% \begin{macro}{\lst@ReenterModes}
% 如果栈不为空，我们就离开所有模式并弹出最顶层的元素
% （它是 |\lst@modestack| 的最后一个元素）。
%    \begin{macrocode}
\def\lst@ReenterModes{%
    \ifx\lst@modestack\@empty\else
        \lst@LeaveAllModes
        \global\let\@gtempa\lst@modestack
        \global\let\lst@modestack\@empty
        \expandafter\lst@ReenterModes@\@gtempa\relax
    \fi}
\def\lst@ReenterModes@#1#2{%
    \ifx\relax#2\@empty
%    \end{macrocode}
% 如果已到达 |\relax|，我们也就找到了最后一个元素：执行 |#1|
% 并在 |\fi| 之后吃掉 |{#2}|=|{\relax}|。
%    \begin{macrocode}
        \gdef\@gtempa##1{#1}%
        \expandafter\@gtempa
    \else
%    \end{macrocode}
% 否则我们只把该元素加入 |\lst@modestack| 并继续循环。
%    \begin{macrocode}
        \lst@AddTo\lst@modestack{{#1}}%
        \expandafter\lst@ReenterModes@
    \fi
    {#2}}
%    \end{macrocode}
% \end{macro}
%
% \begin{macro}{\lst@LeaveAllModes}
% 离开所有模式意味着不断关闭分组，直到模式等于 |\lst@nomode|。
%    \begin{macrocode}
\def\lst@LeaveAllModes{%
    \ifnum\lst@mode=\lst@nomode
        \expandafter\lsthk@EndGroup
    \else
        \expandafter\egroup\expandafter\lst@LeaveAllModes
    \fi}
%    \end{macrocode}
% 我们需要这个宏来正确地结束列表。
%    \begin{macrocode}
\lst@AddToHook{ExitVars}{\lst@LeaveAllModes}
%    \end{macrocode}
% \end{macro}
%
% \begin{macro}{\lst@Pmode}
% \begin{macro}{\lst@GPmode}
% `处理'模式与通用模式。
%    \begin{macrocode}
\lst@NewMode\lst@Pmode
\lst@NewMode\lst@GPmode
%    \end{macrocode}
% \end{macro}
% \end{macro}
%
% \begin{macro}{\lst@modetrue}
% 赋值给布尔值的常规宏，区别是我们还会执行一个钩子。
%    \begin{macrocode}
\def\lst@modetrue{\let\lst@ifmode\iftrue \lsthk@ModeTrue}
\let\lst@ifmode\iffalse % init
\lst@AddToHook{ModeTrue}{}% init
%    \end{macrocode}
% \end{macro}
%
% \begin{macro}{\lst@ifLmode}
% 注释行使用静态模式。它在行尾终止。
%    \begin{macrocode}
\def\lst@Lmodetrue{\let\lst@ifLmode\iftrue}
\let\lst@ifLmode\iffalse % init
\lst@AddToHook{EOL}{\@whilesw \lst@ifLmode\fi \lst@LeaveMode}
%    \end{macrocode}
% \end{macro}
%
%
% \subsection{各种辅助宏}
%
% \begin{macro}{\lst@NormedDef}
% 像 |\def| 那样工作（没有任何参数！），但通过把所有
% 字符变成小写并去掉空格来规范化替换文本。
%    \begin{macrocode}
\def\lst@NormedDef#1#2{\lowercase{\edef#1{\zap@space#2 \@empty}}}
%    \end{macrocode}
% \end{macro}
%
% \begin{macro}{\lst@NormedNameDef}
% 像 |\global\@namedef| 那样工作（同样没有任何参数！），但会
% 同时规范化宏名和替换文本。
%    \begin{macrocode}
\def\lst@NormedNameDef#1#2{%
    \lowercase{\edef\lst@temp{\zap@space#1 \@empty}%
    \expandafter\xdef\csname\lst@temp\endcsname{\zap@space#2 \@empty}}}
%    \end{macrocode}
% \end{macro}
%
% \begin{macro}{\lst@GetFreeMacro}
% 初始化 |\@tempcnta| 和 |\lst@freemacro|，\ldots
%    \begin{macrocode}
\def\lst@GetFreeMacro#1{%
    \@tempcnta\z@ \def\lst@freemacro{#1\the\@tempcnta}%
    \lst@GFM@}
%    \end{macrocode}
% \ldots\ 然后或者构造控制序列，或者推进计数器并
% 继续。
%    \begin{macrocode}
\def\lst@GFM@{%
    \expandafter\ifx \csname\lst@freemacro\endcsname \relax
        \edef\lst@freemacro{\csname\lst@freemacro\endcsname}%
    \else
        \advance\@tempcnta\@ne
        \expandafter\lst@GFM@
    \fi}
%    \end{macrocode}
% \end{macro}
%
% \begin{macro}{\lst@gtempboxa}
%    \begin{macrocode}
\newbox\lst@gtempboxa
%    \end{macrocode}
%    \begin{macrocode}
%</kernel>
%    \end{macrocode}
% \end{macro}
%
%
% \section{输出}
%
%
% \subsection{基本寄存器与键}
%
%    \begin{macrocode}
%<*kernel>
%    \end{macrocode}
%
% \paragraph{当前字符串}
% 保存在一个记号寄存器中，用一个计数器保存其长度。
% 这里我们定义把字符放入输出队列的宏。
%
% \begin{macro}{\lst@token}
% \begin{macro}{\lst@length}
% 在这里分配。相当有用的一条注释，不是吗？
%    \begin{macrocode}
\newtoks\lst@token \newcount\lst@length
%    \end{macrocode}
% \end{macro}
% \end{macro}
%
% \begin{macro}{\lst@ResetToken}
% \begin{macro}{\lst@lastother}
% 这两个寄存器在每行开头分别清空和清零。
% 在收到 \lsthelper{Claus~Atzenbeck}{1999/11/24}{HTML：
% 输出单元在 > 之后重复}的报告之后——我多次修复过这类缺陷——
% 我决定也在 \hookname{EndGroup} 钩子中重置这两个寄存器。
%    \begin{macrocode}
\def\lst@ResetToken{\lst@token{}\lst@length\z@}
%    \end{macrocode}
%    \begin{macrocode}
\lst@AddToHook{InitVarsBOL}{\lst@ResetToken \let\lst@lastother\@empty}
\lst@AddToHook{EndGroup}{\lst@ResetToken \let\lst@lastother\@empty}
%    \end{macrocode}
% 宏 |\lst@lastother| 将等价于最后一个`其他'字符，
% 这引导我们使用 |\lst@ifletter|。
% \end{macro}
% \end{macro}
%
% \begin{macro}{\lst@ifletter}
% 指示该记号包含的是标识符还是其他字符。
%    \begin{macrocode}
\def\lst@lettertrue{\let\lst@ifletter\iftrue}
\def\lst@letterfalse{\let\lst@ifletter\iffalse}
\lst@AddToHook{InitVars}{\lst@letterfalse}
%    \end{macrocode}
% \end{macro}
%
% \begin{macro}{\lst@Append}
% 把参数放入输出队列。
%    \begin{macrocode}
\def\lst@Append#1{\advance\lst@length\@ne
                  \lst@token=\expandafter{\the\lst@token#1}}
%    \end{macrocode}
% \end{macro}
%
% \begin{macro}{\lst@AppendOther}
% 根据当前状态，我们先把这个字符串作为标识符输出。
% 然后用 |\futurelet| 保存`参数'，并调用宏
% |\lst@Append| 完成其余工作。
%    \begin{macrocode}
\def\lst@AppendOther{%
    \lst@ifletter \lst@Output\lst@letterfalse \fi
    \futurelet\lst@lastother\lst@Append}
%    \end{macrocode}
% \end{macro}
%
% \begin{macro}{\lst@AppendLetter}
% 必要时我们输出非标识符字符串，并调用 |\lst@Append|。
%    \begin{macrocode}
\def\lst@AppendLetter{%
    \lst@ifletter\else \lst@OutputOther\lst@lettertrue \fi
    \lst@Append}
%    \end{macrocode}
% \end{macro}
%
% \begin{macro}{\lst@SaveToken}
% \begin{macro}{\lst@RestoreToken}
% 如果出现分组结束并破坏了字符串，我们可以用这些
% 宏来保存和恢复内容。|\lst@thestyle| 是当前的
% 打印样式，也必须保存和恢复。
%    \begin{macrocode}
\def\lst@SaveToken{%
    \global\let\lst@gthestyle\lst@thestyle
    \global\let\lst@glastother\lst@lastother
    \xdef\lst@RestoreToken{\noexpand\lst@token{\the\lst@token}%
                           \noexpand\lst@length\the\lst@length\relax
                           \noexpand\let\noexpand\lst@thestyle
                                        \noexpand\lst@gthestyle
                           \noexpand\let\noexpand\lst@lastother
                                        \noexpand\lst@glastother}}
%    \end{macrocode}
% 现在——也就是说在 \lsthelper{Rolf~Niepraschk}
% {2002/04/12}{|\RequirePackage| 在页面顶部附近时缺少 keywordstyle}
% 报告缺陷之后——|\lst@lastother| 也会被保存和恢复。
% \end{macro}
% \end{macro}
%
% \begin{macro}{\lst@IfLastOtherOneOf}
% 最后是这个显而易见的实现。
%    \begin{macrocode}
\def\lst@IfLastOtherOneOf#1{\lst@IfLastOtherOneOf@ #1\relax}
\def\lst@IfLastOtherOneOf@#1{%
    \ifx #1\relax
        \expandafter\@secondoftwo
    \else
        \ifx\lst@lastother#1%
            \lst@IfLastOtherOneOf@t
        \else
            \expandafter\expandafter\expandafter\lst@IfLastOtherOneOf@
        \fi
    \fi}
\def\lst@IfLastOtherOneOf@t#1\fi\fi#2\relax{\fi\fi\@firstoftwo}
%    \end{macrocode}
% \end{macro}
%
%
% \paragraph{当前位置}
% 要么是尺寸 |\lst@currlwidth|，即不考虑当前字符串的
% 水平位置，要么是从 0 开始的当前列号。这是
% |\lst@column| $-$ |\lst@pos| $+$ |\lst@length|。
% 此外还有 |\lst@lostspace|，即当前行宽与期望行宽
% 之差。我们定义宏来插入
% 这段丢失的空间。
%
% \begin{macro}{\lst@currlwidth}
% \begin{macro}{\lst@column}
% \begin{macro}{\lst@pos}
% 当前行宽和两个计数器。
%    \begin{macrocode}
\newdimen\lst@currlwidth % \global
\newcount\lst@column \newcount\lst@pos % \global
\lst@AddToHook{InitVarsBOL}
    {\global\lst@currlwidth\z@ \global\lst@pos\z@ \global\lst@column\z@}
%    \end{macrocode}
% \end{macro}
% \end{macro}
% \end{macro}
%
% \begin{macro}{\lst@CalcColumn}
% 把 |\@tempcnta| 设为当前列号。
% 注意 |\lst@pos| 将是非正的。
%    \begin{macrocode}
\def\lst@CalcColumn{%
            \@tempcnta\lst@column
    \advance\@tempcnta\lst@length
    \advance\@tempcnta-\lst@pos}
%    \end{macrocode}
% \end{macro}
%
% \begin{macro}{\lst@lostspace}
% 只要这个尺寸为正，我们就可以插入空格。负的`丢失
% 空间'意味着已打印的行比期望的更宽。
%    \begin{macrocode}
\newdimen\lst@lostspace % \global
\lst@AddToHook{InitVarsBOL}{\global\lst@lostspace\z@}
%    \end{macrocode}
% \end{macro}
%
% \begin{macro}{\lst@UseLostSpace}
% 当且仅当 |\lst@lostspace| 为正时，我们插入空格并重置它。
%    \begin{macrocode}
\def\lst@UseLostSpace{\ifdim\lst@lostspace>\z@ \lst@InsertLostSpace \fi}
%    \end{macrocode}
% \end{macro}
%
% \begin{macro}{\lst@InsertLostSpace}
% \begin{macro}{\lst@InsertHalfLostSpace}
% 同上，但即使为负也插入。|\lst@Kern| 很快就定义。
%    \begin{macrocode}
\def\lst@InsertLostSpace{%
    \lst@Kern\lst@lostspace \global\lst@lostspace\z@}
\def\lst@InsertHalfLostSpace{%
    \global\lst@lostspace.5\lst@lostspace \lst@Kern\lst@lostspace}
%    \end{macrocode}
% \end{macro}
% \end{macro}
%
%
% \paragraph{列宽}
% 这里我们处理单列的宽度，它等于单
% 字符盒的宽度。请记住存在固定和灵活两种列
% 格式。
%
% \begin{macro}{\lst@width}
% \begin{lstkey}{basewidth}
% \keyname{basewidth} 把值赋给宏并测试它们是否
% 为负。
%    \begin{macrocode}
\newdimen\lst@width
\lst@Key{basewidth}{0.6em,0.45em}{\lstKV@CSTwoArg{#1}%
    {\def\lst@widthfixed{##1}\def\lst@widthflexible{##2}%
     \ifx\lst@widthflexible\@empty
         \let\lst@widthflexible\lst@widthfixed
     \fi
     \def\lst@temp{\PackageError{Listings}%
                                {Negative value(s) treated as zero}%
                                \@ehc}%
     \let\lst@error\@empty
     \ifdim \lst@widthfixed<\z@
         \let\lst@error\lst@temp \let\lst@widthfixed\z@
     \fi
     \ifdim \lst@widthflexible<\z@
         \let\lst@error\lst@temp \let\lst@widthflexible\z@
     \fi
     \lst@error}}
%    \end{macrocode}
% 我们在一个特殊钩子中设置该尺寸。
%    \begin{macrocode}
\lst@AddToHook{FontAdjust}
    {\lst@width=\lst@ifflexible\lst@widthflexible
                          \else\lst@widthfixed\fi \relax}
%    \end{macrocode}
% \end{lstkey}
% \end{macro}
%
% \begin{lstkey}{fontadjust}
% \begin{macro}{\lst@FontAdjust}
% 这个钩子由开关控制，并且总是在
% \hookname{InitVars}.
%    \begin{macrocode}
\lst@Key{fontadjust}{false}[t]{\lstKV@SetIf{#1}\lst@iffontadjust}
\def\lst@FontAdjust{\lst@iffontadjust \lsthk@FontAdjust \fi}
%    \end{macrocode}
%    \begin{macrocode}
\lst@AddToHook{InitVars}{\lsthk@FontAdjust}
%    \end{macrocode}
% \end{macro}
% \end{lstkey}
%
%
% \subsection{底层与中层输出}
%
% \paragraph{执行输出}
% 意味着把字符串放入盒子寄存器，更新所有内部
% 数据，并最终把盒子交给 \TeX。
%
% \begin{macro}{\lst@OutputBox}
% \begin{macro}{\lst@alloverstyle}
% 最底层是输出一个盒子寄存器。
% 这里我们以 |\box#1| 作为 |\lst@alloverstyle| 的参数。
%    \begin{macrocode}
\def\lst@OutputBox#1{\lst@alloverstyle{\box#1}}
%    \end{macrocode}
% \begin{ALTERNATIVE}
% 在两个定义 |\lst@Kern| 和 |\lst@CalcLostSpaceAndOutput| 中，我们可以
% 在这里推进该尺寸，而不是用 |\global\advance\lst@currlwidth| |\wd|\meta{盒子号}。
% 但我决定不这样做，因为它简化了对 |\lst@OutputBox| 可能的
% 重新定义：我们不必操心 |\lst@currlwidth|。
% \end{ALTERNATIVE}
%    \begin{macrocode}
\def\lst@alloverstyle#1{#1}% init
%    \end{macrocode}
% \end{macro}
% \end{macro}
%
% \begin{macro}{\lst@Kern}
% 用于插入`丢失空间'。
% 它绝不能使用 |\@tempboxa|，因为那个\ldots
%    \begin{macrocode}
\def\lst@Kern#1{%
    \setbox\z@\hbox{{\lst@currstyle{\kern#1}}}%
    \global\advance\lst@currlwidth \wd\z@
    \lst@OutputBox\z@}
%    \end{macrocode}
% \end{macro}
%
% \begin{macro}{\lst@CalcLostSpaceAndOutput}
% \ldots\ 在这里被使用。
% 我们跟踪 |\lst@lostspace|、|\lst@currlwidth| 和 |\lst@pos|。
%    \begin{macrocode}
\def\lst@CalcLostSpaceAndOutput{%
    \global\advance\lst@lostspace \lst@length\lst@width
    \global\advance\lst@lostspace-\wd\@tempboxa
    \global\advance\lst@currlwidth \wd\@tempboxa
    \global\advance\lst@pos -\lst@length
%    \end{macrocode}
% 在输出 |\@tempboxa| 之前，如果有足够的丢失空间，我们就插入空格。
% 这可能会通过`插入一半丢失空间'调用 |\lst@Kern|，这正是
% 我们上面不能使用 |\@tempboxa| 的原因。通过重新定义，我们防止
% |\lst@OutputBox| 在 |\lst@Kern| 中使用任何特殊样式。
%    \begin{macrocode}
    \setbox\@tempboxa\hbox{\let\lst@OutputBox\box
        \ifdim\lst@lostspace>\z@ \lst@leftinsert \fi
        \box\@tempboxa
        \ifdim\lst@lostspace>\z@ \lst@rightinsert \fi}%
%    \end{macrocode}
% 最后我们可以输出新的盒子。
%    \begin{macrocode}
    \lst@OutputBox\@tempboxa \lsthk@PostOutput}
%    \end{macrocode}
%    \begin{macrocode}
\lst@AddToHook{PostOutput}{}% init
%    \end{macrocode}
% \end{macro}
%
% \begin{macro}{\lst@OutputToken}
% 现在来一个中层定义。
% 这里我们用 |\lst@token| 设置 |\@tempboxa| 并最终输出盒子。
% 我们处理字体调整和特殊输出样式。
% 尚未知道的宏将在下面的小节中定义。
%    \begin{macrocode}
\def\lst@OutputToken{%
    \lst@TrackNewLines \lst@OutputLostSpace
    \lst@ifgobbledws
        \lst@gobbledwhitespacefalse
        \lst@@discretionary
    \fi
    \lst@CheckMerge
    {\lst@thestyle{\lst@FontAdjust
     \setbox\@tempboxa\lst@hbox
        {\lsthk@OutputBox
         \lst@lefthss
         \expandafter\lst@FillOutputBox\the\lst@token\@empty
         \lst@righthss}%
     \lst@CalcLostSpaceAndOutput}}%
    \lst@ResetToken}
%    \end{macrocode}
%    \begin{macrocode}
\lst@AddToHook{OutputBox}{}% init
%    \end{macrocode}
%    \begin{macrocode}
\def\lst@gobbledwhitespacetrue{\global\let\lst@ifgobbledws\iftrue}
\def\lst@gobbledwhitespacefalse{\global\let\lst@ifgobbledws\iffalse}
\lst@AddToHookExe{InitBOL}{\lst@gobbledwhitespacefalse}% init
%    \end{macrocode}
% \end{macro}
%
%
% \paragraph{延迟输出}
% 意味着把字符串保存到某处，并在必要时推回。我们也可以
% 把字符串附加到下一个输出盒上，而不
% 影响样式检测：两者都会以随后输出的样式
% 打印。我们把这称为`合并'。
%
% \begin{macro}{\lst@Delay}
% \begin{macro}{\lst@Merge}
% 为了延迟或合并 |#1|，我们照常处理它，并只把状态
% 保存在宏里。对于延迟字符，我们还需要当前`活动的'
% 输出例程。两个定义都先检查是否已经存在
% 延迟或`合并'的字符。
%    \begin{macrocode}
\def\lst@Delay#1{%
    \lst@CheckDelay
    #1%
    \lst@GetOutputMacro\lst@delayedoutput
    \edef\lst@delayed{\the\lst@token}%
    \edef\lst@delayedlength{\the\lst@length}%
    \lst@ResetToken}
%    \end{macrocode}
%    \begin{macrocode}
\def\lst@Merge#1{%
    \lst@CheckMerge
    #1%
    \edef\lst@merged{\the\lst@token}%
    \edef\lst@mergedlength{\the\lst@length}%
    \lst@ResetToken}
%    \end{macrocode}
% \end{macro}
% \end{macro}
%
% \begin{macro}{\lst@MergeToken}
% 这里我们把各部分重新拼起来。
%    \begin{macrocode}
\def\lst@MergeToken#1#2{%
    \advance\lst@length#2%
    \lst@lExtend#1{\the\lst@token}%
    \expandafter\lst@token\expandafter{#1}%
    \let#1\@empty}
%    \end{macrocode}
% \end{macro}
%
% \begin{macro}{\lst@CheckDelay}
% 我们需要打印延迟字符。模式取决于当前输出
% 宏。如果它等于保存的定义，我们就把延迟字符放在
% 字符串前面（即合并它们），因为没有发生
% 字母与其他之间或其他与字母之间的跳变。否则我们局部重置
% 当前字符串，把这个空字符串与延迟的合并，
% 然后输出它。
%    \begin{macrocode}
\def\lst@CheckDelay{%
    \ifx\lst@delayed\@empty\else
        \lst@GetOutputMacro\@gtempa
        \ifx\lst@delayedoutput\@gtempa
            \lst@MergeToken\lst@delayed\lst@delayedlength
        \else
            {\lst@ResetToken
             \lst@MergeToken\lst@delayed\lst@delayedlength
             \lst@delayedoutput}%
            \let\lst@delayed\@empty
        \fi
    \fi}
%    \end{macrocode}
% \end{macro}
%
% \begin{macro}{\lst@CheckMerge}
% 对 |\lst@merged| 来说这一切更简单。
%    \begin{macrocode}
\def\lst@CheckMerge{%
    \ifx\lst@merged\@empty\else
        \lst@MergeToken\lst@merged\lst@mergedlength
    \fi}
%    \end{macrocode}
%    \begin{macrocode}
\let\lst@delayed\@empty % init
\let\lst@merged\@empty % init
%    \end{macrocode}
% \end{macro}
%
%
% \subsection{列格式}
%
% 现在该处理固定和灵活列模式了。
% 一些尚未展开的定义现在被填上。
%
% \begin{macro}{\lst@column@fixed}
% 切换到固定列格式。这里的定义控制上述定义的
% 输出看起来是什么样。
%    \begin{macrocode}
\def\lst@column@fixed{%
    \lst@flexiblefalse
    \lst@width\lst@widthfixed\relax
    \let\lst@OutputLostSpace\lst@UseLostSpace
    \let\lst@FillOutputBox\lst@FillFixed
    \let\lst@hss\hss
    \def\lst@hbox{\hbox to\lst@length\lst@width}}
%    \end{macrocode}
% \end{macro}
%
% \begin{macro}{\lst@FillFixed}
% 填满一个固定模式盒很容易。
%    \begin{macrocode}
\def\lst@FillFixed#1{#1\lst@FillFixed@}
%    \end{macrocode}
% 在未到达末尾（来自上方的 |\@empty|）时，我们插入动态空格，
% 输出该参数并再次调用子宏。
%    \begin{macrocode}
\def\lst@FillFixed@#1{%
    \ifx\@empty#1\else \lst@hss#1\expandafter\lst@FillFixed@ \fi}
%    \end{macrocode}
% \end{macro}
%
% \begin{macro}{\lst@column@flexible}
% 第一个灵活格式。
%    \begin{macrocode}
\def\lst@column@flexible{%
    \lst@flexibletrue
    \lst@width\lst@widthflexible\relax
    \let\lst@OutputLostSpace\lst@UseLostSpace
    \let\lst@FillOutputBox\@empty
    \let\lst@hss\@empty
    \let\lst@hbox\hbox}
%    \end{macrocode}
% \end{macro}
%
% \begin{macro}{\lst@column@fullflexible}
% 这个列格式除行首之外不插入任何丢失空间。
%    \begin{macrocode}
\def\lst@column@fullflexible{%
    \lst@column@flexible
    \def\lst@OutputLostSpace{\lst@ifnewline \lst@UseLostSpace\fi}%
    \let\lst@leftinsert\@empty
    \let\lst@rightinsert\@empty}
%    \end{macrocode}
% \end{macro}
%
% \begin{macro}{\lst@column@spaceflexible}
% 这个列格式只通过拉伸已有的（不可见）空格来插入丢失
% 空间；在原文没有空格的地方，它不会在标识符
% 与其他字符之间插入丢失空间。这
% 由 \lsthelper{Andrei~Alexandrescu}{2007-02-26}{-}提出。
%    \begin{macrocode}
\def\lst@column@spaceflexible{%
    \lst@column@flexible
    \def\lst@OutputLostSpace{%
      \lst@ifwhitespace
        \ifx\lst@outputspace\lst@visiblespace
        \else
          \lst@UseLostSpace
        \fi
      \else
        \lst@ifnewline \lst@UseLostSpace\fi
      \fi}%
    \let\lst@leftinsert\@empty
    \let\lst@rightinsert\@empty}
%    \end{macrocode}
% \end{macro}
%
% 这样我们就有了这些列格式。现在我们定义使用它们的宏。
%
% \begin{macro}{\lst@outputpos}
% 这个宏设置`输出盒定位'参数（旧键
% \keyname{outputpos}）。我们测试 |l|、|c| 和 |r|。
% 固定格式使用 |\lst@lefthss| 和 |\lst@righthss|，而
% 灵活格式需要 |\lst@leftinsert| 和 |\lst@rightinsert|。
%    \begin{macrocode}
\def\lst@outputpos#1#2\relax{%
    \def\lst@lefthss{\lst@hss}\let\lst@righthss\lst@lefthss
    \let\lst@rightinsert\lst@InsertLostSpace
    \ifx #1c%
        \let\lst@leftinsert\lst@InsertHalfLostSpace
    \else\ifx #1r%
        \let\lst@righthss\@empty
        \let\lst@leftinsert\lst@InsertLostSpace
        \let\lst@rightinsert\@empty
    \else
        \let\lst@lefthss\@empty
        \let\lst@leftinsert\@empty
        \ifx #1l\else \PackageWarning{Listings}%
            {Unknown positioning for output boxes}%
        \fi
    \fi\fi}
%    \end{macrocode}
% \end{macro}
%
% \begin{macro}{\lst@ifflexible}
% 指示列模式，但不区分不同的固定
% 或灵活模式。
%    \begin{macrocode}
\def\lst@flexibletrue{\let\lst@ifflexible\iftrue}
\def\lst@flexiblefalse{\let\lst@ifflexible\iffalse}
%    \end{macrocode}
% \end{macro}
%
% \begin{lstkey}{columns}
% 这在这里完成：检查可选参数，然后构造列格式的
% 控制序列。
%    \begin{macrocode}
\lst@Key{columns}{[c]fixed}{\lstKV@OptArg[]{#1}{%
    \ifx\@empty##1\@empty\else \lst@outputpos##1\relax\relax \fi
    \expandafter\let\expandafter\lst@arg
                                \csname\@lst @column@##2\endcsname
%    \end{macrocode}
% 我们发出警告，或保存定义以备后用。
%    \begin{macrocode}
    \lst@arg
    \ifx\lst@arg\relax
        \PackageWarning{Listings}{Unknown column format `##2'}%
    \else
        \lst@ifflexible
            \let\lst@columnsflexible\lst@arg
        \else
            \let\lst@columnsfixed\lst@arg
        \fi
    \fi}}
%    \end{macrocode}
%    \begin{macrocode}
\let\lst@columnsfixed\lst@column@fixed % init
\let\lst@columnsflexible\lst@column@flexible % init
%    \end{macrocode}
% \end{lstkey}
%
% \begin{lstkey}{flexiblecolumns}
% 不过是切换最后使用的灵活与固定模式的一个键。
%    \begin{macrocode}
\lst@Key{flexiblecolumns}\relax[t]{%
    \lstKV@SetIf{#1}\lst@ifflexible
    \lst@ifflexible \lst@columnsflexible
              \else \lst@columnsfixed \fi}
%    \end{macrocode}
% \end{lstkey}
%
%
% \subsection{新行}
%
% \begin{macro}{\lst@newlines}
% 这个计数器保存我们需要执行的`新行'（cr+lf）数目。
%    \begin{macrocode}
\newcount\lst@newlines
\lst@AddToHook{InitVars}{\global\lst@newlines\z@}
\lst@AddToHook{InitVarsBOL}{\global\advance\lst@newlines\@ne}
%    \end{macrocode}
% \end{macro}
%
% \begin{macro}{\lst@NewLine}
% 我们这样开始新的一行：开始新段落并输出一个空
% 盒子。如果低级定义 |\lst@OutputBox| 只是吃掉盒子，
% 我们就不会开始新行。这用于丢弃整个输出。
%    \begin{macrocode}
\def\lst@NewLine{%
    \ifx\lst@OutputBox\@gobble\else
        \par\noindent \hbox{}%
    \fi
    \global\advance\lst@newlines\m@ne
    \lst@newlinetrue}
%    \end{macrocode}
% 定义 |\lst@newlinetrue| 并在输出之后重置。
%    \begin{macrocode}
\def\lst@newlinetrue{\global\let\lst@ifnewline\iftrue}
\lst@AddToHookExe{PostOutput}{\global\let\lst@ifnewline\iffalse}% init
%    \end{macrocode}
% \end{macro}
%
% \begin{macro}{\lst@TrackNewLines}
% 如果 |\lst@newlines| 为正，我们执行钩子并插入
% 新行。
%    \begin{macrocode}
\def\lst@TrackNewLines{%
    \ifnum\lst@newlines>\z@
        \lsthk@OnNewLine
        \lst@DoNewLines
    \fi}
\lst@AddToHook{OnNewLine}{}% init
%    \end{macrocode}
% \end{macro}
%
% \begin{lstkey}{emptylines}
% \lsthelper{Adam~Prugel-Bennett}{2001/02/19}{空行的间距}请求
% 这样一个键——如果我没有误解他的话。我们检查可选星号
% 并设置 |\lst@maxempty| 和开关。
%    \begin{macrocode}
\lst@Key{emptylines}\maxdimen{%
    \@ifstar{\lst@true\@tempcnta\@gobble#1\relax\lst@GobbleNil}%
            {\lst@false\@tempcnta#1\relax\lst@GobbleNil}#1\@nil
    \advance\@tempcnta\@ne
    \edef\lst@maxempty{\the\@tempcnta\relax}%
    \let\lst@ifpreservenumber\lst@if}
%    \end{macrocode}
% \end{lstkey}
%
% \begin{macro}{\lst@DoNewLines}
% 我们先处理 |\lst@maxempty|，然后处理剩余的空行。
%    \begin{macrocode}
\def\lst@DoNewLines{%
    \@whilenum\lst@newlines>\lst@maxempty \do
        {\lst@ifpreservenumber
            \lsthk@OnEmptyLine
            \global\advance\c@lstnumber\lst@advancelstnum
         \fi
         \global\advance\lst@newlines\m@ne}%
    \@whilenum \lst@newlines>\@ne \do
        {\lsthk@OnEmptyLine \lst@NewLine}%
    \ifnum\lst@newlines>\z@ \lst@NewLine \fi}
\lst@AddToHook{OnEmptyLine}{}% init
%    \end{macrocode}
% \end{macro}
%
%
% \subsection{高层输出}
%
% \begin{lstkey}{identifierstyle}
% 一个简单的键。
%    \begin{macrocode}
\lst@Key{identifierstyle}{}{\def\lst@identifierstyle{#1}}
\lst@AddToHook{EmptyStyle}{\let\lst@identifierstyle\@empty}
%    \end{macrocode}
% \end{lstkey}
%
% \begin{macro}{\lst@GotoTabStop}
% 这里我们检查该行是否已含有已打印的字符。
% 如果为真，我们就输出一个宽度等于空白字符的盒子。
%    \begin{macrocode}
\def\lst@GotoTabStop{%
    \ifnum\lst@newlines=\z@
        \setbox\@tempboxa\hbox{\lst@outputspace}%
        \setbox\@tempboxa\hbox to\wd\@tempboxa{{\lst@currstyle{\hss}}}%
        \lst@CalcLostSpaceAndOutput
%    \end{macrocode}
% 大概不清楚为什么输出单个空格就足以到达
% 下一个制表位。只需注意，这一过程丢掉的空间
% 就是上面意义上的`丢失空间'，因此会在输出下一个
% 字符之前被插入。
%    \begin{macrocode}
    \else
%    \end{macrocode}
% 否则（没有已打印的字符）我们只需推进 |\lst@lostspace|
% （它由上面的 |\lst@OutputToken| 插入），并更新列号。
%    \begin{macrocode}
        \global\advance\lst@lostspace \lst@length\lst@width
        \global\advance\lst@column\lst@length \lst@length\z@
    \fi}
%    \end{macrocode}
% 注意这一版本在灵活列模式下也能工作。
% 事实上，它主要是 \packagename{listings} 0.20 的灵活版本。
% \begin{TODO}
% 改用 |\lst@ifnewline| 而不是 |\ifnum\lst@newlines=\z@|？
% \end{TODO}
% \end{macro}
%
% \begin{macro}{\lst@OutputOther}
% 有了前面的定义就变得很容易。
%    \begin{macrocode}
\def\lst@OutputOther{%
    \lst@CheckDelay
    \ifnum\lst@length=\z@\else
        \let\lst@thestyle\lst@currstyle
        \lsthk@OutputOther
        \lst@OutputToken
    \fi}
%    \end{macrocode}
%    \begin{macrocode}
\lst@AddToHook{OutputOther}{}% init
\let\lst@currstyle\relax % init
%    \end{macrocode}
% \end{macro}
%
% \begin{macro}{\lst@Output}
% 我们也许可以把标识符样式作为默认值。
%    \begin{macrocode}
\def\lst@Output{%
    \lst@CheckDelay
    \ifnum\lst@length=\z@\else
        \ifx\lst@currstyle\relax
            \let\lst@thestyle\lst@identifierstyle
        \else
            \let\lst@thestyle\lst@currstyle
        \fi
        \lsthk@Output
        \lst@OutputToken
    \fi
    \let\lst@lastother\relax}
%    \end{macrocode}
% 注意 |\lst@lastother| 变成等价于 |\relax|，而不是像其他
% 地方那样等价于 |\@empty|。我不知道这在将来是否
% 重要。
%    \begin{macrocode}
\lst@AddToHook{Output}{}% init
%    \end{macrocode}
% \end{macro}
%
% \begin{macro}{\lst@GetOutputMacro}
% 只是保存要使用的输出宏。
%    \begin{macrocode}
\def\lst@GetOutputMacro#1{%
    \lst@ifletter \global\let#1\lst@Output
            \else \global\let#1\lst@OutputOther\fi}
%    \end{macrocode}
% \end{macro}
%
% \begin{macro}{\lst@PrintToken}
% 以字母或非字母模式输出当前字符字符串。
%    \begin{macrocode}
\def\lst@PrintToken{%
    \lst@ifletter \lst@Output \lst@letterfalse
            \else \lst@OutputOther \let\lst@lastother\@empty \fi}
%    \end{macrocode}
% \end{macro}
%
% \begin{macro}{\lst@XPrintToken}
% 是一个特殊定义，也打印合并的字符。
%    \begin{macrocode}
\def\lst@XPrintToken{%
    \lst@PrintToken \lst@CheckMerge
    \ifnum\lst@length=\z@\else \lst@PrintToken \fi}
%    \end{macrocode}
% \end{macro}
%
%
% \subsection{丢弃整个输出}
%
% \begin{macro}{\lst@BeginDropOutput}
% 有时把列表的一部分照常处理、但丢弃其
% 输出是有用的。这个宏做主要工作，取一个参数，即它
% 进入的内部模式。我们保存 |\lst@newlines|，用 |\aftergroup| 恢复它，
% 并重新定义一个宏，即 |\lst@OutputBox|。在收到
% \lsthelper{Gunther~Schmidl}{2002/02/27}{在 printpod=false 时
% 合并空行无效}的缺陷报告之后
%    \begin{macrocode}
\def\lst@BeginDropOutput#1{%
    \xdef\lst@BDOnewlines{\the\lst@newlines}%
    \global\let\lst@BDOifnewline\lst@ifnewline
    \lst@EnterMode{#1}%
        {\lst@modetrue
         \let\lst@OutputBox\@gobble
         \aftergroup\lst@BDORestore}}
%    \end{macrocode}
% 恢复日期很容易：
%    \begin{macrocode}
\def\lst@BDORestore{%
    \global\lst@newlines\lst@BDOnewlines
    \global\let\lst@ifnewline\lst@BDOifnewline}
%    \end{macrocode}
% \end{macro}
%
% \begin{macro}{\lst@EndDropOutput}
% 等价于 |\lst@LeaveMode|。
%    \begin{macrocode}
\let\lst@EndDropOutput\lst@LeaveMode
%    \end{macrocode}
%    \begin{macrocode}
%</kernel>
%    \end{macrocode}
% \end{macro}
%
%
% \subsection{写入外部文件}
%
% \begin{aspect}{writefile}
% 现在最好了解一些字符类的情况，因为我们需要
% 取得真正的输入字符，例如制表符，而不是它`展开'成的
% 空格。
%    \begin{macrocode}
%<*misc>
\lst@BeginAspect{writefile}
%    \end{macrocode}
%
% \begin{macro}{\lst@WF}
% \begin{macro}{\lst@WFtoken}
% 该记号的内容将被写入文件。
%    \begin{macrocode}
\newtoks\lst@WFtoken % global
\lst@AddToHook{InitVarsBOL}{\global\lst@WFtoken{}}
%    \end{macrocode}
%    \begin{macrocode}
\newwrite\lst@WF
\global\let\lst@WFifopen\iffalse % init
%    \end{macrocode}
% \end{macro}
% \end{macro}
%
% \begin{macro}{\lst@WFWriteToFile}
% 为此，我们必须展开内容，然后通过 |\edef| 展开它。
% 空的 |\lst@UM| 确保特殊字符（下划线、美元等）
% 被正确写出。
%    \begin{macrocode}
\gdef\lst@WFWriteToFile{%
  \begingroup
   \let\lst@UM\@empty
   \expandafter\edef\expandafter\lst@temp\expandafter{\the\lst@WFtoken}%
   \immediate\write\lst@WF{\lst@temp}%
  \endgroup
  \global\lst@WFtoken{}}
%    \end{macrocode}
% \end{macro}
%
% \begin{macro}{\lst@WFAppend}
% 类似 |\lst@Append|，但使用 |\lst@WFtoken|。
%    \begin{macrocode}
\gdef\lst@WFAppend#1{%
    \global\lst@WFtoken=\expandafter{\the\lst@WFtoken#1}}
%    \end{macrocode}
% \end{macro}
%
% \begin{macro}{\lst@BeginWriteFile}
% \begin{macro}{\lst@BeginAlsoWriteFile}
% 对 |\lst@OutputBox| 使用不同的宏，以（不）丢弃输出。
%    \begin{macrocode}
\gdef\lst@BeginWriteFile{\lst@WFBegin\@gobble}
\gdef\lst@BeginAlsoWriteFile{\lst@WFBegin\lst@OutputBox}
%    \end{macrocode}
% \end{macro}
% \end{macro}
%
% \begin{macro}{\lst@WFBegin}
% 这里\ldots
%    \begin{macrocode}
\begingroup \catcode`\^^I=11
\gdef\lst@WFBegin#1#2{%
    \begingroup
    \let\lst@OutputBox#1%
%    \end{macrocode}
% \ldots\ 我们必须更新 |\lst@WFtoken| 并\ldots
%    \begin{macrocode}
    \def\lst@Append##1{%
        \advance\lst@length\@ne
        \expandafter\lst@token\expandafter{\the\lst@token##1}%
        \ifx ##1\lst@outputspace \else
            \lst@WFAppend##1%
        \fi}%
    \lst@lAddTo\lst@PreGotoTabStop{\lst@WFAppend{^^I}}%
    \lst@lAddTo\lst@ProcessSpace{\lst@WFAppend{ }}%
%    \end{macrocode}
% \ldots\ 需要不同的`EOL'和`DeInit'定义，以把
% 记号寄存器写入文件。
%    \begin{macrocode}
    \let\lst@DeInit\lst@WFDeInit
    \let\lst@MProcessListing\lst@WFMProcessListing
%    \end{macrocode}
% 最后我们在必要时打开文件。
%    \begin{macrocode}
    \lst@WFifopen\else
        \immediate\openout\lst@WF=#2\relax
        \global\let\lst@WFifopen\iftrue
        \@gobbletwo\fi\fi
    \fi}
\endgroup
%    \end{macrocode}
% \end{macro}
%
% \begin{macro}{\lst@EndWriteFile}
% 关闭文件并恢复原始定义。
%    \begin{macrocode}
\gdef\lst@EndWriteFile{%
    \immediate\closeout\lst@WF \endgroup
    \global\let\lst@WFifopen\iffalse}
%    \end{macrocode}
% \end{macro}
%
% \begin{macro}{\lst@WFMProcessListing}
% \begin{macro}{\lst@WFDeInit}
% 另外把 |\lst@WFtoken| 写入外部文件。
%    \begin{macrocode}
\global\let\lst@WFMProcessListing\lst@MProcessListing
\global\let\lst@WFDeInit\lst@DeInit
\lst@AddToAtTop\lst@WFMProcessListing{\lst@WFWriteToFile}
\lst@AddToAtTop\lst@WFDeInit{%
    \ifnum\lst@length=\z@\else \lst@WFWriteToFile \fi}
%    \end{macrocode}
% \end{macro}
% \end{macro}
%
%    \begin{macrocode}
\lst@EndAspect
%</misc>
%    \end{macrocode}
% \end{aspect}
%
%
% \section{字符类}\label{iCharacterClasses}
%
% 本节我们定义基本字符类的行为，之后
% 再转向字符表的选择以及如何特化
% 字符。
%
%
% \subsection{字母、数字与其他}
%
%    \begin{macrocode}
%<*kernel>
%    \end{macrocode}
%
% \begin{macro}{\lst@ProcessLetter}
% 我们把不是空白的字母放入输出队列。
%    \begin{macrocode}
\def\lst@ProcessLetter{\lst@whitespacefalse \lst@AppendLetter}
%    \end{macrocode}
% \end{macro}
%
% \begin{macro}{\lst@ProcessOther}
% 同上。
%    \begin{macrocode}
\def\lst@ProcessOther{\lst@whitespacefalse \lst@AppendOther}
%    \end{macrocode}
% \end{macro}
%
% \begin{macro}{\lst@ProcessDigit}
% 数字把字符追加到当前字符串。但我们必须
% 使用正确的宏。这允许数字成为标识符或
% 数值常量的一部分。
%    \begin{macrocode}
\def\lst@ProcessDigit{%
    \lst@whitespacefalse
    \lst@ifletter \expandafter\lst@AppendLetter
            \else \expandafter\lst@AppendOther\fi}
%    \end{macrocode}
% \end{macro}
%
% \begin{macro}{\lst@ifwhitespace}
% 指示最后处理过的字符是否为空白。
%    \begin{macrocode}
\def\lst@whitespacetrue{\global\let\lst@ifwhitespace\iftrue}
\def\lst@whitespacefalse{\global\let\lst@ifwhitespace\iffalse}
\lst@AddToHook{InitVarsBOL}{\lst@whitespacetrue}
%    \end{macrocode}
% \end{macro}
%
%
% \subsection{空白字符}
%
% 这里我们必须处理两件事：丢弃列表末尾的空行，
% 以及不同的列格式。两者都使用 |\lst@lostspace|。只含
% 制表符和空格的行应被视为空行。为了
% 做到这一点，行首的制表符和空格不输出
% 任何字符，而是推进 |\lst@lostspace|。只要这个尺寸
% 为正，我们就在输出字符串之前插入这段空间。因此，
% 如果只有制表符和空格，该行就是`空'的，因为我们
% 没有做任何输出。
%
% 对于灵活列，我们还要多做一点。空白可以修正列
% 对齐：如果实际行比期望的更宽，制表符至少
% 有一个空格宽；剩下的空间修正对齐。如果存在两个或
% 更多空格字符，至少打印一个；其余修正列
% 对齐。
%
%
% \paragraph{制表符}
% 分三个阶段处理。你已经见过最后阶段
% |\lst@GotoTabStop|。另外两个计算所需宽度，并处理
% 可见的制表符和空格。
%
% \begin{lstkey}{tabsize}
% 在保存之前我们检查参数是否合法。默认 tabsize 为 8，
% 由 \lsthelper{Rolf~Niepraschk}{1997/04/24}{tabsize=8}提议。
%    \begin{macrocode}
\lst@Key{tabsize}{8}
    {\ifnum#1>\z@ \def\lst@tabsize{#1}\else
         \PackageError{Listings}{Strict positive integer expected}%
         {You can't use `#1' as tabsize. \@ehc}%
     \fi}
%    \end{macrocode}
% \end{lstkey}
%
% \begin{lstkey}{showtabs}
% \begin{lstkey}{tab}
% 另外两个用于制表控制的用户键。
%    \begin{macrocode}
\lst@Key{showtabs}f[t]{\lstKV@SetIf{#1}\lst@ifshowtabs}
\lst@Key{tab}{\kern.06em\hbox{\vrule\@height.3ex}%
              \hrulefill\hbox{\vrule\@height.3ex}}
    {\def\lst@tab{#1}}
%    \end{macrocode}
% \end{lstkey}
% \end{lstkey}
%
% \begin{macro}{\lst@ProcessTabulator}
% 制表符输出前面的字符，这会把 |\lst@pos| 减少
% 已打印字符的数目。
%    \begin{macrocode}
\def\lst@ProcessTabulator{%
    \lst@XPrintToken \lst@whitespacetrue
%    \end{macrocode}
% 然后我们计算到达下一个制表位需要多少列：
% 我们不断加上 |\lst@tabsize|，直到 |\lst@pos| 严格为正。换句话说，
% |\lst@pos| 是列号模 |tabsize|，我们正在寻找一个正的
% 代表值。我们把它赋给 |\lst@length|，并在
% 子宏中重置 |\lst@pos|。
%    \begin{macrocode}
    \global\advance\lst@column -\lst@pos
    \@whilenum \lst@pos<\@ne \do
        {\global\advance\lst@pos\lst@tabsize}%
    \lst@length\lst@pos
    \lst@PreGotoTabStop}
%    \end{macrocode}
% \end{macro}
%
% \begin{macro}{\lst@PreGotoTabStop}
% 可见制表符打印 |\lst@tab|。
%    \begin{macrocode}
\def\lst@PreGotoTabStop{%
    \lst@ifshowtabs
        \lst@TrackNewLines
        \setbox\@tempboxa\hbox to\lst@length\lst@width
            {{\lst@currstyle{\hss\lst@tab}}}%
        \lst@CalcLostSpaceAndOutput
    \else
%    \end{macrocode}
% 如果建议我们保留空格，我们就插入正确数目的空格。
%    \begin{macrocode}
        \lst@ifkeepspaces
            \@tempcnta\lst@length \lst@length\z@
            \@whilenum \@tempcnta>\z@ \do
                {\lst@AppendOther\lst@outputspace
                 \advance\@tempcnta\m@ne}%
            \lst@OutputOther
        \else
            \lst@GotoTabStop
        \fi
    \fi
    \lst@length\z@ \global\lst@pos\z@}
%    \end{macrocode}
% \end{macro}
%
%
% \paragraph{空格}
% 按本小节开头所述实现。但首先
% 我们定义一些用户键。
%
% \begin{macro}{\lst@outputspace}
% \begin{macro}{\lst@visiblespace}
%         \lsthelper{Denis Bitouz\'e}{2020/09/12}{`lst@visiblespace：
%         `verbvisiblespace 而非 `textvisiblespace}指出：
%         在 Lua\LaTeX{} 和某些在 32 号槽位没有合适字形的等宽字体
%         下，使用 |\textvisiblespace| 时 `showspaces' 完全没有效果。
%         所以现在我们使用 |\verbvisiblespace| 作为第一个宏的默认
%         定义，针对字体族的测试
%         |\lst@ttfamily| 被删除，因为在大多数情况下它导致
%         |\char32| 的错误输出。latex.ltx 中 |\verbvisible| 的定义
%         考虑了现代 \TeX{} 引擎的不同行为，并定义了有效的 |\verbvisiblespace|
%    \begin{macrocode}
\def\lst@outputspace{\ }
\def\lst@visiblespace{\verbvisiblespace}
%    \end{macrocode}
% \end{macro}
% \end{macro}
%
% \begin{lstkey}{showspaces}
% \begin{lstkey}{keepspaces}
% \ldots\ 它会按用户要求被修改。
%    \begin{macrocode}
\lst@Key{showspaces}{false}[t]{\lstKV@SetIf{#1}\lst@ifshowspaces}
\lst@Key{keepspaces}{false}[t]{\lstKV@SetIf{#1}\lst@ifkeepspaces}
\lst@AddToHook{Init}
    {\lst@ifshowspaces
         \let\lst@outputspace\lst@visiblespace
         \lst@keepspacestrue
     \fi}
\def\lst@keepspacestrue{\let\lst@ifkeepspaces\iftrue}
%    \end{macrocode}
% \end{lstkey}
% \end{lstkey}
%
% \begin{macro}{\lst@ProcessSpace}
% 我们检查空格是否修正列对齐。在后一种情况下
% 我们追加一个空格；否则\ldots
% \lsthelper{Andrei~Alexandrescu}{2007/02/27}{-}测试了 |spaceflexible|
% 列设置，并发现了一个缺陷，它源于 |\lst@PrintToken| 与
% |\lst@whitespacetrue| 在这里的执行顺序颠倒。
%    \begin{macrocode}
\def\lst@ProcessSpace{%
    \lst@ifkeepspaces
        \lst@PrintToken
        \lst@whitespacetrue
        \lst@AppendOther\lst@outputspace
        \lst@PrintToken
    \else \ifnum\lst@newlines=\z@
%    \end{macrocode}
% \ldots\ 如果该行非空，我们就追加一个`特殊空格'。
%    \begin{macrocode}
        \lst@AppendSpecialSpace
    \else \ifnum\lst@length=\z@
%    \end{macrocode}
% 如果该行是空的，我们就检查输出队列中
% 是否有字符。如果没有字符，我们只推进 |\lst@lostspace|。
% 否则我们追加该空格。
%    \begin{macrocode}
            \global\advance\lst@lostspace\lst@width
            \global\advance\lst@pos\m@ne
            \lst@whitespacetrue
        \else
            \lst@AppendSpecialSpace
        \fi
    \fi \fi}
%    \end{macrocode}
% 注意这一版本对固定和灵活列输出都能工作。
% \end{macro}
%
% \begin{macro}{\lst@AppendSpecialSpace}
% 如果至少有两个空白，我们就输出前面的字符并
% 推进 |\lst@lostspace| 以避免对齐问题。否则我们
% 把空格追加到当前字符串。另外，|\lst@whitespacetrue| 已
% 被移到 |\lst@PrintToken| 之后，以便记号打印器能正确
% 检查它打印的是否是空白；这一点以前妨碍了
% |spaceflexible| 列设置正确工作。
%    \begin{macrocode}
\def\lst@AppendSpecialSpace{%
    \lst@ifwhitespace
        \lst@PrintToken
        \global\advance\lst@lostspace\lst@width
        \global\advance\lst@pos\m@ne
        \lst@gobbledwhitespacetrue
    \else
        \lst@PrintToken
        \lst@whitespacetrue
        \lst@AppendOther\lst@outputspace
        \lst@PrintToken
    \fi}
%    \end{macrocode}
% \end{macro}
%
%
% \paragraph{换页符}
% 是在与以下各方交流之后引入的
% \lsthelper{Jan~Braun}{1998/04/27}{formfeed}.
%
% \begin{lstkey}{formfeed}
% 让用户做调整。
%    \begin{macrocode}
\lst@Key{formfeed}{\bigbreak}{\def\lst@formfeed{#1}}
%    \end{macrocode}
% \end{lstkey}
%
% \begin{macro}{\lst@ProcessFormFeed}
% 这里我们根据是否已经开始新行来执行一些宏。
% else 分支中不再使用 |\lst@EOLUpdate|
% ——\lsthelper{Kalle~Tuulos}{2001/01/14}{换页符吃掉后续
% 输出单元}发来了缺陷报告。
%    \begin{macrocode}
\def\lst@ProcessFormFeed{%
    \lst@XPrintToken
    \ifnum\lst@newlines=\z@
        \lst@EOLUpdate \lsthk@InitVarsBOL
    \fi
    \lst@formfeed
    \lst@whitespacetrue}
%    \end{macrocode}
% \end{macro}
%
%
% \subsection{字符表}\label{iCharacterTables}
%
%
% \subsubsection{标准字符表}
%
% 标准字符表由 |\lst@SelectStdCharTable| 选择，它展开为一个
% 记号序列
%    \ldots|\def| |A{\lst@ProcessLetter| |A}|\ldots\space
% 其中第一个 A 是活动的，第二个的 catcode 为 12。我们用
% 下面的宏来构造字符表。
% \begin{syntax}
%     \item[0.19] |\lst@CCPut|\meta{class
%         macro}\meta{$c_1$}\ldots\meta{$c_k$}|\z@|
%
%         用码为 \meta{$c_1$}\ldots\meta{$c_k$} 的字符扩展标准字符表，
%         使每个字符使用 \meta{class macro}。所有这些字符都必须能通过
%     |\char|\meta{$c_i$}.
%
%   \item[0.20]
%     |\lst@CCPutMacro|\meta{class$_1$}\meta{$c_1$}\meta{definition$_1$}^^A
%     \ldots|\@empty\z@\@empty|
%
%         也扩展标准字符表：字符
%         \meta{$c_i$} 将使用 \meta{class$_i$}，并通过
%         \meta{definition$_i$} 打印。这些定义必须是第
%         \ref{dCharacterTables} 节意义上的 \meta{spec. token}。
% \end{syntax}
%
% \begin{macro}{\lst@Def}
% \begin{macro}{\lst@Let}
% 为了速度，我们不会太频繁使用这些辅助宏。
%    \begin{macrocode}
\def\lst@Def#1{\lccode`\~=#1\lowercase{\def~}}
\def\lst@Let#1{\lccode`\~=#1\lowercase{\let~}}
%    \end{macrocode}
% \end{macro}
% \end{macro}
%
% \begingroup
% 下面空格的定义不会造成任何问题。但其他方面，
% 例如 \aspectname{lineshape} 和 \aspectname{formats}，也会重新定义
% 宏 |\space|。现在，如果 \LaTeX\ 调用 |\try@load@fontshape|，|.log|
% 消息就会显示一些奇怪的东西，因为 \LaTeX\ 在这些
% 消息中使用 |\space|。下面的补充确保 |\space| 展开为一个空格
% 而不是别的什么。这又是一个缺陷，由
% \lsthelper{Denis~Girou}{1999/09/16}{breaklines 时的字体提示信息有误}。
%    \begin{macrocode}
\lst@AddToAtTop{\try@load@fontshape}{\def\space{ }}
%    \end{macrocode}
% \endgroup
%
% \begin{macro}{\lst@SelectStdCharTable}
% 前三个标准字符。|\lst@Let| 已被替换为
% |\lst@Def|，这是在 \lsthelper{Chris~Edwards}{2002/02/15}{firstline>1 时制表符会显示出来}报告缺陷之后。
%    \begin{macrocode}
\def\lst@SelectStdCharTable{%
    \lst@Def{9}{\lst@ProcessTabulator}%
    \lst@Def{12}{\lst@ProcessFormFeed}%
    \lst@Def{32}{\lst@ProcessSpace}}
%    \end{macrocode}
% \end{macro}
%
% \begin{macro}{\lst@CCPut}
% 第一个参数给出字符类，然后是
% 码。
%
% \lsthelper{Joseph~Wright}{2016/04/04}{unguarded setting of |lccode|}
% 指出了 TeX StackExchange 上提出的一个缺陷
% （|http://tex.stackexchange.com/questions/302437/textcase-lstings-and-tilde|）。与
% |\lst@CCPutMacro| 不同，|\lccode| 的设置不是局部的，因而导致了
% 该错误。
%    \begin{macrocode}
\def\lst@CCPut#1#2{%
    \ifnum#2=\z@
        \expandafter\@gobbletwo
    \else
        \begingroup\lccode`\~=#2\lccode`\/=#2\lowercase{\endgroup\lst@CCPut@~{#1/}}%
    \fi
    \lst@CCPut#1}
\def\lst@CCPut@#1#2{\lst@lAddTo\lst@SelectStdCharTable{\def#1{#2}}}
%    \end{macrocode}
% 现在我们插入更多标准字符。
%    \begin{macrocode}
\lst@CCPut \lst@ProcessOther
    {"21}{"22}{"28}{"29}{"2B}{"2C}{"2E}{"2F}
    {"3A}{"3B}{"3D}{"3F}{"5B}{"5D}
    \z@
\lst@CCPut \lst@ProcessDigit
    {"30}{"31}{"32}{"33}{"34}{"35}{"36}{"37}{"38}{"39}
    \z@
\lst@CCPut \lst@ProcessLetter
    {"40}{"41}{"42}{"43}{"44}{"45}{"46}{"47}
    {"48}{"49}{"4A}{"4B}{"4C}{"4D}{"4E}{"4F}
    {"50}{"51}{"52}{"53}{"54}{"55}{"56}{"57}
    {"58}{"59}{"5A}
         {"61}{"62}{"63}{"64}{"65}{"66}{"67}
    {"68}{"69}{"6A}{"6B}{"6C}{"6D}{"6E}{"6F}
    {"70}{"71}{"72}{"73}{"74}{"75}{"76}{"77}
    {"78}{"79}{"7A}
    \z@
%    \end{macrocode}
% \end{macro}
%
% \begin{macro}{\lst@CCPutMacro}
% 现在到了一个微妙的点。如果直接输入这些字符，尚未插入的字符
% 或者不可打印（|_|、|$|、\ldots），或者打印效果不好（|*|、|-|、\ldots）。
% 因此我们使用合适的宏来打印这些字符。
% 效果完美。问题在于当前字符串为了速度已经
% 是可打印的，例如 |_| 已被替换为宏版本，
% 但新的关键词测试需要原始字符。
%
% 解决方案：我们定义 |\def _{\lst@ProcessLetter\lst@um_}|，其中第一个
% 下划线是活动的，第二个属于该控制序列。
% 此外我们还有 |\def\lst@um_{\lst@UM _}|，其中第二个下划线
% 具有通常的含义。现在关键词测试只需把 |\lst@UM| 设为空，
% 就能访问原始字符。默认定义取得下一个
% 记号并构造控制序列 |\lst@um_@|，我们将定义它来
% 打印该字符。容易，不是吗？^^A ;-)
%
% 下面的定义为我们都做好了。第一个参数给出
% 字符类，第二个给出字符码，最后一个给出实际
% 打印该字符的定义。我们构造名字 |\lst@um_| 和
% |\lst@um_@|，并把它们交给一个子宏。
%    \begin{macrocode}
\def\lst@CCPutMacro#1#2#3{%
    \ifnum#2=\z@ \else
        \begingroup\lccode`\~=#2\relax \lccode`\/=#2\relax
        \lowercase{\endgroup\expandafter\lst@CCPutMacro@
            \csname\@lst @um/\expandafter\endcsname
            \csname\@lst @um/@\endcsname /~}#1{#3}%
        \expandafter\lst@CCPutMacro
    \fi}
%    \end{macrocode}
% 现在的参数是 |\lst@um_|、|\lst@um_@|、非活动字符、活动
% 字符、字符类和打印定义。我们把 |\def _{|
% |\lst@ProcessLetter| |\lst@um_}| 加入 |\lst@SelectStdCharTable|（以及其他
% 特殊字符同理），并定义 |\def\lst@um_{\lst@UM _}| 和 |\lst@um_@|。
%    \begin{macrocode}
\def\lst@CCPutMacro@#1#2#3#4#5#6{%
    \lst@lAddTo\lst@SelectStdCharTable{\def#4{#5#1}}%
    \def#1{\lst@UM#3}%
    \def#2{#6}}
%    \end{macrocode}
% |\lst@UM| 的默认定义：
%    \begin{macrocode}
\def\lst@UM#1{\csname\@lst @um#1@\endcsname}
%    \end{macrocode}
% 以及所有剩余的标准字符。\lsthelper{Peter
% Cock}{2024/09/18}{LaTeX listings 中的直双引号}
% 指出了以下讨论
% \url{https://tex.stackexchange.com/questions/166790/how-can-i-get-straight-double-quotes-in-listings}
% 中关于缺少直双引号支持的问题，现在
% 由下面几行实现。
%    \begin{macrocode}
\lst@CCPutMacro
    \lst@ProcessOther {"22}{\lst@ifupquote \textquotedbl
                                     \else \char34\relax \fi}
    \lst@ProcessOther {"23}\#
    \lst@ProcessLetter{"24}\textdollar
    \lst@ProcessOther {"25}\%
    \lst@ProcessOther {"26}\&
    \lst@ProcessOther {"27}{\lst@ifupquote \textquotesingle
                                     \else \char39\relax \fi}
    \lst@ProcessOther {"2A}{\lst@ttfamily*\textasteriskcentered}
%    \end{macrocode}
% \lsthelper{Ulrike~Fischer}{2020/02/19}{breqn/flexisym 与 listings
% 不兼容}指出 \packagename{flexisym} 与
% \packagename{listings} 之间的不兼容：\packagename{flexisym}
% 改变数学码，而 \packagename{listings} 改变
% 含义。于是减号消失了。把原来的 |$-$| 替换为
% |\textminus| 应当能补救这个问题。
%    \begin{macrocode}
    \lst@ProcessOther {"2D}{\lst@ttfamily{-{}}{\textminus}}
    \lst@ProcessOther {"3C}{\lst@ttfamily<\textless}
    \lst@ProcessOther {"3E}{\lst@ttfamily>\textgreater}
    \lst@ProcessOther {"5C}{\lst@ttfamily{\char92}\textbackslash}
    \lst@ProcessOther {"5E}\textasciicircum
    \lst@ProcessLetter{"5F}{\lst@ttfamily{\char95}\textunderscore}
                                     % or \char"5F
    \lst@ProcessOther {"60}{\lst@ifupquote \textasciigrave
                                     \else \char96\relax \fi}
    \lst@ProcessOther {"7B}{\lst@ttfamily{\char123}\textbraceleft}
    \lst@ProcessOther {"7C}{\lst@ttfamily|\textbar}
    \lst@ProcessOther {"7D}{\lst@ttfamily{\char125}\textbraceright}
    \lst@ProcessOther {"7E}\textasciitilde
    \lst@ProcessOther {"7F}-
    \@empty\z@\@empty
%    \end{macrocode}
% \end{macro}
%
% \begin{macro}{\lst@ttfamily}
% 这个神秘的宏是什么？它打印第一个或第二个
% 参数。在 |\ttfamily| 中它确保 |----| 被排成 |----| 而不是
% 像 0.17 版那样排成 $-$$-$$-$$-$。缺陷由以下人员遇到：
% \lsthelper{Dr.~Jobst~Hoffmann}{1998/03/30}{|\lst@minus| 与 |\ttfamily|}。
% 此外，在收到来自
% \lsthelper{Magnus~Lewis-Smith}{1999/08/06}{! Bad character code (920).}
%    \begin{macrocode}
\def\lst@ttfamily#1#2{\ifx\f@family\ttdefault#1\relax\else#2\fi}
%    \end{macrocode}
% |\ttdefault| 被定义为 |\long|，所以 |\ifx| 不起作用，因为 |\f@family|
% 不是 |\long|！我们通过局部重新定义 |\ttdefault| 绕过这个问题：
%    \begin{macrocode}
\lst@AddToHook{Init}{\edef\ttdefault{\ttdefault}}
%    \end{macrocode}
% \end{macro}
%
% \begin{lstkey}{upquote}
% 在上面用于决定打印哪个引号。如果
% 所需的 \packagename{textcomp} 命令不可用，我们就打印错误消息。这个键
% 是在 \lsthelper{Frank~Mittelbach}{2003/06/18}{listings 与 upquote}的邮件之后添加的。
%    \begin{macrocode}
\lst@Key{upquote}{false}[t]{\lstKV@SetIf{#1}\lst@ifupquote
    \lst@ifupquote
       \@ifundefined{textasciigrave}%
          {\let\KV@lst@upquote\@gobble
           \lstKV@SetIf f\lst@ifupquote \@gobble\fi
           \PackageError{Listings}{Option `upquote' requires `textcomp'
            package.\MessageBreak The option has been disabled}%
          {Add \string\usepackage{textcomp} to your preamble.}}%
          {}%
    \fi}
%    \end{macrocode}
% 如果加载了 \packagename{upquote} 宏包，\keyname{upquote} 选项
% 就默认启用。
%    \begin{macrocode}
\AtBeginDocument{%
  \@ifpackageloaded{upquote}{\RequirePackage{textcomp}%
                             \lstset{upquote}}{}%
  \@ifpackageloaded{upquote2}{\lstset{upquote}}{}}
%    \end{macrocode}
% \end{lstkey}
%
% \begin{macro}{\lst@ifactivechars}
% 一个简单的开关。
%    \begin{macrocode}
\def\lst@activecharstrue{\let\lst@ifactivechars\iftrue}
\def\lst@activecharsfalse{\let\lst@ifactivechars\iffalse}
\lst@activecharstrue
%    \end{macrocode}
% \end{macro}
%
% \begin{macro}{\lst@SelectCharTable}
% 我们选择标准字符表并切换到活动 catcode。
%    \begin{macrocode}
\def\lst@SelectCharTable{%
    \lst@SelectStdCharTable
    \lst@ifactivechars
        \catcode9\active \catcode12\active \catcode13\active
        \@tempcnta=32\relax
        \@whilenum\@tempcnta<128\do
            {\catcode\@tempcnta\active\advance\@tempcnta\@ne}%
    \fi
    \lst@ifec \lst@DefEC \fi
%    \end{macrocode}
% 下面这一行以及下面的相应宏，是在
% \lsthelper{Fr\'ed\'eric~Boulanger}{2001/02/27}{连字}报告缺陷之后添加的。
% 对 |\do@noligs| 的赋值改成 |\do|，是在
% \lsthelper{Peter~Ruckdeschel}{2002/04/12}{同时使用
% seminar.sty 和 listings.sty 的问题}报告缺陷之后。这个缺陷修正由以下人员友情提供：
% \lsthelper{Timothy~Van~Zandt}{2002/04/13}{Re: ...}.
%    \begin{macrocode}
    \let\do\lst@do@noligs \verbatim@nolig@list
%    \end{macrocode}
% 有两种方式调整标准表：在钩子内部，或通过
% |\lst@DeveloperSCT|。我们使用这些宏，并在必要时
% 初始化反斜杠。|\lst@DefRange| 被移到钩子之外，这是在
% \lsthelper{Michael~Bachmann}{2004/07/21}{无法进行 label 引用
% \ldots}\phantomsection\label{rLst@DefRange-moved}报告缺陷之后。
% ^^A \phantomsection is used to correct the hyperref link
%    \begin{macrocode}
    \lsthk@SelectCharTable
    \lst@DeveloperSCT\lst@DefRange
    \ifx\lst@Backslash\relax\else
        \lst@LetSaveDef{"5C}\lsts@backslash\lst@Backslash
    \fi}
%    \end{macrocode}
% \end{macro}
%
% \begin{lstkey}{SelectCharTable}
% \begin{lstkey}{MoreSelectCharTable}
% 用于调整 |\lst@DeveloperSCT| 的键。
%    \begin{macrocode}
\lst@Key{SelectCharTable}{}{\def\lst@DeveloperSCT{#1}}
\lst@Key{MoreSelectCharTable}\relax{\lst@lAddTo\lst@DeveloperSCT{#1}}
%    \end{macrocode}
%    \begin{macrocode}
\lst@AddToHook{SetLanguage}{\let\lst@DeveloperSCT\@empty}
%    \end{macrocode}
% \end{lstkey}
% \end{lstkey}
%
% \begin{macro}{\lst@do@noligs}
% 为了阻止连字，这个宏把记号 |\lst@NoLig| 插到
% |\lst@Process|\meta{whatever}\meta{spec.~token} 前面。对某些字符，这由
% |\verbatim@nolig@list| 完成。注意该子宏是
% 一种特殊的局部 |\lst@AddToAtTop|。子宏定义
% 得感谢 \lsthelper{Peter~Bartke}{2002/04/10}{糟糕的`noligs'处理}得以修正。
%    \begin{macrocode}
\def\lst@do@noligs#1{%
    \begingroup \lccode`\~=`#1\lowercase{\endgroup
    \lst@do@noligs@~}}
\def\lst@do@noligs@#1{%
    \expandafter\expandafter\expandafter\def
    \expandafter\expandafter\expandafter#1%
    \expandafter\expandafter\expandafter{\expandafter\lst@NoLig#1}}
%    \end{macrocode}
% \end{macro}
%
% \begin{macro}{\lst@NoLig}
% 当这个额外的宏被处理时，它把 |\lst@nolig| 加入输出队列
% 而不增加其长度。对于关键词检测，如果 |\lst@UM| 为空，
% 它必须展开为无。
%    \begin{macrocode}
\def\lst@NoLig{\advance\lst@length\m@ne \lst@Append\lst@nolig}
\def\lst@nolig{\lst@UM\@empty}%
%    \end{macrocode}
% 但 |\lst@UM| 的通常含义会构造下面的控制序列，
% 它以 \LaTeX 的 |\do@noligs| 的方式阻止连字。
%    \begin{macrocode}
\@namedef{\@lst @um@}{\leavevmode\kern\z@}
%    \end{macrocode}
% \end{macro}
%
% \begin{macro}{\lst@SaveOutputDef}
% 为了得到字符 |#1| 的 \meta{spec.~token} 含义，我们在
% |\lst@SelectStdCharTable| 中查找 |\def| `活动字符 |#1|'，取得替换
% 文本，用 |\@gobble| 去掉字符类，然后赋予该含义。
% 注意如果使用非法的 \meta{character code}=|#1|，你会得到一个``runaway argument''
% 错误。
%    \begin{macrocode}
\def\lst@SaveOutputDef#1#2{%
    \begingroup \lccode`\~=#1\relax \lowercase{\endgroup
    \def\lst@temp##1\def~##2##3\relax}{%
        \global\expandafter\let\expandafter#2\@gobble##2\relax}%
    \expandafter\lst@temp\lst@SelectStdCharTable\relax}
%    \end{macrocode}
% \end{macro}
%
% \begin{macro}{\lstum@backslash}
% 一个常用的字符。
%    \begin{macrocode}
\lst@SaveOutputDef{"5C}\lstum@backslash
%    \end{macrocode}
% \end{macro}
%
%
% \subsubsection{国家或地区字符}
%
% \begin{lstkey}{extendedchars}
% 用于激活 128--255 扩展字符的用户键。
%    \begin{macrocode}
\lst@Key{extendedchars}{true}[t]{\lstKV@SetIf{#1}\lst@ifec}
%    \end{macrocode}
% \end{lstkey}
%
% \begin{macro}{\lst@DefEC}
% 目前 128--255 范围内的每个字符都被当作字母。
%    \begin{macrocode}
\def\lst@DefEC{%
    \lst@CCECUse \lst@ProcessLetter
      ^^80^^81^^82^^83^^84^^85^^86^^87^^88^^89^^8a^^8b^^8c^^8d^^8e^^8f%
      ^^90^^91^^92^^93^^94^^95^^96^^97^^98^^99^^9a^^9b^^9c^^9d^^9e^^9f%
      ^^a0^^a1^^a2^^a3^^a4^^a5^^a6^^a7^^a8^^a9^^aa^^ab^^ac^^ad^^ae^^af%
      ^^b0^^b1^^b2^^b3^^b4^^b5^^b6^^b7^^b8^^b9^^ba^^bb^^bc^^bd^^be^^bf%
      ^^c0^^c1^^c2^^c3^^c4^^c5^^c6^^c7^^c8^^c9^^ca^^cb^^cc^^cd^^ce^^cf%
      ^^d0^^d1^^d2^^d3^^d4^^d5^^d6^^d7^^d8^^d9^^da^^db^^dc^^dd^^de^^df%
      ^^e0^^e1^^e2^^e3^^e4^^e5^^e6^^e7^^e8^^e9^^ea^^eb^^ec^^ed^^ee^^ef%
      ^^f0^^f1^^f2^^f3^^f4^^f5^^f6^^f7^^f8^^f9^^fa^^fb^^fc^^fd^^fe^^ff%
      ^^00}
%    \end{macrocode}
% \end{macro}
%
% \begin{macro}{\lst@CCECUse}
% 到达列表末尾（|^^00|）时我们终止循环。
% 否则，如果字符不是活动的，我们就做与 |\lst@CCPut| 相同的事。
% 但如果字符是活动的，我们在重新定义之前保存其含义。
%    \begin{macrocode}
\def\lst@CCECUse#1#2{%
    \ifnum`#2=\z@
        \expandafter\@gobbletwo
    \else
        \ifnum\catcode`#2=\active
            \lccode`\~=`#2\lccode`\/=`#2\lowercase{\lst@CCECUse@#1~/}%
        \else
            \lst@ifactivechars \catcode`#2=\active \fi
            \lccode`\~=`#2\lccode`\/=`#2\lowercase{\def~{#1/}}%
        \fi
    \fi
    \lst@CCECUse#1}
%    \end{macrocode}
% 我们如前所述保存含义。这里我们还必须使用`|\lst@UM|
% 构造'，因为扩展字符经常出现在单词中 =
% 标识符中。缺陷由 \lsthelper{Denis~Girou}{1999/07/26}{与 inputenc 不兼容}报告。
%    \begin{macrocode}
\def\lst@CCECUse@#1#2#3{%
    \expandafter\def\csname\@lst @EC#3\endcsname{\lst@UM#3}%
    \expandafter\let\csname\@lst @um#3@\endcsname #2%
    \edef#2{\noexpand#1%
            \expandafter\noexpand\csname\@lst @EC#3\endcsname}}
%    \end{macrocode}
% \lsthelper{Daniel~Gerigk}{2001/10/25}{extendedchars 不起作用}和
% \lsthelper{Heiko~Oberdiek}{2001/10/26}{extendedchars 不起作用：um@\#3@
% 应为 @um\#3@}分别报告了一个错误和一个解决方案。
% \end{macro}
%
%
% \subsubsection{catcode 的问题}
%
% \begin{macro}{\lst@nfss@catcodes}
% \lsthelper{Anders~Edenbrandt}{1997/04/22}{.fd 文件的预加载}发现了一个与
% \texttt{.fd} 文件有关的缺陷。由于我们改动 catcode 而这些文件是按需
% 读入的，我们必须在输入这些文件之前重置 catcode。我们
% 对 |\nfss@catcodes| 做局部重新定义。
%    \begin{macrocode}
\lst@AddToHook{Init}
    {\let\lsts@nfss@catcodes\nfss@catcodes
     \let\nfss@catcodes\lst@nfss@catcodes}
%    \end{macrocode}
% 在 \lsthelper
% {David~Aspinall}{2003/07/17}{在 tabular 内加载 .fd 文件产生
% 错误}报告缺陷之后，|&| 字符变成了 |\&|。
%    \begin{macrocode}
\def\lst@nfss@catcodes{%
    \lst@makeletter
        ABCDEFGHIJKLMNOPQRSTUVWXYZabcdefghijklmnopqrstuvwxyz\relax
    \@makeother (\@makeother )\@makeother ,\@makeother :\@makeother\&%
    \@makeother 0\@makeother 1\@makeother 2\@makeother 3\@makeother 4%
    \@makeother 5\@makeother 6\@makeother 7\@makeother 8\@makeother 9%
    \@makeother =\lsts@nfss@catcodes}
%    \end{macrocode}
% 对 \lsthelper{Christian~Gudrian}
% {2000/11/16}{mathpple 的问题}所报告缺陷的调查表明，等号需要
% 具有`其他'catcode，如上所赋值。
% \lsthelper{Svend~Tollak~Munkejord}{2002/04/17}{宏包与
%   Lucida .fd 文件不兼容}报告了 Lucida .fd 文件的问题，而
% \lsthelper{Heiko~Oberdiek}{2002/04/17}{回复：listings 在 Lucida
% 字体下失败}分析了该缺陷，这导致上面以
% |\@makeaother (|.
%
% |\lst@makeletter| 这个名字是模仿 \LaTeX 的 |\@makeother|。
%    \begin{macrocode}
\def\lst@makeletter#1{%
    \ifx\relax#1\else\catcode`#111\relax \expandafter\lst@makeletter\fi}
%    \end{macrocode}
% \end{macro}
%
% \begin{lstkey}{useoutput}
% \begin{macro}{\output}
% 另一个问题最初由 \lsthelper{Marcin~Kasperski}
% {1999/04/28}{listings 破坏目录}报告。它也与 catcode 有关，
% \lsthelper{Donald~Arseneau}{1999/05/13}{comp.text.tex 回复：延迟写与
% catcode 改动}让我理解了它。要点是：当 \TeX\ 把非 |\immediate| 的
% |\write| 写入文件时，它似乎使用\emph{当前}活动的 catcode 表，
% 而不是\emph{读入}字符时所涉及的 catcode。
% 因此，如果一个列表跨两页，节标题 |\L a| 会被写成 |\La|：
% 因为在把 |\La| 写入文件时使用了非标准的 catcode 表，
% 先前附着的 catcode 不起作用。还有一个缺陷是，当列表跨
% 两页时，页眉或页脚中的重音符号会丢失。
% \lsthelper{Denis~Girou}{1999/08/03}{列表跨两页时
% 页眉中的重音丢失}发现了后一个缺陷。与
% 波浪号类似的另一个问题由 \lsthelper{Thorsten~Vitt}{2001/06/25}{fancyhdr +
% listings 跨页 ==> 页眉中 ~ 而非空格}报告。
%
% 我们可以在三种可能性之间选择。
% \lsthelper{Donald~Arseneau}{2006/09/14}{cannot select output routine 1}
% 在这里的 |\ifcase| 参数中发现了一个缺陷。
%    \begin{macrocode}
\lst@Key{useoutput}{2}{\edef\lst@useoutput{\ifcase0#1 0\or 1\else 2\fi}}
%    \end{macrocode}
% 第一种不改动已有的输出例程。
%    \begin{macrocode}
\lst@AddToHook{Init}
{\edef\lst@OrgOutput{\the\output}%
\ifcase\lst@useoutput\relax
\or
%    \end{macrocode}
% 第二种可能性如下：我们中断当前模式——
% 特别是带修改过的 catcode 表的 |\lst@Pmode|——调用原始输出
% 例程，然后重新进入该模式。这必须小心处理。
% 首先我们必须关闭 \TeX\ 在输出例程开头
% 打开的分组。单个 |\egroup| 会给出`输出例程不平衡'
% 错误。但 |\expandafter\egroup| 有效。又是
% \lsthelper{Donald~Arseneau}{2001/01/10}{comp.text.tex Re: \output puzzle}
% 给出了解释：|\expandafter| 把 |\bgroup| 的记号类型设为
% |backed_up|，这阻止 \TeX 从输出例程不平衡的状态恢复。
% \lsthelper{Heiko~Oberdiek}{2001/01/05}{comp.text.tex 回复：
% \output 之谜}报告说 |\csname| |egroup||\endcsname| 也能
% 达到目的。
%
% 不过，由于我们关闭分组时 \TeX\ 会检查 |\box| 255 的内容
% （`output routine didn't use all of |\box| 255'），我们
% 必须临时保存它。
%    \begin{macrocode}
 \output{\global\setbox\lst@gtempboxa\box\@cclv
         \expandafter\egroup
%    \end{macrocode}
% 现在我们可以中断模式了，但我们必须保存当前字符串
% 和当前样式。
%    \begin{macrocode}
         \lst@SaveToken
     \lst@InterruptModes
%    \end{macrocode}
% 我们恢复内容，使用原始输出例程，然后\ldots
%    \begin{macrocode}
     \setbox\@cclv\box\lst@gtempboxa
     \bgroup\lst@OrgOutput\egroup
%    \end{macrocode}
% \ldots\ 打开一个与 \TeX\ 在输出例程末尾插入的 |}| 相匹配的分组。
% 我们重新进入模式，并用 |\aftergroup| 恢复字符串和
% 样式。此外我们需要重置 |\pagegoal|——这是在
% \lsthelper{Jochen~Schneider}{2002/03/09}{de.comp.text.tex：
% Listings 宏包 1.0-Beta 的问题；前面有浮动体时
% 无故分页}报告缺陷之后添加的。
%    \begin{macrocode}
     \bgroup
     \aftergroup\pagegoal\aftergroup\vsize
     \aftergroup\lst@ReenterModes\aftergroup\lst@RestoreToken}%
\else
%    \end{macrocode}
% 第三种选择是在修改过的输出例程内恢复所有 catcode 和含义，
% 然后再调用原始例程。
%    \begin{macrocode}
 \output{\lst@RestoreOrigCatcodes
         \lst@ifec \lst@RestoreOrigExtendedCatcodes \fi
         \lst@OrgOutput}%
\fi}
%    \end{macrocode}
% 注意这个输出例程不会经常被使用。只有在列表
% 可能跨两页时才执行它：列表在页面底部结束或
% 从页面顶部开始，或者列表确实被分割。
% \end{macro}
% \end{lstkey}
%
% \begin{macro}{\lst@GetChars}
% \begin{macro}{\lst@ScanChars}
% \begin{lstkey}{rescanchars}
% 为了让第三种 |\output| 选择工作，我们必须扫描 catcode
% 以及活动字符的含义：
%    \begin{macrocode}
\def\lst@GetChars#1#2#3{%
    \let#1\@empty
    \@tempcnta#2\relax \@tempcntb#3\relax
    \loop \ifnum\@tempcnta<\@tempcntb\relax
        \lst@lExtend#1{\expandafter\catcode\the\@tempcnta=}%
        \lst@lExtend#1{\the\catcode\@tempcnta\relax}%
        \ifnum\the\catcode\@tempcnta=\active
            \begingroup\lccode`\~=\@tempcnta
            \lowercase{\endgroup
            \lst@lExtend#1{\expandafter\let\expandafter~\csname
                                    lstecs@\the\@tempcnta\endcsname}%
            \expandafter\let\csname lstecs@\the\@tempcnta\endcsname~}%
        \fi
        \advance\@tempcnta\@ne
    \repeat}
%    \end{macrocode}
% \label{catcode12}根据 \lsthelper{Benjamin~Lings}{2004/10/15}%
% {\usepackage{xy,listings} 产生："Forbidden control sequence...."}的缺陷报告，我们
% 临时停用 |^^L| 的 |\outer| 定义（在 |\lst@ScanChars| 内部和外部）
% 并在宏包末尾通过 |\lst@RestoreCatcodes| 命令恢复
% catcode。
%    \begin{macrocode}
\begingroup \catcode12=\active\let^^L\@empty
\gdef\lst@ScanChars{%
  \let\lsts@ssL^^L%
  \def^^L{\par}%
    \lst@GetChars\lst@RestoreOrigCatcodes\@ne {128}%
  \let^^L\lsts@ssL
    \lst@GetChars\lst@RestoreOrigExtendedCatcodes{128}{256}}
\endgroup
%    \end{macrocode}
% 该扫描可以手工触发，也可以在文档开头触发。
%    \begin{macrocode}
\lst@Key{rescanchars}\relax{\lst@ScanChars}
\AtBeginDocument{\lst@ScanChars}
%    \end{macrocode}
% \end{lstkey}
% \end{macro}
% \end{macro}
%
%
% \subsubsection{调整字符表}
%
% 我们从基本字符类的修饰符开始。
%
% \begin{lstkey}{alsoletter}
% \begin{lstkey}{alsodigit}
% \begin{lstkey}{alsoother}
% 宏 |\lst@also|\ldots\ 将保存 |\def|\meta{char}|{|\ldots|}| 序列，
% 用于调整标准字符表。
%    \begin{macrocode}
\lst@Key{alsoletter}\relax{%
    \lst@DoAlso{#1}\lst@alsoletter\lst@ProcessLetter}
\lst@Key{alsodigit}\relax{%
    \lst@DoAlso{#1}\lst@alsodigit\lst@ProcessDigit}
\lst@Key{alsoother}\relax{%
    \lst@DoAlso{#1}\lst@alsoother\lst@ProcessOther}
%    \end{macrocode}
% 这在 \hookname{SelectCharTable} 处完成，每次选择语言时
% 这些宏都会变为空。
%    \begin{macrocode}
\lst@AddToHook{SelectCharTable}
    {\lst@alsoother \lst@alsodigit \lst@alsoletter}
\lst@AddToHookExe{SetLanguage}% init
    {\let\lst@alsoletter\@empty
     \let\lst@alsodigit\@empty
     \let\lst@alsoother\@empty}
%    \end{macrocode}
% 服务宏开始一个循环，并\ldots
%    \begin{macrocode}
\def\lst@DoAlso#1#2#3{%
    \lst@DefOther\lst@arg{#1}\let#2\@empty
    \expandafter\lst@DoAlso@\expandafter#2\expandafter#3\lst@arg\relax}
\def\lst@DoAlso@#1#2#3{%
    \ifx\relax#3\expandafter\@gobblethree \else
%    \end{macrocode}
% \ldots\ 在未到达 |\relax| 时，我们采用来自
% |\lst@SaveOutputDef| 的 \TeX 技巧，把类替换为 |#2|。最后我们把
% 新定义追加到 |#1|。
%    \begin{macrocode}
        \begingroup \lccode`\~=`#3\relax \lowercase{\endgroup
        \def\lst@temp##1\def~##2##3\relax{%
            \edef\lst@arg{\def\noexpand~{\noexpand#2\expandafter
                                         \noexpand\@gobble##2}}}}%
        \expandafter\lst@temp\lst@SelectStdCharTable\relax
        \lst@lExtend#1{\lst@arg}%
    \fi
    \lst@DoAlso@#1#2}
%    \end{macrocode}
% \end{lstkey}
% \end{lstkey}
% \end{lstkey}
%
% \begin{macro}{\lst@SaveDef}
% \begin{macro}{\lst@DefSaveDef}
% \begin{macro}{\lst@LetSaveDef}
% 这些宏可用于语言定义中做特殊改动。
% 它们保存定义，并定义或赋予新的定义。
%    \begin{macrocode}
\def\lst@SaveDef#1#2{%
    \begingroup \lccode`\~=#1\relax \lowercase{\endgroup\let#2~}}
\def\lst@DefSaveDef#1#2{%
    \begingroup \lccode`\~=#1\relax \lowercase{\endgroup\let#2~\def~}}
\def\lst@LetSaveDef#1#2{%
    \begingroup \lccode`\~=#1\relax \lowercase{\endgroup\let#2~\let~}}
%    \end{macrocode}
% \end{macro}
% \end{macro}
% \end{macro}
%
% 现在来看更强大的定义。
%
% \begin{macro}{\lst@CDef}
% 这里我们展开第一个参数 \meta{1st}\marg{2nd}\marg{rest}，并说
% 这个输入字符串被`替换'为 \meta{save 1st}\marg{2nd}^^A
% \marg{rest}——外加 \meta{execute}、\meta{pre} 和 \meta{post}。这项主要
% 工作由 |\lst@CDefIt| 完成。
%    \begin{macrocode}
\def\lst@CDef#1{\lst@CDef@#1}
\def\lst@CDef@#1#2#3#4{\lst@CDefIt#1{#2}{#3}{#4#2#3}#4}
%    \end{macrocode}
% \end{macro}
%
% \begin{macro}{\lst@CDefX}
% 丢弃输入字符串。
%    \begin{macrocode}
\def\lst@CDefX#1{\lst@CDefX@#1}
\def\lst@CDefX@#1#2#3{\lst@CDefIt#1{#2}{#3}{}}
%    \end{macrocode}
% \end{macro}
%
% \begin{macro}{\lst@CDefIt}
% 是前面几个宏的主工作过程。它重新定义字符
% 序列 |#1#2#3|。至少 |#1| 必须是活动的；另外两个
% 参数可能为空，但不是等价于空！
%    \begin{macrocode}
\def\lst@CDefIt#1#2#3#4#5#6#7#8{%
    \ifx\@empty#2\@empty
%    \end{macrocode}
% 对于单个字符，我们只按正确顺序执行这些参数。
% 你可能想回到第~\ref{dCharacterTablesManipulated} 节去
% 查阅它们。
%    \begin{macrocode}
        \def#1{#6\def\lst@next{#7#4#8}\lst@next}%
    \else \ifx\@empty#3\@empty
%    \end{macrocode}
% 对于两个字符的序列，我们测试是否必须执行 \meta{pre} 和 \meta{post}。
%    \begin{macrocode}
        \def#1##1{%
            #6%
            \ifx##1#2\def\lst@next{#7#4#8}\else
                     \def\lst@next{#5##1}\fi
            \lst@next}%
    \else
%    \end{macrocode}
% 对任意字符序列我们做同样的事——只是我们必须
% 使用 |\lst@IfNextCharsArg| 而不是 |\ifx|\ldots|\fi|。
%    \begin{macrocode}
        \def#1{%
            #6%
            \lst@IfNextCharsArg{#2#3}{#7#4#8}%
                                     {\expandafter#5\lst@eaten}}%
    \fi \fi}
%    \end{macrocode}
% \end{macro}
%
% \begin{macro}{\lst@CArgX}
% 我们把 |#1#2| 变成活动的，并调用 |\lst@CArg|。
%    \begin{macrocode}
\def\lst@CArgX#1#2\relax{%
    \lst@DefActive\lst@arg{#1#2}%
    \expandafter\lst@CArg\lst@arg\relax}
%    \end{macrocode}
% \end{macro}
%
% \begin{macro}{\lst@CArg}
% 为 |\lst@CDef|[|X|] 安排前两个参数。我们得到一个未定义的
% 宏，并使用 |\@empty\@empty\relax| 作为子宏的定界符。
%    \begin{macrocode}
\def\lst@CArg#1#2\relax{%
    \lccode`\/=`#1\lowercase{\def\lst@temp{/}}%
    \lst@GetFreeMacro{lst@c\lst@temp}%
    \expandafter\lst@CArg@\lst@freemacro#1#2\@empty\@empty\relax}
%    \end{macrocode}
% 把 \meta{1st}=|#2| 的含义保存在 \meta{save 1st}=|#1| 中，并用正确的
% 参数调用宏 |#6|。从 1.0 版起，|#2|、|#3| 和 |#4|
% （或空参数）用分组花括号绑定在一起。
% 这允许我们在其他定义中保存两个参数，例如在
% |\lst@DefDelimB|.
%    \begin{macrocode}
\def\lst@CArg@#1#2#3#4\@empty#5\relax#6{%
    \let#1#2%
    \ifx\@empty#3\@empty
        \def\lst@next{#6{#2{}{}}}%
    \else
        \def\lst@next{#6{#2#3{#4}}}%
    \fi
    \lst@next #1}
%    \end{macrocode}
% \end{macro}
%
% \begin{macro}{\lst@CArgEmpty}
% `执行'一个 |\@empty| 定界的参数。我们将把它用于定界符。
%    \begin{macrocode}
\def\lst@CArgEmpty#1\@empty{#1}
%    \end{macrocode}
% \end{macro}
%
%
% \subsection{定界符}
%
% 这里我们从所有定界符共有的通用定义开始。
%
% \begin{lstkey}{excludedelims}
% 控制哪些定界符不按 \meta{whatever}样式打印。我们只需
% 把 |\lst@ifex|\meta{whatever} 定义为真。这样的开关在
% \hookname{ExcludeDelims} 钩子中被置为假，并由各个
% 定界符处理。
%    \begin{macrocode}
\lst@Key{excludedelims}\relax
    {\lsthk@ExcludeDelims \lst@NormedDef\lst@temp{#1}%
     \expandafter\lst@for\lst@temp\do
     {\expandafter\let\csname\@lst @ifex##1\endcsname\iftrue}}
%    \end{macrocode}
% \end{lstkey}
%
% \begin{macro}{\lst@DelimPrint}
% 而这个宏可能有助于做到这一点。|#1| 是 |\lst@ifex|\meta{whatever}
% （外加 |\else|）或只是 |\iffalse|，而 |#2| 将是该定界符。
% 临时模式改动确保这些字符既不能结束当前
% 定界符，也不能开始新的定界符。
%    \begin{macrocode}
\def\lst@DelimPrint#1#2{%
    #1%
      \begingroup
        \lst@mode\lst@nomode \lst@modetrue
        #2\lst@XPrintToken
      \endgroup
      \lst@ResetToken
    \fi}
%    \end{macrocode}
% \end{macro}
%
% \begin{macro}{\lst@DelimOpen}
% 我们打印前面的字符和该定界符，进入合适的模式，
% 再次打印该定界符，并执行 |#3|。事实上，参数 |#1| 和
% |#2| 确保定界符只被打印一次。
%    \begin{macrocode}
\def\lst@DelimOpen#1#2#3#4#5#6\@empty{%
    \lst@TrackNewLines \lst@XPrintToken
    \lst@DelimPrint#1{#6}%
    \lst@EnterMode{#4}{\def\lst@currstyle#5}%
    \lst@DelimPrint{#1#2}{#6}%
    #3}
%    \end{macrocode}
% \end{macro}
%
% \begin{macro}{\lst@DelimClose}
% 是顺序相反的同一件事。
%    \begin{macrocode}
\def\lst@DelimClose#1#2#3\@empty{%
    \lst@TrackNewLines \lst@XPrintToken
    \lst@DelimPrint{#1#2}{#3}%
    \lst@LeaveMode
    \lst@DelimPrint{#1}{#3}}
%    \end{macrocode}
% \end{macro}
%
% \begin{macro}{\lst@BeginDelim}
% \begin{macro}{\lst@EndDelim}
% 这些定义是 |\lst@DelimOpen| 和 |\lst@DelimClose| 的应用：
% 定界符与被定界的文本样式相同。
%    \begin{macrocode}
\def\lst@BeginDelim{\lst@DelimOpen\iffalse\else{}}
\def\lst@EndDelim{\lst@DelimClose\iffalse\else}
%    \end{macrocode}
% \end{macro}
% \end{macro}
%
% \begin{macro}{\lst@BeginIDelim}
% \begin{macro}{\lst@EndIDelim}
% 另一个应用：不打印定界符。
%    \begin{macrocode}
\def\lst@BeginIDelim{\lst@DelimOpen\iffalse{}{}}
\def\lst@EndIDelim{\lst@DelimClose\iffalse{}}
%    \end{macrocode}
% \end{macro}
% \end{macro}
%
% \begin{macro}{\lst@DefDelims}
% 这个宏定义所有定界符，因此每次选择语言时
% 都被重置。
%    \begin{macrocode}
\lst@AddToHook{SelectCharTable}{\lst@DefDelims}
\lst@AddToHookExe{SetLanguage}{\let\lst@DefDelims\@empty}
%    \end{macrocode}
% \end{macro}
%
% \begin{macro}{\lst@Delim}
% 我们先设默认值：无 |\lst@modetrue|、累积样式，且
% |\lst@Delim|[|DM|]|@|\meta{type} 不带参数。
%    \begin{macrocode}
\def\lst@Delim#1{%
    \lst@false \let\lst@cumulative\@empty \let\lst@arg\@empty
%    \end{macrocode}
% 这些是双星形式（double-star-form）的正确设置，所以这种情况下我们
% 立即调用子宏。否则我们或者只是抑制累积
% 样式，或者甚至用 |\lst@true| 表明使用了 |\lst@modetrue|。
%    \begin{macrocode}
    \@ifstar{\@ifstar{\lst@Delim@{#1}}%
                     {\let\lst@cumulative\relax
                      \lst@Delim@{#1}}}%
            {\lst@true\lst@Delim@{#1}}}
%    \end{macrocode}
% 类型参数被保存以备后用。我们以 |#1| 为默认值
% 与可选的 \meta{style} 参数比较，定义 |\lst@delimstyle|，
% 并查找可选的 \meta{type option}，它只是被保存在 |\lst@arg| 中。
%    \begin{macrocode}
\def\lst@Delim@#1[#2]{%
    \gdef\lst@delimtype{#2}%
    \@ifnextchar[\lst@Delim@sty
                 {\lst@Delim@sty[#1]}}
\def\lst@Delim@sty[#1]{%
    \def\lst@delimstyle{#1}%
    \ifx\@empty#1\@empty\else
        \lst@Delim@sty@ #1\@nil
    \fi
    \@ifnextchar[\lst@Delim@option
                 \lst@Delim@delim}
\def\lst@Delim@option[#1]{\def\lst@arg{[#1]}\lst@Delim@delim}
%    \end{macrocode}
% 上面替换文本中的 |[| 和 |]| 是在
% \lsthelper{Stephen~Reindl}{2002/05/28}{使用 Cobol 时无法使用}报告缺陷之后添加的。
%
% |\lst@delimstyle| 的定义取决于第一个记号是否为
% 控制序列。这里我们可能构造 |\lst@|\meta{style}。
%    \begin{macrocode}
\def\lst@Delim@sty@#1#2\@nil{%
    \if\relax\noexpand#1\else
        \edef\lst@delimstyle{\expandafter\noexpand
                             \csname\@lst @\lst@delimstyle\endcsname}%
    \fi}
%    \end{macrocode}
% \end{macro}
%
% \begin{macro}{\lst@Delim@delim}
% 最终调用这个宏。首先我们可能需要删除一批
% 定界符。如果没有定界符，我们可能删除一个子类。
%    \begin{macrocode}
\def\lst@Delim@delim#1\relax#2#3#4#5#6#7#8{%
    \ifx #4\@empty \lst@Delim@delall{#2}\fi
    \ifx\@empty#1\@empty
        \ifx #4\@nil
            \@ifundefined{\@lst @#2DM@\lst@delimtype}%
                {\lst@Delim@delall{#2@\lst@delimtype}}%
                {\lst@Delim@delall{#2DM@\lst@delimtype}}%
        \fi
    \else
%    \end{macrocode}
% 如果定界符不为空，我们就转换它并追加到
% |\lst@arg|。|\lst@Begin|\ldots、|\lst@End|\ldots，以及样式和
% 模式选择同理。
%    \begin{macrocode}
        \expandafter\lst@Delim@args\expandafter
            {\lst@delimtype}{#1}{#5}#6{#7}{#8}#4%
%    \end{macrocode}
% 如果类型已知，我们或者选择动态模式或者选择静态模式，并把
% |\lst@arg| 的内容用作参数。所有这些都被放入 |\lst@delim|。
%    \begin{macrocode}
        \let\lst@delim\@empty
        \expandafter\lst@IfOneOf\lst@delimtype\relax#3%
        {\@ifundefined{\@lst @#2DM@\lst@delimtype}%
             {\lst@lExtend\lst@delim{\csname\@lst @#2@\lst@delimtype
                                     \expandafter\endcsname\lst@arg}}%
             {\lst@lExtend\lst@delim{\expandafter\lst@UseDynamicMode
                                     \csname\@lst @#2DM@\lst@delimtype
                                     \expandafter\endcsname\lst@arg}}%
%    \end{macrocode}
% 现在，根据模式 |#4|，我们或者删除这个特定定界符，
% 或者把它追加到当前所有定界符。
%    \begin{macrocode}
         \ifx #4\@nil
             \let\lst@temp\lst@DefDelims \let\lst@DefDelims\@empty
             \expandafter\lst@Delim@del\lst@temp\@empty\@nil\@nil\@nil
         \else
             \lst@lExtend\lst@DefDelims\lst@delim
         \fi}%
%    \end{macrocode}
% 未知类型会引发错误。
%    \begin{macrocode}
        {\PackageError{Listings}{Illegal type `\lst@delimtype'}%
                                {#2 types are #3.}}%
     \fi}
%    \end{macrocode}
% \end{macro}
%
% \begin{macro}{\lst@Delim@args}
% 现在我们看看如何把参数加入 |\lst@arg|。首先我们初始化
% 转换，只为让所有字符变活动。但如果类型的第一个字符
% 等于 |#4|，\ldots
%    \begin{macrocode}
\def\lst@Delim@args#1#2#3#4#5#6#7{%
    \begingroup
    \lst@false \let\lst@next\lst@XConvert
%    \end{macrocode}
% \ldots\ 我们从 |\lst@delimtype| 中删除该字符，而 |#5| 可能选择
% 不同的转换设置或宏。
%    \begin{macrocode}
    \@ifnextchar #4{\xdef\lst@delimtype{\expandafter\@gobble
                                        \lst@delimtype}%
                    #5\lst@next#2\@nil
                    \lst@lAddTo\lst@arg{\@empty#6}%
                    \lst@GobbleNil}%
%    \end{macrocode}
% 由于我们处在上面`特殊'的情况，我们还把特殊的
% |\lst@Begin|\ldots\ 和 |\lst@End|\ldots\ 宏加入了 |\lst@arg|
% （以及作为定界符刹车的 |\@empty|）。如果字符不相等，
% 就不需要做任何特殊处理。
%    \begin{macrocode}
                   {\lst@next#2\@nil
                    \lst@lAddTo\lst@arg{\@empty#3}%
                    \lst@GobbleNil}%
                 #1\@nil
%    \end{macrocode}
% 我们总是把参数转移到分组之外，并当且仅当我们不是在删除
% 定界符时追加样式和模式选择。
% 为此我们展开定界符样式。
%    \begin{macrocode}
    \global\let\@gtempa\lst@arg
    \endgroup
    \let\lst@arg\@gtempa
    \ifx #7\@nil\else
        \expandafter\lst@Delim@args@\expandafter{\lst@delimstyle}%
    \fi}
%    \end{macrocode}
% 回想样式是由 |\def\lst@currstyle#5| `选定'的，而这个
% `参数' |#5| 现在要被加入。根据开头的
% 设置，我们或者使用 |{\meta{style}}\lst@modetrue|——它选定
% 样式并停用关键词检测——，或者使用 |{}\meta{style}|——它
% 定义一个空样式宏并执行累积样式的样式——，
% 或者使用 |{\meta{style}|——它只定义样式宏。注意我们必须
% 在下面额外使用两个分组层级：一个由 |\lst@lAddTo| 直接丢弃，
% 另一个由 |\lst@Delim|[|DM|]|@|\meta{type} 丢弃。
%    \begin{macrocode}
\def\lst@Delim@args@#1{%
    \lst@if
        \lst@lAddTo\lst@arg{{{#1}\lst@modetrue}}%
    \else
        \ifx\lst@cumulative\@empty
            \lst@lAddTo\lst@arg{{{}#1}}%
        \else
            \lst@lAddTo\lst@arg{{{#1}}}%
        \fi
    \fi}
%    \end{macrocode}
% \end{macro}
%
% \begin{macro}{\lst@Delim@del}
% 为了删除某个特定定界符，我们沿定界符列表向下迭代，
% 把当前项与用户提供的那一个比较。
%    \begin{macrocode}
\def\lst@Delim@del#1\@empty#2#3#4{%
    \ifx #2\@nil\else
        \def\lst@temp{#1\@empty#2#3}%
        \ifx\lst@temp\lst@delim\else
            \lst@lAddTo\lst@DefDelims{#1\@empty#2#3{#4}}%
        \fi
        \expandafter\lst@Delim@del
    \fi}
%    \end{macrocode}
% \end{macro}
%
% \begin{macro}{\lst@Delim@delall}
% 为了删除整类定界符，我们首先展开控制序列
% 名，初始化其他数据，并调用一个子宏来做这项工作。
%    \begin{macrocode}
\def\lst@Delim@delall#1{%
    \begingroup
    \edef\lst@delim{\expandafter\string\csname\@lst @#1\endcsname}%
    \lst@false \global\let\@gtempa\@empty
    \expandafter\lst@Delim@delall@\lst@DefDelims\@empty
    \endgroup
    \let\lst@DefDelims\@gtempa}
%    \end{macrocode}
% 我们首先丢弃前面的 |\lst@UseDynamicMode|。
%    \begin{macrocode}
\def\lst@Delim@delall@#1{%
    \ifx #1\@empty\else
        \ifx #1\lst@UseDynamicMode
            \lst@true
            \let\lst@next\lst@Delim@delall@do
        \else
            \def\lst@next{\lst@Delim@delall@do#1}%
        \fi
        \expandafter\lst@next
    \fi}
%    \end{macrocode}
% 然后我们就可以检查（后面的）|\lst@|\meta{delimiter name}\ldots\
% 是否与 |\lst@delim| 给出的定界符类匹配。
%    \begin{macrocode}
\def\lst@Delim@delall@do#1#2\@empty#3#4#5{%
    \expandafter\lst@IfSubstring\expandafter{\lst@delim}{\string#1}%
      {}%
      {\lst@if \lst@AddTo\@gtempa\lst@UseDynamicMode \fi
       \lst@AddTo\@gtempa{#1#2\@empty#3#4{#5}}}%
    \lst@false \lst@Delim@delall@}
%    \end{macrocode}
% \end{macro}
%
% \begin{macro}{\lst@DefDelimB}
% 这里我们把参数拼在一起，以适应 |\lst@CDef|。注意传给
% |\lst@CDef| 的最后一个参数 |\@empty| 是 |\lst@CArgEmpty|
% 和 |\lst@DelimOpen| 的刹车。
%    \begin{macrocode}
\gdef\lst@DefDelimB#1#2#3#4#5#6#7#8{%
    \lst@CDef{#1}#2%
        {#3}%
        {\let\lst@bnext\lst@CArgEmpty
         \lst@ifmode #4\else
             #5%
             \def\lst@bnext{#6{#7}{#8}}%
         \fi
         \lst@bnext}%
        \@empty}
%    \end{macrocode}
% 在 \lsthelper{Vespe~Savikko}{2000/11/06}{同时加载 HTML 和 Python 时
% doc-string 输出异常}报告缺陷之后，我在 |#7| 周围加了花括号。
% \end{macro}
%
% \begin{macro}{\lst@DefDelimE}
% 第 5 行的 |\ifnum #7=\lst@mode| 确保各定界符
% 彼此匹配。
%    \begin{macrocode}
\gdef\lst@DefDelimE#1#2#3#4#5#6#7{%
    \lst@CDef{#1}#2%
        {#3}%
        {\let\lst@enext\lst@CArgEmpty
         \ifnum #7=\lst@mode%
             #4%
             \let\lst@enext#6%
         \else
             #5%
         \fi
         \lst@enext}%
        \@empty}
%    \end{macrocode}
%    \begin{macrocode}
\lst@AddToHook{Init}{\let\lst@bnext\relax \let\lst@enext\relax}
%    \end{macrocode}
% \end{macro}
%
% \begin{macro}{\lst@DefDelimBE}
% 这个服务宏将实际定义所有字符串定界符。
%    \begin{macrocode}
\gdef\lst@DefDelimBE#1#2#3#4#5#6#7#8#9{%
    \lst@CDef{#1}#2%
        {#3}%
        {\let\lst@bnext\lst@CArgEmpty
         \ifnum #7=\lst@mode
             #4%
             \let\lst@bnext#9%
         \else
             \lst@ifmode\else
                 #5%
                 \def\lst@bnext{#6{#7}{#8}}%
             \fi
         \fi
         \lst@bnext}%
        \@empty}
%    \end{macrocode}
% \end{macro}
%
% \begin{macro}{\lst@delimtypes}
% 是通用定界符类型的列表。
%    \begin{macrocode}
\gdef\lst@delimtypes{s,l}
%    \end{macrocode}
% \end{macro}
%
% \begin{macro}{\lst@DelimKey}
% 我们只是把用于 |\lst@Delim| 的参数拼起来。
%    \begin{macrocode}
\gdef\lst@DelimKey#1#2{%
    \lst@Delim{}#2\relax
        {Delim}\lst@delimtypes #1%
                {\lst@BeginDelim\lst@EndDelim}
        i\@empty{\lst@BeginIDelim\lst@EndIDelim}}
%    \end{macrocode}
% \end{macro}
%
% \begin{lstkey}{delim}
% \begin{lstkey}{moredelim}
% \begin{lstkey}{deletedelim}
% 都使用 |\lst@DelimKey|。
%    \begin{macrocode}
\lst@Key{delim}\relax{\lst@DelimKey\@empty{#1}}
\lst@Key{moredelim}\relax{\lst@DelimKey\relax{#1}}
\lst@Key{deletedelim}\relax{\lst@DelimKey\@nil{#1}}
%    \end{macrocode}
% \end{lstkey}
% \end{lstkey}
% \end{lstkey}
%
% \begin{macro}{\lst@DelimDM@l}
% \begin{macro}{\lst@DelimDM@s}
% 这里没什么特别的。
%    \begin{macrocode}
\gdef\lst@DelimDM@l#1#2\@empty#3#4#5{%
    \lst@CArg #2\relax\lst@DefDelimB{}{}{}#3{#1}{#5\lst@Lmodetrue}}
%    \end{macrocode}
%    \begin{macrocode}
\gdef\lst@DelimDM@s#1#2#3\@empty#4#5#6{%
    \lst@CArg #2\relax\lst@DefDelimB{}{}{}#4{#1}{#6}%
    \lst@CArg #3\relax\lst@DefDelimE{}{}{}#5{#1}}
%    \end{macrocode}
%    \begin{macrocode}
%</kernel>
%    \end{macrocode}
% \end{macro}
% \end{macro}
%
%
% \subsubsection{字符串}
%
% \begin{aspect}{strings}
% 只是开始一个新的方面。
%    \begin{macrocode}
%<*misc>
\lst@BeginAspect{strings}
%    \end{macrocode}
%
% \begin{macro}{\lst@stringtypes}
% 是 \ldots\ 字符串类型的列表？若干有效的字符串
% 类型定义在第~\pageref{uStrings} 页。
%    \begin{macrocode}
\gdef\lst@stringtypes{d,b,m,bd,db,s}
%    \end{macrocode}
% \end{macro}
%
% \begin{macro}{\lst@StringKey}
% 我们只是把用于 |\lst@Delim| 的参数拼起来。
%    \begin{macrocode}
\gdef\lst@StringKey#1#2{%
    \lst@Delim\lst@stringstyle #2\relax
        {String}\lst@stringtypes #1%
                     {\lst@BeginString\lst@EndString}%
        \@@end\@empty{}}
%    \end{macrocode}
% \end{macro}
%
% \begin{lstkey}{string}
% \begin{lstkey}{morestring}
% \begin{lstkey}{deletestring}
% 都使用 |\lst@StringKey|。
%    \begin{macrocode}
\lst@Key{string}\relax{\lst@StringKey\@empty{#1}}
\lst@Key{morestring}\relax{\lst@StringKey\relax{#1}}
\lst@Key{deletestring}\relax{\lst@StringKey\@nil{#1}}
%    \end{macrocode}
% \end{lstkey}
% \end{lstkey}
% \end{lstkey}
%
% \begin{lstkey}{stringstyle}
% 你不需要对下面两行加注释，是吧？
%    \begin{macrocode}
\lst@Key{stringstyle}{}{\def\lst@stringstyle{#1}}
\lst@AddToHook{EmptyStyle}{\let\lst@stringstyle\@empty}
%    \end{macrocode}
% \end{lstkey}
%
% \begin{lstkey}{showstringspaces}
% 感谢 \lsthelper{Knut~M\"uller}{1997/04/28}{\blankstringtrue}
% 报告了 |\blankstringtrue|（现为 |showstringspaces=false|）的问题。
% 该问题已消失。
%    \begin{macrocode}
\lst@Key{showstringspaces}t[t]{\lstKV@SetIf{#1}\lst@ifshowstringspaces}
%    \end{macrocode}
% \end{lstkey}
%
% \begin{macro}{\lst@BeginString}
% 注意 |\lst@DelimOpen| 之后的记号是参数！这里唯一特别的
% 是，如有必要，我们在开始字符串之后切换到`保留空格'。
% \lsthelper{Vespe~Savikko}{2000/09/27}{stringstyle 也被用于
% 前面的其他字符}报告的缺陷由于使用 |\lst@DelimOpen| 而消失。
%    \begin{macrocode}
\gdef\lst@BeginString{%
    \lst@DelimOpen
        \lst@ifexstrings\else
        {\lst@ifshowstringspaces
             \lst@keepspacestrue
             \let\lst@outputspace\lst@visiblespace
         \fi}}
%    \end{macrocode}
%    \begin{macrocode}
\lst@AddToHookExe{ExcludeDelims}{\let\lst@ifexstrings\iffalse}
%    \end{macrocode}
% \end{macro}
%
% \begin{macro}{\lst@EndString}
% 同样，|\lst@DelimClose| 之后的两个记号是参数。
%    \begin{macrocode}
\gdef\lst@EndString{\lst@DelimClose\lst@ifexstrings\else}
%    \end{macrocode}
% \end{macro}
%
% 现在是所有 |\lst@StringDM@|\meta{type} 定义。
%
% \begin{macro}{\lst@StringDM@d}
% `d' 表示没有额外工作；|\lst@DefDelimBE| 之后的头三个参数
% 留空。其余用于开始和结束字符串。
%    \begin{macrocode}
\gdef\lst@StringDM@d#1#2\@empty#3#4#5{%
    \lst@CArg #2\relax\lst@DefDelimBE{}{}{}#3{#1}{#5}#4}
%    \end{macrocode}
% \end{macro}
%
% \begin{macro}{\lst@StringDM@b}
% |\lst@ifletter|\ldots|\fi| 是在以下缺陷报告之后插入的：
% \lsthelper{Daniel~Gerigk}{2001/10/25}{C++ 中不正确的字符串}和
% \lsthelper{Peter~Bartke}{2001/11/01}{C++ 中不正确的字符串}。如果最后
% 一个其他字符是反斜杠（第 4 行），我们就吃掉`结束字符串'记号
% 序列。
%    \begin{macrocode}
\gdef\lst@StringDM@b#1#2\@empty#3#4#5{%
    \let\lst@ifbstring\iftrue
    \lst@CArg #2\relax\lst@DefDelimBE
       {\lst@ifletter \lst@Output \lst@letterfalse \fi}%
       {\ifx\lst@lastother\lstum@backslash
            \expandafter\@gobblethree
        \fi}{}#3{#1}{#5}#4}
%    \end{macrocode}
%    \begin{macrocode}
\global\let\lst@ifbstring\iffalse % init
%    \end{macrocode}
% \lsthelper{Heiko~Heil}{2002/02/08}{字符串 '\\' 在定界符后
% 没有结束}报告了双反斜杠的问题。所以：
%    \begin{macrocode}
\lst@AddToHook{SelectCharTable}{%
    \lst@ifbstring
        \lst@CArgX \\\\\relax \lst@CDefX{}%
           {\lst@ProcessOther\lstum@backslash
            \lst@ProcessOther\lstum@backslash
            \let\lst@lastother\relax}%
           {}%
    \fi}
%    \end{macrocode}
% 对 |\lst@lastother| 的重置是在以下缺陷报告之后添加的：
% \lsthelper{Hermann~H\"uttler}{2002/10/05}{C++ 字符串 "..." \\ 不以第二个
% 双引号结束}以及 \lsthelper{Dan~Luecking}{2003/01/15}{字符串 "\\" 在第二个引号之后并未结束}。
% \end{macro}
%
% \begin{macro}{\lst@StringDM@bd}
% \begin{macro}{\lst@StringDM@db}
% 只是相同，并且与 |\lst@StringDM@b| 相同。
%    \begin{macrocode}
\global\let\lst@StringDM@bd\lst@StringDM@b
\global\let\lst@StringDM@db\lst@StringDM@bd
%    \end{macrocode}
% \end{macro}\end{macro}
%
% \begin{macro}{\lst@StringDM@m}
% 用于 Matlab。只有当最后一个字符不在下面这个例外字符
% 列表中时，我们才进入字符串模式：字母、数字、句点、
% 引号、右圆括号、右方括号和右花括号。第一个列表
% 是在以下缺陷报告之后扩充的：\lsthelper{Christian~Kindinger}
% {2002/03/??}{]' starts a string in Matlab}, \lsthelper{Benjamin~Schubert}
% {2003/02/05}{.' 在 Matlab 中开始字符串}，以及 \lsthelper{Stefan~Stoll}
% {2003/02/18}{0123456789\} 中任一个加引号都会在 Matlab 中开始字符串}。
%    \begin{macrocode}
\gdef\lst@StringDM@m#1#2\@empty#3#4#5{%
    \lst@CArg #2\relax\lst@DefDelimBE{}{}%
        {\let\lst@next\@gobblethree
         \lst@ifletter\else
             \lst@IfLastOtherOneOf{)].0123456789\lstum@rbrace'}%
                 {}%
                 {\let\lst@next\@empty}%
         \fi
         \lst@next}#3{#1}{#5}#4}
%    \end{macrocode}
% \end{macro}
%
% \begin{macro}{\lst@StringDM@s}
% 用于以字符串定界的字符串，就像注释那样。这需要
% 用于 Ruby，可能还有其他语言。
%    \begin{macrocode}
\gdef\lst@StringDM@s#1#2#3\@empty#4#5#6{%
    \lst@CArg #2\relax\lst@DefDelimB{}{}{}#4{#1}{#6}%
    \lst@CArg #3\relax\lst@DefDelimE{}{}{}#5{#1}}
%    \end{macrocode}
% \end{macro}
%
% \begin{macro}{\lstum@rbrace}
% 这在上面已用过。
%    \begin{macrocode}
\lst@SaveOutputDef{"7D}\lstum@rbrace
%    \end{macrocode}
% \end{macro}
%
%    \begin{macrocode}
\lst@EndAspect
%</misc>
%    \end{macrocode}
% \end{aspect}
%
%
% \begin{aspect}{mf}
% 对 MetaFont 和 MetaPost，我们现在定义以 stringstyle 打印
% 输入文件名的宏。
%    \begin{macrocode}
%<*misc>
\lst@BeginAspect{mf}
%    \end{macrocode}
%
% \begin{macro}{\lst@mfinputmode}
% \begin{macro}{\lst@String@mf}
%    \begin{macrocode}
\lst@AddTo\lst@stringtypes{,mf}
\lst@NewMode\lst@mfinputmode
%    \end{macrocode}
%    \begin{macrocode}
\gdef\lst@String@mf#1\@empty#2#3#4{%
  \lst@CArg #1\relax\lst@DefDelimB
       {}{}{\lst@ifletter \expandafter\@gobblethree \fi}%
       \lst@BeginStringMFinput\lst@mfinputmode{#4\lst@Lmodetrue}%
  \@ifundefined{lsts@semicolon}%
  {\lst@DefSaveDef{`\;}\lsts@semicolon{% ; and space end the filename
      \ifnum\lst@mode=\lst@mfinputmode
          \lst@XPrintToken
          \expandafter\lst@LeaveMode
      \fi
      \lsts@semicolon}%
   \lst@DefSaveDef{`\ }\lsts@space{%
      \ifnum\lst@mode=\lst@mfinputmode
          \lst@XPrintToken
          \expandafter\lst@LeaveMode
      \fi
      \lsts@space}%
  }{}}
%    \end{macrocode}
% \end{macro}
% \end{macro}
%
% \begin{macro}{\lst@BeginStringMFinput}
% 还剩下定义这个宏。与 |\lst@PrintDelim| 不同，我们不
% 使用 |\lst@modetrue|，以便这里允许关键词检测。
%    \begin{macrocode}
\gdef\lst@BeginStringMFinput#1#2#3\@empty{%
    \lst@TrackNewLines \lst@XPrintToken
      \begingroup
        \lst@mode\lst@nomode
        #3\lst@XPrintToken
      \endgroup
      \lst@ResetToken
    \lst@EnterMode{#1}{\def\lst@currstyle#2}%
    \lst@ifshowstringspaces
         \lst@keepspacestrue
         \let\lst@outputspace\lst@visiblespace
    \fi}
%    \end{macrocode}
% \end{macro}
%
%    \begin{macrocode}
\lst@EndAspect
%</misc>
%    \end{macrocode}
% \end{aspect}
%
%
% \subsubsection{注释}
%
% \begin{aspect}{comments}
% 这正是我们正在做的事。
%    \begin{macrocode}
%<*misc>
\lst@BeginAspect{comments}
%    \end{macrocode}
%
% \begin{macro}{\lst@commentmode}
% 是注释的通用模式。
%    \begin{macrocode}
\lst@NewMode\lst@commentmode
%    \end{macrocode}
% \end{macro}
%
% \begin{macro}{\lst@commenttypes}
% 通过 \keyname{comment} 可用的注释类型：\textbf{行式}、\textbf{固定
% 列}、\textbf{单行}，以及 \textbf{嵌套}，并且都允许带
% 前置的 \textbf{i} 表示不可见注释。
%    \begin{macrocode}
\gdef\lst@commenttypes{l,f,s,n}
%    \end{macrocode}
% \end{macro}
%
% \begin{macro}{\lst@CommentKey}
% 我们只是把用于 |\lst@Delim| 的参数拼起来。
%    \begin{macrocode}
\gdef\lst@CommentKey#1#2{%
    \lst@Delim\lst@commentstyle #2\relax
        {Comment}\lst@commenttypes #1%
                {\lst@BeginComment\lst@EndComment}%
        i\@empty{\lst@BeginInvisible\lst@EndInvisible}}
%    \end{macrocode}
% \end{macro}
%
% \begin{lstkey}{comment}
% \begin{lstkey}{morecomment}
% \begin{lstkey}{deletecomment}
% 这些键很容易，因为它们用 |\lst@CommentKey| 定义。
%    \begin{macrocode}
\lst@Key{comment}\relax{\lst@CommentKey\@empty{#1}}
\lst@Key{morecomment}\relax{\lst@CommentKey\relax{#1}}
\lst@Key{deletecomment}\relax{\lst@CommentKey\@nil{#1}}
%    \end{macrocode}
% \end{lstkey}
% \end{lstkey}
% \end{lstkey}
%
% \begin{lstkey}{commentstyle}
% 还需要什么提示吗？
%    \begin{macrocode}
\lst@Key{commentstyle}{}{\def\lst@commentstyle{#1}}
\lst@AddToHook{EmptyStyle}{\let\lst@commentstyle\itshape}
%    \end{macrocode}
% \end{lstkey}
%
% \begin{macro}{\lst@BeginComment}
% \begin{macro}{\lst@EndComment}
% 再一次，|\lst@DelimOpen| 之后的三个记号是参数。
%    \begin{macrocode}
\gdef\lst@BeginComment{%
    \lst@DelimOpen
        \lst@ifexcomments\else
        \lsthk@AfterBeginComment}
%    \end{macrocode}
% 同上。
%    \begin{macrocode}
\gdef\lst@EndComment{\lst@DelimClose\lst@ifexcomments\else}
%    \end{macrocode}
%    \begin{macrocode}
\lst@AddToHook{AfterBeginComment}{}
\lst@AddToHookExe{ExcludeDelims}{\let\lst@ifexcomments\iffalse}
%    \end{macrocode}
% \end{macro}
% \end{macro}
%
% \begin{macro}{\lst@BeginInvisible}
% \begin{macro}{\lst@EndInvisible}
% 打印前面的字符并开始丢弃输出。
%    \begin{macrocode}
\gdef\lst@BeginInvisible#1#2#3\@empty{%
    \lst@TrackNewLines \lst@XPrintToken
    \lst@BeginDropOutput{#1}}
%    \end{macrocode}
% 不打印定界符并结束丢弃输出。
%    \begin{macrocode}
\gdef\lst@EndInvisible#1\@empty{\lst@EndDropOutput}
%    \end{macrocode}
% \end{macro}
% \end{macro}
%
% 现在我们提供所有 |\lst@Comment|[|DM|]|@|\meta{type} 宏。
%
% \begin{macro}{\lst@CommentDM@l}
% 很容易——感谢 |\lst@CArg| 和 |\lst@DefDelimB|。注意
% `结束注释'参数 |#4| 在这里没有使用。
%    \begin{macrocode}
\gdef\lst@CommentDM@l#1#2\@empty#3#4#5{%
    \lst@CArg #2\relax\lst@DefDelimB{}{}{}#3{#1}{#5\lst@Lmodetrue}}
%    \end{macrocode}
% \end{macro}
%
% \begin{macro}{\lst@CommentDM@f}
% 要多做一点工作。首先我们提供前面列数。
%    \begin{macrocode}
\gdef\lst@CommentDM@f#1{%
    \@ifnextchar[{\lst@Comment@@f{#1}}%
                 {\lst@Comment@@f{#1}[0]}}
%    \end{macrocode}
% 我们按上面同样的方式定义注释，但我们进入注释模式
% 当且仅当该字符位于第 |#2| 列（从零开始计数）。
%    \begin{macrocode}
\gdef\lst@Comment@@f#1[#2]#3\@empty#4#5#6{%
    \lst@CArg #3\relax\lst@DefDelimB{}{}%
        {\lst@CalcColumn
         \ifnum #2=\@tempcnta\else
             \expandafter\@gobblethree
         \fi}%
        #4{#1}{#6\lst@Lmodetrue}}
%    \end{macrocode}
% \end{macro}
%
% \begin{macro}{\lst@CommentDM@s}
% 这里没什么特别的。
%    \begin{macrocode}
\gdef\lst@CommentDM@s#1#2#3\@empty#4#5#6{%
    \lst@CArg #2\relax\lst@DefDelimB{}{}{}#4{#1}{#6}%
    \lst@CArg #3\relax\lst@DefDelimE{}{}{}#5{#1}}
%    \end{macrocode}
% \end{macro}
%
% \begin{macro}{\lst@CommentDM@n}
% 我们或者给出错误消息，或者定义嵌套注释。
%    \begin{macrocode}
\gdef\lst@CommentDM@n#1#2#3\@empty#4#5#6{%
    \ifx\@empty#3\@empty\else
        \def\@tempa{#2}\def\@tempb{#3}%
        \ifx\@tempa\@tempb
            \PackageError{Listings}{Identical delimiters}%
            {These delimiters make no sense with nested comments.}%
        \else
            \lst@CArg #2\relax\lst@DefDelimB
                {}%
%    \end{macrocode}
% 注意下面的 |\@gobble| 会吃掉来自 |\lst@DefDelimB| 的一个 |\else|。
%    \begin{macrocode}
                {\ifnum\lst@mode=#1\relax \expandafter\@gobble \fi}%
                {}#4{#1}{#6}%
            \lst@CArg #3\relax\lst@DefDelimE{}{}{}#5{#1}%
        \fi
    \fi}
%    \end{macrocode}
% \end{macro}
%
%    \begin{macrocode}
\lst@EndAspect
%</misc>
%    \end{macrocode}
% \end{aspect}
%
%
% \subsubsection{POD 注释}
%
% \begin{aspect}{pod}
% POD 被定义为一个单独的方面。
%    \begin{macrocode}
%<*misc>
\lst@BeginAspect{pod}
%    \end{macrocode}
%
% \begin{lstkey}{printpod}
% \begin{lstkey}{podcomment}
% 我们从用户键开始，这些键是我在与以下各方交流之后引入的
% \lsthelper{Michael~Piotrowski}{1997/11/11}{printpod}.
%    \begin{macrocode}
\lst@Key{printpod}{false}[t]{\lstKV@SetIf{#1}\lst@ifprintpod}
\lst@Key{podcomment}{false}[t]{\lstKV@SetIf{#1}\lst@ifpodcomment}
\lst@AddToHookExe{SetLanguage}{\let\lst@ifpodcomment\iffalse}
%    \end{macrocode}
% \end{lstkey}
% \end{lstkey}
%
% \begin{macro}{\lst@PODmode}
% 是 POD 的静态模式。
%    \begin{macrocode}
\lst@NewMode\lst@PODmode
%    \end{macrocode}
% \end{macro}
%
% 如果用户选择了 |podcomment=true|，我们就调整某些字符。
%    \begin{macrocode}
\lst@AddToHook{SelectCharTable}
    {\lst@ifpodcomment
         \lst@CArgX =\relax\lst@DefDelimB{}{}%
%    \end{macrocode}
% 如果找到等号并且尚未进入模式（事实上是允许模式改动时），
% 就执行下面的代码：我们`开始丢弃输出'，并
% 在 POD 不打印时吃掉通常的注释开始序列（通过 |\@gobblethree|）。
% 此外，如果当前列号不是零，我们也吃掉它
% ——|\@tempcnta| 在下面被赋值。
%    \begin{macrocode}
           {\ifnum\@tempcnta=\z@
                \lst@ifprintpod\else
                    \def\lst@bnext{\lst@BeginDropOutput\lst@PODmode}%
                    \expandafter\expandafter\expandafter\@gobblethree
                \fi
            \else
               \expandafter\@gobblethree
            \fi}%
           \lst@BeginComment\lst@PODmode{{\lst@commentstyle}}%
%    \end{macrocode}
% 如果遇到 |=|，我们就计算当前列号（从零起算）。
%    \begin{macrocode}
         \lst@CArgX =cut\^^M\relax\lst@DefDelimE
           {\lst@CalcColumn}%
%    \end{macrocode}
% 如果另外还有 |cut|+EOL，并且我们处在 |\lst@PODmode| 但不在
% 第一列，我们就必须吃掉`结束注释序列'。
%    \begin{macrocode}
           {\ifnum\@tempcnta=\z@\else
                \expandafter\@gobblethree
            \fi}%
           {}%
           \lst@EndComment\lst@PODmode
     \fi}
%    \end{macrocode}
%
%    \begin{macrocode}
\lst@EndAspect
%</misc>
%    \end{macrocode}
% \end{aspect}
%
%
% \subsubsection{标签}
%
% \begin{aspect}{html}
% 对 HTML 和其他`标记语言'的支持。
%    \begin{macrocode}
%<*misc>
\lst@BeginAspect[keywords]{html}
%    \end{macrocode}
%
% \begin{macro}{\lst@tagtypes}
% 我们再次从标签类型列表开始。它相当短。
%    \begin{macrocode}
\gdef\lst@tagtypes{s}
%    \end{macrocode}
% \end{macro}
%
% \begin{macro}{\lst@TagKey}
% 我们再次只把用于 |\lst@Delim| 的参数拼起来，并\ldots
%    \begin{macrocode}
\gdef\lst@TagKey#1#2{%
    \lst@Delim\lst@tagstyle #2\relax
        {Tag}\lst@tagtypes #1%
                     {\lst@BeginTag\lst@EndTag}%
        \@@end\@empty{}}
%    \end{macrocode}
% \end{macro}
%
% \begin{lstkey}{tag}
% \ldots\ 我们使用这里的定义。
%    \begin{macrocode}
\lst@Key{tag}\relax{\lst@TagKey\@empty{#1}}
%    \end{macrocode}
% \end{lstkey}
%
% \begin{lstkey}{tagstyle}
% 你不需要对下面两行加注释，是吧？
%    \begin{macrocode}
\lst@Key{tagstyle}{}{\def\lst@tagstyle{#1}}
\lst@AddToHook{EmptyStyle}{\let\lst@tagstyle\@empty}
%    \end{macrocode}
% \end{lstkey}
%
% \begin{macro}{\lst@BeginTag}
% 这里的特别之处是：(1) 我们在标签内部激活关键词检测，
% (2) 我们在必要时初始化开关 |\lst@iffirstintag|。
%    \begin{macrocode}
\gdef\lst@BeginTag{%
    \lst@DelimOpen
        \lst@ifextags\else
        {\let\lst@ifkeywords\iftrue
         \lst@ifmarkfirstintag \lst@firstintagtrue \fi}}
%    \end{macrocode}
%    \begin{macrocode}
\lst@AddToHookExe{ExcludeDelims}{\let\lst@ifextags\iffalse}
%    \end{macrocode}
% \end{macro}
%
% \begin{macro}{\lst@EndTag}
% 就和其他 |\lst@End|\meta{whatever} 定义一样。
%    \begin{macrocode}
\gdef\lst@EndTag{\lst@DelimClose\lst@ifextags\else}
%    \end{macrocode}
% \end{macro}
%
% \begin{lstkey}{usekeywordsintag}
% \begin{lstkey}{markfirstintag}
% 第二个键已被`使用'过。
%    \begin{macrocode}
\lst@Key{usekeywordsintag}t[t]{\lstKV@SetIf{#1}\lst@ifusekeysintag}
\lst@Key{markfirstintag}f[t]{\lstKV@SetIf{#1}\lst@ifmarkfirstintag}
%    \end{macrocode}
% 为此，我们安装一个（全局）开关，\ldots
%    \begin{macrocode}
\gdef\lst@firstintagtrue{\global\let\lst@iffirstintag\iftrue}
\global\let\lst@iffirstintag\iffalse
%    \end{macrocode}
% \ldots\ 它由标识符的输出重置，但不被其他
% 输出重置。
%    \begin{macrocode}
\lst@AddToHook{PostOutput}{\lst@tagresetfirst}
\lst@AddToHook{Output}
    {\gdef\lst@tagresetfirst{\global\let\lst@iffirstintag\iffalse}}
\lst@AddToHook{OutputOther}{\gdef\lst@tagresetfirst{}}
%    \end{macrocode}
% 现在我们只需在 \hookname{Output} 钩子中对该开关做测试。
%    \begin{macrocode}
\lst@AddToHook{Output}
    {\ifnum\lst@mode=\lst@tagmode
         \lst@iffirstintag \let\lst@thestyle\lst@gkeywords@sty \fi
%    \end{macrocode}
% 此外我们在这里检查关键词样式是否总要使用。
%    \begin{macrocode}
         \lst@ifusekeysintag\else \let\lst@thestyle\lst@gkeywords@sty\fi
     \fi}
%    \end{macrocode}
% \end{lstkey}
% \end{lstkey}
%
% \begin{macro}{\lst@tagmode}
% 我们分配模式和\ldots
%    \begin{macrocode}
\lst@NewMode\lst@tagmode
%    \end{macrocode}
% 如果定义了任何标签定界符（见下），就停用关键词检测。
%    \begin{macrocode}
\lst@AddToHook{Init}{\global\let\lst@ifnotag\iftrue}
\lst@AddToHook{SelectCharTable}{\let\lst@ifkeywords\lst@ifnotag}
%    \end{macrocode}
% \end{macro}
%
% \begin{macro}{\lst@Tag@s}
% 唯一一种定界符类型的定义不那么有趣。
% 与其他相比，我们设置 |\lst@ifnotag|，并且仅当
% 我们不在标签模式时才进入标签模式。
%    \begin{macrocode}
\gdef\lst@Tag@s#1#2\@empty#3#4#5{%
    \global\let\lst@ifnotag\iffalse
    \lst@CArg #1\relax\lst@DefDelimB {}{}%
        {\ifnum\lst@mode=\lst@tagmode \expandafter\@gobblethree \fi}%
        #3\lst@tagmode{#5}%
    \lst@CArg #2\relax\lst@DefDelimE {}{}{}#4\lst@tagmode}%
%    \end{macrocode}
% \end{macro}
%
% \begin{macro}{\lst@BeginCDATA}
% 这个宏由 XML 语言定义使用。
%    \begin{macrocode}
\gdef\lst@BeginCDATA#1\@empty{%
    \lst@TrackNewLines \lst@PrintToken
    \lst@EnterMode\lst@GPmode{}\let\lst@ifmode\iffalse
    \lst@mode\lst@tagmode #1\lst@mode\lst@GPmode\relax\lst@modetrue}
%    \end{macrocode}
% \end{macro}
%
%    \begin{macrocode}
\lst@EndAspect
%</misc>
%    \end{macrocode}
% \end{aspect}
%
%
% \subsection{替换输入}
%
% \begingroup
%    \begin{macrocode}
%<*kernel>
%    \end{macrocode}
% \endgroup
%
% \begin{macro}{\lst@ReplaceInput}
% 用 |\lst@CArgX| 和 |\lst@CDefX| 定义。
%    \begin{macrocode}
\def\lst@ReplaceInput#1{\lst@CArgX #1\relax\lst@CDefX{}{}}
%    \end{macrocode}
% \end{macro}
%
% \begin{lstkey}{literate}
% \lsthelper{Jason~Alexander}{1999/03/10}{文学化编程}请求
% 类似的功能。该键查找星号并保存参数。
%    \begin{macrocode}
\def\lst@Literatekey#1\@nil@{\let\lst@ifxliterate\lst@if
                             \def\lst@literate{#1}}
\lst@Key{literate}{}{\@ifstar{\lst@true \lst@Literatekey}
                             {\lst@false\lst@Literatekey}#1\@nil@}
\lst@AddToHook{SelectCharTable}
    {\ifx\lst@literate\@empty\else
         \expandafter\lst@Literate\lst@literate{}\relax\z@
     \fi}
%    \end{macrocode}
% 在内部我们不再使用`替换输入'特性。
%^^A We print the preceding text, assign token and length, and output it.
%    \begin{macrocode}
\def\lst@Literate#1#2#3{%
    \ifx\relax#2\@empty\else
        \lst@CArgX #1\relax\lst@CDef
            {}
            {\let\lst@next\@empty
             \lst@ifxliterate
                \lst@ifmode \let\lst@next\lst@CArgEmpty \fi
             \fi
             \ifx\lst@next\@empty
                 \ifx\lst@OutputBox\@gobble\else
                   \lst@XPrintToken \let\lst@scanmode\lst@scan@m
                   \lst@token{#2}\lst@length#3\relax
                   \lst@XPrintToken
                 \fi
                 \let\lst@next\lst@CArgEmptyGobble
             \fi
             \lst@next}%
            \@empty
        \expandafter\lst@Literate
    \fi}
\def\lst@CArgEmptyGobble#1\@empty{}
%    \end{macrocode}
% 注意我们检查 |\lst@OutputBox| 是否为 |\@gobble|。这是因为
% \lsthelper{Jared~Warren}{2003/07/10}{literate 替换
% 产生"幽灵"}的缺陷报告。
% \end{lstkey}
%
% \begin{macro}{\lst@BeginDropInput}
% 我们停用所有`处理'宏。|\lst@modetrue| 对所有
% 将要出现的字符串定界符、注释等做到这一点。
%    \begin{macrocode}
\def\lst@BeginDropInput#1{%
    \lst@EnterMode{#1}%
    {\lst@modetrue
     \let\lst@OutputBox\@gobble
     \let\lst@ifdropinput\iftrue
     \let\lst@ProcessLetter\@gobble
     \let\lst@ProcessDigit\@gobble
     \let\lst@ProcessOther\@gobble
     \let\lst@ProcessSpace\@empty
     \let\lst@ProcessTabulator\@empty
     \let\lst@ProcessFormFeed\@empty}}
\let\lst@ifdropinput\iffalse % init
%    \end{macrocode}
% \end{macro}
%
% \begingroup
%    \begin{macrocode}
%</kernel>
%    \end{macrocode}
% \endgroup
%
%
% \subsection{转义到 \LaTeX}
%
% \begin{aspect}{escape}
% 现在我们定义 \ldots\ 该死的 \ldots\ 方面跑了！
%    \begin{macrocode}
%<*misc>
\lst@BeginAspect{escape}
%    \end{macrocode}
%
% \begin{lstkey}{texcl}
% 与 \lsthelper{J\"orn~Wilms}{1997/07/07}{\TeX\ 注释}的交流
% 是这个键的由来。定义和前几个钩子很容易。
%    \begin{macrocode}
\lst@Key{texcl}{false}[t]{\lstKV@SetIf{#1}\lst@iftexcl}
\lst@AddToHook{TextStyle}{\let\lst@iftexcl\iffalse}
\lst@AddToHook{EOL}
    {\ifnum\lst@mode=\lst@TeXLmode
         \expandafter\lst@escapeend
         \expandafter\lst@LeaveAllModes
         \expandafter\lst@ReenterModes
     \fi}
%    \end{macrocode}
% 如果用户想要 \TeX\ 注释行，我们就打印注释分隔符并
% 中断正常处理。
%    \begin{macrocode}
\lst@AddToHook{AfterBeginComment}
    {\lst@iftexcl \lst@ifLmode \lst@ifdropinput\else
         \lst@PrintToken
         \lst@LeaveMode \lst@InterruptModes
         \lst@EnterMode{\lst@TeXLmode}{\lst@modetrue\lst@commentstyle}%
         \expandafter\expandafter\expandafter\lst@escapebegin
     \fi \fi \fi}
%    \end{macrocode}
%    \begin{macrocode}
\lst@NewMode\lst@TeXLmode
%    \end{macrocode}
% \end{lstkey}
%
% \begin{macro}{\lst@ActiveCDefX}
% 同 |\lst@CDefX|，但我们既把 |#1| 变成活动的，又给它赋新的 catcode。
%    \begin{macrocode}
\gdef\lst@ActiveCDefX#1{\lst@ActiveCDefX@#1}
\gdef\lst@ActiveCDefX@#1#2#3{%
    \catcode`#1\active\lccode`\~=`#1%
    \lowercase{\lst@CDefIt~}{#2}{#3}{}}
%    \end{macrocode}
% \end{macro}
%
% \begin{macro}{\lst@Escape}
% 总共得到四个参数。第一个和第二个是`开始'和
% `结束'转义序列，第三个在转义开始时执行，
% 第四个在结束转义之前执行。我们使用与 \TeX\ 注释
% 行相同的机制。|\lst@ifdropinput| 测试是在
% \lsthelper{Michael~Weber}{2002/03/26}{在 firstline 之前的行上转义会破坏
% 输出}报告缺陷之后添加的。|\lst@newlines\z@| 则是由于
% \lsthelper{Frank~Atanassow}{2004/10/07}{mathescape 后的空格未被
% 保留}的缺陷报告而添加的。
%    \begin{macrocode}
\gdef\lst@Escape#1#2#3#4{%
    \lst@CArgX #1\relax\lst@CDefX
        {}%
        {\lst@ifdropinput\else
         \lst@TrackNewLines\lst@OutputLostSpace \lst@XPrintToken
         \lst@InterruptModes
         \lst@EnterMode{\lst@TeXmode}{\lst@modetrue}%
%    \end{macrocode}
% 现在我们必须定义结束转义的字符序列。
%    \begin{macrocode}
         \ifx\^^M#2%
             \lst@CArg #2\relax\lst@ActiveCDefX
                 {}%
                 {\lst@escapeend #4\lst@LeaveAllModes\lst@ReenterModes}%
                 {\lst@MProcessListing}%
         \else
             \lst@CArg #2\relax\lst@ActiveCDefX
                 {}%
                 {\lst@escapeend #4\lst@LeaveAllModes\lst@ReenterModes
                  \lst@newlines\z@ \lst@whitespacefalse}%
                 {}%
         \fi
         #3\lst@escapebegin
         \fi}%
        {}}
%    \end{macrocode}
% 上面的 |\lst@whitespacefalse| 是在以下缺陷报告之后添加的
% \lsthelper{Martin~Steffen}{2001/04/07}{mathescape drops subsequent space}.
%    \begin{macrocode}
\lst@NewMode\lst@TeXmode
%    \end{macrocode}
% \end{macro}
%
% \begin{lstkey}{escapebegin}
% \begin{lstkey}{escapeend}
% 这些键只是保存参数。
%    \begin{macrocode}
\lst@Key{escapebegin}{}{\def\lst@escapebegin{#1}}
\lst@Key{escapeend}{}{\def\lst@escapeend{#1}}
%    \end{macrocode}
% \end{lstkey}
% \end{lstkey}
%
% \begin{lstkey}{escapechar}
% 这个键的引入源于与 \lsthelper
% {Rui~Oliveira}{1998/06/05}{转义字符}的交流。我们定义 |\lst@DefEsc|，
% 并在选择标准字符表之后执行它。
%    \begin{macrocode}
\lst@Key{escapechar}{}
    {\ifx\@empty#1\@empty
         \let\lst@DefEsc\relax
     \else
         \def\lst@DefEsc{\lst@Escape{#1}{#1}{}{}}%
     \fi}
\lst@AddToHook{TextStyle}{\let\lst@DefEsc\@empty}
\lst@AddToHook{SelectCharTable}{\lst@DefEsc}
%    \end{macrocode}
% \end{lstkey}
%
% \begin{lstkey}{escapeinside}
% 几乎相同。
%    \begin{macrocode}
\lst@Key{escapeinside}{}{\lstKV@TwoArg{#1}%
    {\let\lst@DefEsc\@empty
     \ifx\@empty##1@empty\else \ifx\@empty##2\@empty\else
         \def\lst@DefEsc{\lst@Escape{##1}{##2}{}{}}%
     \fi\fi}}
%    \end{macrocode}
% \end{lstkey}
%
% \begin{lstkey}{mathescape}
% 这是一个开关，在选择字符表之后检查。我们用
% |\lst@Escape| 以数学移位为参数，但全部放在 |\hbox| 内
% 以确定正确的宽度。
%    \begin{macrocode}
\lst@Key{mathescape}{false}[t]{\lstKV@SetIf{#1}\lst@ifmathescape}
\lst@AddToHook{SelectCharTable}
    {\lst@ifmathescape \lst@Escape{\$}{\$}%
        {\setbox\@tempboxa=\hbox\bgroup$}%
        {$\egroup \lst@CalcLostSpaceAndOutput}\fi}
%    \end{macrocode}
% \end{lstkey}
%
%    \begin{macrocode}
\lst@EndAspect
%</misc>
%    \end{macrocode}
% \end{aspect}
%
%
% \section{关键词}
%
%
% \subsection{构造测试}\label{iMakingTests}
%
% \begin{aspect}{keywords}
% 我们开始一个崭新而且非常重要的方面。
% 首先我们需要初始化一些变量，以绕开
% \lsthelper{Beat~Birkhofer}{2001/06/15}{savemem 不起作用}报告的缺陷。
%    \begin{macrocode}
%<*misc>
\lst@BeginAspect{keywords}
%    \end{macrocode}
%    \begin{macrocode}
\global\let\lst@ifsensitive\iftrue % init
\global\let\lst@ifsensitivedefed\iffalse % init % \global
%    \end{macrocode}
% 所有关键词测试都取下面三个参数。
% \begin{macroargs}
% \item \meta{prefix}
% \item |\lst@|\meta{name}|@list|（一个含关键词的宏的列表）
% \item |\lst@g|\meta{name}|@sty|（全局样式宏）
% \end{macroargs}
% 我们从非节省内存的测试开始。
% \begingroup
%    \begin{macrocode}
\lst@ifsavemem\else
%    \end{macrocode}
% \endgroup
%
% \begin{macro}{\lst@KeywordTest}
% 快速关键词测试利用了第~\ref{iCharacterTables} 节中的 |\lst@UM| 构造。
% 如果 |\lst@UM| 为空，所有`使用宏'字符
% 都展开为它们的原始字符。由于 |\lst|\meta{prefix}|@|\meta{keyword}
% 将等价于合适的样式，我们只需构造控制
% 序列 |\lst|\meta{prefix}|@|\meta{current token} 并把它赋给
% |\lst@thestyle|.
%    \begin{macrocode}
\gdef\lst@KeywordTest#1#2#3{%
    \begingroup \let\lst@UM\@empty
    \global\expandafter\let\expandafter\@gtempa
        \csname\@lst#1@\the\lst@token\endcsname
    \endgroup
    \ifx\@gtempa\relax\else
        \let\lst@thestyle\@gtempa
    \fi}
%    \end{macrocode}
% 注意我们这里既不需要 |#2| 也不需要 |#3|。
% \end{macro}
%
% \begin{macro}{\lst@KEYWORDTEST}
% 不区分大小写的测试把当前字符串变成大写，并
% 把它交给一个类似 |\lst@KeywordTest| 的子宏。
%    \begin{macrocode}
\gdef\lst@KEYWORDTEST{%
    \uppercase\expandafter{\expandafter
        \lst@KEYWORDTEST@\the\lst@token}\relax}
\gdef\lst@KEYWORDTEST@#1\relax#2#3#4{%
    \begingroup \let\lst@UM\@empty
    \global\expandafter\let\expandafter\@gtempa
        \csname\@lst#2@#1\endcsname
    \endgroup
    \ifx\@gtempa\relax\else
        \let\lst@thestyle\@gtempa
    \fi}
%    \end{macrocode}
% \end{macro}
%
% \begin{macro}{\lst@WorkingTest}
% \begin{macro}{\lst@WORKINGTEST}
% 相同，只是 |\lst|\meta{prefix}|@|\meta{current token} 可能
% 是一个工作过程；它被执行。
%    \begin{macrocode}
\gdef\lst@WorkingTest#1#2#3{%
    \begingroup \let\lst@UM\@empty
    \global\expandafter\let\expandafter\@gtempa
        \csname\@lst#1@\the\lst@token\endcsname
    \endgroup
    \@gtempa}
%    \end{macrocode}
%    \begin{macrocode}
\gdef\lst@WORKINGTEST{%
    \uppercase\expandafter{\expandafter
        \lst@WORKINGTEST@\the\lst@token}\relax}
\gdef\lst@WORKINGTEST@#1\relax#2#3#4{%
    \begingroup \let\lst@UM\@empty
    \global\expandafter\let\expandafter\@gtempa
        \csname\@lst#2@#1\endcsname
    \endgroup
    \@gtempa}
%    \end{macrocode}
% \end{macro}
% \end{macro}
%
% \begin{macro}{\lst@DefineKeywords}
% 最后我们需要定义和取消定义
% |\lst|\meta{prefix}|@|\meta{keyword} 的宏。这里是参数：
% \begin{macroargs}
% \item \meta{prefix}
% \item |\lst@|\meta{name}（一个关键词列表）
% \item |\lst@g|\meta{name}|@sty|
% \end{macroargs}
% 我们在必要时把关键词变成大写，\ldots
%    \begin{macrocode}
\gdef\lst@DefineKeywords#1#2#3{%
    \lst@ifsensitive
        \def\lst@next{\lst@for#2}%
    \else
        \def\lst@next{\uppercase\expandafter{\expandafter\lst@for#2}}%
    \fi
    \lst@next\do
%    \end{macrocode}
% \ldots\ 遍历列表，并（若未定义）把
% |\lst|\meta{prefix}|@|\meta{keyword} 设为等价于
% |\lst@g|\meta{name}|@sty|，它可能是一个工作宏。
%    \begin{macrocode}
    {\expandafter\ifx\csname\@lst#1@##1\endcsname\relax
        \global\expandafter\let\csname\@lst#1@##1\endcsname#3%
     \fi}}
%    \end{macrocode}
% \end{macro}
%
% \begin{macro}{\lst@UndefineKeywords}
% 我们在必要时把关键词变成大写，\ldots
%    \begin{macrocode}
\gdef\lst@UndefineKeywords#1#2#3{%
    \lst@ifsensitivedefed
        \def\lst@next{\lst@for#2}%
    \else
        \def\lst@next{\uppercase\expandafter{\expandafter\lst@for#2}}%
    \fi
    \lst@next\do
%    \end{macrocode}
% \ldots\ 遍历列表，并`取消定义'
% |\lst|\meta{prefix}|@|\meta{keyword}，如果它等价于
% |\lst@g|\meta{name}|@sty|。
%    \begin{macrocode}
    {\expandafter\ifx\csname\@lst#1@##1\endcsname#3%
        \global\expandafter\let\csname\@lst#1@##1\endcsname\relax
     \fi}}
%    \end{macrocode}
% 感谢 \lsthelper{Magnus~Lewis-Smith}{1999/09/08}{关键词
% 无法取消定义}，替换文本中错误的 |#2| 得以改成 |#3|。
% \end{macro}
%
% \begingroup
% 现在来看节省内存的测试。
%    \begin{macrocode}
\fi
\lst@ifsavemem
%    \end{macrocode}
% \endgroup
%
% \begin{macro}{\lst@IfOneOutOf}
% 这里的定义类似 |\lst@IfOneOf|，但它的第二个参数
% 是一个 |\lst@|\meta{name}|@list|。因此我们这里测试宏的列表。
%    \begin{macrocode}
\gdef\lst@IfOneOutOf#1\relax#2{%
    \def\lst@temp##1,#1,##2##3\relax{%
        \ifx\@empty##2\else \expandafter\lst@IOOOfirst \fi}%
    \def\lst@next{\lst@IfOneOutOf@#1\relax}%
    \expandafter\lst@next#2\relax\relax}
%    \end{macrocode}
% 我们或者执行 \meta{else} 部分，或者进行下一个测试。
%    \begin{macrocode}
\gdef\lst@IfOneOutOf@#1\relax#2#3{%
    \ifx#2\relax
        \expandafter\@secondoftwo
    \else
        \expandafter\lst@temp\expandafter,#2,#1,\@empty\relax
        \expandafter\lst@next
    \fi}
\ifx\iffalse\else\fi
\gdef\lst@IOOOfirst#1\relax#2#3{\fi#2}
%    \end{macrocode}
% 行 |\ifx\iffalse\else\fi| 平衡 |\lst@IOOOfirst| 内的 |\fi|。
% \end{macro}
%
% \begin{macro}{\lst@IFONEOUTOF}
% 如 |\lst@IFONEOF| 中那样，我们这里需要两个 |\uppercase|。
%    \begin{macrocode}
\gdef\lst@IFONEOUTOF#1\relax#2{%
    \uppercase{\def\lst@temp##1,#1},##2##3\relax{%
        \ifx\@empty##2\else \expandafter\lst@IOOOfirst \fi}%
    \def\lst@next{\lst@IFONEOUTOF@#1\relax}%
    \expandafter\lst@next#2\relax}
\gdef\lst@IFONEOUTOF@#1\relax#2#3{%
    \ifx#2\relax
        \expandafter\@secondoftwo
    \else
        \uppercase
            {\expandafter\lst@temp\expandafter,#2,#1,\@empty\relax}%
        \expandafter\lst@next
    \fi}
%    \end{macrocode}
% 注意：倒数第三行利用了这样一个事实：如果关键词不区分大小写，
% 关键词列表（不是关键词列表的列表）已经被转成大写。
% \end{macro}
%
% \begin{macro}{\lst@KWTest}
% 是关键词和工作标识符测试的辅助宏。我们展开
% 记号并调用 |\lst@IfOneOf|。下面的测试会追加合适的
% \meta{then} 和 \meta{else} 参数。
%    \begin{macrocode}
\gdef\lst@KWTest{%
    \begingroup \let\lst@UM\@empty
    \expandafter\xdef\expandafter\@gtempa\expandafter{\the\lst@token}%
    \endgroup
    \expandafter\lst@IfOneOutOf\@gtempa\relax}
%    \end{macrocode}
% \end{macro}
%
% \begin{macro}{\lst@KeywordTest}
% \begin{macro}{\lst@KEYWORDTEST}
% 现在相当容易了。注意我们这里不需要 |#1|=\meta{prefix}。
%    \begin{macrocode}
\gdef\lst@KeywordTest#1#2#3{\lst@KWTest #2{\let\lst@thestyle#3}{}}
\global\let\lst@KEYWORDTEST\lst@KeywordTest
%    \end{macrocode}
% 对于不区分大小写的测试，我们把不区分大小写的版本赋给
% |\lst@IfOneOutOf|。因此这里不需要额外定义。
% \end{macro}
% \end{macro}
%
% \begin{macro}{\lst@WorkingTest}
% \begin{macro}{\lst@WORKINGTEST}
% 同上。
%    \begin{macrocode}
\gdef\lst@WorkingTest#1#2#3{\lst@KWTest #2#3{}}
\global\let\lst@WORKINGTEST\lst@WorkingTest
%    \end{macrocode}
% \end{macro}
% \end{macro}
%
% \begingroup
%    \begin{macrocode}
\fi
%    \end{macrocode}
% \endgroup
%
% \begin{lstkey}{sensitive}
% 是一个开关，每次选择语言时预置为 \texttt{true}。
%    \begin{macrocode}
\lst@Key{sensitive}\relax[t]{\lstKV@SetIf{#1}\lst@ifsensitive}
\lst@AddToHook{SetLanguage}{\let\lst@ifsensitive\iftrue}
%    \end{macrocode}
% 我们在必要时选择不区分大小写的定义。
%    \begin{macrocode}
\lst@AddToHook{Init}
    {\lst@ifsensitive\else
         \let\lst@KeywordTest\lst@KEYWORDTEST
         \let\lst@WorkingTest\lst@WORKINGTEST
         \let\lst@IfOneOutOf\lst@IFONEOUTOF
     \fi}
%    \end{macrocode}
% \end{lstkey}
%
% \begin{macro}{\lst@MakeMacroUppercase}
% （若已定义）把 |#1| 的内容变成大写。
%    \begin{macrocode}
\gdef\lst@MakeMacroUppercase#1{%
    \ifx\@undefined#1\else \uppercase\expandafter
        {\expandafter\def\expandafter#1\expandafter{#1}}%
    \fi}
%    \end{macrocode}
% \end{macro}
%
%
% \subsection{安装测试}
%
% \begin{macro}{\lst@InstallTest}
% 参数是
% \begin{macroargs}
% \item \meta{prefix}
% \item |\lst@|\meta{name}|@list|
% \item |\lst@|\meta{name}
% \item |\lst@g|\meta{name}|@list|
% \item |\lst@g|\meta{name}
% \item |\lst@g|\meta{name}|@sty|
% \item \alternative{w,s}（工作过程或样式）
% \item \alternative{d,o}（\hookname{DetectKeywords} 或 \hookname{Output} 钩子）
% \end{macroargs}
% 我们只插入钩子材料。测试会按需插入。
%    \begin{macrocode}
\gdef\lst@InstallTest#1#2#3#4#5#6#7#8{%
    \lst@AddToHook{TrackKeywords}{\lst@TrackKeywords{#1}#2#4#6#7#8}%
    \lst@AddToHook{PostTrackKeywords}{\lst@PostTrackKeywords#2#3#4#5}}
%    \end{macrocode}
%    \begin{macrocode}
\lst@AddToHook{Init}{\lsthk@TrackKeywords\lsthk@PostTrackKeywords}
\lst@AddToHook{TrackKeywords}
    {\global\let\lst@DoDefineKeywords\@empty}% init
\lst@AddToHook{PostTrackKeywords}
    {\lst@DoDefineKeywords
     \global\let\lst@DoDefineKeywords\@empty}% init
%    \end{macrocode}
% 我们必须某处检测关键词。
%    \begin{macrocode}
\lst@AddToHook{Output}{\lst@ifkeywords \lsthk@DetectKeywords \fi}
\lst@AddToHook{DetectKeywords}{}% init
\lst@AddToHook{ModeTrue}{\let\lst@ifkeywords\iffalse}
\lst@AddToHookExe{Init}{\let\lst@ifkeywords\iftrue}
%    \end{macrocode}
% \end{macro}
%
% \begin{macro}{\lst@InstallTestNow}
% 实际插入一个测试。
% \begin{macroargs}
% \item \meta{prefix}
% \item |\lst@|\meta{name}|@list|
% \item |\lst@g|\meta{name}|@sty|
% \item \alternative{w,s}（工作过程或样式）
% \item \alternative{d,o}（\hookname{DetectKeywords} 或 \hookname{Output} 钩子）
% \end{macroargs}
% 例如，|#4#5|=|sd| 会把
%    |\lst@KeywordTest{|\meta{prefix}|}|
%             |\lst@|\meta{name}|@list| |\lst@g|\meta{name}|@sty|
% 加入 \hookname{DetectKeywords} 钩子。
%    \begin{macrocode}
\gdef\lst@InstallTestNow#1#2#3#4#5{%
    \@ifundefined{\string#2#1}%
    {\global\@namedef{\string#2#1}{}%
     \edef\@tempa{%
         \noexpand\lst@AddToHook{\ifx#5dDetectKeywords\else Output\fi}%
         {\ifx #4w\noexpand\lst@WorkingTest
             \else\noexpand\lst@KeywordTest \fi
          {#1}\noexpand#2\noexpand#3}}%
%    \end{macrocode}
% 如果有人建议我们节省内存，我们就为每个 \meta{name} 插入一个测试。
% 否则我们按 \meta{prefix} 安装测试。
%    \begin{macrocode}
     \lst@ifsavemem
         \@tempa
     \else
         \@ifundefined{\@lst#1@if@ins}%
             {\@tempa \global\@namedef{\@lst#1@if@ins}{}}%
             {}%
     \fi}
    {}}
%    \end{macrocode}
% \end{macro}
%
% \begin{macro}{\lst@TrackKeywords}
% 现在有点棘手。我们在 |\lst@TK@{|\meta{prefix}|}||\lst@g|\meta{name}|@sty| 之后
% 展开类列表 |\lst@|\meta{name}|@list|，并使用两个
% |\relax| 作为终止符。这会把所有类的关键词
% 定义为类型 \meta{prefix} 的关键词。更多细节稍后给出。
%    \begin{macrocode}
\gdef\lst@TrackKeywords#1#2#3#4#5#6{%
    \lst@false
    \def\lst@arg{{#1}#4}%
    \expandafter\expandafter\expandafter\lst@TK@
        \expandafter\lst@arg#2\relax\relax
%    \end{macrocode}
% 取消所有过时关键词也几乎相同，这仅在
% 我们不节省内存时才需要。
%    \begin{macrocode}
    \lst@ifsavemem\else
        \def\lst@arg{{#1}#4#2}%
        \expandafter\expandafter\expandafter\lst@TK@@
            \expandafter\lst@arg#3\relax\relax
    \fi
%    \end{macrocode}
% 最后，如果关键词改变（特别是它们第一次被定义时），
% 我们就安装关键词测试。注意 |\lst@InstallTestNow| 只
% 插入一次测试。
%    \begin{macrocode}
    \lst@if \lst@InstallTestNow{#1}#2#4#5#6\fi}
%    \end{macrocode}
% 回到当前的关键词。全局宏 |\lst@g|\meta{id} 包含
% 全局定义的关键词，而 |\lst@|\meta{id} 包含真正的
% 关键词。这样我们就能跟踪关键词：如果关键词或
% \keyname{sensitive} 改变了，我们取消定义旧（= 全局定义的）
% 关键词并定义真正的关键词。|\lst@TK@| 的参数是
% \begin{macroargs}
% \item \meta{prefix}
% \item |\lst@g|\meta{name}|@sty|
% \item |\lst@|\meta{id}
% \item |\lst@g|\meta{id}
% \end{macroargs}
% 感谢 \lsthelper{Holger~Arndt}{2004/05/27}{关键词在两种语言中
% 顺序交叉时 \lst@UndefineKeywords \lst@DefineKeywords 序列出错}
% 关键词的定义现在通过 |\lst@DoDefineKeywords| 延迟执行。
%    \begin{macrocode}
\gdef\lst@TK@#1#2#3#4{%
  \ifx\lst@ifsensitive\lst@ifsensitivedefed
    \ifx#3#4\else
      \lst@true
      \lst@ifsavemem\else
          \lst@UndefineKeywords{#1}#4#2%
          \lst@AddTo\lst@DoDefineKeywords{\lst@DefineKeywords{#1}#3#2}%
      \fi
    \fi
  \else
    \ifx#3\relax\else
      \lst@true
      \lst@ifsavemem\else
          \lst@UndefineKeywords{#1}#4#2%
          \lst@AddTo\lst@DoDefineKeywords{\lst@DefineKeywords{#1}#3#2}%
      \fi
    \fi
  \fi
%    \end{macrocode}
% 如果我们试图节省内存，就不定义也不取消定义关键词。但我们
% 可能需要把它们变成大写，这又会浪费一些内存。
%    \begin{macrocode}
  \lst@ifsavemem \ifx#3\relax\else
      \lst@ifsensitive\else \lst@MakeMacroUppercase#3\fi
  \fi \fi
%    \end{macrocode}
% 到达类列表末尾时，我们结束循环。
%    \begin{macrocode}
  \ifx#3\relax
      \expandafter\@gobblethree
  \fi
  \lst@TK@{#1}#2}
%    \end{macrocode}
% 现在我们取消定义过时的关键词。在未到达全局列表
% 末尾时，我们检查关键词类 |#4#5| 是否仍在使用，或
% 是否需要取消定义。我们的参数是
% \begin{macroargs}
% \item \meta{prefix}
% \item |\lst@g|\meta{name}|@sty|
% \item |\lst@|\meta{name}|@list|
% \item |\lst@|\meta{id}
% \item |\lst@g|\meta{id}
% \end{macroargs}
%    \begin{macrocode}
\gdef\lst@TK@@#1#2#3#4#5{%
    \ifx#4\relax
        \expandafter\@gobblefour
    \else
        \lst@IfSubstring{#4#5}#3{}{\lst@UndefineKeywords{#1}#5#2}%
    \fi
    \lst@TK@@{#1}#2#3}
%    \end{macrocode}
% 关键词在 \hookname{InitVars} 之后是最新的。
%    \begin{macrocode}
\lst@AddToHook{InitVars}
    {\global\let\lst@ifsensitivedefed\lst@ifsensitive}
%    \end{macrocode}
% \end{macro}
%
% \begin{macro}{\lst@PostTrackKeywords}
% 更新所有关键词之后，全局关键词和全局列表
% 变得等价于局部的那些。
%    \begin{macrocode}
\gdef\lst@PostTrackKeywords#1#2#3#4{%
    \lst@ifsavemem\else
        \global\let#3#1%
        \global\let#4#2%
    \fi}
%    \end{macrocode}
% \end{macro}
%
%
% \subsection{类与族}
%
% \begin{lstkey}{classoffset}
% 只是把参数存入一个宏。
%    \begin{macrocode}
\lst@Key{classoffset}\z@{\def\lst@classoffset{#1}}
%    \end{macrocode}
% \end{lstkey}
%
% \begin{macro}{\lst@InstallFamily}
% 回想参数
% \begin{macroargs}
% \item \meta{prefix}
% \item \meta{name}
% \item \meta{style name}
% \item \meta{style init}
% \item \meta{default style name}
% \item \meta{working procedure}
% \item \alternative{l,o}（语言或其他键）
% \item \alternative{d,o}（\hookname{DetectKeywords} 或 \hookname{Output} 钩子）
% \end{macroargs}
% 首先我们定义各个键以及样式键 \meta{style name}，当且仅当
% 名字非空。
%    \begin{macrocode}
\gdef\lst@InstallFamily#1#2#3#4#5{%
    \lst@Key{#2}\relax{\lst@UseFamily{#2}##1\relax\lst@MakeKeywords}%
    \lst@Key{more#2}\relax
        {\lst@UseFamily{#2}##1\relax\lst@MakeMoreKeywords}%
    \lst@Key{delete#2}\relax
        {\lst@UseFamily{#2}##1\relax\lst@DeleteKeywords}%
    \ifx\@empty#3\@empty\else
        \lst@Key{#3}{#4}{\lstKV@OptArg[\@ne]{##1}%
            {\@tempcnta\lst@classoffset \advance\@tempcnta####1\relax
             \@namedef{lst@#3\ifnum\@tempcnta=\@ne\else \the\@tempcnta
                             \fi}{####2}}}%
    \fi
    \expandafter\lst@InstallFamily@
        \csname\@lst @#2@data\expandafter\endcsname
        \csname\@lst @#5\endcsname {#1}{#2}{#3}}
%    \end{macrocode}
% 现在我们检查 \meta{working procedure} 是否为空。相应地我们在
% `数据'定义中使用 \texttt{工作过程} 或 \texttt{样式}。
% 工作过程在必要时就在这里定义。
%    \begin{macrocode}
\gdef\lst@InstallFamily@#1#2#3#4#5#6#7#8{%
    \gdef#1{{#3}{#4}{#5}#2#7}%
    \long\def\lst@temp##1{#6}%
    \ifx\lst@temp\@gobble
        \lst@AddTo#1{s#8}%
    \else
        \lst@AddTo#1{w#8}%
        \global\@namedef{lst@g#4@wp}##1{#6}%
    \fi}
%    \end{macrocode}
% 这里不再定义别的，其余全部按需完成。
% \end{macro}
%
% \begin{macro}{\lst@UseFamily}
% 我们查找可选的类号，提供这个成员，\ldots
%    \begin{macrocode}
\gdef\lst@UseFamily#1{%
    \def\lst@family{#1}%
    \@ifnextchar[\lst@UseFamily@{\lst@UseFamily@[\@ne]}}
\gdef\lst@UseFamily@[#1]{%
    \@tempcnta\lst@classoffset \advance\@tempcnta#1\relax
    \lst@ProvideFamily\lst@family
%    \end{macrocode}
% \ldots\ 并构造控制序列\ldots
%    \begin{macrocode}
    \lst@UseFamily@a
        {\lst@family\ifnum\@tempcnta=\@ne\else \the\@tempcnta \fi}}
\gdef\lst@UseFamily@a#1{%
    \expandafter\lst@UseFamily@b
       \csname\@lst @#1@list\expandafter\endcsname
       \csname\@lst @#1\expandafter\endcsname
       \csname\@lst @#1@also\expandafter\endcsname
       \csname\@lst @g#1\endcsname}
%    \end{macrocode}
% \ldots\ 它们是 |\lst@MakeKeywords| 和 |#6| 所需的。
%    \begin{macrocode}
\gdef\lst@UseFamily@b#1#2#3#4#5\relax#6{\lstKV@XOptArg[]{#5}#6#1#2#3#4}
%    \end{macrocode}
% \end{macro}
%
% \begin{macro}{\lst@ProvideFamily}
% 提供族 |#1| 的成员`|\the\@tempcnta|'。如果该成员
% 已存在，我们什么都不做。否则我们展开上面定义的数据宏。
% 注意如果计数器等于 1 我们就不使用它。因为有个缺陷报告
% \lsthelper{Kris~Luyten}{2002/08/03}{Undefined control sequence \lst@thestyle}
% 关键词族使用前缀 |lstfam| 而不是 |lst|。标记
% |\lstfam@#1|\oarg{number} 被全局定义，因为
% \lsthelper{Edsko~de~Vries}{2003/07/20}{仅在可选参数中
% 选择语言时关键词出错}报告了缺陷。
%    \begin{macrocode}
\gdef\lst@ProvideFamily#1{%
    \@ifundefined{lstfam@#1\ifnum\@tempcnta=\@ne\else\the\@tempcnta\fi}%
    {\global\@namedef{lstfam@#1\ifnum\@tempcnta=\@ne\else
                                        \the\@tempcnta\fi}{}%
     \expandafter\expandafter\expandafter\lst@ProvideFamily@
         \csname\@lst @#1@data\endcsname
         {\ifnum\@tempcnta=\@ne\else \the\@tempcnta \fi}}%
    {}}%
%    \end{macrocode}
% 现在我们有以下参数
% \begin{macroargs}
% \item \meta{prefix}
% \item \meta{name}
% \item \meta{style name}
% \item \meta{default style name}
% \item \alternative{l,o}（语言或其他键）
% \item \alternative{w,s}（工作过程或样式）
% \item \alternative{d,o}（\hookname{DetectKeywords} 或 \hookname{Output} 钩子）
% \item |\ifnum\@tempcnta=\@ne\else \the\@tempcnta \fi|
% \end{macroargs}
% 我们定义 |\lst@g|\meta{name}\meta{number}|@sty|，用来调用
% |\lst@g|\meta{name}|@wp|（以该编号为参数）或
% |\lst@|\meta{style name}\meta{number}，其中编号属于该控制
% 序列。
%    \begin{macrocode}
\gdef\lst@ProvideFamily@#1#2#3#4#5#6#7#8{%
    \expandafter\xdef\csname\@lst @g#2#8@sty\endcsname
    {\if #6w%
         \expandafter\noexpand\csname\@lst @g#2@wp\endcsname{#8}%
     \else
         \expandafter\noexpand\csname\@lst @#3#8\endcsname
     \fi}%
%    \end{macrocode}
% 我们确保样式宏存在。这在 \hookname{Init} 钩子中完成
% ——必要时赋以默认样式。
%    \begin{macrocode}
    \ifx\@empty#3\@empty\else
        \edef\lst@temp{\noexpand\lst@AddToHook{Init}{%
            \noexpand\lst@ProvideStyle\expandafter\noexpand
                \csname\@lst @#3#8\endcsname\noexpand#4}}%
        \lst@temp
    \fi
%    \end{macrocode}
% 我们调用一个子宏完成其余工作。它需要一些控制序列。
%    \begin{macrocode}
    \expandafter\lst@ProvideFamily@@
         \csname\@lst @#2#8@list\expandafter\endcsname
         \csname\@lst @#2#8\expandafter\endcsname
         \csname\@lst @#2#8@also\expandafter\endcsname
         \csname\@lst @g#2#8@list\expandafter\endcsname
         \csname\@lst @g#2#8\expandafter\endcsname
         \csname\@lst @g#2#8@sty\expandafter\endcsname
         {#1}#5#6#7}
%    \end{macrocode}
% 现在我们有了（除了 \meta{number} 可能总是缺失）
% \begin{macroargs}
% \item |\lst@|\meta{name}\meta{number}|@list|
% \item |\lst@|\meta{name}\meta{number}
% \item |\lst@|\meta{name}\meta{number}|@also|
% \item |\lst@g|\meta{name}\meta{number}|@list|
% \item |\lst@g|\meta{name}\meta{number}
% \item |\lst@g|\meta{name}\meta{number}|@sty|
% \item \meta{prefix}
% \item \alternative{l,o}（语言或其他键）
% \item \alternative{w,s}（工作过程或样式）
% \item \alternative{d,o}（\hookname{DetectKeywords} 或 \hookname{Output} 钩子）
% \end{macroargs}
% 注意 |#9| 和`|#10|'由 |\lst@InstallTest| 读取。我们（在 \hookname{SetLanguage} 处）
% 初始化所有需要的`变量'，并安装测试（该
% 定义事实上也是延迟的）。
%    \begin{macrocode}
\gdef\lst@ProvideFamily@@#1#2#3#4#5#6#7#8{%
    \gdef#1{#2#5}\global\let#2\@empty \global\let#3\@empty % init
    \gdef#4{#2#5}\global\let#5\@empty % init
    \if #8l\relax
        \lst@AddToHook{SetLanguage}{\def#1{#2#5}\let#2\@empty}%
    \fi
    \lst@InstallTest{#7}#1#2#4#5#6}
%    \end{macrocode}
% \end{macro}
%
% \begin{macro}{\lst@InstallKeywords}
% 现在我们利用上面可选参数的构造。因此我们
% 只在键的定义中插入 |[\@ne]| 作为 \meta{number}。
%    \begin{macrocode}
\gdef\lst@InstallKeywords#1#2#3#4#5{%
    \lst@Key{#2}\relax
        {\lst@UseFamily{#2}[\@ne]##1\relax\lst@MakeKeywords}%
    \lst@Key{more#2}\relax
        {\lst@UseFamily{#2}[\@ne]##1\relax\lst@MakeMoreKeywords}%
    \lst@Key{delete#2}\relax
        {\lst@UseFamily{#2}[\@ne]##1\relax\lst@DeleteKeywords}%
    \ifx\@empty#3\@empty\else
        \lst@Key{#3}{#4}{\@namedef{lst@#3}{##1}}%
    \fi
    \expandafter\lst@InstallFamily@
        \csname\@lst @#2@data\expandafter\endcsname
        \csname\@lst @#5\endcsname {#1}{#2}{#3}}
%    \end{macrocode}
% \end{macro}
%
% \begin{macro}{\lst@ProvideStyle}
% 如果样式宏 |#1| 未定义，它就变为等价于 |#2|。
%    \begin{macrocode}
\gdef\lst@ProvideStyle#1#2{%
    \ifx#1\@undefined \let#1#2%
    \else\ifx#1\relax \let#1#2\fi\fi}
%    \end{macrocode}
% \end{macro}
%
% 最后我们定义 |\lst@MakeKeywords|、\ldots、|\lst@DeleteKeywords|。
% 我们从两个辅助宏开始。
%
% \begin{macro}{\lst@BuildClassList}
% |#1| 之后跟一个以 |,\relax,| 结尾的逗号分隔关键词类列表，
% 例如 |keywords2,emph1,\relax,|。对这个列表中每个 \meta{item}，
% 我们\emph{追加}两个宏 |\lst@|\meta{item}|\lst@g|\meta{item}
% 到 |#1|。
%    \begin{macrocode}
\gdef\lst@BuildClassList#1#2,{%
    \ifx\relax#2\@empty\else
        \ifx\@empty#2\@empty\else
            \lst@lExtend#1{\csname\@lst @#2\expandafter\endcsname
                           \csname\@lst @g#2\endcsname}%
        \fi
        \expandafter\lst@BuildClassList\expandafter#1
    \fi}
%    \end{macrocode}
% \end{macro}
%
% \begin{macro}{\lst@DeleteClassesIn}
% 删除成对的记号，即传给子宏的参数 |#2#3|。
%    \begin{macrocode}
\gdef\lst@DeleteClassesIn#1#2{%
    \expandafter\lst@DCI@\expandafter#1#2\relax\relax}
\gdef\lst@DCI@#1#2#3{%
    \ifx#2\relax
        \expandafter\@gobbletwo
    \else
%    \end{macrocode}
% 如果尚未到达类列表末尾，我们就定义一个临时宏
% 来删除所有出现。
%    \begin{macrocode}
        \def\lst@temp##1#2#3##2{%
            \lst@lAddTo#1{##1}%
            \ifx ##2\relax\else
                \expandafter\lst@temp
            \fi ##2}%
        \let\@tempa#1\let#1\@empty
        \expandafter\lst@temp\@tempa#2#3\relax
    \fi
    \lst@DCI@#1}
%    \end{macrocode}
% \end{macro}
%
% \begin{macro}{\lst@MakeKeywords}
% 我们清空一些宏，并利用 |\lst@MakeMoreKeywords|。
% 注意这个定义和接下来的两个定义都有下列参数：
% \begin{macroargs}
% \item 类列表（在方括号中）
% \item 关键词列表
% \item |\lst@|\meta{name}|@list|
% \item |\lst@|\meta{name}
% \item |\lst@|\meta{name}|@also|
% \item |\lst@g|\meta{name}
% \end{macroargs}
%    \begin{macrocode}
\gdef\lst@MakeKeywords[#1]#2#3#4#5#6{%
    \def#3{#4#6}\let#4\@empty \let#5\@empty
    \lst@MakeMoreKeywords[#1]{#2}#3#4#5#6}
%    \end{macrocode}
% \end{macro}
%
% \begin{macro}{\lst@MakeMoreKeywords}
% 我们追加类和关键词。
%    \begin{macrocode}
\gdef\lst@MakeMoreKeywords[#1]#2#3#4#5#6{%
    \lst@BuildClassList#3#1,\relax,%
    \lst@DefOther\lst@temp{,#2}\lst@lExtend#4\lst@temp}
%    \end{macrocode}
% \end{macro}
%
% \begin{macro}{\lst@DeleteKeywords}
% 我们通过 |\lst@MakeKeywords| 转换关键词参数，并删除
% 类和关键词。
%    \begin{macrocode}
\gdef\lst@DeleteKeywords[#1]#2#3#4#5#6{%
    \lst@MakeKeywords[#1]{#2}\@tempa\@tempb#5#6%
    \lst@DeleteClassesIn#3\@tempa
    \lst@DeleteKeysIn#4\@tempb}
%    \end{macrocode}
% \end{macro}
%
%
% \subsection{主要的族与类}
%
%
% \paragraph{关键词}
%
% \begin{lstkey}{keywords}
% 定义关键词族变得非常、非常容易。
%    \begin{macrocode}
\lst@InstallFamily k{keywords}{keywordstyle}\bfseries{keywordstyle}{}ld
%    \end{macrocode}
% 下面的宏设置一个 keywordstyle，它\ldots
%    \begin{macrocode}
\gdef\lst@DefKeywordstyle#1#2\@nil@{%
   \@namedef{lst@keywordstyle\ifnum\@tempcnta=\@ne\else\the\@tempcnta
                             \fi}{#1#2}}%
%    \end{macrocode}
% \ldots\ 在这里拼起来。如果我们在类号之后检测到星号，
% 我们就插入让关键词变成大写的代码。
%    \begin{macrocode}
\lst@Key{keywordstyle}{\bfseries}{\lstKV@OptArg[\@ne]{#1}%
  {\@tempcnta\lst@classoffset \advance\@tempcnta##1\relax
   \@ifstar{\lst@DefKeywordstyle{\uppercase\expandafter{%
                                 \expandafter\lst@token
                                 \expandafter{\the\lst@token}}}}%
           {\lst@DefKeywordstyle{}}##2\@nil@}}
%    \end{macrocode}
% \end{lstkey}
%
% \begin{lstkey}{ndkeywords}
% 二阶关键词使用与 |\lst@InstallKeywords| 相同的技巧。
%    \begin{macrocode}
\lst@Key{ndkeywords}\relax
    {\lst@UseFamily{keywords}[\tw@]#1\relax\lst@MakeKeywords}%
\lst@Key{morendkeywords}\relax
    {\lst@UseFamily{keywords}[\tw@]#1\relax\lst@MakeMoreKeywords}%
\lst@Key{deletendkeywords}\relax
    {\lst@UseFamily{keywords}[\tw@]#1\relax\lst@DeleteKeywords}%
\lst@Key{ndkeywordstyle}\relax{\@namedef{lst@keywordstyle2}{#1}}%
%    \end{macrocode}
% \lsthelper{Dr.~Peter~Leibner}{1999/11/05}{未定义的 \lst@UseKeywords、
% 非法参数号（##1）}报告了两个缺陷：|\lst@UseKeywords| 和
% |##1| 变成了 |\lst@UseFamily| 和 |#1|。
% \end{lstkey}
%
% \begin{lstkey}{keywordsprefix}
% 是实验性实现的。唯一的前缀通过把 |\lst@prefixkeyword| 置空
% 来表明它的存在。我们可以在 \keyname{Output} 钩子中
% 捕获这一信息。
%    \begin{macrocode}
\lst@Key{keywordsprefix}\relax{\lst@DefActive\lst@keywordsprefix{#1}}
\global\let\lst@keywordsprefix\@empty
\lst@AddToHook{SelectCharTable}
    {\ifx\lst@keywordsprefix\@empty\else
         \expandafter\lst@CArg\lst@keywordsprefix\relax
             \lst@CDef{}%
                      {\lst@ifletter\else
                           \global\let\lst@prefixkeyword\@empty
                       \fi}%
                      {}%
     \fi}
\lst@AddToHook{Init}{\global\let\lst@prefixkeyword\relax}
\lst@AddToHook{Output}
    {\ifx\lst@prefixkeyword\@empty
         \let\lst@thestyle\lst@gkeywords@sty
         \global\let\lst@prefixkeyword\relax
     \fi}%
%    \end{macrocode}
% \end{lstkey}
%
% \begin{lstkey}{otherkeywords}
% 感谢 \lsthelper{Bradford~Chamberlain}{2001/07/07}{otherkeywords={@,@^}
% 不起作用}，我们现在沿`其他关键词'列表迭代，并逐个
% 变成活动的——而不是把整个参数变成活动的。我们把活动的
% 记号序列追加到 |\lst@otherkeywords|，以定义每个`其他'关键词。
%    \begin{macrocode}
\lst@Key{otherkeywords}{}{%
    \let\lst@otherkeywords\@empty
    \lst@for{#1}\do{%
      \lst@MakeActive{##1}%
      \lst@lExtend\lst@otherkeywords{%
          \expandafter\lst@CArg\lst@temp\relax\lst@CDef
              {}\lst@PrintOtherKeyword\@empty}}}
\lst@AddToHook{SelectCharTable}{\lst@otherkeywords}
%    \end{macrocode}
% 在 \lsthelper{Peter~Bartke}{2001/11/06}{未定义的控制序列
% \lst@PrintOtherkeyword}.
% \end{lstkey}
%
% \begin{macro}{\lst@PrintOtherKeyword}
% 打印前面的字符、准备输出并以关键词样式排印参数。
% \lsthelper{James~Willans}{2004/07/23}{问题：otherkeywords}
% 报告了在 |\begingroup| 和 |\endgroup| 内调用输出例程时的问题。
% 现在该定义已重构。
%    \begin{macrocode}
\gdef\lst@PrintOtherKeyword#1\@empty{%
    \lst@XPrintToken
    \begingroup
      \lst@modetrue \lsthk@TextStyle
      \let\lst@ProcessDigit\lst@ProcessLetter
      \let\lst@ProcessOther\lst@ProcessLetter
      \lst@lettertrue
      #1%
	  \lst@SaveToken
    \endgroup
	\lst@RestoreToken
	\global\let\lst@savedcurrstyle\lst@currstyle
	\let\lst@currstyle\lst@gkeywords@sty
    \lst@Output
	\let\lst@currstyle\lst@savedcurrstyle}
%    \end{macrocode}
% \begin{TODO}
% \hookname{TextStyle} 钩子的哪一部分是必需的？在重构之后
% 它还是必需的吗？需要移到别处吗？
% \end{TODO}
% \end{macro}
%
%    \begin{macrocode}
\lst@EndAspect
%</misc>
%    \end{macrocode}
% \end{aspect}
%
%
% \paragraph{强调族}
%
% \begin{aspect}{emph}
% 这里只是一次宏调用。
%    \begin{macrocode}
%<*misc>
\lst@BeginAspect[keywords]{emph}
\lst@InstallFamily e{emph}{emphstyle}{}{emphstyle}{}od
\lst@EndAspect
%</misc>
%    \end{macrocode}
% \end{aspect}
%
%
% \paragraph{\TeX\ 控制序列}
%
% \begin{aspect}{tex}
% 这里我们检查最后处理过的`其他'记号。
%    \begin{macrocode}
%<*misc>
\lst@BeginAspect[keywords]{tex}
%    \end{macrocode}
%    \begin{macrocode}
\lst@InstallFamily {cs}{texcs}{texcsstyle}\relax{keywordstyle}
    {\ifx\lst@lastother\lstum@backslash
         \expandafter\let\expandafter\lst@thestyle
                         \csname lst@texcsstyle#1\endcsname
     \fi}
    ld
%    \end{macrocode}
% 样式键检查可选星号（它必须在
% 可选类参数之前）。
%    \begin{macrocode}
\lst@Key{texcsstyle}\relax
  {\@ifstar{\lst@true\lst@DefTexcsstyle}%
           {\lst@false\lst@DefTexcsstyle}#1\@nil@}
\gdef\lst@DefTexcsstyle#1\@nil@{%
    \let\lst@iftexcsincludebs\lst@if
    \lstKV@OptArg[\@ne]{#1}%
    {\@tempcnta\lst@classoffset \advance\@tempcnta##1\relax
     \@namedef{lst@texcsstyle\ifnum\@tempcnta=\@ne\else
                                   \the\@tempcnta \fi}{##2}}}%
\global\let\lst@iftexcsincludebs\iffalse
%    \end{macrocode}
% 为了让反斜杠属于控制序列，它与后面的记号合并。
% 这个选项由 \lsthelper{Morten~H\o gholm}提出
% {2004/07/16}{defining new (colored) texcs}.
% \lsthelper{Christian~Schneider}{2006/09/08}{-}指出，当标识符前面有`其他'
% 字符时原始实现是坏的。为修正这一点（以及其他缺陷），
% 我们先输出当前记号中的任何内容，然后再合并。
%    \begin{macrocode}
\let\lst@iftexcsincludebs\iffalse
\lst@AddToHook{SelectCharTable}
{\lst@iftexcsincludebs \ifx\@empty\lst@texcs\else
     \lst@DefSaveDef{`\\}\lsts@texcsbs
      {\lst@ifletter
           \lst@Output
       \else
           \lst@OutputOther
       \fi
       \lst@Merge\lsts@texcsbs}%
 \fi \fi}
%    \end{macrocode}
%    \begin{macrocode}
\lst@EndAspect
%</misc>
%    \end{macrocode}
% \end{aspect}
%
%
% \paragraph{编译器指令}
%
% \begin{aspect}{directives}
% \begin{lstkey}{directives}
% 首先是一些通常的内容。
%    \begin{macrocode}
%<*misc>
\lst@BeginAspect[keywords]{directives}
%    \end{macrocode}
% |\lst@directives| 的初始化是在
% \lsthelper{Kris~Luyten}{2002/07/30}{加载 C 之后
% |\lst@directives| 未定义导致未定义的控制序列 |\lst@thestyle|}缺陷报告之后添加的。
%    \begin{macrocode}
\lst@NewMode\lst@CDmode
\lst@AddToHook{EOL}{\ifnum\lst@mode=\lst@CDmode \lst@LeaveMode \fi}
\lst@InstallKeywords{d}{directives}{directivestyle}\relax{keywordstyle}
    {\ifnum\lst@mode=\lst@CDmode
         \let\lst@thestyle\lst@directivestyle
     \fi}
    ld
\global\let\lst@directives\@empty % init
%    \end{macrocode}
% 现在我们为指令定义一个新的定界符：我们只在第一列
% 进入`指令模式'。
%    \begin{macrocode}
\lst@AddTo\lst@delimtypes{,directive}
\gdef\lst@Delim@directive#1\@empty#2#3#4{%
    \lst@CArg #1\relax\lst@DefDelimB
        {\lst@CalcColumn}%
        {}%
        {\ifnum\@tempcnta=\z@
             \def\lst@bnext{#2\lst@CDmode{#4\lst@Lmodetrue}%
			                \let\lst@currstyle\lst@directivestyle}%
		 \fi
		 \@gobblethree}%
        #2\lst@CDmode{#4\lst@Lmodetrue}}
%    \end{macrocode}
% 我们引入一个新的字符串类型（感谢 \lsthelper{R.~Isernhagen}
% {1999/11/12}{float 在 #include <float> 中不是关键词}），它\ldots
%    \begin{macrocode}
\lst@AddTo\lst@stringtypes{,directive}
\gdef\lst@StringDM@directive#1#2#3\@empty{%
    \lst@CArg #2\relax\lst@CDef
        {}%
%    \end{macrocode}
% \ldots\ 只在 |\lst@CDmode| 中活动：
%    \begin{macrocode}
        {\let\lst@bnext\lst@CArgEmpty
         \ifnum\lst@mode=\lst@CDmode
             \def\lst@bnext{\lst@BeginString{#1}}%
         \fi
         \lst@bnext}%
        \@empty
    \lst@CArg #3\relax\lst@CDef
        {}%
        {\let\lst@enext\lst@CArgEmpty
         \ifnum #1=\lst@mode
             \let\lst@bnext\lst@EndString
         \fi
         \lst@bnext}%
        \@empty}
%    \end{macrocode}
% \end{lstkey}
%
%    \begin{macrocode}
\lst@EndAspect
%</misc>
%    \end{macrocode}
% \end{aspect}
%
%
% \subsection{关键词注释}
%
% \begin{aspect}{keywordcomments}
% 包含两种注释类型，可能被拆分为这个和 |dkcs|。
%    \begin{macrocode}
%<*misc>
\lst@BeginAspect[keywords,comments]{keywordcomments}
%    \end{macrocode}
%
% \begin{macro}{\lst@BeginKC}
% \begin{macro}{\lst@BeginKCS}
% 开始关键词注释很容易，但是：(1) 各子宏被调用
% 在两个分组层级之外，且\ldots
%    \begin{macrocode}
\lst@NewMode\lst@KCmode \lst@NewMode\lst@KCSmode
\gdef\lst@BeginKC{\aftergroup\aftergroup\aftergroup\lst@BeginKC@}%
\gdef\lst@BeginKC@{%
    \lst@ResetToken
    \lst@BeginComment\lst@KCmode{{\lst@commentstyle}\lst@modetrue}%
                     \@empty}%
\gdef\lst@BeginKCS{\aftergroup\aftergroup\aftergroup\lst@BeginKCS@}%
\gdef\lst@BeginKCS@{%
    \lst@ResetToken
    \lst@BeginComment\lst@KCSmode{{\lst@commentstyle}\lst@modetrue}%
                     \@empty}%
%    \end{macrocode}
% (2) 我们必须确保注释在打印注释定界符
% 之后开始，因为该定界符可能是关键词。我们把 |\lst@BeginKC|[|S|] 赋给
% |\lst@KCpost|，它在 \hookname{PostOutput} 中被执行并重置。
%    \begin{macrocode}
\lst@AddToHook{PostOutput}{\lst@KCpost \global\let\lst@KCpost\@empty}
\global\let\lst@KCpost\@empty % init
%    \end{macrocode}
% \end{macro}
% \end{macro}
%
% \begin{macro}{\lst@EndKC}
% 在打印（临时保存的）注释定界符之前
% 离开注释模式。
%    \begin{macrocode}
\gdef\lst@EndKC{\lst@SaveToken \lst@LeaveMode \lst@RestoreToken
    \let\lst@thestyle\lst@identifierstyle \lsthk@Output}
%    \end{macrocode}
% \end{macro}
%
% \begin{lstkey}{keywordcomment}
% 这里定界符必须相同，因此我们使用 |\lst@KCmatch|。注意传给
% |\lst@InstallKeywords| 的最后一个参数 |o|：工作测试被安装在
% \hookname{Output} 钩子中，而不是 \hookname{DetectKeywords} 中。
% 否则我们无法检测结束定界符，因为关键词检测当且
% 仅当允许模式改动时才进行。
%    \begin{macrocode}
\lst@InstallKeywords{kc}{keywordcomment}{}\relax{}
    {\ifnum\lst@mode=\lst@KCmode
         \edef\lst@temp{\the\lst@token}%
         \ifx\lst@temp\lst@KCmatch
             \lst@EndKC
         \fi
     \else
         \lst@ifmode\else
             \xdef\lst@KCmatch{\the\lst@token}%
             \global\let\lst@KCpost\lst@BeginKC
         \fi
     \fi}
    lo
%    \end{macrocode}
% \end{lstkey}
%
% \begin{lstkey}{keywordcommentsemicolon}
% 这个键只是保存关键词。在 \lsthelper
% {Norbert~Eisinger}{2002/11/26}{keywordcommentsemicolon 在
% 语言改变之后仍活动}报告缺陷之后，\hookname{SetLanguage} 中的初始化
% 被添加。
%    \begin{macrocode}
\lst@Key{keywordcommentsemicolon}{}{\lstKV@ThreeArg{#1}%
    {\def\lst@KCAkeywordsB{##1}%
     \def\lst@KCAkeywordsE{##2}%
     \def\lst@KCBkeywordsB{##3}%
     \def\lst@KCkeywords{##1##2##3}}}
\lst@AddToHook{SetLanguage}{%
    \let\lst@KCAkeywordsB\@empty \let\lst@KCAkeywordsE\@empty
    \let\lst@KCBkeywordsB\@empty \let\lst@KCkeywords\@empty}
%    \end{macrocode}
% 如果定义了这种关键词注释类型，我们就定义一个合适的
% 分号。合适意味着：若关键词注释模式处于活动状态，我们就离开它。
% \lsthelper{Oldrich~Jedlicka}{2001/12/12}{keywordcomment(semicolon) fails}
% 报告了缺陷并提供了修正，即那两个 |\@empty|。
%    \begin{macrocode}
\lst@AddToHook{SelectCharTable}
    {\ifx\lst@KCkeywords\@empty\else
        \lst@DefSaveDef{`\;}\lsts@EKC
            {\lst@XPrintToken
             \ifnum\lst@mode=\lst@KCmode \lst@EndComment\@empty \else
             \ifnum\lst@mode=\lst@KCSmode \lst@EndComment\@empty
             \fi \fi
             \lsts@EKC}%
     \fi}
%    \end{macrocode}
% `工作标识符'宏进入或离开注释模式。
%    \begin{macrocode}
\gdef\lst@KCAWorkB{%
    \lst@ifmode\else \global\let\lst@KCpost\lst@BeginKC \fi}
\gdef\lst@KCBWorkB{%
    \lst@ifmode\else \global\let\lst@KCpost\lst@BeginKCS \fi}
\gdef\lst@KCAWorkE{\ifnum\lst@mode=\lst@KCmode \lst@EndKC \fi}
%    \end{macrocode}
% 现在我们安装测试并初始化给定的宏。
%    \begin{macrocode}
\lst@ProvideFamily@@
    \lst@KCAkeywordsB@list\lst@KCAkeywordsB \lst@KC@also
    \lst@gKCAkeywordsB@list\lst@gKCAkeywordsB \lst@KCAWorkB
    {kcb}owo % prefix, other key, working procedure, Output hook
\lst@ProvideFamily@@
    \lst@KCAkeywordsE@list\lst@KCAkeywordsE \lst@KC@also
    \lst@gKCAkeywordsE@list\lst@gKCAkeywordsE \lst@KCAWorkE
    {kce}owo
\lst@ProvideFamily@@
    \lst@KCBkeywordsB@list\lst@KCBkeywordsB \lst@KC@also
    \lst@gKCBkeywordsB@list\lst@gKCBkeywordsB \lst@KCBWorkB
    {kcs}owo
%    \end{macrocode}
% \end{lstkey}
%
%    \begin{macrocode}
\lst@EndAspect
%</misc>
%    \end{macrocode}
% \end{aspect}
%
%
% \subsection{标识符的导出}
%
% \begin{aspect}{index}
% \begin{macro}{\lstindexmacro}
% 又一个`关键词'类。
%    \begin{macrocode}
%<*misc>
\lst@BeginAspect[keywords]{index}
\lst@InstallFamily w{index}{indexstyle}\lstindexmacro{indexstyle}
    {\csname\@lst @indexstyle#1\expandafter\endcsname
         \expandafter{\the\lst@token}}
    od
\lst@UserCommand\lstindexmacro#1{\index{{\ttfamily#1}\encapchar hdpindex{main}}}
\lst@EndAspect
%</misc>
%    \end{macrocode}
% \end{macro}
% \end{aspect}
%
% \begin{aspect}{procnames}
% \begin{lstkey}{procnamestyle}
% \begin{lstkey}{procnamekeys}
% \begin{lstkey}{indexprocnames}
% 这里的`想法'是使用全局的 |\lst@ifprocname|，它表示前面
% 有`过程关键词'。其余都是已知的东西。
%    \begin{macrocode}
%<*misc>
\lst@BeginAspect[keywords]{procnames}
\gdef\lst@procnametrue{\global\let\lst@ifprocname\iftrue}
\gdef\lst@procnamefalse{\global\let\lst@ifprocname\iffalse}
\lst@AddToHook{Init}{\lst@procnamefalse}
\lst@AddToHook{DetectKeywords}
    {\lst@ifprocname
         \let\lst@thestyle\lst@procnamestyle
         \lst@ifindexproc \csname\@lst @gindex@sty\endcsname \fi
         \lst@procnamefalse
     \fi}
%    \end{macrocode}
% 这两个就是实现的键：
%    \begin{macrocode}
\lst@Key{procnamestyle}{}{\def\lst@procnamestyle{#1}}
\lst@Key{indexprocnames}{false}[t]{\lstKV@SetIf{#1}\lst@ifindexproc}
\lst@AddToHook{Init}{\lst@ifindexproc \lst@indexproc \fi}
\gdef\lst@indexproc{%
    \@ifundefined{lst@indexstyle1}%
        {\@namedef{lst@indexstyle1}##1{}}%
        {}}
%    \end{macrocode}
% 上面 |\lst@indexstyle| 的默认定义被移到钩子之外，
% 这是在 \lsthelper{Ulrich~G.~Wortmann}{2002/01/22}{procnames 不工作}报告缺陷之后。
%    \begin{macrocode}
\lst@InstallKeywords w{procnamekeys}{}\relax{}
    {\global\let\lst@PNpost\lst@procnametrue}
    od
\lst@AddToHook{PostOutput}{\lst@PNpost\global\let\lst@PNpost\@empty}
\global\let\lst@PNpost\@empty % init
\lst@EndAspect
%</misc>
%    \end{macrocode}
% \end{lstkey}
% \end{lstkey}
% \end{lstkey}
% \end{aspect}
%
%
% \section{更多 aspect 和键}
%
% \begin{lstkey}{basicstyle}
% \begin{lstkey}{inputencoding}
% 我认为没有更好的地方来定义这些键了。
%    \begin{macrocode}
%<*kernel>
\lst@Key{basicstyle}\relax{\def\lst@basicstyle{#1}}
\lst@Key{inputencoding}\relax{\def\lst@inputenc{#1}}
\lst@AddToHook{Init}
    {\lst@basicstyle
     \ifx\lst@inputenc\@empty\else
         \@ifundefined{inputencoding}{}%
            {\inputencoding\lst@inputenc}%
     \fi}
\lst@AddToHookExe{EmptyStyle}
    {\let\lst@basicstyle\@empty
     \let\lst@inputenc\@empty}
\lst@Key{multicols}{}{\@tempcnta=0#1\relax\def\lst@multicols{#1}}
%</kernel>
%    \end{macrocode}
% Michael Niedermair 请求一个像 \keyname{inputencoding} 这样的键。
% \end{lstkey}
% \end{lstkey}
%
%
% \subsection{样式与语言}
%
% \begin{aspect}{style}
% 我们从样式定义和选择开始。
%    \begin{macrocode}
%<*misc>
\lst@BeginAspect{style}
%    \end{macrocode}
%
% \begin{macro}{\lststylefiles}
% 该宏被定义当且仅当它尚未定义。
%    \begin{macrocode}
\@ifundefined{lststylefiles}
    {\lst@UserCommand\lststylefiles{lststy0.sty}}{}
%    \end{macrocode}
% \end{macro}
%
% \begin{macro}{\lstdefinestyle}
% \begin{macro}{\lst@definestyle}
% \begin{macro}{\lst@DefStyle}
% 用 |\lst@DefStyle| 定义，而后者通过
% |\lst@DefDriver|.
%    \begin{macrocode}
\lst@UserCommand\lstdefinestyle{\lst@DefStyle\iftrue}
\lst@UserCommand\lst@definestyle{\lst@DefStyle\iffalse}
\gdef\lst@DefStyle{\lst@DefDriver{style}{sty}\lstset}
%    \end{macrocode}
% `空'样式调用初始的空钩子 \hookname{EmptyStyle}。
%    \begin{macrocode}
\global\@namedef{lststy@$}{\lsthk@EmptyStyle}
\lst@AddToHook{EmptyStyle}{}% init
%    \end{macrocode}
% \end{macro}
% \end{macro}
% \end{macro}
%
% \begin{lstkey}{style}
% 是 |\lst@LAS| 的一个应用。我们只需指定钩子，并把空
% 参数作为`前'和`后'代码。
%    \begin{macrocode}
\lst@Key{style}\relax{%
    \lst@LAS{style}{sty}{[]{#1}}\lst@NoAlias\lststylefiles
        \lsthk@SetStyle
        {}}
%    \end{macrocode}
%    \begin{macrocode}
\lst@AddToHook{SetStyle}{}% init
%    \end{macrocode}
% \end{lstkey}
%
%    \begin{macrocode}
\lst@EndAspect
%</misc>
%    \end{macrocode}
% \end{aspect}
%
% \begin{aspect}{language}
% 现在我们处理用于定义和选择编程语言的命令，
% 特别是别名。
%    \begin{macrocode}
%<*misc>
\lst@BeginAspect{language}
%    \end{macrocode}
%
% \begin{macro}{\lstlanguagefiles}
% 该宏被定义当且仅当它尚未定义。
%    \begin{macrocode}
\@ifundefined{lstdriverfiles}
    {\lst@UserCommand\lstlanguagefiles{lstlang0.sty}}{}
%    \end{macrocode}
% \end{macro}
%
% \begin{macro}{\lstdefinelanguage}
% \begin{macro}{\lst@definelanguage}
% \begin{macro}{\lst@DefLang}
% 用 |\lst@DefLang| 定义，而后者通过
% |\lst@DefDriver|.
%    \begin{macrocode}
\lst@UserCommand\lstdefinelanguage{\lst@DefLang\iftrue}
\lst@UserCommand\lst@definelanguage{\lst@DefLang\iffalse}
\gdef\lst@DefLang{\lst@DefDriver{language}{lang}\lstset}
%    \end{macrocode}
% 现在我们可以提供`空'语言。
%    \begin{macrocode}
\lstdefinelanguage{}{}
%    \end{macrocode}
% \end{macro}
% \end{macro}
% \end{macro}
%
% \begin{lstkey}{language}
% \begin{lstkey}{alsolanguage}
% 主要是 |\lst@LAS| 的一个应用。
%    \begin{macrocode}
\lst@Key{language}\relax{\lstKV@OptArg[]{#1}%
    {\lst@LAS{language}{lang}{[##1]{##2}}\lst@FindAlias\lstlanguagefiles
         \lsthk@SetLanguage
         {\lst@FindAlias[##1]{##2}%
          \let\lst@language\lst@malias
          \let\lst@dialect\lst@oalias}}}
%    \end{macrocode}
% 同上，我们只是不执行 |\lsthk@SetLanguage|。
%    \begin{macrocode}
\lst@Key{alsolanguage}\relax{\lstKV@OptArg[]{#1}%
    {\lst@LAS{language}{lang}{[##1]{##2}}\lst@FindAlias\lstlanguagefiles
         {}%
         {\lst@FindAlias[##1]{##2}%
          \let\lst@language\lst@malias
          \let\lst@dialect\lst@oalias}}}
%    \end{macrocode}
%    \begin{macrocode}
\lst@AddToHook{SetLanguage}{}% init
%    \end{macrocode}
% \end{lstkey}
% \end{lstkey}
%
% \begin{macro}{\lstalias}
% 现在我们关注别名和默认方言。
% |\lsta@|\meta{language}|$|\meta{dialect} 和 |\lsta@|\meta{language} 分别包含
% 某个特定方言和整个语言的别名。
% 我们将使用 |$| 字符分隔语言名与其方言。
% 感谢 \lsthelper{Walter~E.~Brown}{2004/02/25}{\lstalias
% （+ \lstdefinelanguage）失败}报告了 |\lstalias@| 先前定义中
% 参数定界符`['的问题。
%    \begin{macrocode}
\lst@UserCommand\lstalias{\@ifnextchar[\lstalias@\lstalias@@}
\gdef\lstalias@[#1]#2{\lstalias@b #2$#1}
\gdef\lstalias@b#1[#2]#3{\lst@NormedNameDef{lsta@#1}{#3$#2}}
\gdef\lstalias@@#1#2{\lst@NormedNameDef{lsta@#1}{#2}}
%    \end{macrocode}
% \end{macro}
%
% \begin{lstkey}{defaultdialect}
% 我们只保存方言。
%    \begin{macrocode}
\lst@Key{defaultdialect}\relax
    {\lstKV@OptArg[]{#1}{\lst@NormedNameDef{lstdd@##2}{##1}}}
%    \end{macrocode}
% \end{lstkey}
%
% \begin{macro}{\lst@FindAlias}
% 现在我们必须找一种语言。先测试整个语言的别名，
% 然后在必要时设置默认方言。
%    \begin{macrocode}
\gdef\lst@FindAlias[#1]#2{%
    \lst@NormedDef\lst@oalias{#1}%
    \lst@NormedDef\lst@malias{#2}%
    \@ifundefined{lsta@\lst@malias}{}%
        {\edef\lst@malias{\csname\@lst a@\lst@malias\endcsname}}%
%    \end{macrocode}
%    \begin{macrocode}
    \ifx\@empty\lst@oalias \@ifundefined{lstdd@\lst@malias}{}%
        {\edef\lst@oalias{\csname\@lst dd@\lst@malias\endcsname}}%
    \fi
%    \end{macrocode}
% 现在我们可以处理单个方言的别名了。
%    \begin{macrocode}
    \edef\lst@temp{\lst@malias $\lst@oalias}%
    \@ifundefined{lsta@\lst@temp}{}%
        {\edef\lst@temp{\csname\@lst a@\lst@temp\endcsname}}%
%    \end{macrocode}
% 最后我们再次设置默认方言——用于方言别名的情况。
%    \begin{macrocode}
    \expandafter\lst@FindAlias@\lst@temp $}
\gdef\lst@FindAlias@#1$#2${%
    \def\lst@malias{#1}\def\lst@oalias{#2}%
    \ifx\@empty\lst@oalias \@ifundefined{lstdd@\lst@malias}{}%
        {\edef\lst@oalias{\csname\@lst dd@\lst@malias\endcsname}}%
    \fi}
%    \end{macrocode}
% \end{macro}
%
% \begin{macro}{\lst@RequireLanguages}
% 这个定义将等价于 |\lstloadlanguages|。我们请求
% 给定的语言列表，并加载额外需要的方面。
%    \begin{macrocode}
\gdef\lst@RequireLanguages#1{%
    \lst@Require{language}{lang}{#1}\lst@FindAlias\lstlanguagefiles
    \ifx\lst@loadaspects\@empty\else
        \lst@RequireAspects\lst@loadaspects
    \fi}
%    \end{macrocode}
% \end{macro}
%
% \begin{macro}{\lstloadlanguages}
% 与 |\lst@RequireLanguages| 相同。
%    \begin{macrocode}
\global\let\lstloadlanguages\lst@RequireLanguages
%    \end{macrocode}
% \end{macro}
%
%    \begin{macrocode}
\lst@EndAspect
%</misc>
%    \end{macrocode}
% \end{aspect}
%
%
% \subsection{格式定义*}
%
% \begin{aspect}{formats}
%    \begin{macrocode}
%<*misc>
\lst@BeginAspect{formats}
%    \end{macrocode}
%
% \begin{macro}{\lstformatfiles}
% 该宏被定义当且仅当它尚未定义。
%    \begin{macrocode}
\@ifundefined{lstformatfiles}
    {\lst@UserCommand\lstformatfiles{lstfmt0.sty}}{}
%    \end{macrocode}
% \end{macro}
%
% \begin{macro}{\lstdefineformat}
% \begin{macro}{\lst@defineformat}
% \begin{macro}{\lst@DefFormat}
% 用 |\lst@DefFormat| 定义，而后者通过
% |\lst@DefDriver|.
%    \begin{macrocode}
\lst@UserCommand\lstdefineformat{\lst@DefFormat\iftrue}
\lst@UserCommand\lst@defineformat{\lst@DefFormat\iffalse}
\gdef\lst@DefFormat{\lst@DefDriver{format}{fmt}\lst@UseFormat}
%    \end{macrocode}
% 我们提供`空'格式。
%    \begin{macrocode}
\lstdefineformat{}{}
%    \end{macrocode}
% \end{macro}
% \end{macro}
% \end{macro}
%
% \begin{lstkey}{format}
% 是 |\lst@LAS| 的一个应用。我们只需把钩子指定为`前'，
% 把空参数指定为`后'代码。
%    \begin{macrocode}
\lst@Key{format}\relax{%
    \lst@LAS{format}{fmt}{[]{#1}}\lst@NoAlias\lstformatfiles
        \lsthk@SetFormat
        {}}
%    \end{macrocode}
%    \begin{macrocode}
\lst@AddToHook{SetFormat}{\let\lst@fmtformat\@empty}% init
%    \end{macrocode}
% \end{lstkey}
%
%
% \paragraph{辅助宏}
% 我们的目标是定义未知的 |\lst@UseFormat|。这个定义
% 将解析用户提供的格式。我们从一些通用宏开始。
%
% \begin{macro}{\lst@fmtSplit}
% 把宏 |#1| 的内容在 |#2| 处分割为前面的字符
% |\lst@fmta| 和后面的字符 |\lst@fmtb|。|\lst@if| 为假当且
% 仅当 |#1| 不含 |#2|。
%    \begin{macrocode}
\gdef\lst@fmtSplit#1#2{%
    \def\lst@temp##1#2##2\relax##3{%
        \ifnum##3=\z@
            \ifx\@empty##2\@empty
                \lst@false
                \let\lst@fmta#1%
                \let\lst@fmtb\@empty
            \else
                \expandafter\lst@temp#1\relax\@ne
            \fi
        \else
            \def\lst@fmta{##1}\def\lst@fmtb{##2}%
        \fi}%
    \lst@true
    \expandafter\lst@temp#1#2\relax\z@}
%    \end{macrocode}
% \end{macro}
%
% \begin{macro}{\lst@IfNextCharWhitespace}
% 用 |\lst@IfSubstring| 定义。
%    \begin{macrocode}
\gdef\lst@IfNextCharWhitespace#1#2#3{%
    \lst@IfSubstring#3\lst@whitespaces{#1}{#2}#3}
%    \end{macrocode}
% 下面是所有空白字符。
%    \begin{macrocode}
\begingroup
\catcode`\^^I=12\catcode`\^^J=12\catcode`\^^M=12\catcode`\^^L=12\relax%
\lst@DefActive\lst@whitespaces{\ ^^I^^J^^M}% add ^^L
\global\let\lst@whitespaces\lst@whitespaces%
\endgroup
%    \end{macrocode}
% \end{macro}
%
% \begin{macro}{\lst@fmtIfIdentifier}
% 测试 |#1| 的第一个字符
%    \begin{macrocode}
\gdef\lst@fmtIfIdentifier#1{%
    \ifx\relax#1\@empty
        \expandafter\@secondoftwo
    \else
        \expandafter\lst@fmtIfIdentifier@\expandafter#1%
    \fi}
%    \end{macrocode}
% 针对`字母' |_|、|@|、|A|,\ldots,|Z| 和 |a|,\ldots,|z|。
%    \begin{macrocode}
\gdef\lst@fmtIfIdentifier@#1#2\relax{%
    \let\lst@next\@secondoftwo
    \ifnum`#1=`_\else
    \ifnum`#1<64\else
    \ifnum`#1<91\let\lst@next\@firstoftwo\else
    \ifnum`#1<97\else
    \ifnum`#1<123\let\lst@next\@firstoftwo\else
    \fi \fi \fi \fi \fi
    \lst@next}
%    \end{macrocode}
% \end{macro}
%
% \begin{macro}{\lst@fmtIfNextCharIn}
% 是可选 \meta{exceptional characters} 所需的。
% 实现很容易——参见第~\ref{iSubstringTests} 节。
%    \begin{macrocode}
\gdef\lst@fmtIfNextCharIn#1{%
    \ifx\@empty#1\@empty \expandafter\@secondoftwo \else
                         \def\lst@next{\lst@fmtIfNextCharIn@{#1}}%
                         \expandafter\lst@next\fi}
\gdef\lst@fmtIfNextCharIn@#1#2#3#4{%
    \def\lst@temp##1#4##2##3\relax{%
        \ifx \@empty##2\expandafter\@secondoftwo
                 \else \expandafter\@firstoftwo \fi}%
    \lst@temp#1#4\@empty\relax{#2}{#3}#4}
%    \end{macrocode}
% \end{macro}
%
% \begin{macro}{\lst@fmtCDef}
% 我们需要 |\lst@CDef| 和 |\lst@CDefX| 的派生版本：我们必须测试
% 下一个字符是否属于例外字符序列 |#5|。
% 这些测试在这里插入。
%    \begin{macrocode}
\gdef\lst@fmtCDef#1{\lst@fmtCDef@#1}
\gdef\lst@fmtCDef@#1#2#3#4#5#6#7{%
    \lst@CDefIt#1{#2}{#3}%
               {\lst@fmtIfNextCharIn{#5}{#4#2#3}{#6#4#2#3#7}}%
               #4%
               {}{}{}}
%    \end{macrocode}
% \end{macro}
%
% \begin{macro}{\lst@fmtCDefX}
% 相同，但`丢弃输入'。
%    \begin{macrocode}
\gdef\lst@fmtCDefX#1{\lst@fmtCDefX@#1}
\gdef\lst@fmtCDefX@#1#2#3#4#5#6#7{%
    \let#4#1%
    \ifx\@empty#2\@empty
        \def#1{\lst@fmtIfNextCharIn{#5}{#4}{#6#7}}%
    \else \ifx\@empty#3\@empty
        \def#1##1{%
            \ifx##1#2%
                \def\lst@next{\lst@fmtIfNextCharIn{#5}{#4##1}%
                                                      {#6#7}}%
            \else
                 \def\lst@next{#4##1}%
            \fi
            \lst@next}%
    \else
        \def#1{%
            \lst@IfNextCharsArg{#2#3}%
                {\lst@fmtIfNextCharIn{#5}{\expandafter#4\lst@eaten}%
                                         {#6#7}}%
                {\expandafter#4\lst@eaten}}%
    \fi \fi}
%    \end{macrocode}
% \end{macro}
%
%
% \paragraph{解析器}
% 把 |\lst@fmtSplit| 应用于把格式定义切成条目、把条目切成
% `输入'和`输出'，并把`输出'切成`前'和`后'。如果你接触过格式定义，
% 这应当很清楚。
%
% \begin{macro}{\lst@UseFormat}
% 现在我们可以开始写解析器了。
%    \begin{macrocode}
\gdef\lst@UseFormat#1{%
    \def\lst@fmtwhole{#1}%
    \lst@UseFormat@}
\gdef\lst@UseFormat@{%
    \lst@fmtSplit\lst@fmtwhole,%
%    \end{macrocode}
% 我们赋予格式定义的其余部分，\ldots
%    \begin{macrocode}
    \let\lst@fmtwhole\lst@fmtb
    \ifx\lst@fmta\@empty\else
%    \end{macrocode}
% \ldots\ 在等号处分割该条目，并处理该条目。
%    \begin{macrocode}
        \lst@fmtSplit\lst@fmta=%
        \ifx\@empty\lst@fmta\else
%    \end{macrocode}
% \begin{TODO}
% 插入 |\let\lst@arg\@empty| |\expandafter\lst@XConvert\lst@fmtb\@nil|
% |\let\lst@fmtb\lst@arg|.
% \end{TODO}
%    \begin{macrocode}
            \expandafter\lstKV@XOptArg\expandafter[\expandafter]%
                \expandafter{\lst@fmtb}\lst@UseFormat@b
        \fi
    \fi
%    \end{macrocode}
% 最后，如果其余部分非空，我们处理下一个条目。
%    \begin{macrocode}
    \ifx\lst@fmtwhole\@empty\else
        \expandafter\lst@UseFormat@
    \fi}
%    \end{macrocode}
% 我们让 |\lst@fmtc| 以前面的字符（作为带花括号的参数）为内容。
% 为了添加更多参数，我们首先在控制序列
% |\string| 处分割替换记号。
%    \begin{macrocode}
\gdef\lst@UseFormat@b[#1]#2{%
    \def\lst@fmtc{{#1}}\lst@lExtend\lst@fmtc{\expandafter{\lst@fmta}}%
    \def\lst@fmtb{#2}%
    \lst@fmtSplit\lst@fmtb\string
%    \end{macrocode}
% 我们追加一个空参数，或把带`|\string| 前置'记号的 |\lst@fmtPre| 作为参数。
% 对 |\string| 之后的记号我们做同样的事。
%    \begin{macrocode}
    \ifx\@empty\lst@fmta
        \lst@lAddTo\lst@fmtc{{}}%
    \else
        \lst@lExtend\lst@fmtc{\expandafter
            {\expandafter\lst@fmtPre\expandafter{\lst@fmta}}}%
    \fi
    \ifx\@empty\lst@fmtb
        \lst@lAddTo\lst@fmtc{{}}%
    \else
        \lst@lExtend\lst@fmtc{\expandafter
            {\expandafter\lst@fmtPost\expandafter{\lst@fmtb}}}%
    \fi
%    \end{macrocode}
% 最后我们适当地扩展 |\lst@fmtformat|。注意 |\lst@if|
% 仍指示替换记号是否含有 |\string|。
%    \begin{macrocode}
    \expandafter\lst@UseFormat@c\lst@fmtc}
%    \end{macrocode}
%    \begin{macrocode}
\gdef\lst@UseFormat@c#1#2#3#4{%
    \lst@fmtIfIdentifier#2\relax
    {\lst@fmtIdentifier{#2}%
     \lst@if\else \PackageWarning{Listings}%
         {Cannot drop identifier in format definition}%
     \fi}%
    {\lst@if
         \lst@lAddTo\lst@fmtformat{\lst@CArgX#2\relax\lst@fmtCDef}%
     \else
         \lst@lAddTo\lst@fmtformat{\lst@CArgX#2\relax\lst@fmtCDefX}%
     \fi
     \lst@DefActive\lst@fmtc{#1}%
     \lst@lExtend\lst@fmtformat{\expandafter{\lst@fmtc}{#3}{#4}}}}
%    \end{macrocode}
%    \begin{macrocode}
\lst@AddToHook{SelectCharTable}{\lst@fmtformat}
\global\let\lst@fmtformat\@empty
%    \end{macrocode}
% \end{macro}
%
%
% \paragraph{格式化}
%
% \begin{macro}{\lst@fmtPre}
%    \begin{macrocode}
\gdef\lst@fmtPre#1{%
    \lst@PrintToken
    \begingroup
    \let\newline\lst@fmtEnsureNewLine
    \let\space\lst@fmtEnsureSpace
    \let\indent\lst@fmtIndent
    \let\noindent\lst@fmtNoindent
    #1%
    \endgroup}
%    \end{macrocode}
% \end{macro}
%
% \begin{macro}{\lst@fmtPost}
%    \begin{macrocode}
\gdef\lst@fmtPost#1{%
    \global\let\lst@fmtPostOutput\@empty
    \begingroup
    \def\newline{\lst@AddTo\lst@fmtPostOutput\lst@fmtEnsureNewLine}%
    \def\space{\aftergroup\lst@fmtEnsurePostSpace}%
    \def\indent{\lst@AddTo\lst@fmtPostOutput\lst@fmtIndent}%
    \def\noindent{\lst@AddTo\lst@fmtPostOutput\lst@fmtNoindent}%
    \aftergroup\lst@PrintToken
    #1%
    \endgroup}
%    \end{macrocode}
%    \begin{macrocode}
\lst@AddToHook{Init}{\global\let\lst@fmtPostOutput\@empty}
\lst@AddToHook{PostOutput}
    {\lst@fmtPostOutput \global\let\lst@fmtPostOutput\@empty}
%    \end{macrocode}
% \end{macro}
%
% \begin{macro}{\lst@fmtEnsureSpace}
% \begin{macro}{\lst@fmtEnsurePostSpace}
%    \begin{macrocode}
\gdef\lst@fmtEnsureSpace{%
    \lst@ifwhitespace\else \expandafter\lst@ProcessSpace \fi}
\gdef\lst@fmtEnsurePostSpace{%
    \lst@IfNextCharWhitespace{}{\lst@ProcessSpace}}
%    \end{macrocode}
% \end{macro}
% \end{macro}
%
% \begin{lstkey}{fmtindent}
% \begin{macro}{\lst@fmtIndent}
% \begin{macro}{\lst@fmtNoindent}
%    \begin{macrocode}
\lst@Key{fmtindent}{20pt}{\def\lst@fmtindent{#1}}
\newdimen\lst@fmtcurrindent
\lst@AddToHook{InitVars}{\global\lst@fmtcurrindent\z@}
\gdef\lst@fmtIndent{\global\advance\lst@fmtcurrindent\lst@fmtindent}
\gdef\lst@fmtNoindent{\global\advance\lst@fmtcurrindent-\lst@fmtindent}
%    \end{macrocode}
% \end{macro}
% \end{macro}
% \end{lstkey}
%
% \begin{macro}{\lst@fmtEnsureNewLine}
%    \begin{macrocode}
\gdef\lst@fmtEnsureNewLine{%
    \global\advance\lst@newlines\@ne
    \global\advance\lst@newlinesensured\@ne
    \lst@fmtignoretrue}
%    \end{macrocode}
%    \begin{macrocode}
\lst@AddToAtTop\lst@DoNewLines{%
    \ifnum\lst@newlines>\lst@newlinesensured
        \global\advance\lst@newlines-\lst@newlinesensured
    \fi
    \global\lst@newlinesensured\z@}
\newcount\lst@newlinesensured % global
\lst@AddToHook{Init}{\global\lst@newlinesensured\z@}
%    \end{macrocode}
%    \begin{macrocode}
\gdef\lst@fmtignoretrue{\let\lst@fmtifignore\iftrue}
\gdef\lst@fmtignorefalse{\let\lst@fmtifignore\iffalse}
\lst@AddToHook{InitVars}{\lst@fmtignorefalse}
\lst@AddToHook{Output}{\lst@fmtignorefalse}
%    \end{macrocode}
% \end{macro}
%
% \begin{macro}{\lst@fmtUseLostSpace}
%    \begin{macrocode}
\gdef\lst@fmtUseLostSpace{%
    \lst@ifnewline \kern\lst@fmtcurrindent \global\lst@lostspace\z@
    \else
        \lst@OldOLS
    \fi}
\lst@AddToHook{Init}
    {\lst@true
     \ifx\lst@fmtformat\@empty \ifx\lst@fmt\@empty \lst@false \fi\fi
     \lst@if
        \let\lst@OldOLS\lst@OutputLostSpace
        \let\lst@OutputLostSpace\lst@fmtUseLostSpace
        \let\lst@ProcessSpace\lst@fmtProcessSpace
     \fi}
%    \end{macrocode}
% \begin{TODO}
% 这段`丢失空间'还没有使用 |\lst@alloverstyle|！
% \end{TODO}
% \end{macro}
%
% \begin{macro}{\lst@fmtProcessSpace}
%    \begin{macrocode}
\gdef\lst@fmtProcessSpace{%
    \lst@ifletter
        \lst@Output
        \lst@fmtifignore\else
            \lst@AppendOther\lst@outputspace
        \fi
    \else \lst@ifkeepspaces
        \lst@AppendOther\lst@outputspace
    \else \ifnum\lst@newlines=\z@
        \lst@AppendSpecialSpace
    \else \ifnum\lst@length=\z@
            \global\advance\lst@lostspace\lst@width
            \global\advance\lst@pos\m@ne
        \else
            \lst@AppendSpecialSpace
        \fi
    \fi \fi \fi
    \lst@whitespacetrue}
%    \end{macrocode}
% \end{macro}
%
%
% \paragraph{格式化标识符}
%
% \begin{macro}{\lst@fmtIdentifier}
% 我们为`格式标识符'安装一个（关键词）测试。
%    \begin{macrocode}
\lst@InstallTest{f}
    \lst@fmt@list\lst@fmt \lst@gfmt@list\lst@gfmt
    \lst@gfmt@wp
    wd
\gdef\lst@fmt@list{\lst@fmt\lst@gfmt}\global\let\lst@fmt\@empty
\gdef\lst@gfmt@list{\lst@fmt\lst@gfmt}\global\let\lst@gfmt\@empty
%    \end{macrocode}
% 工作过程展开 |\lst@fmt$|\meta{string}（并定义
% |\lst@PrintToken| 为不做任何事）。
%    \begin{macrocode}
\gdef\lst@gfmt@wp{%
    \begingroup \let\lst@UM\@empty
    \let\lst@PrintToken\@empty
    \csname\@lst @fmt$\the\lst@token\endcsname
    \endgroup}
%    \end{macrocode}
% 这个控制序列大概被定义为`工作标识符'。
%    \begin{macrocode}
\gdef\lst@fmtIdentifier#1#2#3#4{%
    \lst@DefOther\lst@fmta{#2}\edef\lst@fmt{\lst@fmt,\lst@fmta}%
    \@namedef{\@lst @fmt$\lst@fmta}{#3#4}}
%    \end{macrocode}
% |\lst@fmt$|\meta{identifier} 展开为由 |#2| 和 |#3| 定义的
% |\lst@fmtPre|/|\lst@fmtPost| 序列。
% \end{macro}
%
%    \begin{macrocode}
\lst@EndAspect
%</misc>
%    \end{macrocode}
% \end{aspect}
%
%
%
% \subsection{行号}
%
% \begin{aspect}{labels}
% \lsthelper{Rolf~Niepraschk}{1997/04/24}{行号}请求行
% 号，相关的键和`\textit{data}'命令在第~\ref{rLineNumbers} 节
% 描述，这里是它们的实现。
%    \begin{macrocode}
%<*misc>
\lst@BeginAspect{labels}
%    \end{macrocode}
%
% \begin{lstkey}{numbers}
% 根据参数，我们定义 |\lst@PlaceNumber| 来打印行
% 号。
%    \begin{macrocode}
\lst@Key{numbers}{none}{%
    \let\lst@PlaceNumber\@empty
    \lstKV@SwitchCases{#1}%
    {none:\\%
     left:\def\lst@PlaceNumber{\llap{\normalfont
                \lst@numberstyle{\thelstnumber}\kern\lst@numbersep}}\\%
     right:\def\lst@PlaceNumber{\rlap{\normalfont
                \kern\linewidth \kern\lst@numbersep
                \lst@numberstyle{\thelstnumber}}}%
    }{\PackageError{Listings}{value `#1' for keyword `numbers' not supported}\@ehc}}
%    \end{macrocode}
% \end{lstkey}
%
% \begin{lstkey}{stepnumber}
% \begin{lstkey}{numberfirstline}
% \begin{lstkey}{numberstyle}
% \begin{lstkey}{numbersep}
% \begin{lstkey}{numberblanklines}
% 各键的定义。
%    \begin{macrocode}
\lst@Key{stepnumber}{1}{\def\lst@stepnumber{#1\relax}}
\lst@Key{numberfirstline}{f}[t]{\lstKV@SetIf{#1}\lst@ifnumberfirstline}
\lst@Key{numberstyle}{}{\def\lst@numberstyle{#1}}
\lst@Key{numbersep}{10pt}{\def\lst@numbersep{#1}}
\lst@Key{numberblanklines}{true}[t]
    {\lstKV@SetIf{#1}\lst@ifnumberblanklines}
%    \end{macrocode}
% 如果我们使用`空'样式，列表必须从第 1 行开始，所以：
%    \begin{macrocode}
\lst@AddToHook{EmptyStyle}{\let\lst@stepnumber\@ne}
%    \end{macrocode}
% 稍后我们用这一点来防止第一行被第二次
% 编号。
%    \begin{macrocode}
\gdef\lst@numberfirstlinefalse{\let\lst@ifnumberfirstline\iffalse}
%    \end{macrocode}
% \end{lstkey}
% \end{lstkey}
% \end{lstkey}
% \end{lstkey}
% \end{lstkey}
%
% \begin{lstkey}{firstnumber}
% 我们根据参数选择第一个编号。
%    \begin{macrocode}
\lst@Key{firstnumber}{auto}{%
    \lstKV@SwitchCases{#1}%
    {auto:\let\lst@firstnumber\@undefined\\%
     last:\let\lst@firstnumber\c@lstnumber
    }{\def\lst@firstnumber{#1\relax}}}
\lst@AddToHook{PreSet}{\let\lst@advancenumber\z@}
%    \end{macrocode}
% |\lst@firstnumber| 现在被设为 |\lst@lineno| 而不是 |\lst@firstline|，
% 这是根据 |lstpatch.sty| 自 1.3b 起关于行范围标记的改动。
%    \begin{macrocode}
\lst@AddToHook{PreInit}
    {\ifx\lst@firstnumber\@undefined
         \def\lst@firstnumber{\lst@lineno}%
     \fi}
%    \end{macrocode}
% \end{lstkey}
%
% 键 |name| 在第~\ref{sec:list-listings} 节定义，数据
% 命令 |\thelstnumber| 见第~\pageref{iTheLstNumber} 页。
%
% \begin{macro}{\lst@SetFirstNumber}
% \begin{macro}{\lst@SaveFirstNumber}
% \lsthelper{Boris~Veytsman}{1998/03/25}{continue line numbering: a.c b.c a.c}
% 提议按列表名继续行号。我们在这里定义第一个打印行
% 的标签号。\lsthelper{Jens~Schwarzer}{2001/05/29}{first>1 时
% lstlisting 行号错误}报告的缺陷已通过把 |\@ne| 替换为 |\lst@firstline| 而消除。
%    \begin{macrocode}
\gdef\lst@SetFirstNumber{%
    \ifx\lst@firstnumber\@undefined
        \@tempcnta 0\csname\@lst no@\lst@intname\endcsname\relax
        \ifnum\@tempcnta=\z@ \@tempcnta\lst@firstline
                       \else \lst@nololtrue \fi
        \advance\@tempcnta\lst@advancenumber
        \edef\lst@firstnumber{\the\@tempcnta\relax}%
    \fi}
%    \end{macrocode}
% 当前标签保存在 |\lstno@|\meta{name} 中。如果名字为空，
% 我们就改用空格，这会使 |\lstno@| 保持未定义。
%    \begin{macrocode}
\gdef\lst@SaveFirstNumber{%
    \expandafter\xdef
        \csname\@lst no\ifx\lst@intname\@empty @ \else @\lst@intname\fi
        \endcsname{\the\c@lstnumber}}
%    \end{macrocode}
% \end{macro}
% \end{macro}
%
% \begin{macro}{\thelstnumber}
% \begin{macro}{\c@lstnumber}
% \label{iTheLstNumber}%
% 这个计数器保存当前标签号。我们把它用作当前标签，以
% 让行号可以被 |\ref| 引用。这由
% \lsthelper{Boris~Veytsman}{1998/03/25}{让行号可通过
% \label 和 \ref 引用}提出。我们现在用 |\refstepcounter| 完成这件事——感谢
% \lsthelper{Christian~Gudrian}{2000/11/13}{\ref{lst:line}
% 跳到列表顶部而不是该行}的缺陷报告。
%    \begin{macrocode}
\newcounter{lstnumber}% \global
\global\c@lstnumber\@ne % init
\renewcommand*\thelstnumber{\@arabic\c@lstnumber}
\lst@AddToHook{EveryPar}
    {\global\advance\c@lstnumber\lst@advancelstnum
     \global\advance\c@lstnumber\m@ne \refstepcounter{lstnumber}%
     \lst@SkipOrPrintLabel}%
\global\let\lst@advancelstnum\@ne
%    \end{macrocode}
% 注意计数器在标签被打印\emph{之前}推进，而不是
% 之后。否则我们会有错误的引用——由以下人员报告：
% \lsthelper{Gregory~Van~Vooren}{1999/06/04}{引用大了一个单位}。
%    \begin{macrocode}
\lst@AddToHook{Init}{\def\@currentlabel{\thelstnumber}}
%    \end{macrocode}
% \lsthelper{Felix Freiberger}{2025/03/12}{LaTeX listings 宏包：cleveref
% 无法再引用单个行}提到了文章
% \url{https://tex.stackexchange.com/q/738349/30156}，标题为
% ``\packagename{cleveref} 无法再引用
% |lstlisting| 中的行''。\lsthelper{Ulrike Fischer}{2025/03/05}{cleveref 无法再
% 引用 lstlisting 中的行}用下面一行提供了有效的解决方案。
%    \begin{macrocode}
\lst@AddToHook{Init}{\def\@currentcounter{lstnumber}}
%    \end{macrocode}
% 标签号被初始化，我们确保续接列表的行号正确。一个
% 表面多余的、在指定 \texttt{firstnumber=last} 时把行号
% 推进 \verb|-\lst@advancelstnum| 的操作已被删除，
% 这是根据 \lsthelper{Joachim~Breitner}%
% {2006/05/14}{failure to continue counting correctly}.
%    \begin{macrocode}
\lst@AddToHook{InitVars}
    {\global\c@lstnumber\lst@firstnumber
     \global\advance\c@lstnumber\lst@advancenumber
     \global\advance\c@lstnumber-\lst@advancelstnum}
\lst@AddToHook{ExitVars}
    {\global\advance\c@lstnumber\lst@advancelstnum}
%    \end{macrocode}
% \lsthelper{Walter~E.~Brown}{2001/05/22}{pdftex 3.14159-14f 警告：
% 具有相同标识符的目标}报告了 pdftex 与
% \packagename{hyperref} 的问题。原因是 |\theHlstlabel| 的一个糟糕默认值。
% \lsthelper{Heiko~Oberdiek}{2001/11/08}{pdftex 警告：具有
% 相同标识符的目标}发现了另一个缺陷，它源于
% |\lst@neglisting| 的本地化。他也提供了下面的修正，用
% |\ifx| \ldots\ |\fi| 构造替换 |\thelstlisting|。
% \lsthelper{Ivo~Pletikosi\'c}{2015/08/18}{listings + hyperref, lstmisc.sty}
% 报告了 |\thelstlisting| 重新定义方面的另一个问题。
% \lsthelper{Heiko~Oberdiek}{2001/11/08}%
% {http://tex.stackexchange.com/questions/261437/renewing-the-thelstlisting^^A
% -command-when-hyperref-is-on}
% 再次提供了修正：|\thelstlisting| 必须替换为 |\theHlstlisting|。
%    \begin{macrocode}
\AtBeginDocument{%
    \def\theHlstnumber{\ifx\lst@@caption\@empty \lst@neglisting
                                          \else \theHlstlisting \fi
                       .\thelstnumber}}
%    \end{macrocode}
% \end{macro}
% \end{macro}
%
% \begin{macro}{\lst@skipnumbers}
% 还有更多事情要做。我们计算有多少行必须跳过其
% 标签。公式是
%	$$|\lst@skipnumbers|=
%		\textrm{\emph{first printing line}}\bmod|\lst@stepnumber|.$$
% 注意我们对 |\lst@skipnumbers| 使用非正代表值。
%    \begin{macrocode}
\newcount\lst@skipnumbers % \global
\lst@AddToHook{Init}
    {\ifnum \z@>\lst@stepnumber
         \let\lst@advancelstnum\m@ne
         \edef\lst@stepnumber{-\lst@stepnumber}%
     \fi
     \ifnum \z@<\lst@stepnumber
         \global\lst@skipnumbers\lst@firstnumber
         \global\divide\lst@skipnumbers\lst@stepnumber
         \global\multiply\lst@skipnumbers-\lst@stepnumber
         \global\advance\lst@skipnumbers\lst@firstnumber
         \ifnum\lst@skipnumbers>\z@
             \global\advance\lst@skipnumbers -\lst@stepnumber
         \fi
%    \end{macrocode}
% 如果 |\lst@stepnumber| 为零，就不打印任何行号：
%    \begin{macrocode}
     \else
         \let\lst@SkipOrPrintLabel\relax
     \fi}
%    \end{macrocode}
% \end{macro}
%
% \begin{macro}{\lst@SkipOrPrintLabel}
% 但这是默认情况。我们利用 |\lst@skipnumbers| 非正这一事实。
% 每行都推进计数器，如果该计数器为零，我们就打印一个行
% 号并把计数器减去 |\lst@stepnumber|。
%    \begin{macrocode}
\gdef\lst@SkipOrPrintLabel{%
    \ifnum\lst@skipnumbers=\z@
        \global\advance\lst@skipnumbers-\lst@stepnumber\relax
        \lst@PlaceNumber
        \lst@numberfirstlinefalse
    \else
%    \end{macrocode}
% 如果列表的第一行应该得到编号，它就在这里得到。
%    \begin{macrocode}
        \lst@ifnumberfirstline
            \lst@PlaceNumber
            \lst@numberfirstlinefalse
        \fi
    \fi
    \global\advance\lst@skipnumbers\@ne}%
%    \end{macrocode}
%    \begin{macrocode}
\lst@AddToHook{OnEmptyLine}{%
    \lst@ifnumberblanklines\else \ifnum\lst@skipnumbers=\z@
        \global\advance\lst@skipnumbers-\lst@stepnumber\relax
    \fi\fi}
%    \end{macrocode}
% \end{macro}
%
%    \begin{macrocode}
\lst@EndAspect
%</misc>
%    \end{macrocode}
% \end{aspect}
%
%
% \subsection{行形状与断行}
%
% \begin{macro}{\lst@parshape}
% 我们为未加载 \aspectname{lineshape} 方面的情况定义
% |\lst@parshape| 的默认版本。我们每行（事实上每个段落）
% 都使用这个 parshape。此外，如果我们关闭一个分组
% 层级，就必须重复该 parshape——否则形状会被忘记。
%    \begin{macrocode}
%<*kernel>
\def\lst@parshape{\parshape\@ne \z@ \linewidth}
\lst@AddToHookAtTop{EveryLine}{\lst@parshape}
\lst@AddToHookAtTop{EndGroup}{\lst@parshape}
%</kernel>
%    \end{macrocode}
% \end{macro}
%
% \begin{aspect}{lineshape}
% 本节我们的第一个方面。
%    \begin{macrocode}
%<*misc>
\lst@BeginAspect{lineshape}
%    \end{macrocode}
%
% \begin{lstkey}{xleftmargin}
% \begin{lstkey}{xrightmargin}
% \begin{lstkey}{resetmargins}
% \begin{lstkey}{linewidth}
% 通常的内容。
%    \begin{macrocode}
\lst@Key{xleftmargin}{\z@}{\def\lst@xleftmargin{#1}}
\lst@Key{xrightmargin}{\z@}{\def\lst@xrightmargin{#1}}
\lst@Key{resetmargins}{false}[t]{\lstKV@SetIf{#1}\lst@ifresetmargins}
%    \end{macrocode}
% 如果我们围绕列表做一个精确的盒子，边距就变为零。
%    \begin{macrocode}
\lst@AddToHook{BoxUnsafe}{\let\lst@xleftmargin\z@
                          \let\lst@xrightmargin\z@}
\lst@AddToHook{TextStyle}{%
    \let\lst@xleftmargin\z@ \let\lst@xrightmargin\z@
    \let\lst@ifresetmargins\iftrue}
%    \end{macrocode}
% 上面添加的钩子是在 \lsthelper{Magnus~Lewis-Smith}
%{1999/08/06}{|\lstinline| 被缩进}和 \lsthelper{Jos\'e~Romildo~Malaquias}
%{2000/08/22}{|\lstinline| 被缩进（resetmargins）}的缺陷报告之后。
%    \begin{macrocode}
\lst@Key{linewidth}\linewidth{\def\lst@linewidth{#1}}
\lst@AddToHook{PreInit}{\linewidth\lst@linewidth\relax}
%    \end{macrocode}
% \end{lstkey}
% \end{lstkey}
% \end{lstkey}
% \end{lstkey}
%
% \begin{macro}{\lst@parshape}
% 定义本身很容易。
%    \begin{macrocode}
\gdef\lst@parshape{%
    \parshape\@ne \@totalleftmargin \linewidth}
%    \end{macrocode}
% 我们计算列表的行宽和（内/外）缩进。
%    \begin{macrocode}
\lst@AddToHook{Init}
    {\lst@ifresetmargins
         \advance\linewidth\@totalleftmargin
         \advance\linewidth\rightmargin
         \@totalleftmargin\z@
     \fi
     \advance\linewidth-\lst@xleftmargin
     \advance\linewidth-\lst@xrightmargin
     \advance\@totalleftmargin\lst@xleftmargin\relax}
%    \end{macrocode}
% \end{macro}
%
% \begin{lstkey}{lineskip}
% 这个键的引入源于与
% \lsthelper{Andreas~Bartelt}{1997/09/11}{重新定义 \parskip 的问题；
% 引入 \lstlineskip}的交流。1.0 版通过重新定义
% |\baselinestretch| 实现该特性。
%    \begin{macrocode}
\lst@Key{lineskip}{\z@}{\def\lst@lineskip{#1\relax}}
\lst@AddToHook{Init}
    {\parskip\z@
     \ifdim\z@=\lst@lineskip\else
         \@tempdima\baselineskip
         \advance\@tempdima\lst@lineskip
%    \end{macrocode}
% 下面三行模拟`糟糕的' |\divide| |\@tempdima|
% |\strip@pt| |\baselineskip| |\relax|。感谢 \lsthelper{Peter~Bartke}
% {2002/04/10}{糟糕地使用 \strip@pt}的缺陷报告。
%    \begin{macrocode}
         \multiply\@tempdima\@cclvi
         \divide\@tempdima\baselineskip\relax
         \multiply\@tempdima\@cclvi
%    \end{macrocode}
%    \begin{macrocode}
         \edef\baselinestretch{\strip@pt\@tempdima}%
         \selectfont
     \fi}
%    \end{macrocode}
% \end{lstkey}
%
% \begin{lstkey}{breaklines}
% \begin{lstkey}{breakindent}
% \begin{lstkey}{breakautoindent}
% \begin{lstkey}{breakatwhitespace}
% \begin{lstkey}{prebreak}
% \begin{lstkey}{postbreak}
% 像往常一样，宣布更多键没有问题。
% \keyname{breakatwhitespace} 源自 \lsthelper{Javier~Bezos}{2003/09/23}
% {breaklines 在奇怪的位置断行}。不幸的是，该键先前的定义
% 是错误的，正如 \lsthelper{Franz~Rinnerthaler}{2004/03/12}
% {breakatwhitespace 没有效果}和 \lsthelper{Ulrike~Fischer}{2004/07/11}
% {breakatwhitespace 没有效果}所报告的那样。
%    \begin{macrocode}
\lst@Key{breaklines}{false}[t]{\lstKV@SetIf{#1}\lst@ifbreaklines}
\lst@Key{breakindent}{20pt}{\def\lst@breakindent{#1}}
\lst@Key{breakautoindent}{t}[t]{\lstKV@SetIf{#1}\lst@ifbreakautoindent}
\lst@Key{breakatwhitespace}{false}[t]%
    {\lstKV@SetIf{#1}\lst@ifbreakatwhitespace}
\lst@Key{prebreak}{}{\def\lst@prebreak{#1}}
\lst@Key{postbreak}{}{\def\lst@postbreak{#1}}
%    \end{macrocode}
% 我们赋一些不同的宏，（必要时）抑制``underfull
% |\hbox|''消息（并使用不同的 pretolerance）：
%    \begin{macrocode}
\lst@AddToHook{Init}
    {\lst@ifbreaklines
         \hbadness\@M \pretolerance\@M
         \@rightskip\@flushglue \rightskip\@rightskip % \raggedright
         \leftskip\z@skip \parindent\z@
%    \end{macrocode}
% 上面的 |\raggedright| 已被手工设值取代，这是在
% \lsthelper{Morten~H\o gholm}{2004/09/06}{ltugboat.cls 与
% listings}的缺陷报告之后。
%
% 我们使用正常的 parshape 和计算出的 |\lst@breakshape|（见下）。
%    \begin{macrocode}
         \def\lst@parshape{\parshape\tw@ \@totalleftmargin\linewidth
                           \lst@breakshape}%
     \else
         \let\lst@discretionary\@empty
     \fi}
\lst@AddToHook{OnNewLine}
    {\lst@ifbreaklines \lst@breakNewLine \fi}
%    \end{macrocode}
% \end{lstkey}\end{lstkey}\end{lstkey}\end{lstkey}
% \end{lstkey}\end{lstkey}
%
% \begin{macro}{\lst@discretionary}
% \begin{macro}{\lst@spacekern}
% 现在来了全部魔法：我们在每个`输出单元'之后放置一个
% 可断点。不过我们重新定义 |\space| 以便在 |\discretionary| 内使用，
% 并使用 \hookname{EveryLine} 钩子。在 \lsthelper{Carsten~Hamm}
% {2002/04/19}{breaklines 且 xleftmargin>0pt 时框线错误}的缺陷报告之后，我
% 加入了 |\kern-\lst@xleftmargin|，它后来变成 |\kern-\@totalleftmargin|，
% 这是在 \lsthelper{Christian~Kaiser}{2002/12/13}{breaklines=true 时
% itemize 内的框线错误}的缺陷报告之后。
%    \begin{macrocode}
\gdef\lst@discretionary{%
    \lst@ifbreakatwhitespace
        \lst@ifwhitespace \lst@@discretionary \fi
    \else
        \lst@@discretionary
    \fi}%
\gdef\lst@@discretionary{%
    \discretionary{\let\space\lst@spacekern\lst@prebreak}%
                  {\llap{\lsthk@EveryLine
                   \kern\lst@breakcurrindent \kern-\@totalleftmargin}%
                   \let\space\lst@spacekern\lst@postbreak}{}}
\lst@AddToHook{PostOutput}{\lst@discretionary}
\gdef\lst@spacekern{\kern\lst@width}
%    \end{macrocode}
% \begin{ALTERNATIVE}
% 在每个`输出单元'\emph{之前}（即在输出宏中的 |\hbox{...}| 之前）
% 使用 |\penalty\@M \hskip\z@ plus 1fil \penalty0\hskip\z@ plus-1fil|
% 也能按需要断行。但那样我们就没有 |prebreak| 和 |postbreak| 了。
% \end{ALTERNATIVE}
% \end{macro}\end{macro}
%
% \begin{macro}{\lst@breakNewLine}
% 我们使用 \keyname{breakindent}，另外如果`自动缩进'开启，
% 还使用当前行缩进（来自行首的空白）。
%    \begin{macrocode}
\gdef\lst@breakNewLine{%
    \@tempdima\lst@breakindent\relax
    \lst@ifbreakautoindent \advance\@tempdima\lst@lostspace \fi
%    \end{macrocode}
% 现在我们计算折行部分的边距和行宽\ldots
%    \begin{macrocode}
    \@tempdimc-\@tempdima \advance\@tempdimc\linewidth
                          \advance\@tempdima\@totalleftmargin
%    \end{macrocode}
% \ldots\ 并把它存入 |\lst@breakshape|。
%    \begin{macrocode}
    \xdef\lst@breakshape{\noexpand\lst@breakcurrindent \the\@tempdimc}%
    \xdef\lst@breakcurrindent{\the\@tempdima}}
\global\let\lst@breakcurrindent\z@ % init
%    \end{macrocode}
% |\lst@breakcurrindent| 的初始化是在
% \lsthelper{Alvaro~Herrera}{2002/12/09}{fancyvrb 与 breaklines 时
% 出现`未定义的控制序列 \lst@breakcurrindent'}的缺陷报告之后添加的。
% \begin{TODO}
% 我们可以通过分配两个全局尺寸来加速。
% \end{TODO}
% \end{macro}
%
% \begin{macro}{\lst@breakshape}
% \lsthelper{Andreas~Deininger}{2000/08/25}{`breaklines,first>1'导致
% ``未定义的控制序列''错误}报告了一个问题，通过提供
% 默认的断行形状得以解决。
%    \begin{macrocode}
\gdef\lst@breakshape{\@totalleftmargin \linewidth}
%    \end{macrocode}
% \end{macro}
%
% \begin{macro}{\lst@breakProcessOther}
% 与 |\lst@ProcessOther| 相同，只是它还会输出当前
% 记号字符串。这会插入一个可能的断行点。
% 目前只有右圆括号使用这个宏。
%    \begin{macrocode}
\gdef\lst@breakProcessOther#1{\lst@ProcessOther#1\lst@OutputOther}
\lst@AddToHook{SelectCharTable}
    {\lst@ifbreaklines \lst@Def{`)}{\lst@breakProcessOther)}\fi}
%    \end{macrocode}
% \lsthelper{Gabriel~Tauro}{2001/04/18}{字符出现在第一个打印行
% 之前时出现意外的`)'}报告的缺陷，已通过使用
% |\lst@ProcessOther| 而不是 |\lst@AppendOther| 而消除。
% \end{macro}
%
%    \begin{macrocode}
\lst@EndAspect
%</misc>
%    \end{macrocode}
% \end{aspect}
%
%
% \subsection{边框}
%
% \begin{aspect}{frames}
% 另一个方面。
%    \begin{macrocode}
%<*misc>
\lst@BeginAspect[lineshape]{frames}
%    \end{macrocode}
%
% \begin{lstkey}{framexleftmargin}
% \begin{lstkey}{framexrightmargin}
% \begin{lstkey}{framextopmargin}
% \begin{lstkey}{framexbottommargin}
% 这些键只是保存参数。
%    \begin{macrocode}
\lst@Key{framexleftmargin}{\z@}{\def\lst@framexleftmargin{#1}}
\lst@Key{framexrightmargin}{\z@}{\def\lst@framexrightmargin{#1}}
\lst@Key{framextopmargin}{\z@}{\def\lst@framextopmargin{#1}}
\lst@Key{framexbottommargin}{\z@}{\def\lst@framexbottommargin{#1}}
%    \end{macrocode}
% \end{lstkey}
% \end{lstkey}
% \end{lstkey}
% \end{lstkey}
%
% \begin{lstkey}{backgroundcolor}
% \lsthelper{Ralf~Imh\"auser}{2000/01/08}{彩色背景}启发了
% 键 \keyname{backgroundcolor}。所有键都保存参数，并且\ldots
%    \begin{macrocode}
\lst@Key{backgroundcolor}{}{\def\lst@bkgcolor{#1}}
\lst@Key{fillcolor}{}{\def\lst@fillcolor{#1}}
\lst@Key{rulecolor}{}{\def\lst@rulecolor{#1}}
\lst@Key{rulesepcolor}{}{\def\lst@rulesepcolor{#1}}
%    \end{macrocode}
% \ldots\ 其中一些若为空则有默认设置。
%    \begin{macrocode}
\lst@AddToHook{Init}{%
    \ifx\lst@fillcolor\@empty
        \let\lst@fillcolor\lst@bkgcolor
    \fi
    \ifx\lst@rulesepcolor\@empty
        \let\lst@rulesepcolor\lst@fillcolor
    \fi}
%    \end{macrocode}
% \end{lstkey}
%
% \begin{lstkey}{rulesep}
% \begin{lstkey}{framerule}
% \begin{lstkey}{framesep}
% \begin{lstkey}{frameshape}
% 另一组键，它们主要保存各自的参数。
% \keyname{frameshape} 把所有字母变成大写，并检查是否至少有
% 一个圆角。最后我们定义 |\lst@frame| 为空当且
% 仅当没有 frameshape。
%    \begin{macrocode}
\lst@Key{rulesep}{2pt}{\def\lst@rulesep{#1}}
\lst@Key{framerule}{.4pt}{\def\lst@framerulewidth{#1}}
\lst@Key{framesep}{3pt}{\def\lst@frametextsep{#1}}
\lst@Key{frameshape}{}{%
    \let\lst@xrulecolor\@empty
    \lstKV@FourArg{#1}%
    {\uppercase{\def\lst@frametshape{##1}}%
     \uppercase{\def\lst@framelshape{##2}}%
     \uppercase{\def\lst@framershape{##3}}%
     \uppercase{\def\lst@framebshape{##4}}%
     \let\lst@ifframeround\iffalse
     \lst@IfSubstring R\lst@frametshape{\let\lst@ifframeround\iftrue}{}%
     \lst@IfSubstring R\lst@framebshape{\let\lst@ifframeround\iftrue}{}%
     \def\lst@frame{##1##2##3##4}}}
%    \end{macrocode}
% \end{lstkey}
% \end{lstkey}
% \end{lstkey}
% \end{lstkey}
%
% \begin{lstkey}{frameround}
% \begin{lstkey}{frame}
% 我们这里必须做一些转换。
%    \begin{macrocode}
\lst@Key{frameround}\relax
    {\uppercase{\def\lst@frameround{#1}}%
     \expandafter\lstframe@\lst@frameround ffff\relax}
\global\let\lst@frameround\@empty
%    \end{macrocode}
% 如果是 verbose 参数，我们就使用 |trbl| 子集替换。
%    \begin{macrocode}
\lst@Key{frame}\relax{%
    \let\lst@xrulecolor\@empty
    \lstKV@SwitchCases{#1}%
    {none:\let\lst@frame\@empty\\%
     leftline:\def\lst@frame{l}\\%
     topline:\def\lst@frame{t}\\%
     bottomline:\def\lst@frame{b}\\%
     lines:\def\lst@frame{tb}\\%
     single:\def\lst@frame{trbl}\\%
     shadowbox:\def\lst@frame{tRBl}%
            \def\lst@xrulecolor{\lst@rulesepcolor}%
            \def\lst@rulesep{\lst@frametextsep}%
    }{\def\lst@frame{#1}}%
    \expandafter\lstframe@\lst@frameround ffff\relax}
%    \end{macrocode}
% 在其大写版本的情况下加入 |t|、|r|、|b| 和 |l| 会让
% 后面的测试更容易。
%    \begin{macrocode}
\gdef\lstframe@#1#2#3#4#5\relax{%
    \lst@IfSubstring T\lst@frame{\edef\lst@frame{t\lst@frame}}{}%
    \lst@IfSubstring R\lst@frame{\edef\lst@frame{r\lst@frame}}{}%
    \lst@IfSubstring B\lst@frame{\edef\lst@frame{b\lst@frame}}{}%
    \lst@IfSubstring L\lst@frame{\edef\lst@frame{l\lst@frame}}{}%
%    \end{macrocode}
% 现在我们检查上下框线，\ldots
%    \begin{macrocode}
    \let\lst@frametshape\@empty \let\lst@framebshape\@empty
    \lst@frameCheck
        ltr\lst@framelshape\lst@frametshape\lst@framershape #4#1%
    \lst@frameCheck
        LTR\lst@framelshape\lst@frametshape\lst@framershape #4#1%
    \lst@frameCheck
        lbr\lst@framelshape\lst@framebshape\lst@framershape #3#2%
    \lst@frameCheck
        LBR\lst@framelshape\lst@framebshape\lst@framershape #3#2%
%    \end{macrocode}
% \ldots\ 查找圆角\ldots
%    \begin{macrocode}
    \let\lst@ifframeround\iffalse
    \lst@IfSubstring R\lst@frametshape{\let\lst@ifframeround\iftrue}{}%
    \lst@IfSubstring R\lst@framebshape{\let\lst@ifframeround\iftrue}{}%
%    \end{macrocode}
% 并定义左右框形状。
%    \begin{macrocode}
    \let\lst@framelshape\@empty \let\lst@framershape\@empty
    \lst@IfSubstring L\lst@frame
        {\def\lst@framelshape{YY}}%
        {\lst@IfSubstring l\lst@frame{\def\lst@framelshape{Y}}{}}%
    \lst@IfSubstring R\lst@frame
        {\def\lst@framershape{YY}}%
        {\lst@IfSubstring r\lst@frame{\def\lst@framershape{Y}}{}}}
%    \end{macrocode}
% 现在来看用于定义上下框形状的宏。
% 它扩展宏 |#5|。最后两个参数表明左右角是否圆角。
% |#4| 和 |#6| 是临时宏。
% |#1#2#3| 是我们测试的三个字符。
%    \begin{macrocode}
\gdef\lst@frameCheck#1#2#3#4#5#6#7#8{%
    \lst@IfSubstring #1\lst@frame
        {\if #7T\def#4{R}\else \def#4{Y}\fi}%
        {\def#4{N}}%
    \lst@IfSubstring #3\lst@frame
        {\if #8T\def#6{R}\else \def#6{Y}\fi}%
        {\def#6{N}}%
    \lst@IfSubstring #2\lst@frame{\edef#5{#5#4Y#6}}{}}
%    \end{macrocode}
% 对于文本样式列表，所有框线和背景色都被
% 停用——这是在 \lsthelper{Stephen~Reindl}%
% {2002/06/04}{文本样式列表未停用框线}和
% \lsthelper{Thomas~ten~Cate}{2006/07/14}{行内列表在断行后
% 得到背景色}的缺陷报告之后添加的。
%    \begin{macrocode}
\lst@AddToHook{TextStyle}
   {\let\lst@frame\@empty
    \let\lst@frametshape\@empty
    \let\lst@framershape\@empty
    \let\lst@framebshape\@empty
    \let\lst@framelshape\@empty
    \let\lst@bkgcolor\@empty}
%    \end{macrocode}
% \end{lstkey}
% \end{lstkey}
%
% 根据 \lsthelper{Ignacio~Fern\'andez~Galv\'an}{2006/07/26}%
% {带背景色的框线在左侧有小洞}的缺陷报告，边距左侧
% 那一小块背景色现在在绘制框线左侧
% 之前绘制，以便在 Acrobat 中正确重叠。
%
% \begin{macro}{\lst@frameMakeBoxV}
%    \begin{macrocode}
\gdef\lst@frameMakeBoxV#1#2#3{%
    \setbox#1\hbox{%
      \color@begingroup \lst@rulecolor
      \ifx\lst@framelshape\@empty
      \else
            \llap{%
                \lst@frameBlock\lst@fillcolor\lst@frametextsep{#2}{#3}%
                \kern\lst@framexleftmargin}%
      \fi
      \llap{\setbox\z@\hbox{\vrule\@width\z@\@height#2\@depth#3%
                            \lst@frameL}%
            \rlap{\lst@frameBlock\lst@rulesepcolor{\wd\z@}%
                                                  {\ht\z@}{\dp\z@}}%
            \box\z@
            \kern\lst@frametextsep\relax
            \kern\lst@framexleftmargin}%
      \rlap{\kern-\lst@framexleftmargin
                    \@tempdima\linewidth
            \advance\@tempdima\lst@framexleftmargin
            \advance\@tempdima\lst@framexrightmargin
            \lst@frameBlock\lst@bkgcolor\@tempdima{#2}{#3}%
            \ifx\lst@framershape\@empty
                \kern\lst@frametextsep\relax
            \else
                \lst@frameBlock\lst@fillcolor\lst@frametextsep{#2}{#3}%
            \fi
            \setbox\z@\hbox{\vrule\@width\z@\@height#2\@depth#3%
                            \lst@frameR}%
            \rlap{\lst@frameBlock\lst@rulesepcolor{\wd\z@}%
                                                  {\ht\z@}{\dp\z@}}%
            \box\z@}%
      \color@endgroup}}
%    \end{macrocode}
% \end{macro}
%
% \begin{macro}{\lst@frameBlock}
%    \begin{macrocode}
\gdef\lst@frameBlock#1#2#3#4{%
    \color@begingroup
      #1%
      \setbox\z@\hbox{\vrule\@height#3\@depth#4%
                      \ifx#1\@empty \@width\z@ \kern#2\relax
                              \else \@width#2\relax \fi}%
      \box\z@
    \color@endgroup}
%    \end{macrocode}
% \end{macro}
%
% \begin{macro}{\lst@frameR}
% 排印右框线。
% 我们只需遍历 |\lst@framershape|。
%    \begin{macrocode}
\gdef\lst@frameR{%
    \expandafter\lst@frameR@\lst@framershape\relax
    \kern-\lst@rulesep}
\gdef\lst@frameR@#1{%
    \ifx\relax#1\@empty\else
        \if #1Y\lst@framevrule \else \kern\lst@framerulewidth \fi
        \kern\lst@rulesep
        \expandafter\lst@frameR@b
    \fi}
\gdef\lst@frameR@b#1{%
    \ifx\relax#1\@empty
    \else
        \if #1Y\color@begingroup
               \lst@xrulecolor
               \lst@framevrule
               \color@endgroup
        \else
               \kern\lst@framerulewidth
        \fi
        \kern\lst@rulesep
        \expandafter\lst@frameR@
    \fi}
%    \end{macrocode}
% \end{macro}
%
% \begin{macro}{\lst@frameL}
% 左框线同理。
%    \begin{macrocode}
\gdef\lst@frameL{%
    \kern-\lst@rulesep
    \expandafter\lst@frameL@\lst@framelshape\relax}
\gdef\lst@frameL@#1{%
    \ifx\relax#1\@empty\else
        \kern\lst@rulesep
        \if#1Y\lst@framevrule \else \kern\lst@framerulewidth \fi
        \expandafter\lst@frameL@
    \fi}
%    \end{macrocode}
% \end{macro}
%
% \begin{macro}{\lst@frameH}
% 这是用于绘制上下框线的核心宏。
% 第一个参数是 |T| 或 |B|，第二个包含形状。
% 我们用 |\@tempcntb| 作为尺码计数器。
%    \begin{macrocode}
\gdef\lst@frameH#1#2{%
    \global\let\lst@framediml\z@ \global\let\lst@framedimr\z@
    \setbox\z@\hbox{}\@tempcntb\z@
    \expandafter\lst@frameH@\expandafter#1#2\relax\relax\relax
            \@tempdimb\lst@frametextsep\relax
    \advance\@tempdimb\lst@framerulewidth\relax
            \@tempdimc-\@tempdimb
    \advance\@tempdimc\ht\z@
    \advance\@tempdimc\dp\z@
    \setbox\z@=\hbox{%
      \lst@frameHBkg\lst@fillcolor\@tempdimb\@firstoftwo
      \if#1T\rlap{\raise\dp\@tempboxa\box\@tempboxa}%
       \else\rlap{\lower\ht\@tempboxa\box\@tempboxa}\fi
      \lst@frameHBkg\lst@rulesepcolor\@tempdimc\@secondoftwo
      \advance\@tempdimb\ht\@tempboxa
      \if#1T\rlap{\raise\lst@frametextsep\box\@tempboxa}%
       \else\rlap{\lower\@tempdimb\box\@tempboxa}\fi
      \rlap{\box\z@}%
    }}
\gdef\lst@frameH@#1#2#3#4{%
    \ifx\relax#4\@empty\else
        \lst@frameh \@tempcntb#1#2#3#4%
        \advance\@tempcntb\@ne
        \expandafter\lst@frameH@\expandafter#1%
    \fi}
\gdef\lst@frameHBkg#1#2#3{%
    \setbox\@tempboxa\hbox{%
        \kern-\lst@framexleftmargin
        #3{\kern-\lst@framediml\relax}{\@tempdima\z@}%
        \ifdim\lst@framediml>\@tempdimb
            #3{\@tempdima\lst@framediml \advance\@tempdima-\@tempdimb
               \lst@frameBlock\lst@rulesepcolor\@tempdima\@tempdimb\z@}%
              {\kern-\lst@framediml
               \advance\@tempdima\lst@framediml\relax}%
        \fi
        #3{\@tempdima\z@
           \ifx\lst@framelshape\@empty\else
               \advance\@tempdima\@tempdimb
           \fi
           \ifx\lst@framershape\@empty\else
               \advance\@tempdima\@tempdimb
           \fi}%
          {\ifdim\lst@framedimr>\@tempdimb
              \advance\@tempdima\lst@framedimr\relax
           \fi}%
        \advance\@tempdima\linewidth
        \advance\@tempdima\lst@framexleftmargin
        \advance\@tempdima\lst@framexrightmargin
        \lst@frameBlock#1\@tempdima#2\z@
        #3{\ifdim\lst@framedimr>\@tempdimb
               \@tempdima-\@tempdimb
               \advance\@tempdima\lst@framedimr\relax
               \lst@frameBlock\lst@rulesepcolor\@tempdima\@tempdimb\z@
           \fi}{}%
        }}
%    \end{macrocode}
% \end{macro}
%
% \begin{macro}{\lst@frameh}
% 这是用于绘制上下框线的低级宏。
% 它向盒子 0 \emph{添加}一条线和角。
% 第一个参数给出角的大小，第二个是
% |T| 或 |B|。
% |#3#4#5| 是对框线从左到右的描述，位置在
% $\{$\texttt{Y,N,R}$\}\times\{$\texttt{Y,N}$\}\times\{$\texttt{Y,N,R}$\}$.
% 我们移到正确的水平位置，设置左角、
% 水平线，以及右角。
%    \begin{macrocode}
\gdef\lst@frameh#1#2#3#4#5{%
    \lst@frameCalcDimA#1%
    \lst@ifframeround \@getcirc\@tempdima \fi
%    \end{macrocode}
%    \begin{macrocode}
    \setbox\z@\hbox{%
      \begingroup
      \setbox\z@\hbox{%
        \kern-\lst@framexleftmargin
        \color@begingroup
        \ifnum#1=\z@ \lst@rulecolor \else \lst@xrulecolor \fi
%    \end{macrocode}
% |\lst@frameCorner| 取四个参数：
% |\llap|、|TL| 或 |BL|、角类型 $\in\{$\texttt{Y,N,R}$\}$，以及
% 尺寸 |#1|。
%    \begin{macrocode}
        \lst@frameCornerX\llap{#2L}#3#1%
        \ifdim\lst@framediml<\@tempdimb
            \xdef\lst@framediml{\the\@tempdimb}%
        \fi
        \begingroup
        \if#4Y\else \let\lst@framerulewidth\z@ \fi
                \@tempdima\lst@framexleftmargin
        \advance\@tempdima\lst@framexrightmargin
        \advance\@tempdima\linewidth
        \vrule\@width\@tempdima\@height\lst@framerulewidth \@depth\z@
        \endgroup
        \lst@frameCornerX\rlap{#2R}#5#1%
        \ifdim\lst@framedimr<\@tempdimb
            \xdef\lst@framedimr{\the\@tempdimb}%
        \fi
        \color@endgroup}%
%    \end{macrocode}
%    \begin{macrocode}
      \if#2T\rlap{\raise\dp\z@\box\z@}%
       \else\rlap{\lower\ht\z@\box\z@}\fi
      \endgroup
      \box\z@}}
%    \end{macrocode}
% \end{macro}
%
% \begin{macro}{\lst@frameCornerX}
% 排印单个角并返回 |\@tempdimb|，即该角的宽度。
%    \begin{macrocode}
\gdef\lst@frameCornerX#1#2#3#4{%
    \setbox\@tempboxa\hbox{\csname\@lst @frame\if#3RR\fi #2\endcsname}%
    \@tempdimb\wd\@tempboxa
    \if #3R%
        #1{\box\@tempboxa}%
    \else
        \if #3Y\expandafter#1\else
               \@tempdimb\z@ \expandafter\vphantom \fi
        {\box\@tempboxa}%
    \fi}
%    \end{macrocode}
% \end{macro}
%
% \begin{macro}{\lst@frameCalcDimA}
% 计算总宽度；由 |\lst@frameh| 和 |\lst@frameInit| 使用。
%    \begin{macrocode}
\gdef\lst@frameCalcDimA#1{%
            \@tempdima\lst@rulesep
    \advance\@tempdima\lst@framerulewidth
    \multiply\@tempdima#1\relax
    \advance\@tempdima\lst@frametextsep
    \advance\@tempdima\lst@framerulewidth
    \multiply\@tempdima\tw@}
%    \end{macrocode}
% \end{macro}
%
% \begin{macro}{\lst@frameInit}
% 首先我们查看左侧和右侧有哪些框线类型。
% 如果没有竖线，我们就加速处理。
%    \begin{macrocode}
\lst@AddToHook{Init}{\lst@frameInit}
\newbox\lst@framebox
\gdef\lst@frameInit{%
    \ifx\lst@framelshape\@empty \let\lst@frameL\@empty \fi
    \ifx\lst@framershape\@empty \let\lst@frameR\@empty \fi
    \def\lst@framevrule{\vrule\@width\lst@framerulewidth\relax}%
%    \end{macrocode}
% 必要时我们调整数值以处理圆角。
%    \begin{macrocode}
    \lst@ifframeround
        \lst@frameCalcDimA\z@ \@getcirc\@tempdima
        \@tempdimb\@tempdima \divide\@tempdimb\tw@
        \advance\@tempdimb -\@wholewidth
        \edef\lst@frametextsep{\the\@tempdimb}%
        \edef\lst@framerulewidth{\the\@wholewidth}%
%    \end{macrocode}
%    \begin{macrocode}
        \lst@frameCalcDimA\@ne \@getcirc\@tempdima
        \@tempdimb\@tempdima \divide\@tempdimb\tw@
        \advance\@tempdimb -\tw@\@wholewidth
        \advance\@tempdimb -\lst@frametextsep
        \edef\lst@rulesep{\the\@tempdimb}%
    \fi
%    \end{macrocode}
%    \begin{macrocode}
    \lst@frameMakeBoxV\lst@framebox{\ht\strutbox}{\dp\strutbox}%
    \def\lst@framelr{\copy\lst@framebox}%
%    \end{macrocode}
% 最后我们排印各条线（+ 角）。
% 我们可能需要插入负 |\vskip| 以去掉前面文本
% 与上框线之间的空间。
% \begin{TODO}
% 改用 |\vspace| 而不是 |\vskip|？
% \end{TODO}
%    \begin{macrocode}
    \ifx\lst@frametshape\@empty\else
        \lst@frameH T\lst@frametshape
        \ifvoid\z@\else
            \par\lst@parshape
            \@tempdima-\baselineskip \advance\@tempdima\ht\z@
            \ifdim\prevdepth<\@cclvi\p@\else
                \advance\@tempdima\prevdepth
            \fi
            \ifdim\@tempdima<\z@
                \vskip\@tempdima\vskip\lineskip
            \fi
            \noindent\box\z@\par
            \lineskiplimit\maxdimen \lineskip\z@
        \fi
        \lst@frameSpreadV\lst@framextopmargin
    \fi}
%    \end{macrocode}
% |\parshape\lst@parshape| 确保上框线正确缩进。
% 该缺陷由 \lsthelper{Marcin~Kasperski}{1999/04/28}{itemize 内的
% 上框线向右缩进}报告。
%
% 我们每行都排印左右框线。
%    \begin{macrocode}
\lst@AddToHook{EveryLine}{\lst@framelr}
\global\let\lst@framelr\@empty
%    \end{macrocode}
% \end{macro}
%
% \begin{macro}{\lst@frameExit}
% 底部的线。
%    \begin{macrocode}
\lst@AddToHook{DeInit}
    {\ifx\lst@framebshape\@empty\else \lst@frameExit \fi}
\gdef\lst@frameExit{%
    \lst@frameSpreadV\lst@framexbottommargin
    \lst@frameH B\lst@framebshape
    \ifvoid\z@\else
        \everypar{}\par\lst@parshape\nointerlineskip\noindent\box\z@
    \fi}
%    \end{macrocode}
% \end{macro}
%
% \begin{macro}{\lst@frameSpreadV}
% 为垂直延展设置框线。
%    \begin{macrocode}
\gdef\lst@frameSpreadV#1{%
    \ifdim\z@=#1\else
        \everypar{}\par\lst@parshape\nointerlineskip\noindent
        \lst@frameMakeBoxV\z@{#1}{\z@}%
        \box\z@
    \fi}
%    \end{macrocode}
% \end{macro}
%
% \begin{macro}{\lst@frameTR}
% \begin{macro}{\lst@frameBR}
% \begin{macro}{\lst@frameBL}
% \begin{macro}{\lst@frameTL}
% 这些宏绘制竖线和横线。
% 隐式参数 |\@tempdima| 给出两个角的大小，
% 由 |\lst@frameh| 提供。
%    \begin{macrocode}
\gdef\lst@frameTR{%
    \vrule\@width.5\@tempdima\@height\lst@framerulewidth\@depth\z@
    \kern-\lst@framerulewidth
    \raise\lst@framerulewidth\hbox{%
        \vrule\@width\lst@framerulewidth\@height\z@\@depth.5\@tempdima}}
\gdef\lst@frameBR{%
    \vrule\@width.5\@tempdima\@height\lst@framerulewidth\@depth\z@
    \kern-\lst@framerulewidth
    \vrule\@width\lst@framerulewidth\@height.5\@tempdima\@depth\z@}
\gdef\lst@frameBL{%
    \vrule\@width\lst@framerulewidth\@height.5\@tempdima\@depth\z@
    \kern-\lst@framerulewidth
    \vrule\@width.5\@tempdima\@height\lst@framerulewidth\@depth\z@}
\gdef\lst@frameTL{%
    \raise\lst@framerulewidth\hbox{%
        \vrule\@width\lst@framerulewidth\@height\z@\@depth.5\@tempdima}%
    \kern-\lst@framerulewidth
    \vrule\@width.5\@tempdima\@height\lst@framerulewidth\@depth\z@}
%    \end{macrocode}
% \end{macro}\end{macro}\end{macro}\end{macro}
%
% \begin{macro}{\lst@frameRoundT}
% \begin{macro}{\lst@frameRoundB}
% 是排印圆角的辅助宏。我们把高度和深度设为
% 圆形字体的可见部分。
%    \begin{macrocode}
\gdef\lst@frameRoundT{%
    \setbox\@tempboxa\hbox{\@circlefnt\char\@tempcnta}%
    \ht\@tempboxa\lst@framerulewidth
    \box\@tempboxa}
\gdef\lst@frameRoundB{%
    \setbox\@tempboxa\hbox{\@circlefnt\char\@tempcnta}%
    \dp\@tempboxa\z@
    \box\@tempboxa}
%    \end{macrocode}
% \end{macro}
% \end{macro}
%
% \begin{macro}{\lst@frameRTR}
% \begin{macro}{\lst@frameRBR}
% \begin{macro}{\lst@frameRBL}
% \begin{macro}{\lst@frameRTL}
% 圆角。
%    \begin{macrocode}
\gdef\lst@frameRTR{%
    \hb@xt@.5\@tempdima{\kern-\lst@framerulewidth
                           \kern.5\@tempdima \lst@frameRoundT \hss}}
\gdef\lst@frameRBR{%
    \hb@xt@.5\@tempdima{\kern-\lst@framerulewidth
    \advance\@tempcnta\@ne \kern.5\@tempdima \lst@frameRoundB \hss}}
\gdef\lst@frameRBL{%
    \advance\@tempcnta\tw@ \lst@frameRoundB
    \kern-.5\@tempdima}
\gdef\lst@frameRTL{%
    \advance\@tempcnta\thr@@\lst@frameRoundT
    \kern-.5\@tempdima}
%    \end{macrocode}
% \end{macro}\end{macro}\end{macro}\end{macro}
%
%    \begin{macrocode}
\lst@EndAspect
%</misc>
%    \end{macrocode}
% \end{aspect}
%
%
% \subsection{make 的宏用法}
%
% \begin{aspect}{make}
% \begin{macro}{\lst@makemode}
% \begin{macro}{\lst@ifmakekey}
% 如果我们进入了 Make 的特殊模式，我们就保存最后一个
% 标识符是否是一阶关键词。
%    \begin{macrocode}
%<*misc>
\lst@BeginAspect[keywords]{make}
%    \end{macrocode}
%    \begin{macrocode}
\lst@NewMode\lst@makemode
\lst@AddToHook{Output}{%
    \ifnum\lst@mode=\lst@makemode
        \ifx\lst@thestyle\lst@gkeywords@sty
            \lst@makekeytrue
        \fi
    \fi}
%    \end{macrocode}
%    \begin{macrocode}
\gdef\lst@makekeytrue{\let\lst@ifmakekey\iftrue}
\gdef\lst@makekeyfalse{\let\lst@ifmakekey\iffalse}
\global\lst@makekeyfalse % init
%    \end{macrocode}
% \end{macro}\end{macro}
%
% \begin{lstkey}{makemacrouse}
% 必要时调整字符表
%    \begin{macrocode}
\lst@Key{makemacrouse}f[t]{\lstKV@SetIf{#1}\lst@ifmakemacrouse}
%    \end{macrocode}
% \end{lstkey}
%
% \begin{macro}{\lst@MakeSCT}
% 如果`宏使用'开启，开 |$(| 打印前面的字符、进入
% 特殊模式，并把这两个字符与随后的输出合并。
%
%    \begin{macrocode}
\gdef\lst@MakeSCT{%
    \lst@ifmakemacrouse
        \lst@ReplaceInput{$(}{%
            \lst@PrintToken
            \lst@EnterMode\lst@makemode{\lst@makekeyfalse}%
            \lst@Merge{\lst@ProcessOther\$\lst@ProcessOther(}}%
%    \end{macrocode} ^^A $
% 闭圆括号测试模式，并或者照常处理 |)|，
% 或者就在这里输出它（若 |$(| 和 |)| 之间有关键词，则用关键词样式）。
% |)|).
%    \begin{macrocode}
        \lst@ReplaceInput{)}{%
            \ifnum\lst@mode=\lst@makemode
                \lst@PrintToken
                \begingroup
                    \lst@ProcessOther)%
                    \lst@ifmakekey
                        \let\lst@currstyle\lst@gkeywords@sty
                    \fi
                    \lst@OutputOther
                \endgroup
                \lst@LeaveMode
            \else
                \expandafter\lst@ProcessOther\expandafter)%
            \fi}%
%    \end{macrocode}
% 如果 \keyname{makemacrouse} 关闭，那么两个 |$(| 都只是`其他'。
%    \begin{macrocode}
    \else
        \lst@ReplaceInput{$(}{\lst@ProcessOther\$\lst@ProcessOther(}%
    \fi}
%    \end{macrocode} ^^A $
% \end{macro}
%
%    \begin{macrocode}
\lst@EndAspect
%</misc>
%    \end{macrocode}
% \end{aspect}
%
%
% \section{排印列表}
% \label{sec:typesetting-listing}
%
% \begingroup
%    \begin{macrocode}
%<*kernel>
%    \end{macrocode}
% \endgroup
% \begin{macro}{\lst@lineno}
% \begin{lstkey}{print}
% \begin{lstkey}{firstline}
% \begin{lstkey}{lastline}
% `当前行'计数器和三个选项键。
%    \begin{macrocode}
\newcount\lst@lineno % \global
\lst@AddToHook{InitVars}{\global\lst@lineno\@ne}
%    \end{macrocode}
%    \begin{macrocode}
\lst@Key{print}{true}[t]{\lstKV@SetIf{#1}\lst@ifprint}
\lst@Key{firstline}\relax{\def\lst@firstline{#1\relax}}
\lst@Key{lastline}\relax{\def\lst@lastline{#1\relax}}
%    \end{macrocode}
% 用合理的值初始化行范围
%    \begin{macrocode}
\lst@AddToHook{PreSet}
{%
  \def\lst@firstline{1\relax}
  \def\lst@lastline{9999999\relax}
  \let\lst@linerange\@empty
}
%    \end{macrocode}
% \end{lstkey}
% \end{lstkey}
% \end{lstkey}
%
% \begin{lstkey}{nolol}
%     在这里只是又一个含义明显的键。我们当然会在下面
%     用到它。
%    \begin{macrocode}
\lst@Key{nolol}{false}[t]{\lstKV@SetIf{#1}\lst@ifnolol}
\def\lst@nololtrue{\let\lst@ifnolol\iftrue}
\let\lst@ifnolol\iffalse % init
%    \end{macrocode}
% \end{lstkey}
% \end{macro}
%
%
% \subsection{处理行范围}
% \label{sec:lineranges}
%
% 下面的代码只是从当前开发版本以及 Carsten 在 1.3b 版中
% 留下的 |lstpatch.sty| 文件里复制来的，用于带编号和
% 范围标记的行范围。它使用了更多键：
%
% \begin{lstkey}{linerange}
% \begin{lstkey}{consecutivenumbers}
% \begin{lstkey}{rangeprefix}
% \begin{lstkey}{rangesuffix}
% \begin{lstkey}{rangebeginprefix}
% \begin{lstkey}{rangebeginsuffix}
% \begin{lstkey}{rangeendprefix}
% \begin{lstkey}{rangeendsuffix}
% \begin{lstkey}{includerangemarker}
% 首先是控制行范围处理的选项。
%    \begin{macrocode}
\lst@Key{linerange}\relax{\lstKV@OptArg[]{#1}{%
    \def\lst@interrange{##1}\def\lst@linerange{##2,}}}
\lst@Key{consecutivenumbers}{true}[t]{%
  \lstKV@SetIf{#1}\lst@ifconsecutivenumbers}
%    \end{macrocode}
% 接下来的选项用于方便地描述任意行范围
% 标记：
%    \begin{macrocode}
\lst@Key{rangeprefix}\relax{\def\lst@rangebeginprefix{#1}%
                            \def\lst@rangeendprefix{#1}}
\lst@Key{rangesuffix}\relax{\def\lst@rangebeginsuffix{#1}%
                            \def\lst@rangeendsuffix{#1}}
\lst@Key{rangebeginprefix}{}{\def\lst@rangebeginprefix{#1}}
\lst@Key{rangebeginsuffix}{}{\def\lst@rangebeginsuffix{#1}}
\lst@Key{rangeendprefix}{}{\def\lst@rangeendprefix{#1}}
\lst@Key{rangeendsuffix}{}{\def\lst@rangeendsuffix{#1}}
\lst@Key{includerangemarker}{true}[t]{\lstKV@SetIf{#1}%
  \lst@ifincluderangemarker}
%    \end{macrocode}
% 行范围是已知的——它或者被显式设置（通过编号或
% 任意行范围标记），或者由默认值（1 和
% 999999）隐式给出。|\lst@GetLineInterval| 递归解析已知的行范围，
% 因为可能存在逗号分隔的成对列表。
%    \begin{macrocode}
\lst@AddToHook{Init}
{%
  \ifx\lst@linerange\@empty
      \edef\lst@linerange{{\lst@firstline}-{\lst@lastline},}%
  \fi%
  \lst@GetLineInterval%
}%
\def\lst@GetLineInterval{\expandafter\lst@GLI\lst@linerange\@nil}
%    \end{macrocode}
% 逗号分隔的成对列表的分割由以下完成
%    \begin{macrocode}
\def\lst@GLI#1,#2\@nil{% GLI: get line interval
%    \end{macrocode}
% 列表被分成两部分：列表的第一个元素（\#1）和
% 列表的其余部分（\#2），后者被存入 |\lst@linerange| 以便
% 稍后处理。
%    \begin{macrocode}
  \def\lst@linerange{#2}\lst@GLI@#1--\@nil%
}
%    \end{macrocode}
% 进一步解析由 |lst@GLI@| 完成，参数是 |#1--\@nil|，
% 即列表第一个元素加上一个哨兵。这个
% 参数用模板 |#1-#2-#3\@nil| 解析。
%
%    \begin{macrocode}
\def\lst@GLI@#1-#2-#3\@nil{%
%    \end{macrocode}
% 例如以`|12-21|'给出的行范围导致参数
% |12-21--\@nil|，因此被解析为 \#1\(\leftarrow\) |12|、
% \#2\(\leftarrow\)|21|、\#3\(\leftarrow\)|-|。
%
% 否则也可以用任意行范围标记定义范围，
% 如`|start-end|'，这导致参数
% |12-21--\@nil|，因此被解析为 \#1\(\leftarrow\) |12|、
% \#2\(\leftarrow\)|21|、\#3\(\leftarrow\)|-|。


% 首先我们测试行范围是否以数字（由 -、1、
% 2、3、\ldots 组成）或任意行范围标记开头。如果行范围
% 错误地以`|-|'字符开头，则参数 \#1 不被设置，于是
% 范围从 1 开始，给出也许符合预期的结果，但
% 行范围的定义必须总是至少由三部分组成：
% \meta{first}、`|-|'和 \meta{last}。
%    \begin{macrocode}
  \lst@IfNumber{#1}%
  {%
    \ifx\@empty#1\@empty
        \let\lst@firstline\@ne
    \else
        \def\lst@firstline{#1\relax}
    \fi
%    \end{macrocode}
% 现在我们知道行范围的起始编号。对于只由一个
% 编号组成的行范围，\#3 不被设置；在其他所有情况下 \#2 定义
% 最后一行（显式或隐式）。
%    \begin{macrocode}
    \ifx\@empty#3\@empty
        \def\lst@lastline{9999999\relax}
        \ifx\@empty#2\@empty
            \let\lst@lastline\lst@firstline
        \fi
    \else%
        \ifx\@empty#2\@empty
            \def\lst@lastline{9999999\relax}
        \else % doesn't happen(?)
            \def\lst@lastline{#2\relax}
        \fi
    \fi
  }%
%    \end{macrocode}
%        如果我们找到了任意行范围标记，就把 firstline 和 lastline 设为
%        9999999。这（几乎）阻止现在打印任何内容。
%    \begin{macrocode}
  {%
    \def\lst@firstline{9999999\relax}
    \let\lst@lastline\lst@firstline
%    \end{macrocode}
%        我们把前缀和后缀加到标记上。
%    \begin{macrocode}
    \let\lst@rangebegin\lst@rangebeginprefix
    \lst@AddTo\lst@rangebegin{#1}%
    \lst@Extend\lst@rangebegin\lst@rangebeginsuffix
    \ifx\@empty#3\@empty
        \let\lst@rangeend\lst@rangeendprefix
        \lst@AddTo\lst@rangeend{#1}
        \lst@Extend\lst@rangeend\lst@rangeendsuffix
    \else
        \ifx\@empty#2\@empty
            \let\lst@rangeend\@empty
        \else
            \let\lst@rangeend\lst@rangeendprefix
            \lst@AddTo\lst@rangeend{#2}%
            \lst@Extend\lst@rangeend\lst@rangeendsuffix
        \fi
    \fi
%    \end{macrocode}
% 下面的定义将在 \hookname{SelectCharTable} 钩子中执行，
% 并且如果我们已经在处理列表，就立即在这里执行。
%    \begin{macrocode}
    \global\def\lst@DefRange{%
      \expandafter\lst@CArgX\lst@rangebegin\relax\lst@DefRangeB}%
    \ifnum\lst@mode=\lst@Pmode \expandafter\lst@DefRange \fi%
  }%
}
%    \end{macrocode}
% \end{lstkey}
% \end{lstkey}
% \end{lstkey}
% \end{lstkey}
% \end{lstkey}
% \end{lstkey}
% \end{lstkey}
% \end{lstkey}
% \end{lstkey}
%
% \begin{macro}{\lst@DefRange}
%        |\lst@DefRange| 不再通过钩子插入。现在它由
%        |\lst@SelectCharTable| 直接调用。为了摆脱与
%        转义到 LaTeX 特性的干扰，这是必要的。该缺陷由
%        \lsthelper{Michael~Bachmann}{2004/07/21}{无法进行 label 引用
%        \ldots}报告。另一处改动也源于同一缺陷：当找到代码开头时，
%        |\lst@DefRange| 被全局重新定义，见下。该缺陷
%        由 \lsthelper{Tobias~Rapp}{2004/04/06}{列表跨分页时
%        未检测到范围末尾}和 \lsthelper{Markus~Luisser}{2004/08/13}{折行列表中的
%        `linerangemarker'缺陷}报告。
%    \begin{macrocode}
\lst@AddToHookExe{DeInit}{\global\let\lst@DefRange\@empty}
%    \end{macrocode}
%
% 实际上定义标记（通过 |\lst@GLI@|、|\lst@DefRange|、
% |\lst@CArgX|，如上看）类似 |\lst@DefDelimB|——只是
% 我们展开第一个参数并使用不同的 \meta{execute}、
% \meta{pre} 和 \meta{post} 语句。
%    \begin{macrocode}
\def\lst@DefRangeB#1#2{\lst@DefRangeB@#1#2}
\def\lst@DefRangeB@#1#2#3#4{%
  \lst@CDef{#1{#2}{#3}}#4{}{%
    \lst@ifincluderangemarker
        \lst@LeaveMode
        \let#1#4%
        \lst@DefRangeEnd
        \lst@InitLstNumber
    \else
        \@tempcnta\lst@lineno \advance\@tempcnta\@ne
        \edef\lst@firstline{\the\@tempcnta\relax}%
        \gdef\lst@OnceAtEOL{\let#1#4\lst@DefRangeEnd}%
        \lst@InitLstNumber
    \fi
    \global\let\lst@DefRange\lst@DefRangeEnd
    \lst@CArgEmpty%
  }%
  \@empty%
}
%    \end{macrocode}
% \end{macro}
%
% \begin{macro}{\lst@SetFirstNumber}
% 根据
% \lsthelper{Omair-Inam~Abdul-Matin}{2004/05/10}{实验性行范围
% 特性与 firstnumber 不兼容}报告的错误，修改标签并定义上面用到的 |\lst@InitLstNumber|。
%    \begin{macrocode}
\def\lstpatch@labels{%
  \gdef\lst@SetFirstNumber{%
     \ifx\lst@firstnumber\@undefined
         \@tempcnta 0\csname\@lst no@\lst@intname\endcsname\relax
         \ifnum\@tempcnta=\z@ \else
             \lst@nololtrue
             \advance\@tempcnta\lst@advancenumber
             \edef\lst@firstnumber{\the\@tempcnta\relax}
         \fi
    \fi
  }
}
%    \end{macrocode}
% \end{macro}
%
% \begin{macro}{\lst@InitLstNumber}
%    \begin{macrocode}
\def\lst@InitLstNumber{%
  \global\c@lstnumber\lst@firstnumber
  \global\advance\c@lstnumber\lst@advancenumber
  \global\advance\c@lstnumber-\lst@advancelstnum
  \ifx \lst@firstnumber\c@lstnumber
      \global\advance\c@lstnumber-\lst@advancelstnum
  \fi%
%    \end{macrocode}
% \lsthelper{Byron~K.~Boulton}{2013/11/21}{line numbers off by one}
% 报告说，当行范围由模式给出且设置了
% |includerangemarker=false| 时，显示的行号会偏移一位。
% 加入这个测试修正了这一行为。
%    \begin{macrocode}
  \lst@ifincluderangemarker\else%
    \global\advance\c@lstnumber by 1%
  \fi%
}
%    \end{macrocode}
% \end{macro}
%
%        结束标记被定义当且仅当它非空。该定义
%        类似 |\lst@DefDelimE|——带有上面的例外，并且除了我们
%        定义重入点 |\lst@DefRangeE@@| 的方式与上面
%        |\lst@MProcessListing| 新版本中相同。
%    \begin{macrocode}
\def\lst@DefRangeEnd{%
  \ifx\lst@rangeend\@empty\else
      \expandafter\lst@CArgX\lst@rangeend\relax\lst@DefRangeE
  \fi}
\def\lst@DefRangeE#1#2{\lst@DefRangeE@#1#2}
\def\lst@DefRangeE@#1#2#3#4{%
  \lst@CDef{#1#2{#3}}#4{}%
  {\let#1#4%
    \edef\lst@lastline{\the\lst@lineno\relax}%
    \lst@DefRangeE@@}%
  \@empty}
\def\lst@DefRangeE@@#1\@empty{%
  \lst@ifincluderangemarker
    #1\lst@XPrintToken
  \fi
  \lst@LeaveModeToPmode
  \lst@BeginDropInput{\lst@Pmode}}
\def\lst@LeaveModeToPmode{%
  \ifnum\lst@mode=\lst@Pmode
      \expandafter\lsthk@EndGroup
  \else
      \expandafter\egroup\expandafter\lst@LeaveModeToPmode
  \fi}
%    \end{macrocode}
%
% 有时把一个程序的多个片段连续编号是有好处的，
% 即 1--11 而不是 3--5、12--17、20--21。因此我们
% 引入 |\lst@DisplayConsecutiveNumbersOrNot|，它把
% 显示的行号修正为总是等于输入行号
%    \begin{macrocode}
\def\lst@DisplayConsecutiveNumbersOrNot{%
    \lst@ifconsecutivenumbers\else%
        \c@lstnumber=\numexpr-1+\lst@lineno
    \fi%
}
%    \end{macrocode}
%
%        最后我们不应忘记安装 |\lst@OnceAtEOL|，它也必须
%        在 |\lst@MSkipToFirst| 中被调用。
%    \begin{macrocode}
\lst@AddToHook{EOL}{\lst@OnceAtEOL\global\let\lst@OnceAtEOL\@empty}
\gdef\lst@OnceAtEOL{}% Init
%    \end{macrocode}
% 下面的代码在 listings 1.4 版中引入。该代码
% 类似 1.3c 版的第 14031--14061 行，但包含循环：
% \newline|\def\lst@next{...}\expandafter\lst@next|.\newline
% 实际上那段代码永远不会被展开，因为第~\ref{iTheInputCommand} 节下面
% 有两个类似的定义，它们会覆盖
% 这两条命令。
%    \begin{macrocode}
\def\lst@MSkipToFirst{%
    \global\advance\lst@lineno\@ne
    \ifnum \lst@lineno=\lst@firstline
        \def\lst@next{\lst@LeaveMode \global\lst@newlines\z@
        \lst@OnceAtEOL \global\let\lst@OnceAtEOL\@empty
        \lst@InitLstNumber % Added to work with modified \lsthk@PreInit.
        \lsthk@InitVarsBOL
        \lst@BOLGobble}%
        \expandafter\lst@next
    \fi}
\def\lst@SkipToFirst{%
    \ifnum \lst@lineno<\lst@firstline
        \def\lst@next{\lst@BeginDropInput\lst@Pmode
        \lst@Let{13}\lst@MSkipToFirst
        \lst@Let{10}\lst@MSkipToFirst}%
        \expandafter\lst@next
    \else
        \expandafter\lst@BOLGobble
    \fi}
%    \end{macrocode}
%
%    最后是服务宏 |\lst@IfNumber|：
%    \begin{macrocode}
\def\lst@IfNumber#1{%
    \ifx\@empty#1\@empty
        \let\lst@next\@firstoftwo
    \else
        \lst@IfNumber@#1\@nil
    \fi
    \lst@next}
\def\lst@IfNumber@#1#2\@nil{%
    \let\lst@next\@secondoftwo
    \ifnum`#1>47\relax \ifnum`#1>57\relax\else
        \let\lst@next\@firstoftwo
    \fi\fi}
%    \end{macrocode}
%
%
% \subsection{浮动体、盒子与题注}
%
% \begin{lstkey}{captionpos}
% \begin{lstkey}{abovecaptionskip}
% \begin{lstkey}{belowcaptionskip}
% \begin{lstkey}{label}
% \begin{lstkey}{title}
% \begin{lstkey}{caption}
% 一些键和\ldots
%    \begin{macrocode}
\lst@Key{captionpos}{t}{\def\lst@captionpos{#1}}
\lst@Key{abovecaptionskip}\smallskipamount{%
  \def\lst@abovecaption{#1}}
\lst@Key{belowcaptionskip}\smallskipamount{%
  \def\lst@belowcaption{#1}}
%    \end{macrocode}
% \lsthelper{Rolf~Niepraschk}{2000/01/10}{键：title}提出
% \keyname{title}.
%    \begin{macrocode}
\lst@Key{label}\relax{\def\lst@label{#1}}
\lst@Key{title}\relax{\def\lst@title{#1}\let\lst@caption\relax}
\lst@Key{caption}\relax{\lstKV@OptArg[{#1}]{#1}%
    {\def\lst@caption{##2}\def\lst@@caption{##1}}%
     \let\lst@title\@empty}
\lst@AddToHookExe{TextStyle}
    {\let\lst@caption\@empty \let\lst@@caption\@empty
     \let\lst@title\@empty \let\lst@label\@empty}
%    \end{macrocode}
% \end{lstkey}
% \end{lstkey}
% \end{lstkey}
% \end{lstkey}
% \end{lstkey}
% \end{lstkey}
%
% \begin{macro}{\thelstlisting}
% \begin{macro}{\lstlistingname}
% \begin{macro}{\lstlistingnamestyle}
% \begin{lstkey}{numberbychapter}
% \ldots\ 以及题注编号的样子。我把 |\ifx| |\@undefined| 改为
% |\@ifundefined|，这是在
% \lsthelper{Denis~Girou}{1999/07/26}{先加载 hyperref 再加载
% listings 时不兼容}的错误报告之后。
%
% 这个设置在 |\AtBeginDocument| 中完成，以便用户可以在计数器定义之前
% 用 |numberbychapter| 键指定是否在每个章处重置
% 该计数器。
%    \begin{macrocode}
\AtBeginDocument{%
  \@ifundefined{thechapter}{\let\lst@ifnumberbychapter\iffalse}{}
  \lst@ifnumberbychapter
      \newcounter{lstlisting}[chapter]
      \gdef\thelstlisting%
           {\ifnum \c@chapter>\z@ \thechapter.\fi \@arabic\c@lstlisting}
  \else
      \newcounter{lstlisting}
      \gdef\thelstlisting{\@arabic\c@lstlisting}
  \fi}
%    \end{macrocode}
%    \begin{macrocode}
\lst@UserCommand\lstlistingname{Listing}
\lst@UserCommand\lstlistingnamestyle{}
%    \end{macrocode}
%    \begin{macrocode}
\lst@Key{numberbychapter}{true}[t]{%
  \lstKV@SetIf{#1}\lst@ifnumberbychapter}
%    \end{macrocode}
% \end{lstkey}
% \end{macro}
% \end{macro}
% \end{macro}
%
% \begin{macro}{\lst@MakeCaption}
% 在定义这个宏之前，我们确保其他一些控制序列
% 存在——\lsthelper{Adam~Prugel-Bennett}{2001/02/19}{\abovecaptionskip
% 在 slides.cls 中未定义}报告了 slides 文档类的问题。
% 特别地，我们分配上题注和下题注间距寄存器，并定义
% |\@makecaption|，它是 article 类中该定义的确切副本。
% 为遵守 LPPL：你应当在你的 \TeX\ 系统上有该类的副本，或者从 CTAN
% 取得副本，例如从 ftp 服务器
% \texttt{ftp.dante.de}.
%
% Axel Sommerfeldt 提出了若干关于题注和
% 标题的改进。第一项是把间距寄存器的定义与
% |\@makecaption|.
%    \begin{macrocode}
\@ifundefined{abovecaptionskip}{%
  \newskip\abovecaptionskip%
  \newskip\belowcaptionskip%
}{}
\@ifundefined{@makecaption}{%
  \long\def\@makecaption#1#2{%
    \vskip\abovecaptionskip%
    \sbox\@tempboxa{#1: #2}%
    \ifdim \wd\@tempboxa >\hsize%
        #1: #2\par%
    \else%
        \global \@minipagefalse%
        \hb@xt@\hsize{\hfil\box\@tempboxa\hfil}%
    \fi%
   \vskip\belowcaptionskip}%
}{}
%    \end{macrocode}
% |\fnum@lstlisting| 的引入也归功于 Axel。此前替换文本
% 直接在 |\lst@MakeCaption| 中使用。|\noindent| 已
% 被移到别处，并变成 |\@parboxrestore|，这是在
% \lsthelper{Frank~Mittelbach}{2004/02/13}{回复：信息：caption 与
% listings 之间的不兼容}的缺陷报告之后。\lsthelper{Karl Berry}{2019/08/18}{加粗的
%   "listings 1"？}请求能够自定义标签。于是
% |\lstlistingnamestyle| 被引入到 |\lstlistingname| 前面。
%    \begin{macrocode}
\def\fnum@lstlisting{%
  {\lstlistingnamestyle\lstlistingname
    \ifx\lst@@caption\@empty\else~\thelstlisting\fi}%
}
%    \end{macrocode}
% 硬编码扩展名难以使用不同的扩展名，例如用于
% 附录。\lsthelper{Markus Kohm}{2023/09/05}{扩展 listings 使其
% 也主动支持 tocbasic}建议改为定义并使用 |\ext@lstlisting|。
%    \begin{macrocode}
\def\ext@lstlisting{lol}
%    \end{macrocode}
% 题注只为显示样式列表设置——感谢
% \lsthelper{Peter~L\"offler}{2004/04/24}{pdfTeX 警告（dest）：名称 figure.1
% 已被引用但不存在}报告该缺陷，并感谢
% \lsthelper{Axel~Sommerfeldt}{2004/02/27}{Re: caption + listings + hyperref}
% 分析了该缺陷。
% 我们 |\refstepcounter| 列表计数器当且仅当 |\lst@@caption|
% 非空。否则我们确保正确的超引用，
% 见下面的 |\lst@HRefStepCounter|。我们每个列表做一次，即在
% 顶部。
%    \begin{macrocode}
\def\lst@MakeCaption#1{%
  \lst@ifdisplaystyle
    \ifx #1t \allowbreak%
        \ifx\lst@@caption\@empty\expandafter\lst@HRefStepCounter \else
                                \expandafter\refstepcounter
        \fi {lstlisting}%
        \ifx\lst@label\@empty\else \label{\lst@label}\fi
%    \end{macrocode}
% 下面的代码在 \lsthelper{Rolf~Niepraschk}{2003/06/11}{题注与
% 列表之间的分页}的缺陷报告之后，从 \hookname{Init} 钩子移到这里。此外 |\lst@name| 等的初始化
% 是在 \lsthelper{Werner~Struckmann}
% 的缺陷报告之后插入这里的。
% {2003/06/25}{undefined control sequence \lst@name}.
% 如果名字既不为空也不是单个空格，我们就做一个`lol'条目。但
% 我们先测试 |\lst@|(|@|)|caption| 和 |\lst@ifnolol|。
%    \begin{macrocode}
        \let\lst@arg\lst@intname \lst@ReplaceIn\lst@arg\lst@filenamerpl
        \global\let\lst@name\lst@arg \global\let\lstname\lst@name
        \lst@ifnolol\else
            \ifx\lst@@caption\@empty
                \ifx\lst@caption\@empty
                    \ifx\lst@intname\@empty \else \def\lst@temp{ }%
                    \ifx\lst@intname\lst@temp \else
                        \addcontentsline{\ext@lstlisting}{lstlisting}\lst@name
                    \fi\fi
                \fi
            \else
                \addcontentsline{\ext@lstlisting}{lstlisting}%
                    {\protect\numberline{\thelstlisting}\lst@@caption}%
            \fi
        \fi
    \fi
%    \end{macrocode}
% 我们当且仅当题注非空且用户在 |#1| $\in\{\mathtt t,\mathtt b\}$ 处
% 请求了题注时才制作题注。为了禁止题注（或标题）与
% 列表之间的分页，我们局部重新定义原语 |\vskip| 以插入 |\nobreak|。
% 注意我们允许在`上题注'之前和`下题注'之后
% 分页。另外，|\@makecaption| 调用中的 |\ignorespaces| 是为了与
% \LaTeX\ 在 |\@caption| 中的做法一致；AMSbook 类（也许还有其他类）
% 假定它存在，并在测试空题注时试图去掉它，从而引起了
% \lsthelper{Xiaobo~Peng}{2006/06/29}{使用 amsbook 类时不显示题注}
% 注意到的缺陷。
% \begin{TODO}
% 这次重新定义是一个暴力方法。有没有更好的办法？
% \end{TODO}
%    \begin{macrocode}
    \ifx\lst@caption\@empty\else
        \lst@IfSubstring #1\lst@captionpos
            {\begingroup \let\@@vskip\vskip
             \def\vskip{\afterassignment\lst@vskip \@tempskipa}%
             \def\lst@vskip{\nobreak\@@vskip\@tempskipa\nobreak}%
             \par\@parboxrestore\normalsize\normalfont % \noindent (AS)
             %%\ifx #1t\allowbreak \fi
             \ifx\lst@title\@empty
                 \lst@makecaption\fnum@lstlisting{%
                   \ignorespaces \lst@caption}
             \else
                 \lst@maketitle\lst@title % (AS)
             \fi
             \ifx #1b\allowbreak \fi
             \endgroup}{}%
    \fi
  \fi}
%    \end{macrocode}
% 我在 \lsthelper{Andreas~Matthias}{2000/01/04}{题注受 basicstyle 影响}
% 的缺陷报告之后插入了 |\normalsize|，
% 并在收到 \lsthelper{Sonja~Weidmann}{2000/02/01}{listings 与 caption 宏包
% 不兼容}的另一个报告之后把它移到 |\@makecaption| 之前。
% \end{macro}
%
% \begin{macro}{\lst@makecaption}
% \begin{macro}{\lst@maketitle}
% Axel 提出了第一个定义。另外两个是默认定义。
% 可以调整它们以使 \packagename{listings} 与其他
% 宏包和文档类兼容。
% Markus 建议也定义 |\@captype|，这样他的
% |\raggedlstlistingcaption| 就可以使用了。
%    \begin{macrocode}
\def\lst@makecaption{\def\@captype{lstlisting}\@makecaption}
\def\lst@maketitle{\@makecaption\lst@title@dropdelim}
\def\lst@title@dropdelim#1{\ignorespaces}
%    \end{macrocode}
% 下面的 \packagename{caption}(\packagename{2}) 支持同样来自
% Axel。
%    \begin{macrocode}
\AtBeginDocument{%
\@ifundefined{captionlabelfalse}{}{%
  \def\lst@maketitle{\captionlabelfalse\@makecaption\@empty}}%
\@ifundefined{caption@startrue}{}{%
  \def\lst@maketitle{\caption@startrue\@makecaption\@empty}}%
}
%    \end{macrocode}
% \end{macro}\end{macro}
%
% \begin{macro}{\lst@HRefStepCounter}
% 这个宏把列表编号设为负值，因为用户
% 不应引用这样的列表。如果存在 \packagename{hyperref} 宏包，
% 我们就用`lstlisting.'（上面的参数）与
% 当前（负）编号拼接后交给 hyperref。
% 加入这些分组是为了防止其他宏包（即
% \packagename{tabularx}）读取局部改动的计数器
% 并把它全局写回。感谢 \lsthelper{Michael~Niedermair}
% {2001/09/18}{列表编号异常}的报告。不幸的是
% 这种局部化又导致了另一个缺陷，见 |\theHlstnumber|。
%    \begin{macrocode}
\def\lst@HRefStepCounter#1{%
  \begingroup
  \c@lstlisting\lst@neglisting
  \advance\c@lstlisting\m@ne \xdef\lst@neglisting{\the\c@lstlisting}%
  \ifx\hyper@refstepcounter\@undefined\else
      \hyper@refstepcounter{#1}%
  \fi
  \endgroup
}
\gdef\lst@neglisting{\z@}% init, \z@ can be used both for 0pt and 0,
                         % remember: _neg_ative number, but no
                         % counter!
%    \end{macrocode}
% \end{macro}
%
% \begin{lstkey}{boxpos}
% \begin{macro}{\lst@boxtrue}
% 设置（可能使用的）盒子的垂直对齐，或者指示
% 使用了盒子。
%    \begin{macrocode}
\lst@Key{boxpos}{c}{\def\lst@boxpos{#1}}
%    \end{macrocode}
%    \begin{macrocode}
\def\lst@boxtrue{\let\lst@ifbox\iftrue}
\let\lst@ifbox\iffalse
%    \end{macrocode}
% \end{macro}\end{lstkey}
%
% \begin{lstkey}{float}
% \begin{lstkey}{floatplacement}
% Matthias Zenger 请求双栏浮动体，所以我插入了一些代码。
% 我们首先检查星号\ldots
%    \begin{macrocode}
\lst@Key{float}\relax[\lst@floatplacement]{%
    \lstKV@SwitchCases{#1}%
    {true:\let\lst@floatdefault\lst@floatplacement
          \let\lst@float\lst@floatdefault\\%
     false:\let\lst@floatdefault\relax
           \let\lst@float\lst@floatdefault
    }{\def\lst@next{\@ifstar{\let\lst@beginfloat\@dblfloat
                             \let\lst@endfloat\end@dblfloat
                             \lst@KFloat}%
                            {\let\lst@beginfloat\@float
                             \let\lst@endfloat\end@float
                             \lst@KFloat}}
      \edef\lst@float{#1}%
      \expandafter\lst@next\lst@float\relax}}
%    \end{macrocode}
% \ldots\ 并定义 |\lst@float|。
%    \begin{macrocode}
\def\lst@KFloat#1\relax{%
    \ifx\@empty#1\@empty
        \let\lst@float\lst@floatplacement
    \else
        \def\lst@float{#1}%
    \fi}
%    \end{macrocode}
% 设置 |\lst@AddToHook{PreSet}{\let\lst@float\relax}| 被
% 应 \lsthelper{Tanguy~Fautr\'e}{2004/02/02}{listings
% 不遵循 float 指令？}的要求而修改。这也导致上面的一些调整。
%    \begin{macrocode}
\lst@Key{floatplacement}{tbp}{\def\lst@floatplacement{#1}}
\lst@AddToHook{PreSet}{\let\lst@float\lst@floatdefault}
\lst@AddToHook{TextStyle}{\let\lst@float\relax}
\let\lst@floatdefault\relax % init
%    \end{macrocode}
% |\lst@doendpe| 根据 |\lst@float| 设置——感谢
% \lsthelper{Andreas~Schmidt}{2004/05/15}{浮动列表
% 跟在 \section 之后时间距错误}和 \lsthelper{Heiko~Oberdiek}{2004/05/18}{同上}。
%    \begin{macrocode}
\lst@AddToHook{DeInit}{%
    \ifx\lst@float\relax
        \global\let\lst@doendpe\@doendpe
    \else
        \global\let\lst@doendpe\@empty
    \fi}
%    \end{macrocode}
% 浮动类型 |\ftype@lstlisting| 根据是否加载
% \packagename{float} 宏包以及是否定义了 \texttt{figure} 和
% \texttt{table} 浮动体来设置。这在 |\begin{document}| 处完成，
% 以使代码与宏包加载顺序无关。
%    \begin{macrocode}
\AtBeginDocument{%
\@ifundefined{c@float@type}%
    {\edef\ftype@lstlisting{\ifx\c@figure\@undefined 1\else 4\fi}}
    {\edef\ftype@lstlisting{\the\c@float@type}%
     \addtocounter{float@type}{\value{float@type}}}%
}
%    \end{macrocode}
% \end{lstkey}
% \end{lstkey}
%
%
% \subsection{初始化与行尾}
%
% \begin{lstkey}{aboveskip}
% \begin{lstkey}{belowskip}
% 我们定义并初始化这些键，并为`行内'
% 列表防止多余间距（特别是当 \packagename{fancyvrb} 接口活动时，
% 问题由 \lsthelper{Denis~Girou}{1999/08/03}{间距错误}报告）。
%    \begin{macrocode}
\lst@Key{aboveskip}\medskipamount{\def\lst@aboveskip{#1}}
\lst@Key{belowskip}\medskipamount{\def\lst@belowskip{#1}}
\lst@AddToHook{TextStyle}
    {\let\lst@aboveskip\z@ \let\lst@belowskip\z@}
%    \end{macrocode}
% \end{lstkey}\end{lstkey}
%
% \begin{lstkey}{everydisplay}
% \begin{macro}{\lst@ifdisplaystyle}
% 有些东西取决于显示样式列表。
%    \begin{macrocode}
\lst@Key{everydisplay}{}{\def\lst@EveryDisplay{#1}}
\lst@AddToHook{TextStyle}{\let\lst@ifdisplaystyle\iffalse}
\lst@AddToHook{DisplayStyle}{\let\lst@ifdisplaystyle\iftrue}
\let\lst@ifdisplaystyle\iffalse
%    \end{macrocode}
% \end{macro}
% \end{lstkey}
%
% \begin{macro}{\lst@Init}
% 如有要求，开始浮动体或多栏环境。
%    \begin{macrocode}
\def\lst@Init#1{%
    \begingroup
    \ifx\lst@float\relax\else
        \edef\@tempa{\noexpand\lst@beginfloat{lstlisting}[\lst@float]}%
        \expandafter\@tempa
    \fi
    \ifx\lst@multicols\@empty\else
        \edef\lst@next{\noexpand\multicols{\lst@multicols}}
        \expandafter\lst@next
    \fi
%    \end{macrocode}
% 在受限的水平 \TeX\ 模式下我们切换到 |\lst@boxtrue|。
% 这种情况下我们在列表周围制作合适的盒子。
%    \begin{macrocode}
    \ifhmode\ifinner \lst@boxtrue \fi\fi
    \lst@ifbox
        \lsthk@BoxUnsafe
        \hbox to\z@\bgroup
             $\if t\lst@boxpos \vtop
        \else \if b\lst@boxpos \vbox
        \else \vcenter \fi\fi
        \bgroup \par\noindent
    \else
        \lst@ifdisplaystyle
            \lst@EveryDisplay
            \par\penalty-50\relax
            \vspace\lst@aboveskip
        \fi
    \fi
%    \end{macrocode} ^^A$ to make emacs font locking happy
% 把 |\vspace| 移到 |\par| 之后——否则列表顶部会出现空行。
% 缺陷由 \lsthelper{Jim~Hefferon}{1999/08/27}{使用
% |\lstinputlisting| 的列表前有空行}报告。
%
% 现在制作上题注。
%    \begin{macrocode}
    \normalbaselines
    \abovecaptionskip\lst@abovecaption\relax
    \belowcaptionskip\lst@belowcaption\relax
    \lst@MakeCaption t%
%    \end{macrocode}
% 一些初始化。
% 我在 \lsthelper{Jim~Hefferon}{1999/08/23}{lstlisting 之后的垂直间距错误}
% 的缺陷报告之后移除了 |\par\nointerlineskip| |\normalbaselines|。
% 他报告了与 Aidan Philip Heerdegen 相同的问题（见下），但我
%    \begin{macrocode}
    \lsthk@PreInit \lsthk@Init
    \lst@ifdisplaystyle
        \global\let\lst@ltxlabel\@empty
        \if@inlabel
            \lst@ifresetmargins
                \leavevmode
            \else
                \xdef\lst@ltxlabel{\the\everypar}%
                \lst@AddTo\lst@ltxlabel{%
                    \global\let\lst@ltxlabel\@empty
                    \everypar{\lsthk@EveryLine\lsthk@EveryPar}}%
            \fi
        \fi
        \everypar\expandafter{\lst@ltxlabel
                              \lsthk@EveryLine\lsthk@EveryPar}%
    \else
        \everypar{}\let\lst@NewLine\@empty
    \fi
    \lsthk@InitVars \lsthk@InitVarsBOL
%    \end{macrocode}
% 行尾字符 |chr(13)|=|^^M| 控制处理过程，见下面
% |\lst@MProcessListing| 的定义。
% 参数 |#1| 是 |\relax| 或 |\lstenv@backslash|。
%    \begin{macrocode}
    \lst@Let{13}\lst@MProcessListing
    \let\lst@Backslash#1%
    \lst@EnterMode{\lst@Pmode}{\lst@SelectCharTable}%
    \lst@InitFinalize}
%    \end{macrocode}
% 注意：从 0.19 版起，`列表处理'被实现为一个内部
% 模式，即具有特殊字符表的模式。因为有个缺陷报告
% \lsthelper{Fermin~Reig}{2002/09/04}{bad top frame inside figure+centering}
% |\rightskip| 及其他通过 \hookname{PreInit} 而不是通过
% \hookname{InitVars}.
%    \begin{macrocode}
\let\lst@InitFinalize\@empty % init
\lst@AddToHook{PreInit}
    {\rightskip\z@ \leftskip\z@ \parfillskip=\z@ plus 1fil
     \let\par\@@par}
\lst@AddToHook{EveryLine}{}% init
\lst@AddToHook{EveryPar}{}% init
%    \end{macrocode}
% \end{macro}
%
% \begin{lstkey}{showlines}
% 让用户控制列表末尾的空行是否
% 被打印。但如果你读过用户指南，这你是知道的。
%    \begin{macrocode}
\lst@Key{showlines}f[t]{\lstKV@SetIf{#1}\lst@ifshowlines}
%    \end{macrocode}
% \end{lstkey}
%
% \begin{macro}{\lst@DeInit}
% 输出剩余的字符并更新所有东西。起初我漏掉了
% 在这里使用 |\lst@ifdisplaystyle|，但后来 \lsthelper{KP~Gores}{2001/07/11}
% {每个 \lstinline 之后的 \csname{par}}报告了问题。
% |\everypar| 在 \lsthelper{Michael~Niedermair}{2002/05/14}{listings.sty 与 caption}
% 的缺陷报告之后被放到 |\lsthk@ExitVars| 后面，而
% 我在 \lsthelper{Georg~Rehm}
% {2002/05/14}{listings.sty 与长题注}的缺陷报告之后加入了 |\normalbaselines|，
% 并在 \lsthelper{Walter~E.~Brown}{2004/03/01}{列表底部的题注
% 继承了 basicstyle 的颜色}的报告之后加入了 |\normalcolor|。
%    \begin{macrocode}
\def\lst@DeInit{%
    \lst@XPrintToken \lst@EOLUpdate
    \global\advance\lst@newlines\m@ne
    \lst@ifshowlines
        \lst@DoNewLines
    \else
        \setbox\@tempboxa\vbox{\lst@DoNewLines}%
    \fi
    \lst@ifdisplaystyle \par\removelastskip \fi
    \lsthk@ExitVars\everypar{}\lsthk@DeInit\normalbaselines\normalcolor
%    \end{macrocode}
% 放置下题注。
%    \begin{macrocode}
    \lst@MakeCaption b%
%    \end{macrocode}
% 必要时关闭盒子，并绘制一条线以取得正确的宽度。
% 我在收到 \lsthelper{Aidan~Philip~Heerdegen}
% {1999/07/23}{垂直间距错误}的缺陷报告之后加入了 |\par\nointerlineskip|（后来又
% 移除了 |\nointerlineskip|）。|\everypar{}| 则源于
% \lsthelper{Sonja~Weidmann}{2000/02/01}{listings 与 caption 宏包
% 不兼容}的缺陷报告。
%    \begin{macrocode}
    \lst@ifbox
        \egroup $\hss \egroup
        \vrule\@width\lst@maxwidth\@height\z@\@depth\z@
    \else
        \lst@ifdisplaystyle
            \par\penalty-50\vspace\lst@belowskip
        \fi
    \fi
%    \end{macrocode}
% 必要时结束多栏环境和/或浮动体。
%    \begin{macrocode}
    \ifx\lst@multicols\@empty\else
        \def\lst@next{\global\let\@checkend\@gobble
                      \endmulticols
                      \global\let\@checkend\lst@@checkend}
        \expandafter\lst@next
    \fi
    \ifx\lst@float\relax\else
        \expandafter\lst@endfloat
    \fi
    \endgroup}
\let\lst@@checkend\@checkend
%    \end{macrocode}
% \end{macro}
%
% \begin{macro}{\lst@maxwidth}
% 将被分配、初始化并更新。
%    \begin{macrocode}
\newdimen\lst@maxwidth % \global
\lst@AddToHook{InitVars}{\global\lst@maxwidth\z@}
\lst@AddToHook{InitVarsEOL}
    {\ifdim\lst@currlwidth>\lst@maxwidth
         \global\lst@maxwidth\lst@currlwidth
     \fi}
%    \end{macrocode}
% \end{macro}
%
% \begin{macro}{\lst@EOLUpdate}
% 你觉得这个宏做什么？
%    \begin{macrocode}
\def\lst@EOLUpdate{\lsthk@EOL \lsthk@InitVarsEOL}
%    \end{macrocode}
% \end{macro}
%
% \begin{macro}{\lst@MProcessListing}
% 这是我们在处理列表时行尾必须做的事。
% 我们输出所有剩余字符并更新变量。
% 如果已到达最后一行，我们就检查是否还有下一个行
% 区间要输入。
%    \begin{macrocode}
\def\lst@MProcessListing{%
  \lst@XPrintToken \lst@EOLUpdate \lsthk@InitVarsBOL
  \global\advance\lst@lineno\@ne
  \ifnum \lst@lineno>\lst@lastline
      \lst@ifdropinput \lst@LeaveMode \fi
      \ifx\lst@linerange\@empty
          \expandafter\expandafter\expandafter\lst@EndProcessListing
      \else
          \lst@interrange
          \lst@GetLineInterval
          \expandafter\expandafter\expandafter\lst@SkipToFirst
      \fi
  \else
      \expandafter\lst@BOLGobble
  \fi}
%    \end{macrocode}
% \end{macro}
%
% \begin{macro}{\lst@EndProcessListing}
% 默认定义是 |\endinput|。
% 这适用于 |\lstinputlisting|。
%    \begin{macrocode}
\let\lst@EndProcessListing\endinput
%    \end{macrocode}
% \end{macro}
%
% \begin{lstkey}{gobble}
% 该键设置每行要吞掉的字符数。
%    \begin{macrocode}
\lst@Key{gobble}{0}{\def\lst@gobble{#1}}
%    \end{macrocode}
% \end{lstkey}
%
% \begin{macro}{\lst@BOLGobble}
% 如果数目为正，我们就设置一个临时计数器并开始循环。
%    \begin{macrocode}
\def\lst@BOLGobble{%
    \ifnum\lst@gobble>\z@
        \@tempcnta\lst@gobble\relax
        \expandafter\lst@BOLGobble@
	\fi}
%    \end{macrocode}
% 非正数终止循环（通过不继续）。
% 注意：这不是 |\lst@BOLGobble| 中刚用到的那个宏。
%    \begin{macrocode}
\def\lst@BOLGobble@@{%
    \ifnum\@tempcnta>\z@
        \expandafter\lst@BOLGobble@
    \fi}
%    \end{macrocode}
% 如果我们吞掉一个反斜杠，就必须查看这个反斜杠是否结束
% 一个环境。至于随后的字符是否等于例如~|end{lstlisting}|，
% 我们或者结束环境，或者把刚吞掉的所有字符插到
% `继续循环'宏之后。
%    \begin{macrocode}
\def\lstenv@BOLGobble@@{%
    \lst@IfNextChars\lstenv@endstring{\lstenv@End}%
    {\advance\@tempcnta\m@ne \expandafter\lst@BOLGobble@@\lst@eaten}}
%    \end{macrocode}
% 现在来了循环：如果读到 |\relax|、EOL 或 FF，下一个操作
% 就是完全相同的记号。注意对 FF（以及下面的制表符）我们测试
% 一个含 |\lst@ProcessFormFeed| 的宏。这是一个缺陷，由
% \lsthelper{Heiko~Oberdiek}{2002/04/16}{Re: first experience ...}.
%    \begin{macrocode}
\def\lst@BOLGobble@#1{%
    \let\lst@next#1%
    \ifx \lst@next\relax\else
    \ifx \lst@next\lst@MProcessListing\else
    \ifx \lst@next\lst@processformfeed\else
%    \end{macrocode}
% 否则我们使用两个子宏之一。
%    \begin{macrocode}
    \ifx \lst@next\lstenv@backslash
        \let\lst@next\lstenv@BOLGobble@@
    \else
        \let\lst@next\lst@BOLGobble@@
%    \end{macrocode}
% 现在我们真正吞掉字符。制表符把临时计数器
% 减去 |\lst@tabsize|（并在必要时处理剩余量），\ldots
%    \begin{macrocode}
        \ifx #1\lst@processtabulator
            \advance\@tempcnta-\lst@tabsize\relax
            \ifnum\@tempcnta<\z@
                \lst@length-\@tempcnta \lst@PreGotoTabStop
            \fi
%    \end{macrocode}
% \ldots\ 而任何其他字符把计数器减一。
%    \begin{macrocode}
        \else
            \advance\@tempcnta\m@ne
        \fi
    \fi \fi \fi \fi
    \lst@next}
%    \end{macrocode}
%    \begin{macrocode}
\def\lst@processformfeed{\lst@ProcessFormFeed}
\def\lst@processtabulator{\lst@ProcessTabulator}
%    \end{macrocode}
% \end{macro}
%
%
% \subsection{列表清单}%
% \label{sec:list-listings}
%
% \begin{lstkey}{name}
% \begin{macro}{\lstname}
% \begin{macro}{\lst@name}
% \begin{macro}{\lst@intname}
% 每个美化排版命令在设置任何键之前都给 |\lst@intname| 赋值。
%    \begin{macrocode}
\lst@Key{name}\relax{\def\lst@intname{#1}}
\lst@AddToHookExe{PreSet}{\global\let\lst@intname\@empty}
\lst@AddToHook{PreInit}{%
    \let\lst@arg\lst@intname \lst@ReplaceIn\lst@arg\lst@filenamerpl
    \global\let\lst@name\lst@arg \global\let\lstname\lst@name}
%    \end{macrocode}
% 使用 |\lst@ReplaceIn| 消除了一个缺陷，它最初由
% \lsthelper{Magne~Rudshaug}{1998/01/09}{下划线与列表清单}。
% 这里是替换列表。
%    \begin{macrocode}
\def\lst@filenamerpl{_\textunderscore $\textdollar -\textendash}
%    \end{macrocode} ^^A $
% \end{macro}
% \end{macro}
% \end{macro}
% \end{lstkey}
%
% \begin{macro}{\l@lstlisting}
% 打印一行`lol'。
%    \begin{macrocode}
\def\l@lstlisting#1#2{\@dottedtocline{1}{1.5em}{2.3em}{#1}{#2}}
%    \end{macrocode}
% \end{macro}
%
% \begin{macro}{\lstlistlistingname}
% 只包含标题名。
%    \begin{macrocode}
\lst@UserCommand\lstlistlistingname{Listings}
%    \end{macrocode}
% \end{macro}
%
% \begin{macro}{\lstlistoflistings}
% 我们做局部调整并调用 |\tableofcontents|。这样，
% 对该宏的重新定义（例如内部不带任何 |\MakeUppercase|）
% 也会对列表清单生效。
%    \begin{macrocode}
\lst@UserCommand\lstlistoflistings{\bgroup
    \let\contentsname\lstlistlistingname
    \let\lst@temp\@starttoc
    \def\@starttoc##1{\lst@temp{\ext@lstlisting}}%
    \tableofcontents \egroup}
%    \end{macrocode}
% 对于 \packagename{KOMA}-script 类，我们按 \packagename{KOMA} 的方式定义它，
% 这要感谢
% \lsthelper{Tino~Langer}{2003/11/01}{koma-script 的 listsleft 选项不影响
% lol}的缺陷报告。\lsthelper{Markus~Kohm}{2006/08/12}{koma-script 支持
% 已损坏}提出了一个大大改进的版本，它也能与
% \packagename{float} 宏包一起工作。下面几行注释来自 Markus。
%
% 如果定义了 |\float@listhead|（例如使用 float），就利用它
%    \begin{macrocode}
\@ifundefined{float@listhead}{}{%
  \renewcommand*{\lstlistoflistings}{%
    \begingroup
%    \end{macrocode}
% 如果存在用于切换的开关，就切换到单栏模式。
%    \begin{macrocode}
      \@ifundefined{@restonecoltrue}{}{%
        \if@twocolumn
          \@restonecoltrue\onecolumn
        \else
          \@restonecolfalse
        \fi
      }%
      \float@listhead{\lstlistlistingname}%
%    \end{macrocode}
% 把 |\parskip| 设为 0pt（应当如此！）、|\parindent| 设为 0pt（更好但不总是
% 需要）、|\parfillskip| 设为 0pt plus 1fil（应当如此！）。
%    \begin{macrocode}
      \parskip\z@\parindent\z@\parfillskip \z@ \@plus 1fil%
      \@starttoc{\ext@lstlisting}%
%    \end{macrocode}
% 切回双栏（见上）。
%    \begin{macrocode}
      \@ifundefined{@restonecoltrue}{}{%
        \if@restonecol\twocolumn\fi
      }%
    \endgroup
  }%
}
%    \end{macrocode}
% 如果加载了 tocbasic 宏包就利用它：
%    \begin{macrocode}
\AtBeginDocument{%
  \@ifpackageloaded{tocbasic}{%
    \addtotoclist[float]{\ext@lstlisting}%
    \renewcommand*{\lstlistoflistings}{\listoftoc[{\lstlistlistingname}]{lol}}%
    \DeclareTOCStyleEntry[level=1,numwidth=2.3em,indent=1.5em]{default}{lstlisting}%
  }{}%
}
%    \end{macrocode}
% \end{macro}
%
% \begin{macro}{\float@addtolists}
% \packagename{float} 宏包定义了一种通用方式，让宏包向除
% 标准 figure 和 table 列表之外的所有浮动体列表添加
% 内容（例如章名）。每个定义浮动体列表的
% 宏包都向 |\float@addtolists| 添加一条命令，然后
% 希望向所有浮动体列表添加内容的宏包（例如
% KOMA-script 文档类）就可以使用它，而不必知道可能存在的
% 所有列表。感谢 \lsthelper{Markus~Kohm}{2007/02/25}{-}
% 的建议。
%
% 不幸的是，\packagename{float} 用 |\newcommand| 定义它；因此
% 为避免冲突，我们必须在其加载之后重新定义它。
% |\AtBeginDocument| 是做到这一点的最简单方式。再次感谢
% Markus 的建议。
%    \begin{macrocode}
\AtBeginDocument{%
  \@ifundefined{KOMAClassName}{%
    \@ifundefined{float@addtolists}{%
      \gdef\float@addtolists#1{\addtocontents{\ext@lstlisting}{#1}}%
    }{%
      \let\orig@float@addtolists\float@addtolists
      \gdef\float@addtolists#1{%
        \addtocontents{\ext@lstlisting}{#1}%
        \orig@float@addtolists{#1}}}%
  }{}%
}%
%    \end{macrocode}
% \end{macro}
%
%
% \subsection{行内列表}\label{iInlineListings}
%
% \subsubsection{处理行内列表}\label{uProcessingInline}
%
% \begin{macro}{\lstinline}
%         除 |\lsthk@PreSet| 之外，我们使用 |boxpos=b| 和
%         flexiblecolumns。我在 \lsthelper{Michael~Weber}
%         {1999/12/16}{列表环境中间距错误}的缺陷报告之后插入了 |\leavevmode|。
%         \lsthelper{Olivier~Lecarme}{2001/07/30}{在题注内使用 \lstinline 时
%         `break'行为不一致}报告了一个问题，
%         它在移除 |\let| |\lst@newlines| |\@empty| 之后消失（现在
%         |\lst@newlines| 是一个计数器！）。不幸的是我不知道
%         早些时候插入这段代码的原因！在宏末尾我们
%         检查定界符。\footnote{这是原作者
%         Carsten Heinz 的文字。}
%
%         然后是一个允许花括号的实验版本，但
%         \lsthelper{Luc~Van~Eycken}{2014/01/22}{\rcmdname\lstinline{} 在
%         对齐环境中}报告说，用花括号代替字符围绕源代码
%         的 |\lstinline| 实验性实现，在 tabular 环境中使用时
%         会导致错误。
%
%         他发现该错误来自 master counter（参见 TeXbook 附录
%         D（Dirty Tricks）第 5 项（Brace hacks）
%         （第 385-386 页））。在测试下一个字符之后（下面被注释掉的
%         片段中的第 7 行）加入下面这一行
%\begin{verbatim}
%\ifnum`{=0}\fi%
%\end{verbatim}
%         可修正错误行为。但 \lsthelper{Qing~Lee}{2014/06/28}{-}
%         指出这会破坏第~\ref{uListingsArguments} 节第~\pageref{uListingsArguments} 页
%         所示的那类代码，并
%         提出了另一种解决方案，而这又破坏了 Luc 提供的代码：
%\begin{verbatim}
% \renewcommand\lstinline[1][]{%
%   \leavevmode\bgroup % \hbox\bgroup --> \bgroup
%   \def\lst@boxpos{b}%
%   \lsthk@PreSet\lstset{flexiblecolumns,#1}%
%   \lsthk@TextStyle
%   \ifnum\iffalse{\fi`}=\z@\fi
%   \@ifnextchar\bgroup{%
%     \ifnum`{=\z@}\fi%
%     \afterassignment\lst@InlineG \let\@let@token}{%
%     \ifnum`{=\z@}\fi\lstinline@}%
%}
%\end{verbatim}
%         所以最后又回到了旧代码，那些需要带花括号的
%         |\lstinline| 的人应当使用第~\ref{uListingsArguments} 节
%         第~\pageref{uListingsArguments} 页的变通办法。
%
%         这个长期存在的缺陷现在由
%     \lsthelper{user5028841}{2024/05/05}{关于 listings 宏包}
%         修复，他提供了一个解决方案：用特殊字符作为分组的
%         开始和结束：
%    \begin{macrocode}
\edef\lst@temp{\the\catcode`\^^@}
\catcode`\^^@=1
\newcommand\lstinline[1][]{%
  \leavevmode\bgroup % \hbox\bgroup --> \bgroup
  \def\lst@boxpos{b}%
  \lsthk@PreSet\lstset{flexiblecolumns,#1}%
  \lsthk@TextStyle
  \@ifnextchar\bgroup{%
    \afterassignment\lst@InlineG \romannumeral`^^@\iffalse}\fi
  \let\@let@token}%
  \lstinline@}
%    \end{macrocode}
% 这里我们恢复 |\^^@| 先前的 catcode：
%    \begin{macrocode}
\catcode`\^^@=\lst@temp
%    \end{macrocode}
% \end{macro}
% \begin{macro}{\lst@inline@}
%         这是处理 |\lstinline| 可选参数之后的
%         标准方法。
%    \begin{macrocode}
\def\lstinline@#1{%
  \lst@Init\relax
  \lst@IfNextCharActive{\lst@InlineM#1}{\lst@InlineJ#1}}
%    \end{macrocode}
%    \begin{macrocode}
\lst@AddToHook{TextStyle}{}% init
\lst@AddToHook{SelectCharTable}{\lst@inlinechars}
\global\let\lst@inlinechars\@empty
%    \end{macrocode}
% \end{macro}
%
% \begin{macro}{\lst@InlineM}
% \begin{macro}{\lst@InlineJ}
%         处理`普通'行内与参数内行内列表这两种
%         情况。第一种情况下给定字符结束行内
%         列表，这样的列表内的 EOL 立即结束它，并
%         产生错误消息。
%    \begin{macrocode}
\def\lst@InlineM#1{%
  \gdef\lst@inlinechars{%
    \lst@Def{`#1}{\lst@DeInit\egroup\global\let\lst@inlinechars\@empty}%
    \lst@Def{13}{\lst@DeInit\egroup \global\let\lst@inlinechars\@empty
      \PackageError{Listings}{lstinline ended by EOL}\@ehc}}%
  \lst@inlinechars}
%    \end{macrocode}
% 另一种情况下我们取得直到 |#1| 的所有字符，把这些字符
% 变成活动的，执行（排印）它们并结束列表（全部通过临时宏）。
% 这就是全部内容。
%    \begin{macrocode}
\def\lst@InlineJ#1{%
  \def\lst@temp##1#1{%
    \let\lst@arg\@empty \lst@InsideConvert{##1}\lst@arg
    \lst@DeInit\egroup}%
  \lst@temp}
%    \end{macrocode}
% \end{macro}
% \end{macro}
%
% \begin{macro}{\lst@InlineG}
%     是实验性的。
%    \begin{macrocode}
\def\lst@InlineG{%
  \lst@Init\relax
  \lst@IfNextCharActive{\lst@InlineM\}}{%
    \let\lst@arg\@empty \lst@InlineGJ}}
%    \end{macrocode}
%     这是关闭分组的时点：
%    \begin{macrocode}
\edef\lst@temp{\the\catcode`\^^@}
\catcode`\^^@=2
\def\lst@InlineGJ{\futurelet\@let@token\lst@InlineGJTest}
\def\lst@InlineGJTest{%
  \ifx\@let@token\egroup
    \iffalse{\fi\romannumeral`^^@
      \afterassignment\lst@InlineGJEnd
      \expandafter\let\expandafter\@let@token
  \else
    \ifx\@let@token\@sptoken
      \let\lst@next\lst@InlineGJReadSp
    \else
      \let\lst@next\lst@InlineGJRead
    \fi
    \expandafter\lst@next
    \fi}
%    \end{macrocode}
%     同前：恢复 |\^^@| 先前的 catcode：
%    \begin{macrocode}
\catcode`\^^@=\lst@temp
\def\lst@InlineGJEnd{\lst@arg\lst@DeInit\egroup}
\def\lst@InlineGJRead#1{%
  \lccode`\~=`#1\lowercase{\lst@lAddTo\lst@arg~}%
  \lst@InlineGJ}
\def\lst@InlineGJReadSp#1{%
  \lccode`\~=`\ \lowercase{\lst@lAddTo\lst@arg~}%
  \lst@InlineGJ#1}
%    \end{macrocode}
% \end{macro}
%
%
% \subsubsection{短行内列表环境}
%
% 本节实现基于 \packagename{shortvrb} 宏包，它属于标准 \LaTeX\ 文档
% 宏包中 |doc.dtx| 的 2006/02/02 v2.1d 版。其中部分内容
% 的版权因此属于 The \LaTeX3 Project 1993--2006 以及 Frank
% Mittelbach 1989--1999。\lsthelper{Denis~Bitouz\'e}{2020/03/21}
% {错误消息中的错字}利用疫情危机查看了
% 错误消息，并发现了一些错字。
%
% \begin{macro}{\lstMakeShortInline}
% \begin{macro}{\lstMakeShortInline@}
% 首先，如果省略了可选参数，我们就提供一个。
%    \begin{macrocode}
\newcommand\lstMakeShortInline[1][]{%
  \def\lst@shortinlinedef{\lstinline[#1]}%
  \lstMakeShortInline@}%
\def\lstMakeShortInline@#1{%
  \expandafter\ifx\csname lst@ShortInlineOldCatcode\string#1\endcsname\relax
    \lst@shortlstinlineinfo{Made }{#1}%
    \lst@add@special{#1}%
%    \end{macrocode}
% 该字符的当前 catcode 被存到
% |\lst@ShortInlineOldCatcode\|\meta{c}.
%    \begin{macrocode}
    \expandafter
    \xdef\csname lst@ShortInlineOldCatcode\string#1\endcsname{%
      \the\catcode`#1}%
%    \end{macrocode}
% 该字符用与 |\verb|（例如）中相同的技巧拼进定义，
% 做法是在一个分组内激活 |~|。
%    \begin{macrocode}
    \begingroup
      \catcode`\~\active  \lccode`\~`#1%
      \lowercase{%
%    \end{macrocode}
% 该字符的旧含义被记录在
% |\lst@ShortInlineOldMeaning\|\meta{c} 中，然后才赋予新含义。
%    \begin{macrocode}
        \global\expandafter\let
          \csname lst@ShortInlineOldMeaning\string#1\endcsname~%
          \expandafter\gdef\expandafter~\expandafter{%
            \lst@shortinlinedef#1}}%
    \endgroup
%    \end{macrocode}
% 最后该字符被变成活动的。
%    \begin{macrocode}
    \global\catcode`#1\active
%    \end{macrocode}
% 如果我们怀疑 \meta{c} 已经是一个短引用，我们就提醒
% 用户。现在如果出了什么问题，他/她要自己负责\ldots
% （\packagename{listings} 中的改动：我们在此处给出恰当的错误。）
%    \begin{macrocode}
  \else
    \PackageError{Listings}%
    {\string\lstMakeShorterInline\ definitions cannot be nested}%
    {Use \string\lstDeleteShortInline first.}%
    {}%
  \fi}
%    \end{macrocode}
% \end{macro}
% \end{macro}
% \begin{macro}{\lstDeleteShortInline}
%    \begin{macrocode}
\def\lstDeleteShortInline#1{%
  \expandafter\ifx%
      \csname lst@ShortInlineOldCatcode\string#1\endcsname\relax%
    \PackageError{Listings}%
    {#1 is not a short reference for \string\lstinline}%
    {Use \string\lstMakeShortInline first.}%
    {}%
  \else
    \lst@shortlstinlineinfo{Deleted }{#1 as}%
    \lst@rem@special{#1}%
    \global\catcode`#1\csname lst@ShortInlineOldCatcode\string#1\endcsname
    \global \expandafter\let%
      \csname lst@ShortInlineOldCatcode\string#1\endcsname \relax
    \ifnum\catcode`#1=\active
      \begingroup
        \catcode`\~\active  \lccode`\~`#1%
        \lowercase{%
          \global\expandafter\let\expandafter~%
          \csname lst@ShortInlineOldMeaning\string#1\endcsname}%
      \endgroup
    \fi
  \fi}
%    \end{macrocode}
% \end{macro}
%
% \begin{macro}{\lst@shortlstinlineinfo}
%    \begin{macrocode}
\def\lst@shortlstinlineinfo#1#2{%
     \PackageInfo{Listings}{%
       #1\string#2 a short reference for \string\lstinline}}
%    \end{macrocode}
%  \end{macro}
%
% \begin{macro}{\lst@add@special}
% 这个辅助宏把它的参数加到 |\dospecials| 宏中，
% 后者按惯例被 verbatim 宏用来
% 改变当前活动字符的 catcode。我们需要
% 在 |\dospecials| 的展开中加入 |\do\|\meta{c}，且要先
% 把可能已存在的该字符移除，以避免在 |\lstMakeShortInline| 未被
% |\lstDeleteShortInline| 平衡时堆积多份副本
% （以防任何使用 |\dospecials| 的东西在意重复）。
%    \begin{macrocode}
\def\lst@add@special#1{%
  \lst@rem@special{#1}%
  \expandafter\gdef\expandafter\dospecials\expandafter
    {\dospecials \do #1}%
%    \end{macrocode}
% 同样我们必须把 |\@makeother\|\meta{c} 加到 |\@sanitize|
% （后者用于像"\index"这样的地方，对所有特殊
% 字符（花括号除外）重新设置 catcode）。
%    \begin{macrocode}
  \expandafter\gdef\expandafter\@sanitize\expandafter
    {\@sanitize \@makeother #1}}
%    \end{macrocode}
% \end{macro}
% \begin{macro}{\lst@rem@special}
% |\lst@add@special| 的反操作稍微更棘手。|\do| 被
% 重新定义为：如果其参数是所关心的字符则展开为空，
% 否则只展开为该参数。然后我们
% 重新定义 |\dospecials| 为其自身的展开。|=`##1| 之后的空格
% 可防止展开成 |\relax|！
%    \begin{macrocode}
\def\lst@rem@special#1{%
  \def\do##1{%
    \ifnum`#1=`##1 \else \noexpand\do\noexpand##1\fi}%
  \xdef\dospecials{\dospecials}%
%    \end{macrocode}
% 修正 |\@sanitize| 相同，只是我们需要重新定义
% |\@makeother|，这显然必须在分组内完成。
%    \begin{macrocode}
  \begingroup
    \def\@makeother##1{%
      \ifnum`#1=`##1 \else \noexpand\@makeother\noexpand##1\fi}%
    \xdef\@sanitize{\@sanitize}%
  \endgroup}
%    \end{macrocode}
% \end{macro}
%
%
% \subsection{input 命令}\label{iTheInputCommand}
%
% \begin{macro}{\lst@MakePath}
% \begin{lstkey}{inputpath}
% 如有必要，该宏在路径末尾追加一个斜杠。
%    \begin{macrocode}
\def\lst@MakePath#1{\ifx\@empty#1\@empty\else\lst@MakePath@#1/\@nil/\fi}
\def\lst@MakePath@#1/{#1/\lst@MakePath@@}
\def\lst@MakePath@@#1/{%
    \ifx\@nil#1\expandafter\@gobble
         \else \ifx\@empty#1\else #1/\fi \fi
    \lst@MakePath@@}
%    \end{macrocode}
% 现在我们可以清空路径，或使用 |\lst@MakePath|。
%    \begin{macrocode}
\lst@Key{inputpath}{}{\edef\lst@inputpath{\lst@MakePath{#1}}}
%    \end{macrocode}
% \end{lstkey}
% \end{macro}
%
% \begin{macro}{\lstinputlisting}
% 输入该列表，或向用户询问新的文件名。
%    \begin{macrocode}
\def\lstinputlisting{%
    \begingroup \lst@setcatcodes \lst@inputlisting}
\newcommand\lst@inputlisting[2][]{%
  \endgroup
%    \end{macrocode}
% |\lst@set| 取局部选项，特别是用于格式化输入文件的
% |inputpath=<path>|。所以必须更新 |\lstset| 以得到
% 正确的 |\lst@inputpath| 值。整个流程必须置于分组内以
% 使改动成为局部的。
%    \begin{macrocode}
  \bgroup\def\lst@set{#1}%
  \expandafter\lstset\expandafter{\lst@set}%
  \IfFileExists{\lst@inputpath#2}%
    {\lst@InputListing{\lst@inputpath#2}}%
    {\filename@parse{\lst@inputpath#2}%
      \edef\reserved@a{\noexpand\lst@MissingFileError
        {\filename@area\filename@base}%
        {\ifx\filename@ext\relax tex\else\filename@ext\fi}}%
      \reserved@a
%    \end{macrocode}
% 在不存在文件的（错误）情况下，我们必须为 |\lst@doendpe| 提供
% 有效值。
%    \begin{macrocode}
  \let\lst@doendpe\@empty}%
  \egroup
  \lst@doendpe \@newlistfalse \ignorespaces%
}
%    \end{macrocode}
% 我们使用 |\lst@doendpe| 去掉下一行开头的缩进
% ——除非 |\lstinputlisting| 之后有空行。缺陷
% 由 \lsthelper{David~John~Evans}{1999/06/08}{列表之后的
% 缩进}报告，\lsthelper{David~Carlisle}{1999/06/08}{LaTeX`显示
% 环境'代码}为我指出了解决方案。
% \end{macro}
%
% \begin{macro}{\lst@MissingFileError}
% 是 \LaTeX 的 |\@missingfileerror| 的派生版本。括号是
% 在 \lsthelper{Heiko~Oberdiek}{2003/01/14}{文件 `Makefile.tex' 未
% 找到}报告 TEX-D-L 上讨论的一个问题之后加入的。
%    \begin{macrocode}
\def\lst@MissingFileError#1#2{%
    \typeout{^^J! Package Listings Error: File `#1(.#2)' not found.^^J%
        ^^JType X to quit or <RETURN> to proceed,^^J%
        or enter new name. (Default extension: #2)^^J}%
    \message{Enter file name: }%
    {\endlinechar\m@ne \global\read\m@ne to\@gtempa}%
%    \end{macrocode}
% 输入 |x| 或 |X| 即退出。
%    \begin{macrocode}
    \ifx\@gtempa\@empty \else
        \def\reserved@a{x}\ifx\reserved@a\@gtempa\batchmode\@@end\fi
        \def\reserved@a{X}\ifx\reserved@a\@gtempa\batchmode\@@end\fi
%    \end{macrocode}
% 在其他所有情况下我们尝试新的文件名。
%    \begin{macrocode}
        \filename@parse\@gtempa
        \edef\filename@ext{%
            \ifx\filename@ext\relax#2\else\filename@ext\fi}%
        \edef\reserved@a{\noexpand\IfFileExists %
                {\filename@area\filename@base.\filename@ext}%
            {\noexpand\lst@InputListing %
                {\filename@area\filename@base.\filename@ext}}%
            {\noexpand\lst@MissingFileError
                {\filename@area\filename@base}{\filename@ext}}}%
        \expandafter\reserved@a %
    \fi}
%    \end{macrocode}
% \end{macro}
%
% \begin{macro}{\lst@ifdraft}
% 利用 |\lst@ifprint|。\lsthelper{Enrico~Straube}{2002/02/12}
% {de.comp.text.tex：listings 与 draft 模式}请求了最后一个选项。
%    \begin{macrocode}
\let\lst@ifdraft\iffalse
\DeclareOption{draft}{\let\lst@ifdraft\iftrue}
\DeclareOption{final}{\let\lst@ifdraft\iffalse}
\lst@AddToHook{PreSet}
    {\lst@ifdraft
         \let\lst@ifprint\iffalse
         \@gobbletwo\fi\fi
     \fi}
%    \end{macrocode}
% \end{macro}
%
% \begin{macro}{\lst@InputListing}
% 唯一的参数是文件名，但我们还有`隐式'
% 参数 |\lst@set|。注意 |\lst@Init| 以 |\relax| 为参数。
%    \begin{macrocode}
\def\lst@InputListing#1{%
    \begingroup
      \lsthk@PreSet \gdef\lst@intname{#1}%
      \expandafter\lstset\expandafter{\lst@set}%
      \lsthk@DisplayStyle
      \catcode\active=\active
      \lst@Init\relax \let\lst@gobble\z@
      \lst@SkipToFirst
      \lst@ifprint \def\lst@next{\input{#1}}%
             \else \let\lst@next\@empty \fi
      \lst@next
      \lst@DeInit
    \endgroup}
%    \end{macrocode}
% 行 |\catcode\active=\active| 把 CR 字符变成活动的，
% 这是在 \lsthelper{Rene~H.~Larsen}{2002/04/15}{lstinputlisting 与 texcl 冲突}
% 的缺陷报告之后添加的。
% \end{macro}
%
% \begin{macro}{\lst@SkipToFirst}
% 行尾字符或者处理列表，或者负责
% 丢弃直到第一个打印行的各行。
%    \begin{macrocode}
\def\lst@SkipToFirst{%
    \ifnum \lst@lineno<\lst@firstline
%    \end{macrocode}
% 我们丢弃输入并重新定义行尾字符。
%    \begin{macrocode}
        \lst@BeginDropInput\lst@Pmode
        \lst@Let{13}\lst@MSkipToFirst
        \lst@Let{10}\lst@MSkipToFirst
    \else
        \expandafter\lst@BOLGobble
    \fi}
%    \end{macrocode}
% \end{macro}
%
% \begin{macro}{\lst@MSkipToFirst}
% 我们只查看是要丢弃更多行，还是要离开该模式（该模式恢复
% |chr(13)| 和 |chr(10)| 的定义）。
%    \begin{macrocode}
\def\lst@MSkipToFirst{%
    \global\advance\lst@lineno\@ne
    \ifnum \lst@lineno=\lst@firstline
        \lst@LeaveMode \global\lst@newlines\z@
        \lsthk@InitVarsBOL
        \lst@DisplayConsecutiveNumbersOrNot
        \expandafter\lst@BOLGobble
    \fi}
%    \end{macrocode}
% \end{macro}
%
%
% \subsection{环境}
%
%
% \subsubsection{底层处理}
%
% \begin{macro}{\lstenv@DroppedWarning}
% 如果有字符被丢弃就发出警告。
%    \begin{macrocode}
\def\lstenv@DroppedWarning{%
    \ifx\lst@dropped\@undefined\else
        \PackageWarning{Listings}{Text dropped after begin of listing}%
    \fi}
\let\lst@dropped\@undefined % init
%    \end{macrocode}
% \end{macro}
%
% \begin{macro}{\lstenv@Process}
% 我们根据找到的是活动 EOL 还是非活动的 |^^J|，
% 执行`|\lstenv@ProcessM|'或 |\lstenv@ProcessJ|。
%    \begin{macrocode}
\begingroup \lccode`\~=`\^^M\lowercase{%
\gdef\lstenv@Process#1{%
    \ifx~#1%
%    \end{macrocode}
% 我们不额外定义 |\lstenv@ProcessM|，因为如果找到活动的 EOL，
% 就根本无事可做。
%    \begin{macrocode}
        \lstenv@DroppedWarning \let\lst@next\lst@SkipToFirst
    \else\ifx^^J#1%
        \lstenv@DroppedWarning \let\lst@next\lstenv@ProcessJ
    \else
        \let\lst@dropped#1\let\lst@next\lstenv@Process
    \fi \fi
    \lst@next}
}\endgroup
%    \end{macrocode}
% \end{macro}
%
% \begin{macro}{\lstenv@ProcessJ}
% 现在来了糟糕的情形：参数内的列表。我们在第~\ref{iApplicationsTo} 节
% 已为此做了工作。这里我们必须取得
% 直到`结束环境'的所有字符。我们区分`命令
% 形式'和`真正的环境'两种情况。
%    \begin{macrocode}
\def\lstenv@ProcessJ{%
    \let\lst@arg\@empty
    \ifx\@currenvir\lstenv@name
        \expandafter\lstenv@ProcessJEnv
    \else
%    \end{macrocode}
% 第一种情况很简单：代码由
% |\end|\meta{name of environment} 终止。于是我们展开该控制序列，
% 然后定义一个临时宏，由它取得列表并完成
% 其余工作。回到 |\lstenv@ProcessJ| 的定义，我们在展开 |\fi| 之后
% 调用该临时宏。
%    \begin{macrocode}
        \expandafter\def\expandafter\lst@temp\expandafter##1%
            \csname end\lstenv@name\endcsname
                {\lst@InsideConvert{##1}\lstenv@ProcessJ@}%
        \expandafter\lst@temp
    \fi}
%    \end{macrocode}
% 我们必须把一个活动的反斜杠和`结束字符串'追加到 |\lst@arg|。这样所有
% （事实上是大多数）其他处理都不会注意到代码曾经在
% 参数内。但现在 EOL 字符是 |chr(10)|=|^^J|，而不是 |chr(13)|。
%    \begin{macrocode}
\begingroup \lccode`\~=`\\\lowercase{%
\gdef\lstenv@ProcessJ@{%
    \lst@lExtend\lst@arg
        {\expandafter\ \expandafter~\lstenv@endstring}%
    \catcode10=\active \lst@Let{10}\lst@MProcessListing
%    \end{macrocode}
% 我们执行 |\lst@arg| 以排印该列表。
%    \begin{macrocode}
    \lst@SkipToFirst \lst@arg}
}\endgroup
%    \end{macrocode}
% \end{macro}
%
% \begin{macro}{\lstenv@ProcessJEnv}
% `真正的环境'情况更复杂。我们取得直到
% 一个 |\end| 及其后参数的所有字符。如果它等于 |\lstenv@name|，我们就
% 找到了环境末尾并开始排印。
%    \begin{macrocode}
\def\lstenv@ProcessJEnv#1\end#2{\def\lst@temp{#2}%
    \ifx\lstenv@name\lst@temp
        \lst@InsideConvert{#1}%
        \expandafter\lstenv@ProcessJ@
    \else
%    \end{macrocode}
% 否则我们把这些字符（包括被吞掉的 |\end| 和被吞掉的
% 参数）追加到当前 |\lst@arg|。然后我们再次查找环境末尾。
%    \begin{macrocode}
        \lst@InsideConvert{#1\\end\{#2\}}%
        \expandafter\lstenv@ProcessJEnv
    \fi}
%    \end{macrocode}
% \end{macro}
%
% \begin{macro}{\lstenv@backslash}
% 遇到反斜杠时，我们或者结束列表，或者处理一个反斜杠
% 并重新插入被吞掉的字符。
%    \begin{macrocode}
\def\lstenv@backslash{%
    \lst@IfNextChars\lstenv@endstring
        {\lstenv@End}%
        {\expandafter\lsts@backslash \lst@eaten}}%
%    \end{macrocode}
% \end{macro}
%
% \begin{macro}{\lstenv@End}
% 这个宏刚被使用过，并终止一个列表环境：
% 我们用 |\end| 或作为命令调用`结束环境'宏。
%    \begin{macrocode}
\def\lstenv@End{%
    \ifx\@currenvir\lstenv@name
        \edef\lst@next{\noexpand\end{\lstenv@name}}%
    \else
        \def\lst@next{\csname end\lstenv@name\endcsname}%
    \fi
    \lst@next}
%    \end{macrocode}
% \end{macro}
%
%
% \subsubsection{定义新环境}
%
% \begin{macro}{\lstnewenvironment}
% 现在来了主命令。我们只定义未定义的环境。关于
% 参数文本 |#1#2#|（特别是最后一个井号），见`The \TeX book'
% 第 204 页例 20.5 之后的段落。
%    \begin{macrocode}
\lst@UserCommand\lstnewenvironment#1#2#{%
    \@ifundefined{#1}%
        {\let\lst@arg\@empty
         \lst@XConvert{#1}\@nil
         \expandafter\lstnewenvironment@\lst@arg{#1}{#2}}%
        {\PackageError{Listings}{Environment `#1' already defined}\@eha
         \@gobbletwo}}
\def\@tempa#1#2#3{%
  \gdef\lstnewenvironment@##1##2##3##4##5{%
    \begingroup
%    \end{macrocode}
% 孤立的`结束环境'会产生错误。
%    \begin{macrocode}
    \global\@namedef{end##2}{\lstenv@Error{##2}}%
%    \end{macrocode}
% `主'环境宏为后续使用定义环境名，并
% 调用一个子宏来取得所有参数。我们打开一个分组并把 EOL 变成活动的。
% 这确保 |\@ifnextchar[| 不会读到列表的字符——它读到的是
% 活动的 EOL。
%    \begin{macrocode}
    \global\@namedef{##2}{\def\lstenv@name{##2}%
        \begingroup \lst@setcatcodes \catcode\active=\active
        \csname##2@\endcsname}%
%    \end{macrocode}
% 该子宏通过 |\new@command| 定义。我们误用 |\l@ngrel@x| 使
% 该定义成为 |\global|，并改进 \LaTeX 的 |\@xargdef|。
%    \begin{macrocode}
    \let\l@ngrel@x\global
%    \end{macrocode}
% 在 2025 年 11 月的 \LaTeX{} 内核中，上述技巧——做一个本应
% 是 |\long| 或 |\relax| 的命令却突然应用 |\global|——不再
% 有效，\url{https://tex.stackexchange.com/a/754359} 说明了原因。
% 下面这一行通过换成一个同样可疑的技巧来补救这个缺陷 :-)
%    \begin{macrocode}
    \let\pr@tectedrel@x\global
    \let\@xargdef\lstenv@xargdef
    \expandafter\new@command\csname##2@\endcsname##3%
%    \end{macrocode}
% 首先我们执行 |##4|=\meta{begin code}。然后定义
% 终止字符串（例如 |end{lstlisting}| 或 |endlstlisting|）：
%    \begin{macrocode}
        {\lsthk@PreSet ##4%
         \ifx\@currenvir\lstenv@name
             \def\lstenv@endstring{#1#2##1#3}%
         \else
             \def\lstenv@endstring{#1##1}%
         \fi
%    \end{macrocode}
% 我们（局部）重新定义`结束环境'，因为现在结束是合法的。
% 注意该重新定义在 \TeX\ 注释行内也有效。
%    \begin{macrocode}
         \@namedef{end##2}{\lst@DeInit ##5\endgroup
                          \lst@doendpe \@ignoretrue}%
%    \end{macrocode}
% |\lst@doendpe| 再次消除缩进问题。
%
% 最后我们开始处理。|\lst@EndProcessListing| 的赋值
% 被移到 |\lst@Init| 之前，这是在
% \lsthelper{Andreas~Deininger}{2002/11/11}{Compiling just stops}.
%    \begin{macrocode}
         \lsthk@DisplayStyle
         \let\lst@EndProcessListing\lstenv@SkipToEnd
         \lst@Init\lstenv@backslash
         \lst@ifprint
             \expandafter\expandafter\expandafter\lstenv@Process
         \else
             \expandafter\lstenv@SkipToEnd
         \fi
         \lst@insertargs}%
    \endgroup}%
}
\let\lst@arg\@empty \lst@XConvert{end}\{\}\@nil
\expandafter\@tempa\lst@arg
\let\lst@insertargs\@empty
%    \end{macrocode}
% \end{macro}
%
% \begin{macro}{\lstenv@xargdef}
% 这是 \LaTeX 的 |\@xargdef| 的派生版本。我们展开子宏的名字，
% 使用 |\gdef| 而不是 |\def|，并硬编码一种 |\@protected@testopt|。
%    \begin{macrocode}
\def\lstenv@xargdef#1{%
  \expandafter\lstenv@xargdef@\csname\string#1\endcsname#1}
\def\lstenv@xargdef@#1#2[#3][#4]#5{%
  \@ifdefinable#2{%
    \gdef#2{%
      \ifx\protect\@typeset@protect
          \expandafter\lstenv@testopt
      \else
          \@x@protect#2%
      \fi
      #1%
      {#4}}%
    \@yargdef%
    #1%
    \tw@%
    {#3}%
    {#5}}%
}
%    \end{macrocode}
% \end{macro}
%
% \begin{macro}{\lstenv@testopt}
% 这个宏与 |\@testopt| 的区别在于我们临时
% 重置 EOL 字符 |^^M| 的 catcode 以读取可选参数。
%    \begin{macrocode}
\long\def\lstenv@testopt#1#2{%
  \@ifnextchar[{\catcode\active5\relax \lstenv@testopt@#1}%
               {#1[{#2}]}}
\def\lstenv@testopt@#1[#2]{%
    \catcode\active\active
    #1[#2]}
%    \end{macrocode}
% \end{macro}
%
% \begin{macro}{\lstenv@SkipToEnd}
% 我们使用临时定义
% \begin{itemize}\item[]
%    |\long\def\lst@temp##1\|\meta{content of \textup{\cs{lstenv@endstring}}}|{\lstenv@End}|
% \end{itemize}
% 它会吞掉直到环境末尾的所有字符并结束它。
%    \begin{macrocode}
\begingroup \lccode`\~=`\\\lowercase{%
\gdef\lstenv@SkipToEnd{%
    \long\expandafter\def\expandafter\lst@temp\expandafter##\expandafter
        1\expandafter~\lstenv@endstring{\lstenv@End}%
    \lst@temp}
}\endgroup
%    \end{macrocode}
% \end{macro}
%
% \begin{macro}{\lstenv@Error}
% 由一个孤立的`结束环境'调用。
%    \begin{macrocode}
\def\lstenv@Error#1{\PackageError{Listings}{Extra \string\end#1}%
    {I'm ignoring this, since I wasn't doing a \csname#1\endcsname.}}
%    \end{macrocode}
% \end{macro}
%
% \begin{macro}{\lst@TestEOLChar}
% 这里我们测试两种可能的 EOL 字符。
%    \begin{macrocode}
\begingroup \lccode`\~=`\^^M\lowercase{%
\gdef\lst@TestEOLChar#1{%
    \def\lst@insertargs{#1}%
    \ifx ~#1\@empty \else
    \ifx^^J#1\@empty \else
        \global\let\lst@intname\lst@insertargs
        \let\lst@insertargs\@empty
    \fi \fi}
}\endgroup
%    \end{macrocode}
% \end{macro}
%
% \begin{environment}{lstlisting}
% 棘手的工作已完成，现在定义很容易。我们测试
% 用户是否给出了名字参数、设置各键，并处理
% 续接的行号。
%    \begin{macrocode}
\lstnewenvironment{lstlisting}[2][]{%
     \lst@TestEOLChar{#2}%
     \lstset{#1}%
     \csname\@lst @SetFirstNumber\endcsname%
   }{%
%    \end{macrocode}
% 正如 \lsthelper{Jean-Yves~Baudais}{2016/03/11}{一些
% 垂直间距问题}指出的，列表下方的垂直间距有问题。
% 当列表从 |\paragraph| 顶部或某个 |example| 环境
% 开头开始时，也会出现类似问题。Jean-Yves 提供了一个
% 解决方案——|\let\if@nobreak\iffalse|——如
% |fr.comp.text.tex| 上所讨论。原先的假设——如果列表开头
% 有上框线，或者由 |\leavevmode| 引入列表，问题就会消失——
% 是错误的，正如
% \lsthelper{Karl~Berry}{2018/10/26}{listings 1.7 \rcmdname\leavevmode -> 空行}和
% \lsthelper{Sven~Schreiber}{2018/10/23}{新版 listings 中行为改变}各自
% 独立报告的那样，所以提出的代码进入环境定义的
% 第二部分。\lsthelper{Enrico Gregorio}{2019/05/04}{这是
% listings 宏包的缺陷吗？}在
% \url{https://tex.stackexchange.com/questions/489121/is-it-a-listings-package-bug}
% 回答：先前的解决方案 |\let\if@nobreak\iffalse| 是错误的，因为它是
% 局部赋值，而需要的是全局设置。
%    \begin{macrocode}
     \@nobreakfalse
     \csname\@lst @SaveFirstNumber\endcsname%
   }
%    \end{macrocode}
%    \begin{macrocode}
%</kernel>
%    \end{macrocode}
% \end{environment}
%
%
% \section{文档支持}
%
% \begin{syntax}
% \item[0.19]
%   |\begin{lstsample}|\oarg{point list}\marg{left}\marg{right}
%
%   \leavevmode\hspace*{-\leftmargini}|\end{lstsample}|
%
%             粗略地说，该环境之间的所有材料在左侧执行
%             并在右侧按 verbatim 排印。\meta{left} 在左侧排印
%             之前执行，\meta{right} 同理在右侧
%             之前执行。
%
%             \meta{point list} 被用作 \keyname{point} 键的参数。
%             这是一个用于在示例中高亮各键的特殊键。
%
% \item[1.0]
%   |\begin{lstxsample}|\marg{point list}
%
%   \leavevmode\hspace*{-\leftmargini}|\end{lstxsample}|
%
%             中间的材料 (a) 被加到下一个 \texttt{lstsample} 环境的左侧，
%             (b) 以整行宽度按 verbatim
%             排印。
%
% \item[0.21] |\newdocenvironment|\marg{name}\marg{short name}\marg{begin code}\marg{end code}
%
%             \meta{name} 环境可以像`宏'那样使用。
%             提供的(!)定义
%                     |\Print|\meta{short name}|Name|
%             和 |\SpecialMain|\meta{short name}|Index|
%             像默认的 |\PrintMacroName| 和 |\SpecialMainIndex| 那样
%             控制页边打印和索引。
%
%       这个命令用于定义`方面'和`lstkey'环境。
%
%\item[0.21] \texttt{macroargs} 环境
%
%             这个`enumerate'环境使用标签`|#1| ='、`|#2| ='，
%             依此类推。
%
% \item \texttt{TODO} 环境
% \item \texttt{ALTERNATIVE} 环境
% \item \texttt{REMOVED} 环境
% \item \texttt{OLDDEF} 环境
%
%             这些环境包围关于`待做'、备选方案以及
%             已删除或旧定义的注释。
%
% \item[0.21] |\lstscanlanguages|\meta{list macro}\marg{input files}\marg{don't input}
%
%             在 \marg{input files}$\setminus$\marg{don't input} 中扫描语言
%             定义。可用的语言以 \meta{language}|(|\meta{dialect}|),| 的
%             形式存入 \meta{list macro}。
%
% \item[0.21] |\lstprintlanguages|\meta{list macro}
%
%             以双栏格式打印语言。
% \end{syntax}
% 以及许多更简单的命令。
%
%
% \subsection{必需的宏包}
%
% 大多数`必需的'宏包都是可选的。
% \lsthelper{Stephan~Hennig}{2006-09-25}{文档与 algorithmic 不兼容}
% 注意到一个缺陷：|\ifalgorithmic| 与 |algorithmic.sty| 的更新冲突，因此
% 这已改为 |\ifalgorithmicpkg|。
%    \begin{macrocode}
%<*doc>
\let\lstdoc@currversion\fileversion
\RequirePackage[writefile]{listings}[2004/09/07]
\newif\iffancyvrb \IfFileExists{fancyvrb.sty}{\fancyvrbtrue}{}
\newif\ifcolor \IfFileExists{color.sty}{\colortrue}{}
\lst@false
\newif\ifhyper
\@ifundefined{pdfoutput}
    {}
    {\ifnum\pdfoutput>\z@ \lst@true \fi}
\@ifundefined{VTeXversion}
    {}
    {\ifnum\OpMode>\z@ \lst@true \fi}
\lst@if \IfFileExists{hyperref.sty}{\hypertrue}{}\fi
\newif\ifalgorithmicpkg \IfFileExists{algorithmic.sty}{\algorithmicpkgtrue}{}
\newif\iflgrind \IfFileExists{lgrind.sty}{\lgrindtrue}{}
\iffancyvrb \RequirePackage{fancyvrb}\fi
\ifhyper \RequirePackage[colorlinks]{hyperref}\else
    \def\href#1{\texttt}\fi
\ifcolor \RequirePackage{color}\fi
\ifalgorithmicpkg \RequirePackage{algorithmic}\fi
\iflgrind \RequirePackage{lgrind}\fi
\RequirePackage{nameref}
\RequirePackage{xurl}           % allow URL breaks
\renewcommand\ref{\protect\T@ref}
\renewcommand\pageref{\protect\T@pageref}
%    \end{macrocode}
%
%
% \subsection{注释用环境}
%
% \begin{macro}{\lst@BeginRemark}
% \begin{macro}{\lst@EndRemark}
% 我们先从两个简单的定义开始\ldots
%    \begin{macrocode}
\def\lst@BeginRemark#1{%
    \begin{quote}\topsep0pt\let\small\footnotesize\small#1:}
\def\lst@EndRemark{\end{quote}}
%    \end{macrocode}
% \end{macro}\end{macro}
%
% \begin{environment}{TODO}
% \begin{environment}{ALTERNATIVE}
% \begin{environment}{REMOVED}
% \begin{environment}{OLDDEF}
% \ldots\ 用于定义一些环境。
%    \begin{macrocode}
\newenvironment{TODO}
    {\lst@BeginRemark{To do}}{\lst@EndRemark}
\newenvironment{ALTERNATIVE}
    {\lst@BeginRemark{Alternative}}{\lst@EndRemark}
\newenvironment{REMOVED}
    {\lst@BeginRemark{Removed}}{\lst@EndRemark}
\newenvironment{OLDDEF}
    {\lst@BeginRemark{Old definition}}{\lst@EndRemark}
%    \end{macrocode}
% \end{environment}\end{environment}\end{environment}\end{environment}
%
% \begin{environment}{advise}
% \begin{macro}{\advisespace}
% 该环境使用 |\@listi|。
%    \begin{macrocode}
\def\advise{\par\list\labeladvise
    {\advance\linewidth\@totalleftmargin
     \@totalleftmargin\z@
     \@listi
     \let\small\footnotesize \small\sffamily
     \parsep \z@ \@plus\z@ \@minus\z@
     \topsep6\p@ \@plus1\p@\@minus2\p@
     \def\makelabel##1{\hss\llap{##1}}}}
\let\endadvise\endlist
%    \end{macrocode}
%    \begin{macrocode}
\def\advisespace{\hbox{}\qquad}
\def\labeladvise{$\to$}
%    \end{macrocode}
% \end{macro}
% \end{environment}
%
% \begin{environment}{syntax}
% \begin{macro}{\syntaxbreak}
% \begin{macro}{\syntaxnewline}
% \begin{macro}{\syntaxor}
% 该环境使用带特殊 |\makelabel| 的 |\list|，\ldots
%    \begin{macrocode}
\newenvironment{syntax}
   {\list{}{\itemindent-\leftmargin
    \def\makelabel##1{\hss\lst@syntaxlabel##1,,,,\relax}}}
   {\endlist}
%    \end{macrocode}
% \ldots\ 它在这里定义。逗号分隔的条目按
% 需要放置。
%    \begin{macrocode}
\def\lst@syntaxlabel#1,#2,#3,#4\relax{%
    \llap{\scriptsize\itshape#3}%
    \def\lst@temp{#2}%
    \expandafter\lst@syntaxlabel@\meaning\lst@temp\relax
    \rlap{\hskip-\itemindent\hskip\itemsep\hskip\linewidth
          \llap{\ttfamily\lst@temp}\hskip\labelwidth
          \def\lst@temp{#1}%
          \ifx\lst@temp\lstdoc@currversion#1\fi}}
\def\lst@syntaxlabel@#1>#2\relax
    {\edef\lst@temp{\zap@space#2 \@empty}}
%    \end{macrocode}
%    \begin{macrocode}
\newcommand*\syntaxnewline{\newline\hbox{}\kern\labelwidth}
\newcommand*\syntaxor{\qquad or\qquad}
\newcommand*\syntaxbreak
    {\hfill\kern0pt\discretionary{}{\kern\labelwidth}{}}
\let\syntaxfill\hfill
%    \end{macrocode}
% \end{macro}
% \end{macro}
% \end{macro}
% \end{environment}
%
% \begin{macro}{\alternative}
% 沿列表迭代并插入竖线。
%    \begin{macrocode}
\def\alternative#1{\lst@true \alternative@#1,\relax,}
\def\alternative@#1,{%
    \ifx\relax#1\@empty
        \expandafter\@gobble
    \else
        \ifx\@empty#1\@empty\else
            \lst@if \lst@false \else $\vert$\fi
            \textup{\texttt{#1}}%
        \fi
    \fi
    \alternative@}
%    \end{macrocode}
% \end{macro}
%
%
% \subsection{对 \packagename{doc} 的扩展}
%
% \begin{macro}{\m@cro@}
% 我们需要略微修改 \packagename{doc} 的内部宏。
% 原先的参数 |#2| 变成了 |#3|。此改动未在下面标出。
% 第二个参数现在是 \meta{short name}。
%    \begin{macrocode}
\long\def\m@cro@#1#2#3{\endgroup \topsep\MacroTopsep \trivlist
  \edef\saved@macroname{\string#3}%
  \def\makelabel##1{\llap{##1}}%
  \if@inlabel
    \let\@tempa\@empty \count@\macro@cnt
    \loop \ifnum\count@>\z@
      \edef\@tempa{\@tempa\hbox{\strut}}\advance\count@\m@ne \repeat
    \edef\makelabel##1{\llap{\vtop to\baselineskip
                               {\@tempa\hbox{##1}\vss}}}%
    \advance \macro@cnt \@ne
  \else  \macro@cnt\@ne  \fi
  \edef\@tempa{\noexpand\item[%
     #1%
       \noexpand\PrintMacroName
     \else
%    \end{macrocode}
% 下一行已被修改。
%    \begin{macrocode}
       \expandafter\noexpand\csname Print#2Name\endcsname % MODIFIED
     \fi
     {\string#3}]}%
  \@tempa
  \global\advance\c@CodelineNo\@ne
   #1%
      \SpecialMainIndex{#3}\nobreak
      \DoNotIndex{#3}%
   \else
%    \end{macrocode}
% 同上。
%    \begin{macrocode}
      \csname SpecialMain#2Index\endcsname{#3}\nobreak % MODIFIED
   \fi
  \global\advance\c@CodelineNo\m@ne
  \ignorespaces}
%    \end{macrocode}
% \end{macro}
%
% \begin{macro}{\macro}
% \begin{macro}{\environment}
% 由于修改过的 |\m@cro@|，这两个定义需要小幅调整。
%    \begin{macrocode}
\def\macro{\begingroup
   \catcode`\\12
   \MakePrivateLetters \m@cro@ \iftrue {Macro}}% MODIFIED
\def\environment{\begingroup
   \catcode`\\12
   \MakePrivateLetters \m@cro@ \iffalse {Env}}% MODIFIED
%    \end{macrocode}
% \end{macro}\end{macro}
%
% \begin{macro}{\newdocenvironment}
% 这个命令简单地制作类似`environment'的定义，并提供
% 打印和索引命令。
%    \begin{macrocode}
\def\newdocenvironment#1#2#3#4{%
    \@namedef{#1}{#3\begingroup \catcode`\\12\relax
                  \MakePrivateLetters \m@cro@ \iffalse {#2}}%
    \@namedef{end#1}{#4\endmacro}%
    \@ifundefined{Print#2Name}{\expandafter
        \let\csname Print#2Name\endcsname\PrintMacroName}{}%
    \@ifundefined{SpecialMain#2Index}{\expandafter
        \let\csname SpecialMain#2Index\endcsname\SpecialMainIndex}{}}
%    \end{macrocode}
% \end{macro}
%
% \begin{environment}{aspect}
% \begin{macro}{\PrintAspectName}
% \begin{macro}{\SpecialMainAspectIndex}
% 该环境及其`print'和`index'命令。
%    \begin{macrocode}
\newdocenvironment{aspect}{Aspect}{}{}
\def\PrintAspectName#1{}
\def\SpecialMainAspectIndex#1{%
    \@bsphack
    \index{aspects:\levelchar\protect\aspectname{#1}}%
    \@esphack}
%    \end{macrocode}
% \end{macro}\end{macro}\end{environment}
%
% \begin{environment}{lstkey}
% \begin{macro}{\PrintKeyName}
% \begin{macro}{\SpecialMainKeyIndex}
% 再来一个环境及其`print'和`index'命令。
%    \begin{macrocode}
\newdocenvironment{lstkey}{Key}{}{}
\def\PrintKeyName#1{\strut\keyname{#1}\ }
\def\SpecialMainKeyIndex#1{%
    \@bsphack
    \index{keys\levelchar\protect\keyname{#1}}%
    \@esphack}
%    \end{macrocode}
% \end{macro}\end{macro}\end{environment}
%
% \begin{macro}{\labelargcount}
% \begin{environment}{macroargs}
% 我们只是分配一个计数器，并使用 \LaTeX 的 |\list| 实现这个
% 环境。
%    \begin{macrocode}
\newcounter{argcount}
\def\labelargcount{\texttt{\#\arabic{argcount}}\hskip\labelsep$=$}
%    \end{macrocode}
%    \begin{macrocode}
\def\macroargs{\list\labelargcount
    {\usecounter{argcount}\leftmargin=2\leftmargin
     \parsep \z@ \@plus\z@ \@minus\z@
     \topsep4\p@ \@plus\p@ \@minus2\p@
     \itemsep\z@ \@plus\z@ \@minus\z@
     \def\makelabel##1{\hss\llap{##1}}}}
\def\endmacroargs{\endlist\@endparenv}
%    \end{macrocode}
% \end{environment}\end{macro}
%
%
% \subsection{\texttt{lstsample} 环境}
%
% \begin{environment}{lstsample}
% 我们存储 verbatim 部分，同时把源代码写入文件。
%    \begin{macrocode}
\lst@RequireAspects{writefile}
%    \end{macrocode}
%    \begin{macrocode}
\newbox\lst@samplebox
\lstnewenvironment{lstsample}[3][]
    {\global\let\lst@intname\@empty
     \gdef\lst@sample{#2}%
     \setbox\lst@samplebox=\hbox\bgroup
         \setkeys{lst}{language={},style={},tabsize=4,gobble=5,%
             basicstyle=\small\ttfamily,basewidth=0.51em,point={#1}}
         #3%
         \lst@BeginAlsoWriteFile{\jobname.tmp}}
    {\lst@EndWriteFile\egroup
%    \end{macrocode}
% 现在 |\lst@samplebox| 包含 verbatim 部分。
% 如果它太宽，我们就用 atop 和 below 代替 left 和 right。
%    \begin{macrocode}
     \ifdim \wd\lst@samplebox>.5\linewidth
         \begin{center}%
             \hbox to\linewidth{\box\lst@samplebox\hss}%
         \end{center}%
         \lst@sampleInput
     \else
         \begin{center}%
         \begin{minipage}{0.45\linewidth}\lst@sampleInput\end{minipage}%
         \qquad
         \begin{minipage}{0.45\linewidth}%
             \hbox to\linewidth{\box\lst@samplebox\hss}%
         \end{minipage}%
         \end{center}%
     \fi}
%    \end{macrocode}
% 新的关键词类 \keyname{point}。
%    \begin{macrocode}
\lst@InstallKeywords{p}{point}{pointstyle}\relax{keywordstyle}{}ld
%    \end{macrocode}
% \end{environment}
%
% \begin{environment}{lstxsample}
% 省略 |\lst@EndWriteFile| 会使文件保持打开。
%    \begin{macrocode}
\lstnewenvironment{lstxsample}[1][]
    {\begingroup
         \setkeys{lst}{belowskip=-\medskipamount,language={},style={},%
             tabsize=4,gobble=5,basicstyle=\small\ttfamily,%
             basewidth=0.51em,point={#1}}
         \lst@BeginAlsoWriteFile{\jobname.tmp}}
    {\endgroup
     \endgroup}
%    \end{macrocode}
% \end{environment}
%
% \begin{macro}{\lst@sampleInput}
% 输入`左侧'。
%    \begin{macrocode}
\def\lst@sampleInput{%
    \MakePercentComment\catcode`\^^M=10\relax
    \small\lst@sample
    {\setkeys{lst}{SelectCharTable=\lst@ReplaceInput{\^\^I}%
                                  {\lst@ProcessTabulator}}%
     \leavevmode \input{\jobname.tmp}}\MakePercentIgnore}
%    \end{macrocode}
% \end{macro}
%
%
% \subsection{其他}
%
% \paragraph{分节与交叉引用}
% 我们以一个重新定义段落开始。
%    \begin{macrocode}
\renewcommand\paragraph{\@startsection{paragraph}{4}{\z@}%
                                      {1.25ex \@plus1ex \@minus.2ex}%
                                      {-1em}%
                                      {\normalfont\normalsize\bfseries}}
%    \end{macrocode}
% 我们引入 |\lstref|，它打印节号及其名称。
%    \begin{macrocode}
\def\lstref#1{\emph{\ref{#1} \nameref{#1}}}
%    \end{macrocode}
% 此外我们调整目录。在加入目录行之前使用 |\phantomsection|
% 为 \packagename{hyperref} 提供合适的目录行链接
% 目标，从而确保目录行位于
% PDF 书签树中的正确层级。
%    \begin{macrocode}
\def\@part[#1]#2{\ifhyper\phantomsection\fi
    \addcontentsline{toc}{part}{#1}%
    {\parindent\z@ \raggedright \interlinepenalty\@M
     \normalfont \huge \bfseries #2\markboth{}{}\par}%
    \nobreak\vskip 3ex\@afterheading}
\renewcommand*\l@section[2]{%
    \addpenalty\@secpenalty
    \addvspace{.25em \@plus\p@}%
    \setlength\@tempdima{1.5em}%
    \begingroup
      \parindent \z@ \rightskip \@pnumwidth
      \parfillskip -\@pnumwidth
      \leavevmode
      \advance\leftskip\@tempdima
      \hskip -\leftskip
      #1\nobreak\hfil \nobreak\hb@xt@\@pnumwidth{\hss #2}\par
    \endgroup}
\renewcommand*\l@subsection{\@dottedtocline{2}{0pt}{2.3em}}
\renewcommand*\l@subsubsection{\@dottedtocline{3}{0pt}{3.2em}}
%    \end{macrocode}
%
% \paragraph{索引}
% `用户'命令。标记一个键有两种不同
% 方式。|\ikeyname| 是\textit{行}内使用键的命令，
% |\rkeyname| 定义键的\textit{参}考（显示为绿色），
% |\rstyle| 在下面定义。
%    \begin{macrocode}
\newcommand\ikeyname[1]{%
    \lstkeyindex{#1}{}%
    \lstaspectindex{#1}{}%
    \keyname{#1}}
\newcommand\rkeyname[1]{%
    \@bsphack
    \lstkeyindex{#1}{}%
    \lstaspectindex{#1}{}%
    \@esphack{\rstyle\keyname{#1}}}
%    \end{macrocode}
% 同样的规则适用于 |\icmdname| 和 |\rcmdname|：
%    \begin{macrocode}
\newcommand\icmdname[1]{%
    \@bsphack
    \lstaspectindex{#1}{}%
    \@esphack\texttt{\string#1}}
\newcommand\rcmdname[1]{%
    \@bsphack
    \lstaspectindex{#1}{}%
    \@esphack\texttt{\rstyle\string#1}}
%    \end{macrocode}
% 两个未知的`index'宏中有一个为空，另一个查找
% 给定参数对应的方面名。
%    \begin{macrocode}
\def\lstaspectindex#1#2{%
    \global\@namedef{lstkandc@\string#1}{}%
    \@ifundefined{lstisaspect@\string#1}
        {\index{unknown\levelchar
                \protect\texttt{\protect\string\string#1}#2}}%
        {\index{\@nameuse{lstisaspect@\string#1}\levelchar
                \protect\texttt{\protect\string\string#1}#2}}%
}
\def\lstkeyindex#1#2{%
%    \index{key\levelchar\protect\keyname{#1}#2}%
}
%    \end{macrocode}
% 键/命令与方面的关系在本文件顶部附近用
% 下面的命令定义。将来宏包应当从
% 方面文件读取这一信息。
%    \begin{macrocode}
\def\lstisaspect[#1]#2{%
    \global\@namedef{lstaspect@#1}{#2}%
    \lst@AddTo\lst@allkeysandcmds{,#2}%
    \@for\lst@temp:=#2\do
    {\ifx\@empty\lst@temp\else
         \global\@namedef{lstisaspect@\lst@temp}{#1}%
     \fi}}
\gdef\lst@allkeysandcmds{}
%    \end{macrocode}
% 这个关系也便于打印某一方面的所有键和命令\ldots
%    \begin{macrocode}
\def\lstprintaspectkeysandcmds#1{%
    \lst@true
    \expandafter\@for\expandafter\lst@temp
    \expandafter:\expandafter=\csname lstaspect@#1\endcsname\do
    {\lst@if\lst@false\else, \fi \texttt{\lst@temp}}}
%    \end{macrocode}
% \ldots\ 或检查参考。注意我们已在
% |\lstaspectindex| 中定义了 |\lstkandc@|\meta{name}。
%    \begin{macrocode}
\def\lstcheckreference{%
   \@for\lst@temp:=\lst@allkeysandcmds\do
   {\ifx\lst@temp\@empty\else
        \@ifundefined{lstkandc@\lst@temp}
        {\typeout{\lst@temp\space not in reference guide?}}{}%
    \fi}}
%    \end{macrocode}
%
%
% \paragraph{唯一样式}
%    \begin{macrocode}
\newcommand*\lst{\texttt{lst}}
\newcommand*\Cpp{C\texttt{++}}
\let\keyname\texttt
\let\keyvalue\texttt
\let\hookname\texttt
\newcommand*\aspectname[1]{{\normalfont\sffamily#1}}
%    \end{macrocode}
%    \begin{macrocode}
\DeclareRobustCommand\packagename[1]{%
    {\leavevmode\text@command{#1}%
     \switchfontfamily\sfdefault\rmdefault
     \check@icl #1\check@icr
     \expandafter}}%
\renewcommand\packagename[1]{{\normalfont\sffamily#1}}
\def\switchfontfamily#1#2{%
    \begingroup\xdef\@gtempa{#1}\endgroup
    \ifx\f@family\@gtempa\fontfamily#2%
                    \else\fontfamily#1\fi
    \selectfont}
%    \end{macrocode}
% 主要用于参考指南中键和命令的颜色
% ——\textit{r} 表示参考。
%    \begin{macrocode}
\ifcolor
    \definecolor{darkgreen}{rgb}{0,0.5,0}
    \def\rstyle{\color{darkgreen}}
\else
    \let\rstyle\empty
\fi
%    \end{macrocode}
%
% \paragraph{致谢与帮助者命令}
% 有两个用于致谢和帮助者的命令：
% \begin{enumerate}
%     \item |\lstthanks| 用于把贡献者的名字放入
%         `结语与致谢'节。它有两个参数：|#1| 为
%         名字，|#2| 为电子邮件地址——该地址不会被显示。
%     \item |\lsthelper| 必须在正文中使用，以显示
%         帮助者的名字（参数 |#1|）、贡献日期（参数 |#2|）
%         以及关于该贡献的简短文字（参数 |#3|）。只
%         打印第一个参数。
% \end{enumerate}
%    \begin{macrocode}
\gdef\lst@emails{}
\newcommand*\lstthanks[2]
    {#1\lst@AddTo\lst@emails{,#1,<#2>}%
     \ifx\@empty#2\@empty\typeout{Missing email for #1}\fi}
\newcommand*\lsthelper[3]
    {{\let~\ #1}%
     \lst@IfOneOf#1\relax\lst@emails
     {}{\typeout{^^JWarning: Unknown helper #1.^^J}}}
%    \end{macrocode}
%
% \paragraph{语言与样式}
%    \begin{macrocode}
\lstdefinelanguage[doc]{Pascal}{%
  morekeywords={alfa,and,array,begin,boolean,byte,case,char,const,div,%
     do,downto,else,end,false,file,for,function,get,goto,if,in,%
     integer,label,maxint,mod,new,not,of,or,pack,packed,page,program,%
     procedure,put,read,readln,real,record,repeat,reset,rewrite,set,%
     text,then,to,true,type,unpack,until,var,while,with,write,writeln},%
  sensitive=false,%
  morecomment=[s]{(*}{*)},%
  morecomment=[s]{\{}{\}},%
  morestring=[d]{'}}
%    \end{macrocode}
%    \begin{macrocode}
\lstdefinestyle{}
    {basicstyle={},%
     keywordstyle=\bfseries,identifierstyle={},%
     commentstyle=\itshape,stringstyle={},%
     numberstyle={},stepnumber=1,%
     pointstyle=\pointstyle}
\def\pointstyle{%
    {\let\lst@um\@empty \xdef\@gtempa{\the\lst@token}}%
    \expandafter\lstkeyindex\expandafter{\@gtempa}{}%
    \expandafter\lstaspectindex\expandafter{\@gtempa}{}%
    \rstyle}
\lstset{defaultdialect=[doc]Pascal,language=Pascal,style={}}
%    \end{macrocode}
%
%
% \subsection{扫描语言}
%
% \begin{macro}{\lstscanlanguages}
% 我们修改一些内部定义并输入各文件。
%    \begin{macrocode}
\def\lstscanlanguages#1#2#3{%
    \begingroup
        \def\lst@DefDriver@##1##2##3##4[##5]##6{%
           \lst@false
           \lst@lAddTo\lst@scan{##6(##5),}%
           \begingroup
           \@ifnextchar[{\lst@XDefDriver{##1}##3}{\lst@DefDriver@@##3}}%
        \def\lst@XXDefDriver[##1]{}%
        \lst@InputCatcodes
        \def\lst@dontinput{#3}%
        \let\lst@scan\@empty
        \lst@for{#2}\do{%
            \lst@IfOneOf##1\relax\lst@dontinput
                {}%
                {\InputIfFileExists{##1}{}{}}}%
        \global\let\@gtempa\lst@scan
    \endgroup
    \let#1\@gtempa}
%    \end{macrocode}
% \end{macro}
%
% \begin{macro}{\lstprintlanguages}
% |\do| 创建一个宽度为 0.5|\linewidth| 或 |\linewidth| 的盒子，
% 取决于参数有多宽。这就导致`双栏'输出。
% 另一件主要的事是对列表排序并开始输出。
%    \begin{macrocode}
\def\lstprintlanguages#1{%
    \def\do##1{\setbox\@tempboxa\hbox{##1\space\space}%
        \ifdim\wd\@tempboxa<.5\linewidth \wd\@tempboxa.5\linewidth
                                   \else \wd\@tempboxa\linewidth \fi
        \box\@tempboxa\allowbreak}%
    \begin{quote}
      \par\noindent
      \hyphenpenalty=\@M \rightskip=\z@\@plus\linewidth\relax
      \lst@BubbleSort#1%
      \expandafter\lst@NextLanguage#1\relax(\relax),%
    \end{quote}}
%    \end{macrocode}
% 我们取得并定义当前语言，并\ldots
%    \begin{macrocode}
\def\lst@NextLanguage#1(#2),{%
    \ifx\relax#1\else
        \def\lst@language{#1}\def\lst@dialects{(#2),}%
        \expandafter\lst@NextLanguage@
    \fi}
%    \end{macrocode}
% \ldots\ 收集该语言的所有可用方言（注意
% 列表已排序）
%    \begin{macrocode}
\def\lst@NextLanguage@#1(#2),{%
    \def\lst@temp{#1}%
    \ifx\lst@temp\lst@language
        \lst@lAddTo\lst@dialects{(#2),}%
        \expandafter\lst@NextLanguage@
    \else
%    \end{macrocode}
% 或者开始打印该语言及其所有方言。为此我们
% 对方言排序
%    \begin{macrocode}
        \do{\lst@language
        \ifx\lst@dialects\lst@emptydialect\else
            \expandafter\lst@NormedDef\expandafter\lst@language
                \expandafter{\lst@language}%
            \space(%
            \lst@BubbleSort\lst@dialects
            \expandafter\lst@PrintDialects\lst@dialects(\relax),%
            )%
        \fi}%
        \def\lst@next{\lst@NextLanguage#1(#2),}%
        \expandafter\lst@next
    \fi}
\def\lst@emptydialect{(),}
%    \end{macrocode}
% 并打印方言，之间用合适的逗号。
%    \begin{macrocode}
\def\lst@PrintDialects(#1),{%
    \ifx\@empty#1\@empty empty\else
        \lst@PrintDialect{#1}%
    \fi
    \lst@PrintDialects@}
\def\lst@PrintDialects@(#1),{%
    \ifx\relax#1\else
        , \lst@PrintDialect{#1}%
        \expandafter\lst@PrintDialects@
    \fi}
%    \end{macrocode}
% 这里我们处理默认方言。
%    \begin{macrocode}
\def\lst@PrintDialect#1{%
    \lst@NormedDef\lst@temp{#1}%
    \expandafter\ifx\csname\@lst dd@\lst@language\endcsname\lst@temp
        \texttt{\underbar{#1}}%
    \else
        \texttt{#1}%
    \fi}
%    \end{macrocode}
% \end{macro}
%
%
% \subsection{冒泡排序}
%
% \begin{macro}{\lst@IfLE}
% \meta{string 1}|\relax\@empty|\meta{string 2}|\relax\@empty|\marg{then}\meta{else}.
% 如果 \meta{string 1} $\leq$ \meta{string 2}，我们就执行 \meta{then}，
% 否则执行 \meta{else}。
% 注意这个比较不区分大小写。
%    \begin{macrocode}
\def\lst@IfLE#1#2\@empty#3#4\@empty{%
    \ifx #1\relax
        \let\lst@next\@firstoftwo
    \else \ifx #3\relax
        \let\lst@next\@secondoftwo
    \else
        \lowercase{\ifx#1#3}%
            \def\lst@next{\lst@IfLE#2\@empty#4\@empty}%
        \else
            \lowercase{\ifnum`#1<`#3}\relax
                \let\lst@next\@firstoftwo
            \else
                \let\lst@next\@secondoftwo
            \fi
        \fi
    \fi \fi
    \lst@next}
%    \end{macrocode}
% \end{macro}
%
% \begin{macro}{\lst@BubbleSort}
% 实际上是冒泡排序的派生。
%    \begin{macrocode}
\def\lst@BubbleSort#1{%
    \ifx\@empty#1\else
        \lst@false
%    \end{macrocode}
% 我们对第一、第二、\ldots\ 个元素`冒泡排序'，然后\ldots
%    \begin{macrocode}
        \expandafter\lst@BubbleSort@#1\relax,\relax,%
%    \end{macrocode}
% \ldots\ 然后对第二、第三、\ldots\ 个元素，直到没有元素
% 被交换。
%    \begin{macrocode}
        \expandafter\lst@BubbleSort@\expandafter,\lst@sorted
                                      \relax,\relax,%
        \let#1\lst@sorted
        \lst@if
            \def\lst@next{\lst@BubbleSort#1}%
            \expandafter\expandafter\expandafter\lst@next
        \fi
    \fi}
\def\lst@BubbleSort@#1,#2,{%
    \ifx\@empty#1\@empty
        \def\lst@sorted{#2,}%
        \def\lst@next{\lst@BubbleSort@@}%
    \else
        \let\lst@sorted\@empty
        \def\lst@next{\lst@BubbleSort@@#1,#2,}%
    \fi
    \lst@next}
%    \end{macrocode}
% 但气泡每次调用只上升一步。把元素放到它们
% 最靠上的位置效率不高（因为这种情况下 \TeX\ 必须读
% 多得多的参数）。
%    \begin{macrocode}
\def\lst@BubbleSort@@#1,#2,{%
    \ifx\relax#1\else
        \ifx\relax#2%
            \lst@lAddTo\lst@sorted{#1,}%
            \expandafter\expandafter\expandafter\lst@BubbleSort@@@
        \else
            \lst@IfLE #1\relax\@empty #2\relax\@empty
                          {\lst@lAddTo\lst@sorted{#1,#2,}}%
                {\lst@true \lst@lAddTo\lst@sorted{#2,#1,}}%
            \expandafter\expandafter\expandafter\lst@BubbleSort@@
        \fi
    \fi}
\def\lst@BubbleSort@@@#1\relax,{}
%    \end{macrocode}
%    \begin{macrocode}
%</doc>
%    \end{macrocode}
% \end{macro}
%
%
% \section{与其他程序的接口}
%
%
% \subsection{0.21 兼容性}
%
% \begin{aspect}{0.21}
% 有些键刚被改名。
%    \begin{macrocode}
%<*0.21>
\lst@BeginAspect{0.21}
%    \end{macrocode}
%
%    \begin{macrocode}
\lst@Key{labelstyle}{}{\def\lst@numberstyle{#1}}
\lst@Key{labelsep}{10pt}{\def\lst@numbersep{#1}}
\lst@Key{labelstep}{0}{%
    \ifnum #1=\z@ \KV@lst@numbers{none}%
            \else \KV@lst@numbers{left}\fi
    \def\lst@stepnumber{#1\relax}}
\lst@Key{firstlabel}\relax{\def\lst@firstnumber{#1\relax}}
\lst@Key{advancelabel}\relax{\def\lst@advancenumber{#1\relax}}
\let\c@lstlabel\c@lstnumber
\lst@AddToHook{Init}{\def\thelstnumber{\thelstlabel}}
\newcommand*\thelstlabel{\@arabic\c@lstlabel}
%    \end{macrocode}
% 倒数第二行的 |\let| 已改为 |\def|，这是在
% \lsthelper{Venkatesh~Prasad~Ranganath}{2002/08/31}{使用 0.21 选项时
% 控制序列 \thelstnumber 未定义}的缺陷报告之后。
%    \begin{macrocode}
\lst@Key{first}\relax{\def\lst@firstline{#1\relax}}
\lst@Key{last}\relax{\def\lst@lastline{#1\relax}}
%    \end{macrocode}
%    \begin{macrocode}
\lst@Key{framerulewidth}{.4pt}{\def\lst@framerulewidth{#1}}
\lst@Key{framerulesep}{2pt}{\def\lst@rulesep{#1}}
\lst@Key{frametextsep}{3pt}{\def\lst@frametextsep{#1}}
\lst@Key{framerulecolor}{}{\lstKV@OptArg[]{#1}%
    {\ifx\@empty##2\@empty
         \let\lst@rulecolor\@empty
     \else
         \ifx\@empty##1\@empty
             \def\lst@rulecolor{\color{##2}}%
         \else
             \def\lst@rulecolor{\color[##1]{##2}}%
         \fi
     \fi}}
\lst@Key{backgroundcolor}{}{\lstKV@OptArg[]{#1}%
    {\ifx\@empty##2\@empty
         \let\lst@bkgcolor\@empty
     \else
         \ifx\@empty##1\@empty
             \def\lst@bkgcolor{\color{##2}}%
         \else
             \def\lst@bkgcolor{\color[##1]{##2}}%
         \fi
     \fi}}
\lst@Key{framespread}{\z@}{\def\lst@framespread{#1}}
\lst@AddToHook{PreInit}
    {\@tempdima\lst@framespread\relax \divide\@tempdima\tw@
     \edef\lst@framextopmargin{\the\@tempdima}%
     \let\lst@framexrightmargin\lst@framextopmargin
     \let\lst@framexbottommargin\lst@framextopmargin
     \advance\@tempdima\lst@xleftmargin\relax
     \edef\lst@framexleftmargin{\the\@tempdima}}
%    \end{macrocode}
% \lsthelper{Harald~Harders}{1998/03/30}{内间距与外间距}提出了
% 两个间距（内与外）的想法。我们或者把尺寸减半，或者
% 把这两个尺寸分配给内间距与外间距。
%    \begin{macrocode}
\newdimen\lst@innerspread \newdimen\lst@outerspread
\lst@Key{spread}{\z@,\z@}{\lstKV@CSTwoArg{#1}%
    {\lst@innerspread##1\relax
     \ifx\@empty##2\@empty
         \divide\lst@innerspread\tw@\relax
         \lst@outerspread\lst@innerspread
     \else
         \lst@outerspread##2\relax
     \fi}}
\lst@AddToHook{BoxUnsafe}{\lst@outerspread\z@ \lst@innerspread\z@}
\lst@Key{wholeline}{false}[t]{\lstKV@SetIf{#1}\lst@ifresetmargins}
\lst@Key{indent}{\z@}{\def\lst@xleftmargin{#1}}
\lst@AddToHook{PreInit}
    {\lst@innerspread=-\lst@innerspread
     \lst@outerspread=-\lst@outerspread
     \ifodd\c@page \advance\lst@innerspread\lst@xleftmargin
             \else \advance\lst@outerspread\lst@xleftmargin \fi
     \ifodd\c@page
         \edef\lst@xleftmargin{\the\lst@innerspread}%
         \edef\lst@xrightmargin{\the\lst@outerspread}%
     \else
         \edef\lst@xleftmargin{\the\lst@outerspread}%
         \edef\lst@xrightmargin{\the\lst@innerspread}%
     \fi}
%    \end{macrocode}
%    \begin{macrocode}
\lst@Key{defaultclass}\relax{\def\lst@classoffset{#1}}
\lst@Key{stringtest}\relax{}% dummy
\lst@Key{outputpos}\relax{\lst@outputpos#1\relax\relax}
%    \end{macrocode}
%    \begin{macrocode}
\lst@Key{stringspaces}\relax[t]{\lstKV@SetIf{#1}\lst@ifshowstringspaces}
\lst@Key{visiblespaces}\relax[t]{\lstKV@SetIf{#1}\lst@ifshowspaces}
\lst@Key{visibletabs}\relax[t]{\lstKV@SetIf{#1}\lst@ifshowtabs}
%    \end{macrocode}
%
%    \begin{macrocode}
\lst@EndAspect
%</0.21>
%    \end{macrocode}
% \end{aspect}
%
%
% \subsection{\packagename{fancyvrb}}
%
% \lsthelper{Denis~Girou}{1998/07/26}{fancyvrb}询问
% \packagename{fancyvrb} 与 \packagename{listings} 能否一起工作。
%
% \begin{lstkey}{fancyvrb}
% 我们设置该布尔量并调用一个子宏。
%    \begin{macrocode}
%<*kernel>
\lst@Key{fancyvrb}\relax[t]{%
    \lstKV@SetIf{#1}\lst@iffancyvrb
    \lstFV@fancyvrb}
\ifx\lstFV@fancyvrb\@undefined
    \gdef\lstFV@fancyvrb{\lst@RequireAspects{fancyvrb}\lstFV@fancyvrb}
\fi
%</kernel>
%    \end{macrocode}
% \end{lstkey}
%
% \begin{aspect}{fancyvrb}
% 如果不存在 \packagename{fancyvrb}，我们就结束任务。
%    \begin{macrocode}
%<*misc>
\lst@BeginAspect{fancyvrb}
%    \end{macrocode}
%    \begin{macrocode}
\@ifundefined{FancyVerbFormatLine}
    {\typeout{^^J%
     ***^^J%
     *** `listings.sty' needs `fancyvrb.sty' right now.^^J%
     *** Please ensure its availability and try again.^^J%
     ***^^J}%
     \batchmode \@@end}{}
%    \end{macrocode}
%
% \begin{macro}{\lstFV@fancyvrb}
% 我们给 |\FancyVerbFormatLine| 指派正确的宏。
%    \begin{macrocode}
\gdef\lstFV@fancyvrb{%
    \lst@iffancyvrb
        \ifx\FancyVerbFormatLine\lstFV@FancyVerbFormatLine\else
            \let\lstFV@FVFL\FancyVerbFormatLine
            \let\FancyVerbFormatLine\lstFV@FancyVerbFormatLine
        \fi
    \else
        \ifx\lstFV@FVFL\@undefined\else
            \let\FancyVerbFormatLine\lstFV@FVFL
            \let\lstFV@FVFL\@undefined
        \fi
    \fi}
%    \end{macrocode}
% \end{macro}
%
% \begin{macro}{\lstFV@VerbatimBegin}
% 必要时我们进行初始化。
%    \begin{macrocode}
\gdef\lstFV@VerbatimBegin{%
    \ifx\FancyVerbFormatLine\lstFV@FancyVerbFormatLine
        \lsthk@TextStyle \lsthk@BoxUnsafe
        \lsthk@PreSet
        \lst@activecharsfalse
        \let\normalbaselines\relax
%    \end{macrocode}
% \begin{TODO}
% 这个 |\let| 不好吗？
% \end{TODO}
% 我在 \lsthelper{Peter~Bartke}{1999/11/18}{fancyvrb 框线错误}
% 的缺陷报告之后插入了 |\lst@ifresetmargins|\ldots|\fi|。行宽
% 被保存并恢复，这是在 \lsthelper{Denis~Girou}{2003/07/04}{fancyvrb=true 时列表环境中的问题}。
%    \begin{macrocode}
\xdef\lstFV@RestoreData{\noexpand\linewidth\the\linewidth\relax}%
        \lst@Init\relax
        \lst@ifresetmargins \advance\linewidth-\@totalleftmargin \fi
\lstFV@RestoreData
        \everypar{}\global\lst@newlines\z@
        \lst@mode\lst@nomode \let\lst@entermodes\@empty
        \lst@InterruptModes
%    \end{macrocode}
% \lsthelper{Rolf~Niepraschk}{1998/11/25}{连字问题}报告了一个
% 关于连字的缺陷，给 \lsthelper{Denis~Girou}{1998/11/27}{使用 |\@noligs|}。
%    \begin{macrocode}
%% D.G. modification begin - Nov. 25, 1998
        \let\@noligs\relax
%% D.G. modification end
    \fi}
%    \end{macrocode}
% \end{macro}
%
% \begin{macro}{\lstFV@VerbatimEnd}
% 在 |\lst@DeInit| 之后必须存在一个盒子和宏。
% 我们全局存储它们。
%    \begin{macrocode}
\gdef\lstFV@VerbatimEnd{%
    \ifx\FancyVerbFormatLine\lstFV@FancyVerbFormatLine
        \global\setbox\lstFV@gtempboxa\box\@tempboxa
        \global\let\@gtempa\FV@ProcessLine
        \lst@mode\lst@Pmode
        \lst@DeInit
        \let\FV@ProcessLine\@gtempa
        \setbox\@tempboxa\box\lstFV@gtempboxa
        \par
    \fi}
%    \end{macrocode}
% |\par| 是在 \lsthelper{Peter~Bartke}{2002/04/10}{离开 "Verbatim" 时 TeX 不在垂直模式} 的缺陷报告之后加入的。
%    \begin{macrocode}
\newbox\lstFV@gtempboxa
%    \end{macrocode}
% \end{macro}
%
% \noindent
% 我们在必要时插入 |\lstFV@VerbatimBegin| 和 |\lstFV@VerbatimEnd|。
%    \begin{macrocode}
\lst@AddTo\FV@VerbatimBegin\lstFV@VerbatimBegin
\lst@AddToAtTop\FV@VerbatimEnd\lstFV@VerbatimEnd
\lst@AddTo\FV@LVerbatimBegin\lstFV@VerbatimBegin
\lst@AddToAtTop\FV@LVerbatimEnd\lstFV@VerbatimEnd
\lst@AddTo\FV@BVerbatimBegin\lstFV@VerbatimBegin
\lst@AddToAtTop\FV@BVerbatimEnd\lstFV@VerbatimEnd
%    \end{macrocode}
%
% \begin{macro}{\lstFV@FancyVerbFormatLine}
% `@'终止 |\lst@FVConvert| 的参数。
% 此外 |\lst@ReenterModes| 和 |\lst@InterruptModes| 包围一些代码。
% 这确保我们在宏的开头和结尾处于相同的分组层级
% ——即使用户开始了但没有结束注释，也就是
% 有一个未关闭的分组。
% 此外我们使用 |\vtop| 并重置 |\lst@newlines| 以允许断行。
%    \begin{macrocode}
\gdef\lstFV@FancyVerbFormatLine#1{%
    \let\lst@arg\@empty \lst@FVConvert#1\@nil
    \global\lst@newlines\z@
    \vtop{\noindent\lst@parshape
          \lst@ReenterModes
          \lst@arg \lst@PrintToken\lst@EOLUpdate\lsthk@InitVarsBOL
          \lst@InterruptModes}}
%    \end{macrocode}
% |\vtop| 内的 |\lst@parshape| 源于 \lsthelper{Peter~Bartke}{1999/11/18}{fancyvrb 中的段首缩进错误}
% 的缺陷报告。
% |\leavevmode| 变成了 |\noindent|。
% \end{macro}
%
% \begin{lstkey}{fvcmdparams}
% \begin{lstkey}{morefvcmdparams}
% 这些键调整 |\lst@FVcmdparams|，后面的转换宏将用到它。
% 命令及其参数数目的基础集合
% 由 \lsthelper{Denis~Girou}{2002/05/31}{fvcmdparams 的初始化}提供。
%    \begin{macrocode}
\lst@Key{fvcmdparams}%
    {\overlay\@ne}%
    {\def\lst@FVcmdparams{,#1}}
\lst@Key{morefvcmdparams}\relax{\lst@lAddTo\lst@FVcmdparams{,#1}}
%    \end{macrocode}
% \end{lstkey}
% \end{lstkey}
%
% \begin{macro}{\lst@FVConvert}
% 我们进行转换或者\ldots
%    \begin{macrocode}
\gdef\lst@FVConvert{\@tempcnta\z@ \lst@FVConvertO@}%
\gdef\lst@FVConvertO@{%
    \ifcase\@tempcnta
        \expandafter\futurelet\expandafter\@let@token
        \expandafter\lst@FVConvert@@
    \else
%    \end{macrocode}
% \ldots\ 我们不经转换地追加参数，逐个参数地，
% 共 |\@tempcnta| 次。
%    \begin{macrocode}
        \expandafter\lst@FVConvertO@a
    \fi}
\gdef\lst@FVConvertO@a#1{%
    \lst@lAddTo\lst@arg{{#1}}\advance\@tempcnta\m@ne
    \lst@FVConvertO@}%
%    \end{macrocode}
% 由于 |\@ifnextchar\bgroup| 可能失败，我们这里必须使用 |\ifcat|。
% 缺陷由 \lsthelper{Denis~Girou}{1999/07/26}{fancyvrb=true 且`第二个
% 命令字符'不是 \{ 时无法工作}报告。
% 不过我们不像 |\@ifnextchar| 那样吞掉空格记号。
%    \begin{macrocode}
\gdef\lst@FVConvert@@{%
    \ifcat\noexpand\@let@token\bgroup \expandafter\lst@FVConvertArg
                                \else \expandafter\lst@FVConvert@ \fi}
%    \end{macrocode}
% 遇到这样的 catcode${}={}$1 字符时，我们转换该参数，并
% 连同分组定界符一起加到 |\lst@arg|。
% 我们还加入 |\lst@PrintToken|，它会在我们
% 忘记之前打印所有收集到的字符。
% 最后我们继续转换。
%    \begin{macrocode}
\gdef\lst@FVConvertArg#1{%
    {\let\lst@arg\@empty
     \lst@FVConvert#1\@nil
     \global\let\@gtempa\lst@arg}%
     \lst@lExtend\lst@arg{\expandafter{\@gtempa\lst@PrintToken}}%
     \lst@FVConvert}
%    \end{macrocode}
%    \begin{macrocode}
\gdef\lst@FVConvert@#1{%
    \ifx \@nil#1\else
       \if\relax\noexpand#1%
          \lst@lAddTo\lst@arg{\lst@OutputLostSpace\lst@PrintToken#1}%
       \else
          \lccode`\~=`#1\lowercase{\lst@lAddTo\lst@arg~}%
       \fi
       \expandafter\lst@FVConvert
    \fi}
%    \end{macrocode}
% 没有 |\bgroup| 时，我们查看是否已找到输入末尾，并
% 转换一个记号（（非）活动字符或控制序列）。
%    \begin{macrocode}
\gdef\lst@FVConvert@#1{%
    \ifx \@nil#1\else
       \if\relax\noexpand#1%
          \lst@lAddTo\lst@arg{\lst@OutputLostSpace\lst@PrintToken#1}%
%    \end{macrocode}
% 这里我们检查已注册的带参数命令，并按需
% 设置 |\@tempcnta| 的值。
%    \begin{macrocode}
          \def\lst@temp##1,#1##2,##3##4\relax{%
              \ifx##3\@empty \else \@tempcnta##2\relax \fi}%
          \expandafter\lst@temp\lst@FVcmdparams,#1\z@,\@empty\relax
       \else
          \lccode`\~=`#1\lowercase{\lst@lAddTo\lst@arg~}%
       \fi
       \expandafter\lst@FVConvertO@
    \fi}
%    \end{macrocode}
% \end{macro}
%
%    \begin{macrocode}
\lst@EndAspect
%</misc>
%    \end{macrocode}
% \end{aspect}
%
%
% \subsection{Omega 支持}
%
% \begingroup
% $\Omega$ 支持看起来很简单——我希望它至少在有些情况下有效。
%    \begin{macrocode}
%<*kernel>
%    \end{macrocode}
%    \begin{macrocode}
\@ifundefined{ocp}{}
    {\lst@AddToHook{OutputBox}%
         {\let\lst@ProcessLetter\@firstofone
          \let\lst@ProcessDigit\@firstofone
          \let\lst@ProcessOther\@firstofone}}
%    \end{macrocode}
%    \begin{macrocode}
%</kernel>
%    \end{macrocode}
% \endgroup
%
%
% \subsection{\packagename{LGrind}}
%
% \begin{aspect}{lgrind}
% \begin{macro}{\lst@LGGetNames}
% 用于从 |\lst@arg|（\packagename{LGrind} 的定义）中
% 提取语言名。
%    \begin{macrocode}
%<*misc>
\lst@BeginAspect[keywords,comments,strings,language]{lgrind}
%    \end{macrocode}
%    \begin{macrocode}
\gdef\lst@LGGetNames#1:#2\relax{%
    \lst@NormedDef\lstlang@{#1}\lst@ReplaceInArg\lstlang@{|,}%
    \def\lst@arg{:#2}}
%    \end{macrocode}
% \end{macro}
%
% \begin{macro}{\lst@LGGetValue}
% 在 |\lst@LGvalue| 中返回列表 |\lst@arg| 所给出的
% 能力 |#1| 的值。如果找不到 |#1|，我们就有 |\lst@if|=|\iffalse|。
% 否则它为真，并把``cap=value''对从列表中移除。
% 首先我们测试 |#1|，并
%    \begin{macrocode}
\gdef\lst@LGGetValue#1{%
    \lst@false
    \def\lst@temp##1:#1##2##3\relax{%
        \ifx\@empty##2\else \lst@LGGetValue@{#1}\fi}
    \expandafter\lst@temp\lst@arg:#1\@empty\relax}
%    \end{macrocode}
% 必要时移除该对。
%    \begin{macrocode}
\gdef\lst@LGGetValue@#1{%
    \lst@true
    \def\lst@temp##1:#1##2:##3\relax{%
        \@ifnextchar=\lst@LGGetValue@@{\lst@LGGetValue@@=}##2\relax
        \def\lst@arg{##1:##3}}%
    \expandafter\lst@temp\lst@arg\relax}
\gdef\lst@LGGetValue@@=#1\relax{\def\lst@LGvalue{#1}}
%    \end{macrocode}
% \end{macro}
%
% \begin{macro}{\lst@LGGetComment}
% 如果存在类型 |#1| 的注释且不是注释行，就把注释定界符
% （括在花括号内）存入 |#2|。否则 |#2| 为空。
%    \begin{macrocode}
\gdef\lst@LGGetComment#1#2{%
    \let#2\@empty
    \lst@LGGetValue{#1b}%
    \lst@if
        \let#2\lst@LGvalue
        \lst@LGGetValue{#1e}%
        \ifx\lst@LGvalue\lst@LGEOL
            \edef\lstlang@{\lstlang@,commentline={#2}}%
            \let#2\@empty
        \else
            \edef#2{{#2}{\lst@LGvalue}}%
        \fi
    \fi}
%    \end{macrocode}
% \end{macro}
%
% \begin{macro}{\lst@LGGetString}
% 对字符串定界符做同样的事，但它不`返回'任何值。
%    \begin{macrocode}
\gdef\lst@LGGetString#1#2{%
    \lst@LGGetValue{#1b}%
    \lst@if
        \let#2\lst@LGvalue
        \lst@LGGetValue{#1e}%
        \ifx\lst@LGvalue\lst@LGEOL
            \edef\lstlang@{\lstlang@,morestringizer=[l]{#2}}%
        \else
%    \end{macrocode}
% 我们必须检查 |\e|，即我们是否要使用 \texttt{双写}或
% \texttt{反斜杠}的字符串化器。
%    \begin{macrocode}
            \ifx #2\lst@LGvalue
                \edef\lstlang@{\lstlang@,morestringizer=[d]{#2}}%
            \else
                \edef\lst@temp{\lst@LGe#2}%
                \ifx \lst@temp\lst@LGvalue
                    \edef\lstlang@{\lstlang@,morestringizer=[b]{#2}}%
                \else
                    \PackageWarning{Listings}%
                    {String #2...\lst@LGvalue\space not supported}%
                \fi
            \fi
        \fi
    \fi}
%    \end{macrocode}
% \end{macro}
%
% \begin{macro}{\lst@LGDefLang}
% 定义由 |\lst@arg| 给出的语言，即定义部分，
% 以及 |\lst@language@|，即语言名。首先我们从
% |\lst@arg| 中移除不要的东西，例如把 |:\ :| 换成 |:|。
%    \begin{macrocode}
\gdef\lst@LGDefLang{%
    \lst@LGReplace
    \let\lstlang@\empty
%    \end{macrocode}
% 取得 friends 的关键词和取值。
%    \begin{macrocode}
    \lst@LGGetValue{kw}%
    \lst@if
        \lst@ReplaceInArg\lst@LGvalue{{ },}%
        \edef\lstlang@{\lstlang@,keywords={\lst@LGvalue}}%
    \fi
%    \end{macrocode}
%    \begin{macrocode}
    \lst@LGGetValue{oc}%
    \lst@if
        \edef\lstlang@{\lstlang@,sensitive=f}%
    \fi
%    \end{macrocode}
%    \begin{macrocode}
    \lst@LGGetValue{id}%
    \lst@if
        \edef\lstlang@{\lstlang@,alsoletter=\lst@LGvalue}%
    \fi
%    \end{macrocode}
% 现在我们取得注释定界符，并根据有两个还是四个
% 定界符把它们用作单注释或双注释。
% 注意 |\lst@LGGetComment| 处理注释行。
%    \begin{macrocode}
    \lst@LGGetComment a\lst@LGa
    \lst@LGGetComment c\lst@LGc
    \ifx\lst@LGa\@empty
        \ifx\lst@LGc\@empty\else
            \edef\lstlang@{\lstlang@,singlecomment=\lst@LGc}%
        \fi
    \else
        \ifx\lst@LGc\@empty
            \edef\lstlang@{\lstlang@,singlecomment=\lst@LGa}%
        \else
            \edef\lstlang@{\lstlang@,doublecomment=\lst@LGc\lst@LGa}%
        \fi
    \fi
%    \end{macrocode}
% 现在我们解析字符串化器。
%    \begin{macrocode}
    \lst@LGGetString s\lst@LGa
    \lst@LGGetString l\lst@LGa
%    \end{macrocode}
% 我们测试续行能力，并
%    \begin{macrocode}
    \lst@LGGetValue{tc}%
    \lst@if
        \edef\lstlang@{\lstlang@,lgrindef=\lst@LGvalue}%
    \fi
%    \end{macrocode}
% 定义该语言。
%    \begin{macrocode}
    \expandafter\xdef\csname\@lst LGlang@\lst@language@\endcsname
        {\noexpand\lstset{\lstlang@}}%
%    \end{macrocode}
% 最后我们告知用户所有被忽略的能力。
%    \begin{macrocode}
    \lst@ReplaceInArg\lst@arg{{: :}:}\let\lst@LGvalue\@empty
    \expandafter\lst@LGDroppedCaps\lst@arg\relax\relax
    \ifx\lst@LGvalue\@empty\else
        \PackageWarningNoLine{Listings}{Ignored capabilities for
            \space `\lst@language@' are\MessageBreak\lst@LGvalue}%
    \fi}
%    \end{macrocode}
% \end{macro}
%
% \begin{macro}{\lst@LGDroppedCaps}
% 只是丢掉先前的值，并把下一个能力名追加到
% |\lst@LGvalue|.
%    \begin{macrocode}
\gdef\lst@LGDroppedCaps#1:#2#3{%
    \ifx#2\relax
        \lst@RemoveCommas\lst@LGvalue
    \else
        \edef\lst@LGvalue{\lst@LGvalue,#2#3}%
        \expandafter\lst@LGDroppedCaps
    \fi}
%    \end{macrocode}
% \end{macro}
%
% \begin{macro}{\lst@LGReplace}
% \begin{macro}{\lst@LGe}
% 我们把`转义的 \verb!:^$|!'换成 catcode 11 版本，把其他字符串
% 换成某种短版本（这对于让上面的
% 定义生效是必要的）。
%    \begin{macrocode}
\begingroup
\catcode`\/=0
\lccode`\z=`\:\lccode`\y=`\^\lccode`\x=`\$\lccode`\v=`\|
\catcode`\\=12\relax
/lowercase{%
/gdef/lst@LGReplace{/lst@ReplaceInArg/lst@arg
    {{\:}{z }{\^}{y}{\$}{x}{\|}{v}{ \ }{ }{:\ :}{:}{\ }{ }{\(}({\)})}}
/gdef/lst@LGe{\e}
}
/endgroup
%    \end{macrocode}
% \end{macro}\end{macro}
%
% \begin{macro}{\lst@LGRead}
% 读取一个语言定义，并在找到正确的那一个时
% 定义该语言。
%    \begin{macrocode}
\gdef\lst@LGRead#1\par{%
    \lst@LGGetNames#1:\relax
    \def\lst@temp{endoflanguagedefinitions}%
    \ifx\lstlang@\lst@temp
        \let\lst@next\endinput
    \else
        \expandafter\lst@IfOneOf\lst@language@\relax\lstlang@
            {\lst@LGDefLang \let\lst@next\endinput}%
            {\let\lst@next\lst@LGRead}%
    \fi
    \lst@next}
%    \end{macrocode}
% \end{macro}
%
% \begin{lstkey}{lgrindef}
% 我们只需请求该语言，并
%    \begin{macrocode}
\lst@Key{lgrindef}\relax{%
    \lst@NormedDef\lst@language@{#1}%
    \begingroup
    \@ifundefined{lstLGlang@\lst@language@}%
        {\everypar{\lst@LGRead}%
         \catcode`\\=12\catcode`\{=12\catcode`\}=12\catcode`\%=12%
         \catcode`\#=14\catcode`\$=12\catcode`\^=12\catcode`\_=12\relax
         \input{\lstlgrindeffile}%
        }{}%
    \endgroup
%    \end{macrocode}
% 选择它或发出错误消息。
%    \begin{macrocode}
    \@ifundefined{lstLGlang@\lst@language@}%
        {\PackageError{Listings}%
         {LGrind language \lst@language@\space undefined}%
         {The language is not loadable. \@ehc}}%
        {\lsthk@SetLanguage
         \csname\@lst LGlang@\lst@language@\endcsname}}
%    \end{macrocode}
% \end{lstkey}
%
% \begin{macro}{\lstlgrindeffile}
% 只包含文件名。
%    \begin{macrocode}
\@ifundefined{lstlgrindeffile}
    {\lst@UserCommand\lstlgrindeffile{lgrindef.}}{}
%    \end{macrocode}
% \end{macro}
%
%    \begin{macrocode}
\lst@EndAspect
%</misc>
%    \end{macrocode}
% \end{aspect}
%
%
% \subsection{\packagename{hyperref}}
%
% \begin{aspect}{hyper}
%    \begin{macrocode}
%<*misc>
\lst@BeginAspect[keywords]{hyper}
%    \end{macrocode}
%
% \begin{lstkey}{hyperanchor}
% \begin{lstkey}{hyperlink}
% 分别确定用于设置锚点和链接的宏。
%    \begin{macrocode}
\lst@Key{hyperanchor}\hyper@@anchor{\let\lst@hyperanchor#1}
\lst@Key{hyperlink}\hyperlink{\let\lst@hyperlink#1}
%    \end{macrocode}
% \end{lstkey}\end{lstkey}
% 同样，主要工作是一个特殊的工作流程。首先我们提取
% |\lst@token| 的内容，并为这个当前字符串取得一个空闲宏名
% （使用前缀 |lstHR@| 和一个数字后缀）。然后我们让这个
% 空闲宏等价于 |\@empty|，这样下次它就不会被用到。
%    \begin{macrocode}
\lst@InstallKeywords{h}{hyperref}{}\relax{}
    {\begingroup
         \let\lst@UM\@empty \xdef\@gtempa{\the\lst@token}%
     \endgroup
     \lst@GetFreeMacro{lstHR@\@gtempa}%
     \global\expandafter\let\lst@freemacro\@empty
%    \end{macrocode}
% |\@tempcnta| 是空闲宏的后缀。我们在这里用它引用
% 同一字符串的最后一次出现。为此我们重新定义输出
% 宏 |\lst@alloverstyle| 以设置一个锚点\ldots
%    \begin{macrocode}
     \@tempcntb\@tempcnta \advance\@tempcntb\m@ne
     \edef\lst@alloverstyle##1{%
         \let\noexpand\lst@alloverstyle\noexpand\@empty
         \noexpand\smash{\raise\baselineskip\hbox
             {\noexpand\lst@hyperanchor{lst.\@gtempa\the\@tempcnta}%
                                       {\relax}}}%
%    \end{macrocode}
% \ldots\ 以及指向最后一次出现（如果有）的链接。
%    \begin{macrocode}
         \ifnum\@tempcnta=\z@ ##1\else
             \noexpand\lst@hyperlink{lst.\@gtempa\the\@tempcntb}{##1}%
         \fi}%
    }
    od
%    \end{macrocode}
%
%    \begin{macrocode}
\lst@EndAspect
%</misc>
%    \end{macrocode}
% \end{aspect}
%
%
% \section{后记}
%
% \begingroup
%    \begin{macrocode}
%<*kernel>
%    \end{macrocode}
% 每个选项都把方面名加到 |\lst@loadaspects| 中，或从该数据宏中移除它。
%    \begin{macrocode}
\DeclareOption*{\expandafter\lst@ProcessOption\CurrentOption\relax}
\def\lst@ProcessOption#1#2\relax{%
    \ifx #1!%
        \lst@DeleteKeysIn\lst@loadaspects{#2}%
    \else
        \lst@lAddTo\lst@loadaspects{,#1#2}%
    \fi}
%    \end{macrocode}
% 以下方面默认加载。
%    \begin{macrocode}
\@ifundefined{lst@loadaspects}
  {\def\lst@loadaspects{strings,comments,escape,style,language,%
      keywords,labels,lineshape,frames,emph,index}%
  }{}
%    \end{macrocode}
% 我们加载补丁文件，\ldots
%    \begin{macrocode}
\InputIfFileExists{lstpatch.sty}{}{}
%    \end{macrocode}
% \ldots\ 处理选项，\ldots
%    \begin{macrocode}
\let\lst@ifsavemem\iffalse
\DeclareOption{savemem}{\let\lst@ifsavemem\iftrue}
\DeclareOption{noaspects}{\let\lst@loadaspects\@empty}
\ProcessOptions
%    \end{macrocode}
% \ldots\ 并加载各方面。
%    \begin{macrocode}
\lst@RequireAspects\lst@loadaspects
\let\lst@loadaspects\@empty
%    \end{macrocode}
% 如果存在，我们就选择空样式和空语言。
%    \begin{macrocode}
\lst@UseHook{SetStyle}\lst@UseHook{EmptyStyle}
\lst@UseHook{SetLanguage}\lst@UseHook{EmptyLanguage}
%    \end{macrocode}
% 最后我们加载配置文件。\lsthelper{Ulrike~Fischer}{2020/11/13}
% {.cfg 文件以错误的 catcode 读取}指出这必须在
% 正确的 catcode 下进行。目前 |^^M| 的 catcode 是 9，这是
% 错误的。所以我们在加载文件之前把 catcode
% 重置为正确值：
%    \begin{macrocode}
%<info>\lst@ReportAllocs
\lst@RestoreCatcodes%
\InputIfFileExists{listings.cfg}{}{}
\InputIfFileExists{lstlocal.cfg}{}{}
%    \end{macrocode}
%    \begin{macrocode}
%</kernel>
%    \end{macrocode}
% \endgroup
%
%
% \section{历史}
% \begingroup\small
% 这里只列出主要改动。命令和键的引入版本号
% 在各指南的源码中，这使本历史记录相当
% 简短。
% \renewcommand\labelitemi{--}
% \begin{itemize}
% \item[0.1] 自 1996/03/09
%     \item 测试版本，看看该宏包是否可行
% \item[0.11] 自 1996/08/19
%\iffalse
%     \item 额外的空白选项（= 语言）
%\fi
%     \item 改进的对齐
% \item[0.12] 自 1997/01/16
%     \item 近乎`完美'的对齐
% \item[0.13] 自 1997/02/11
%\iffalse
%     \item 新增语言：Eiffel、Fortran 90、Modula-2、Pascal XSC
%\fi
%     \item 按需加载：语言相关的宏移到驱动文件
%     \item 注释现在只声明一次，不再为每种语言重复实现
%                 （这使 \TeX\ 源码更易读）
% \item[0.14] 自 1997/02/18
%     \item 重写用户指南，实现指南使用 macro 环境
%     \item 实现（不）区分大小写以及多种字符串类型，
%                 即 Modula-2 同时处理两种字符串类型：单引号与双引号
%\iffalse
%     \item 注释声明对用户开放
%     \item 宏包兼容 \verb!german.sty!
%\fi
% \item[0.15] 自 1997/04/18
%\iffalse
%     \item 新增语言：Java、Turbo Pascal
%\fi
%     \item 宏包从 \packagename{listing} 改名为 \packagename{listings}，
%                 因为前者已被占用
% \item[0.16] 自 1997/06/01
%\iffalse
%     \item 在 |\lst@SkipToFirst| 中把`$<$'改为`$>$'
%     \item 移除缺陷：|\lst@Init| 必须放在 |\lst@SkipToFirst| 之前
%\fi
%     \item 重写列表环境
% \item[0.17] 自 1997/09/29
%\iffalse
%     \item |\spreadlisting| 现在工作正确（例如页码不再右移）
%\fi
%     \item 提速（快速`参数为空'判断、除一个之外全部去掉 |\long|、
%                 更快的 \verb!\lst@GotoNextTabStop! 等）
%     \item 改进宽`其他'字符（例如 $==$）的对齐
%\iffalse
%     \item 许多新语言：Ada、Algol、Cobol、Comal 80、Elan、Fortran 77、
%                 Lisp、Logo、Matlab、Oberon、Perl、PL/I、Simula、SQL、\TeX
%\fi
% \item[0.18 之前] 自 1998/03/24（未发布）
%\iffalse
%     \item 移除关于 |\labelstyle|（变成 \keyname{numberstyle}）的缺陷
%                 （现在 oldstylenum 示例可以工作）
%\fi
%     \item 字符类的实验性实现
% \item[0.19] 自 1998/11/09
%     \item 字符类和新的 \lst-方面看来效果良好
%     \item 用户接口使用 \packagename{keyval} 宏包
%     \item \packagename{fancyvrb} 支持
% \item[0.20] 自 1999/07/12
%     \item 新的关键词检测机制
%     \item 新方面：\aspectname{writefile}、\aspectname{breaklines}、
%                 题注、\aspectname{html}
%\iffalse
%     \item 改进的 \packagename{fancyvrb} 支持
%\fi
%     \item 所有方面位于单个文件，语言驱动目前位于
%                 两个文件
% \item[0.21] 2000/08/23
%     \item 全新的用户指南
%     \item 实验性格式定义
%     \item 关键词类被族（family）取代
%     \item 动态模式
% \item[1.0$\beta$] 2001/09/21
%     \item 键名与 \packagename{fancyvrb} 同步
%     \item 扩展 \aspectname{frames} 方面
%     \item 新的输出概念（延迟与合并）
% \item[1.0] 2002/04/01
%     \item 更新所有文档章节，包括开发者指南
%     \item 统一了定界符
% \item[1.1] 2003/06/21
%     \item 缺陷修复版，含一些新键
% \item[1.2] 2004/02/13
%     \item 缺陷修复版，含两个新键和新节 \ref{rArbitraryLinerangeMarkers}
% \item[1.3] 2004/09/07
%     \item 又一个缺陷修复版，符合 LPPL-1.3
% \item[1.4] 2007/02/26
%     \item 许多缺陷修复，以及新的维护者
%     \item 若干新的和更新的语言定义
%     \item 许多小的文档改进
%     \item 新键、多字符字符串定界符、短行内列表等。
% \item[1.5] 2013/06/27
%     \item 新的维护者
% \item[1.6] 2015/05/05
%     \item 加入关于使用 |\lstinline|\oarg{key=value list}\marg{source code} 的讨论
%     \item 加入``已知错误与变通方法''一节。
% \item[1.7] 2018/09/02
%     \item 一些新的或更新的语言定义
%     \item 若干错误修正
% \item[1.8] 自 2019/02/27 起
%     \item 修正并启用了 |inputpath| 选项
%     \item 一些新的或更新的语言定义
%     \item 若干错误修正
%     \item 引入 |\lstlistingnamestyle|
% \item[1.9] 自 2023/02/27 起
%     \item 有望修正了显示可见空格这一长期存在的错误行为
%     \item 推广了 |linerange| 的用法
%     \item 引入 |consecutivenumbers| 键
%     \item 不再允许（上）题注与源代码之间的分页
%     \item 配置文件 |listings.cfg| 和 |lstlocal.cfg| 以正确的 catcode
%         读取
%     \item 一些文档和排版改进。
% \item[1.10] 自 2024/02/15 起
%     \item 如果加载了 |tocbasic| 就使用它，以改进与
%     KOMA-Script 的兼容性并增强功能。
%     \item 区分 Python 2 与 Python 3
%     \item 加入对 Scala 3.0 的初步支持
%     \item 使 |\lstinline{...}| 在 tabular 单元格中可用
%     \item 实现了直双引号
% \end{itemize}
% \endgroup
%
%
% \Finale
%
%    \begin{macrocode}
%<*dev/null>
%    \end{macrocode}
%%% Local Variables:
%%% ispell-local-dictionary: "english"
%%% End:
%    \begin{macrocode}
%<*dev/null>
%    \end{macrocode}
\endinput
